% Self-contained O--QFT/BRST article with conditional fixed-graph convergence
% results and a separately declared O--QED auxiliary test theory.
% Compile with: latexmk -pdf -interaction=nonstopmode -halt-on-error UV-preview.tex
% The interacting completion, physical reconstruction, and all-order graph
% estimates remain hypotheses wherever not explicitly constructed.
\documentclass[final,5p,times,twocolumn,numbers,nopreprintline]{elsarticle}
\usepackage[T1]{fontenc}
\usepackage[utf8]{inputenc}
\renewcommand{\ttdefault}{lmtt} % Outline font for DOI and URL text.
\usepackage{amsmath,amssymb,amsthm,mathtools,bm}
\usepackage{mathrsfs}
\usepackage{slashed}
\usepackage{microtype}
\usepackage{booktabs,tabularx,array}
\IfFileExists{flushend.sty}{\usepackage[nospread]{flushend}}{}
\usepackage{enumitem}
\usepackage{xcolor}
\usepackage[hidelinks]{hyperref}
\usepackage{bookmark}
\hypersetup{
 pdftitle={Conditional All-Orders UV Finiteness: Non-Abelian BRST Closure, Analyticity, and Crossing in Omniverse Quantum Field Theory},
 pdfauthor={Giovanni Joseph Chiappone},
 pdfsubject={Fixed-plane BRST and ultraviolet analysis with an O-QED application},
 pdfkeywords={BRST, Slavnov-Taylor, ultraviolet finiteness, O-QED, geometric time}}
\numberwithin{equation}{section}
\newtheorem{theorem}{Theorem}[section]
\newtheorem{proposition}[theorem]{Proposition}
\newtheorem{lemma}[theorem]{Lemma}
\newtheorem{corollary}[theorem]{Corollary}
\theoremstyle{definition}
\newtheorem{definition}[theorem]{Definition}
\newtheorem{assumption}[theorem]{Assumption}
\theoremstyle{remark}
\newtheorem{remark}[theorem]{Remark}
\newcommand{\dd}{\mathop{}\!\mathrm d}
\newcommand{\ii}{\mathrm i}
\newcommand{\bb}{\mathfrak b}
\newcommand{\BR}{\mathsf s}
\newcommand{\Ph}{\Phi}
\newcommand{\Tr}{\operatorname{tr}}
\newcommand{\PV}{\operatorname{PV}}
\newcommand{\Disc}{\operatorname{Disc}}
\newcommand{\Cut}{\operatorname{Cut}}
\newcommand{\supp}{\operatorname{supp}}
\newcommand{\diag}{\operatorname{diag}}
\newcommand{\rank}{\operatorname{rank}}
\newcommand{\sgn}{\operatorname{sgn}}
\newcommand{\erfc}{\operatorname{erfc}}
\newcommand{\erfcx}{\operatorname{erfcx}}
\newcommand{\RR}{\mathbb R}
\newcommand{\CC}{\mathbb C}
\newcommand{\one}{\mathbf 1}
\newcommand{\vct}[1]{\boldsymbol{#1}}
\newcommand{\norm}[1]{\left\lVert #1\right\rVert}
\newcommand{\abs}[1]{\left\lvert #1\right\rvert}
\newcommand{\cGT}{c_{\mathrm{GT}}}
\newcommand{\ee}{e_{\bb}}
\newcommand{\epsvec}{\hat\epsilon_{\bb}}
\newcommand{\je}{j_{\bb}}
\newcommand{\embed}{\varphi_{\ee}}
\newcommand{\ExpDir}{\operatorname{Exp}_{\ee}}
\newcommand{\SqrtDir}{\operatorname{Sqrt}_{\ee}}
\newcommand{\LogDir}{\operatorname{Log}_{\ee}}
\newcommand{\EoneDir}{(E_1)_{\ee}}
\newcommand{\Mzero}{M_0}
\newcommand{\bmu}{\boldsymbol{\mu}}
\newcommand{\sMand}{\mathsf{s}}
\newcommand{\tMand}{\mathsf{t}}
\newcommand{\uMand}{\mathsf{u}}
% Match the attached geometric-time article's Elsevier 5p typography.
\setlist{nosep,leftmargin=*}
\setlength{\emergencystretch}{2em}
\linespread{0.936}
\setlength{\dbltextfloatsep}{6pt plus 1pt minus 1pt}
\setlength{\textfloatsep}{7pt plus 1pt minus 1pt}
\setlength{\abovecaptionskip}{4pt}
\allowdisplaybreaks[1]
\interfootnotelinepenalty=10000
\raggedbottom
\AtBeginDocument{%
 \setlength{\abovedisplayskip}{7pt plus 2pt minus 2pt}%
 \setlength{\belowdisplayskip}{7pt plus 2pt minus 2pt}%
 \setlength{\abovedisplayshortskip}{4pt plus 2pt minus 2pt}%
 \setlength{\belowdisplayshortskip}{4pt plus 2pt minus 2pt}%
}
% Numeric references are embedded; no BibTeX or external figures are needed.
% Published articles below cite their versions of record. For books,
% original publication, corrected reissues, and later electronic formats
% are distinguished rather than treated as interchangeable editions.
% DLMF entries pin the publicly identified release and direct equation links.
% arXiv entries identify the latest revisions visible on 6 November 2026.
% For author-supplied manuscripts, the associated Zenodo identifiers are
% given without assuming an unverified DOI-to-file or version match.
% elsarticle already includes Appendix in appendix reference labels.
\journal{}
\widowpenalty=10000
\clubpenalty=10000
\displaywidowpenalty=10000

\begin{document}
\begin{frontmatter}
\title{Conditional All-Orders UV Finiteness:\\
Non-Abelian BRST Closure, Analyticity, and Crossing\\
in Omniverse Quantum Field Theory}
\author[aff1]{Giovanni Joseph Chiappone\fnref{orcid}}
\ead{giovannichiappone252@gmail.com}
\address[aff1]{Independent researcher}
\fntext[orcid]{ORCID: \href{https://orcid.org/0009-0005-1689-463X}{0009-0005-1689-463X}.}
\begin{abstract}
We derive sufficient conditions for ultraviolet convergence of
finite Euclidean graphs in fixed-plane Omniverse Quantum Field Theory.
Positive Laplace--Stieltjes kernels yield a uniform Gaussian estimate
under explicit inverse moments and infrared envelopes; a selected-line
criterion permits distinct moment orders and massless selected
propagators when the remaining factors are controlled. An explicit two-loop comparison has
reproducible real- and complex-momentum stability bounds. Residual
coercivity also permits some nonpolynomial completed vertices.
The phrase ``all orders'' refers to each admitted finite graph, not
convergence of the perturbation series.

An anomaly-free completed BRST/BV functional supports the conditional
Slavnov--Taylor, background Ward, and Nielsen identities. We keep the
vectorial boundary prescription, both carrier characters, and source
reciprocity distinct from physical positivity and scattering unitarity.

For standalone support-diagonal O--QED, the normalized square-root
kernel satisfies every finite inverse-moment test. We quantify
positive-depth derivative bounds, exponential momentum tails,
strong-depth crossovers, and a finite zero-depth limit for a common
subtracted scalar bubble. We also identify the universal logarithmic
finite part of its unsubtracted zero-depth asymptotic. These are
kernel and comparison-model results. A two-Ward-vertex polarization
term fails individual absolute convergence, and an explicitly declared
constant-background contact test fails a grouped absolute bound.
A physical interaction completion therefore still requires a common
vertex/contact estimate, compatible source and BRST structure,
reflection positivity, and cut--discontinuity matching. Conditional
pole and response certificates do not replace those constructions.
\end{abstract}
\begin{keyword}
BRST symmetry \sep Slavnov--Taylor identities \sep ultraviolet finiteness
\sep geometric-time locking \sep Laplace--Stieltjes kernels
\sep reflection positivity \sep O--QED \sep Omniverse quantum field theory
\end{keyword}
\end{frontmatter}

\section{Introduction}
\label{sec:introduction}

A positive Laplace--Stieltjes multiplier can suppress large Euclidean
loop momenta by expressing the damping as an average of Gaussian
factors. In an interacting theory, its effectiveness depends on the
vertices accompanying the modified propagators. This article gives
conditions for that suppression to survive in complete graphs of
Omniverse Quantum Field Theory (O--QFT) on one fixed sector plane.
The central result is a uniform absolute-convergence theorem for
finite loop graphs. We then state the additional conditions under
which the same completed action supports BRST identities and a
physical interpretation of its amplitudes.

The geometric input is a selected \(1+3\)-dimensional Lorentzian plane
in an ambient geometry with one or more geometric-time directions.
Geometric-time locking (GTL) restricts fields to that plane; the label
\(\bb=(s,b_s)\) records its sector and affine placement.
The kernel input is
\begin{equation}
\begin{aligned}
 \Ph(z)&=\int_{[0,\infty)}e^{-\tau z}\rho(\dd\tau),\qquad \rho\geq0,\\
 0&<\Mzero:=\rho([0,\infty))<\infty .
\end{aligned}
 \label{eq:bernstein}
\end{equation}
It acts on the nonnegative Euclidean kinetic operator of an uncut
virtual line. The common-depth benchmark uses fixed \(a>0\); the edgewise
model data are specified in \eqref{eq:virtual-edge-depths}. The
construction does not define physical amplitudes by removing \(a\).
The physical kernel convention of Ref.~\cite{ChiapponeMonograph}
selects \(0<\Ph(0)\leq1\). The comparison estimates also admit
any finite positive measure mass \(M_0\) in \eqref{eq:bernstein}.

The geometric and field-theoretic construction is developed in
Ref.~\cite{ChiapponeMonograph}. Here the fixed-plane estimates and
functional identities are derived with their assumptions stated
explicitly. The kernel representation \eqref{eq:bernstein} supplies
the Euclidean estimates. Existence of a compatible completed action
and the hypotheses of physical reconstruction are separate requirements.
The cited reconstruction and crossing results are invoked only within
their stated domains.

The worked application in Section~\ref{sec:OQED-application} uses the
standalone support-diagonal O--QED action and normalized internal kernel
of Ref.~\cite{ChiapponeIgnition}. We evaluate its inverse moments,
specialize the graph criteria, and test the interaction requirement
rather than assuming it. The kernel certificate is affirmative, while
a two-vertex vacuum term satisfying the dressed Ward identities fails
individual-graph convergence. We formulate a common-routing test for
its contact-completed sum and examine a separately declared
constant-background contact realization. This additional test is not
attributed to the source as its completed theory. The medium and RT
response formulas state what a verified UV certificate would control,
not a calibrated plasma coefficient.

\subsection{Three mathematical questions}

The first question is \emph{convergence}. On a flat Euclidean plane, a scalar
covariance has the form
\begin{equation}
 C_{\Ph,E}(q)=\frac{\Ph(a^2q^2)}{q^2+m^2}.
 \label{eq:covariance-intro}
\end{equation}
Using \eqref{eq:bernstein} expresses the damping as a positive
Gaussian average. Control of the measure near \(\tau=0\) bounds
large loop momenta. Tensor amplitudes are bounded in absolute value;
their numerators need not be positive. Section~\ref{sec:uv} and
Appendix~\ref{sec:appendix-gaussian} give the global estimate and a
conditional forest criterion based on simultaneous sector bounds.

The second question is \emph{gauge symmetry}. In a non-Abelian theory, a
nonconstant kinetic modification generally changes the vertices required
by the master equation. The condition
\(\BR\Ph(a^2\mathscr L_{\bb})=0\), with \(\BR\) the BRST differential,
says that the prescribed background multiplier is inert under quantum
BRST. It does not say that a graph expansion with altered propagators
and arbitrary unchanged vertices satisfies the Slavnov--Taylor identity.
Section~\ref{sec:brst} derives the identity from a completed action,
including its nonminimal BV sector.

The third question is \emph{physical reconstruction}: whether the
Euclidean correlations describe physical states and scattering
amplitudes. Positive heat-kernel weights alone do not establish
Osterwalder--Schrader reflection positivity. A prescribed replacement
of a cut line by an undressed on-shell projector must also be shown
to reproduce the discontinuity of the continued graph.
Section~\ref{sec:analytic} states these additional requirements,
and Section~\ref{sec:examples} tests them in explicit models.

\subsection{Assumptions and scope}
\label{sec:assumptions}

The results use the following hypotheses in the combinations stated
below. When convergence and symmetry are claimed together, both sets
of hypotheses must hold for the same completed action.

\begin{description}[leftmargin=1.6cm,style=nextline]
\item[Geometry.]
One fixed locked support \(\bb\), its label-diagonal
restriction, and the independent A2/GTP momentum binder are used.
No sum over supports or transition between supports is implicit. Its Euclidean comparison
operator is nonnegative on a specified domain.

\item[Kernel.]
Each covered uncut internal covariance has a finite positive
representation \eqref{eq:bernstein}. A gap or quantitative inverse
moments are imposed whenever ultraviolet convergence is claimed.
The convenient normalization \(\Ph(0)=1\) is optional for convergence.

\item[Vertices and IR.]
The same completed vertices used in the symmetry identities satisfy
the growth bounds of the ultraviolet theorem. Positive masses or a
separately specified integrable infrared envelope control the long
proper-time region. External data lie in a compact domain on which
these bounds are uniform.

\item[Quantum symmetry.]
An off-shell nilpotent BRST differential, a completed BV action,
compatible source couplings and covariance exchange symmetry, and an
invariant anomaly-free measure
exist on the stated perturbative domain in both carrier characters.
The identities concern this complete system with simultaneous
projection of all fields, antifields, sources, and coefficients.

\item[Continuation.]
A declared continuation of the kernels, the plane-resolved
vectorial boundary datum, and a common integrable majorant are
available on the complex domain actually used.
Contour deformations, when used, have their own boundary and
no-pinch estimates.

\item[Physical completion.]
Reflection positivity on the observable algebra, the other
reconstruction axioms, isolated asymptotic poles, and the matching of
cuts to discontinuities are additional requirements. Any optical
identity also uses scattering unitarity and a complete physical state
resolution, or the explicitly stated equivalent cut-factorization
hypotheses. They are invoked
only in the results that explicitly assume them.
\end{description}

There is no implication from finiteness of each coefficient to
convergence of the infinite perturbation series. Nor does a fixed-plane
estimate control an unrestricted integration over backgrounds, an
infinite population of supports, or support-changing groupoid kernels.
Those extensions require estimates for the added variables. For an
independent mathematical audit, the essential distinction is between
\emph{proving a convergence implication from supplied graph bounds}
and \emph{deriving those bounds for every graph of one specified
interacting action}. Only the former is established by the abstract
convergence theorems; the latter is an application hypothesis.

\paragraph{Notation and organization}
Table~\ref{tab:notation} collects the symbols used across the three
parts of the argument. Section~\ref{sec:framework} fixes the geometry
and carrier conventions. The convergence proof is in
Section~\ref{sec:uv}, the functional identities in
Section~\ref{sec:brst}, and the analytic and physical conditions in
Section~\ref{sec:analytic}. Section~\ref{sec:examples} provides explicit
tests. Section~\ref{sec:OQED-application} applies those results to the
specified O--QED kernel, vertices, and response interface.
Appendix~\ref{sec:dependence} records the hypotheses needed to combine them.

\begin{table*}[t]
\centering
\small
\renewcommand{\arraystretch}{1.15}
\caption{Principal notation. Each support label is fixed unless a sum
or integral is explicitly written.}
\label{tab:notation}
\begin{tabularx}{\textwidth}{@{}>{\raggedright\arraybackslash}p{0.24\textwidth}X@{}}
\toprule
Symbol & Meaning\\
\midrule
\(K,\ \bb=(s,b_s)\) & Number of ambient geometric-time directions;
selected sector and scalar carrier-axis offset.\\
\(X_0,\ k_\parallel\) & Centred normalized carrier coordinate (the monograph's
\(G_{0\bb}\)) and its conjugate momentum coefficient.\\
\(a,\ \rho,\ \Ph\) & Fixed inverse-momentum scale, finite positive
Laplace measure, and its dimensionless transform.\\
\(a_{\rm ker}^{(p)},\ M_{\rm ker},\ \mu\) & O--QED depth parameter,
auxiliary kernel scale, and their product \(\mu=a_{\rm ker}^{(p)}M_{\rm ker}\);
not the electron pole mass or laboratory response length.\\
\(M_r[\rho]\) & Inverse moment of order \(r>0\); absence of an atom
at zero is stated separately.\\
\(\mathcal K,\ D\) & Compact external-momentum set and total
numerator-growth degree in the graph estimate.\\
\(\alpha_i,\tau_i,w_i\) & Schwinger parameter, dimensionless Laplace
parameter, and combined width \(w_i=\alpha_i+a^2\tau_i\).\\
\(\sigma_s,\sigma_\epsilon,\eta\) & Sector orientation, ordering
branch, and their product \(\eta=\sigma_s\sigma_\epsilon\).\\
\(\ee,\hat u_{\bb},\widehat u_k\) & Fixed pole carrier, branch-action unit
\(\hat u_{\bb}=u_{\epsilon,\bb}\), and running line unit.\\
\(\BR,\Sigma,\Gamma\) & BRST differential, completed action with
antifield sources, and one-particle-irreducible effective action.\\
\bottomrule
\end{tabularx}
\end{table*}

\section{Framework on a fixed sector plane}
\label{sec:framework}

\subsection{Locked geometry and canonical normalization}

Let \(\mathcal G\simeq\RR^K\), \(K\geq1\), be the geometric-time space
with its positive inner product, and let \(\hat G_{\bb}\in\mathcal G\)
be a unit carrier. On the ambient product with three spatial directions,
use the quadratic form
\[
 \dd s^2=\norm{\dd X_{\rm GT}}_{\mathcal G}^2-\norm{\dd\mathbf x}^2.
\]
A fixed support label is \(\bb=(s,b_s)\), where \(b_s\in\RR\)
is an affine displacement along the selected carrier. Its ambient offset is
\(B_{{\rm GT},\bb}=b_s\hat G_{\bb}\). We use \(X_0\) for the
\emph{centred, normalized carrier coordinate}, denoted \(G_{0\bb}\)
in Ref.~\cite{ChiapponeMonograph}; it is not the raw common-component
coordinate used in the monograph's diagonal chart. The fibrewise embedding is
\begin{equation}
 \iota_{\bb}(X_0,\mathbf x)
   =\bigl((X_0+b_s)\hat G_{\bb},\mathbf x\bigr),\qquad
 Q_{\bb}X_{\rm GT}=0,\quad
 Q_{\bb}=\one-\hat G_{\bb}\hat G_{\bb}^{\,T}.
 \label{eq:locking}
% Monograph source: eq:GTL-master-affine-embedding; eq:pre-post-GTL-affine-offset-pushforward.
\end{equation}
The label remains part of the event. In the disjoint union of labelled
fibres, \((\bb,X_0,\mathbf x)\) and
\((\bb',X'_0,\mathbf x')\) belong to distinct fibres when
\(\bb\ne\bb'\), even if their forgetful ambient images coincide.
Indeed, translating a complete carrier line along itself leaves its
unlabelled image unchanged. Equality of those images therefore does not
identify the fields, sources, states, or synchronization origins.

On a fixed fibre the pullback metric and the positive intrinsic measure are
\begin{equation}
 \iota_{\bb}^{*}(\dd s^2)=\dd X_0^2-\dd\mathbf x^2,
 \qquad
 \dd^4x_{\bb}=|\dd X_0|\,\dd^3\mathbf x.
 \label{eq:pullback-metric}
\end{equation}
Here \(X_0\) has dimensions of length. With
\(\sigma_N=+1\), \(\sigma_A=-1\), and the chosen synchronized
clock readout,
\[
 X_0=\sigma_s\cGT t,
 \qquad t=\frac{\sigma_s X_0}{\cGT},
 \qquad
 \vct c_{\bb}=\sigma_s\cGT\hat G_{\bb}.
\]
Thus \(\dd X_0=\sigma_s\cGT\dd t\), whereas
\(\dd^4x_{\bb}=\cGT|\dd t|\,\dd^3\mathbf x\) is positive
and contains no orientation sign. The common clock readout is distinct
from the directed carrier coordinate; the antiverse orientation does not
mean a backward-running physical clock.

In an equal-component chart with
\(\hat G=K^{-1/2}(1,\ldots,1)\) and
\(x^\Gamma=x^0\), the coordinate dictionary is
\[
\begin{aligned}
 X_0&=\sqrt K\,x^0-b_s,\\
 \dd^4x_{\bb}&=\sqrt K\,|\dd x^0|\,\dd^3\mathbf x,\\
 p_\parallel&=K^{-1/2}\sum_\Gamma p_\Gamma.
\end{aligned}
\]
The measure, density, Fourier phase, and external normalization must be
transformed together. In particular, one cannot insert \(Kp_0\)
into a Dirac slash while keeping the unrescaled denominator and measure.
Throughout the loop examples,
\begin{equation}
 k=(k_\parallel,\mathbf k),\qquad
 k^2=k_\parallel^2-\mathbf k^2,\qquad
 \slashed k=\gamma^\mu k_\mu
 \label{eq:intrinsic-momentum}
\end{equation}
are canonically normalized intrinsic quantities. The Clifford algebra
is \(\mathrm{Cl}(1,3)\); no additional gamma matrices or multiplicity
factors arise from this coordinate choice.

We use \(\hbar=1\) and momentum-dimension mass parameter \(m\). With
explicit units the shell parameter is \(m_{\rm phys}\cGT\), and the
length corresponding to the inverse-momentum kernel parameter \(a\)
is \(\ell_{\Ph}=\hbar a\). Proper-time variables
\(\alpha\) and \(a^2\tau\) have inverse-momentum-squared dimensions.
The variable \(z=a^2q_E^2\) in \(\Ph\) is dimensionless.

\subsection{Coordinate restriction and momentum alignment}
\label{sec:fourier-locking}

GTL acts at each spacetime point. It does not identify different points
or turn independent source labels into the same label. It is also
important to distinguish restriction to a coordinate subspace from
support on a momentum subspace. For the Fourier convention with
kernel \(\exp[-\ii\sum_\Gamma\kappa_\Gamma x^\Gamma]\),
\begin{equation}
 \int_{\RR^K}\prod_{\Gamma=0}^{K-1}\dd x^\Gamma\,
 e^{-\ii\sum_\Gamma\kappa_\Gamma x^\Gamma}
 \prod_{\Gamma=1}^{K-1}\delta(x^\Gamma-x^0)
 =2\pi\delta\!\left(\sum_{\Gamma=0}^{K-1}\kappa_\Gamma\right).
 \label{eq:fourier-locking}
\end{equation}
The result is not a product of
\(\delta(\kappa_\Gamma-\kappa_0)\).
With a function \(f(x^0)\) on the diagonal, the transform instead
contains \(\widetilde f(\sum_\Gamma\kappa_\Gamma)\).
Equation~\eqref{eq:fourier-locking} uses the bare coordinate-elimination
product. The invariant carrier distribution is instead
\[
 \Delta_{{\rm GTL},\hat G}(X_{\rm GT})
 :=\delta_{\hat G^\perp}(Q_{\hat G}X_{\rm GT})
 =\sqrt K\prod_{\Gamma=1}^{K-1}\delta(x^\Gamma-x^0)
\]
in the normalized diagonal chart. The factor \(\sqrt K\) converts
the common-component measure to the normalized carrier measure; it
does not add a physical degeneracy. For \(K=1\), the empty product
is one.

The physical post-GTL theory also imposes the independent A2/GTP
momentum binder
\begin{equation}
 Q_{\bb}p_{\rm GT}=0,\qquad
 p_{\rm GT}=p_\parallel\hat G_{\bb},\qquad
 p_\parallel=\hat G_{\bb}\!\cdot p_{\rm GT}.
 \label{eq:independent-momentum-binder}
% Monograph source: eq:A2-carrier-covector-pullback; eq:GTL-covector-pullback.
\end{equation}
It is not inferred from the Fourier transform of the coordinate
distribution. The physical Fourier phase is pulled back along the
carrier, giving
\(p_{\rm GT}\cdot[(X_0+b_s)\hat G_{\bb}]
=p_\parallel(X_0+b_s)\).
The intrinsic source theory and asymptotic states used here are the
result of both the coordinate pullback and the declared momentum binder.

\subsection{Plane binding and functional derivatives}

For a discrete or atomic support ledger define
\begin{equation}
 \mathcal P_{\bb\bb'}=\delta_{ss'}\delta_{b_s b'_{s'}},
 \qquad
 C_{\bb\bb'}(x,y)=\mathcal P_{\bb\bb'}C_{\bb}(x,y).
 \label{eq:plane-projector}
\end{equation}
The action and interactions are assumed support diagonal; canonical
brackets alone would not prohibit a separately introduced mixing
interaction. For a continuous support family the authoritative
statements are fibrewise, for almost every fixed \(b_s\). A physical
support integral uses a shared fibre variable and its declared
measure. The symbols in \eqref{eq:plane-projector} are then algebraic
label notation, not a normalization by a singleton projector.

Fix also the ordering branch \(\epsilon\). With
\(\ee:=\hat p_{0,\bb}\), the post--GTL action unit is
\begin{equation}
 u_{\epsilon,\bb}:=\sigma_\epsilon\ee,\qquad
 u_{\epsilon,\bb}^{\,2}=1,\qquad
 \BR u_{\epsilon,\bb}=0.
 \label{eq:action-branch-unit}
\end{equation}
We use \(\hat u_{\bb}\equiv u_{\epsilon,\bb}\) only after this
branch has been fixed. The regulator-free physical action and density
have the directed ledger
\begin{equation}
 \vct S_{\epsilon,\bb}
   =u_{\epsilon,\bb}S_{{\rm sc},\epsilon,\bb},\qquad
 S_{{\rm sc},\epsilon,\bb}
   =u_{\epsilon,\bb}\vct S_{\epsilon,\bb},\qquad
 \vct{\mathcal L}_{\epsilon,\bb}
   =u_{\epsilon,\bb}\mathcal L_{{\rm sc},\epsilon,\bb}.
 \label{eq:directed-action}
% Monograph source: eq:A5-postGTL-scalar-directed-action-ledger.
\end{equation}
For the standalone O--QFT representative, the scalar action ledger
contains the sector-fold sign \(\sigma_s\) exactly once. The branch
unit \(u_{\epsilon,\bb}\), the fixed pole carrier \(\ee\), and the
running line unit \(\widehat u_k\) have different roles and are not
interchanged. In particular, \(u_{\epsilon,\bb}\) is not the
numerator convention \(\sigma_s\ee\) of a fixed virtual-carrier
benchmark.

The scalar gauge-fixing and ghost coefficients obey
\[
 S_{{\rm gf},{\rm sc}}+S_{{\rm gh},{\rm sc}}
       =\BR\Psi_{\rm sc},\qquad
 \vct\Psi=u_{\epsilon,\bb}\Psi_{\rm sc},\qquad
 \vct S_{\rm gf}+\vct S_{\rm gh}=\BR\vct\Psi.
\]
The coefficient representation is useful for Euclidean estimates.
The factorization in \eqref{eq:directed-action} applies to the
regulator-free physical action terms; it is not a rule for attaching
one carrier factor to every regulated functional, determinant,
connected graph, or loop-corrected effective action. Those objects
are reconstructed from their carrier characters as described below.

The variational prescription is pullback first. At fixed carrier and
fixed support, the action, connections, inverse metrics, and other
derived geometric objects are formed intrinsically from the pulled-back
fields before Euler--Lagrange variation or integration by parts.
Holding the support fixed means
\(\delta\hat G_{\bb}=0\) and
\(\delta\Delta_{{\rm GTL},\hat G_{\bb}}=0\); it does not mean
\(\partial_M\Delta_{{\rm GTL},\hat G_{\bb}}=0\).
An ambient integration by parts would generally differentiate the
locking distribution. Ambient source and contraction bookkeeping may
precede restriction, but it does not change this intrinsic
variational prescription.

On a fixed carrier, whenever the corresponding
canonical or variational objects exist on the admitted domain,
they inherit the same carrier factor:
\[
 \vct\pi_{\bb}=\hat u_{\bb}\pi_{{\rm sc},\bb},
 \qquad \vct j_{\bb}^{\,\mu}
          =\hat u_{\bb}j_{{\rm sc},\bb}^{\,\mu}.
\]
The carrier-directed translation charge is geo-momentum. Its scalar
coefficient is useful in the analytic estimates, while the charge
itself retains the direction and sector sign inherited from the action.

On an admitted canonically normalized bosonic free-field domain,
at equal \(X_0\),
\begin{align}
 [\phi_{\bb}(X_0,\mathbf x),
     \pi_{{\rm sc},\bb'}(X_0,\mathbf y)]
   &=\ii\mathcal P_{\bb\bb'}\delta^{(3)}(\mathbf x-\mathbf y),
 \label{eq:ccr}\\
 [\phi_{\bb},\vct\pi_{\bb'}]
   &=\ii\hat u_{\bb}\mathcal P_{\bb\bb'}
       \delta^{(3)}(\mathbf x-\mathbf y).
 \label{eq:vector-ccr}
\end{align}
Fermions have the corresponding canonically normalized CAR with
the same label identity. Positivity and stability of that
quantization remain properties of the admitted physical completion.

Independent intrinsic sources obey the usual identity
\begin{equation}
 \frac{\delta J^A_{\bb}(x)}{\delta J^B_{\bb'}(y)}
   =\delta^A{}_B\,\mathcal P_{\bb\bb'}\delta^{(4)}(x-y).
 \label{eq:source-identity}
\end{equation}
Ambient transport phases, if present before restriction, belong to the
transport kernel. They do not modify the identity operator on this
source space.

The identity kernel in \eqref{eq:source-identity} does not absorb
the coefficient of the physical source coupling. On this fixed
support the latter is
\begin{equation}
\begin{aligned}
 a_{{\rm src},\epsilon,\bb}
   &:=\varsigma_{\rm rep}(s)u_{\epsilon,\bb},\\
 \varsigma_{\rm OQFT}(s)&=\sigma_s,\\
 a_{{\rm src},\epsilon,\bb}&=\eta\ee\quad\text{in standalone O--QFT}.
\end{aligned}
 \label{eq:physical-source-carrier}
% Monograph source: eq:A5-postGTL-connected-cumulant-rule; def:W-Gamma-corrected.
\end{equation}
The metric-origin representative of HOUQFT uses
\(\varsigma_{\rm HOU}(s)=1\); that different ledger is not obtained
by inserting another O--QFT sector fold. We denote physical source
coordinates by \(\mathcal J_A\), reserving \(J_A\) in the
Euclidean BRST calculation for canonically normalized coefficient
sources. Any map from auxiliary Wick fields to physical source
fields is part of the completed source pairing. Neither the label
identity nor the carrier factor determines that map.

Repeated support labels never imply a sum. For an
\(n\)-leg amplitude in an atomic ledger, the external binder is
\begin{equation}
 \mathcal P^{(n)}_{\bb\mid{\rm ext}}
   =\prod_{\ell=1}^n\delta_{s s_\ell}
                         \delta_{b_s b_{s_\ell}}.
 \label{eq:external-binder}
\end{equation}
Support diagonality yields a selection rule, not a proof of unitarity.

\subsection{Virtual-line depths and placement}
In the monograph's Postulate~\(\Omega\), the positive virtual-lapse
law and its depth assignment are additional model data. Coordinate
locking, a unit-modulus endpoint phase, or an angular or Gaussian
support average does not uniquely determine that law.
For a graph \(G\), the nondegenerate assignment is
\begin{equation}
 e\in E_{\rm u}(G)\longmapsto a_{e,\bb_e}>0,
 \qquad \ell_{\Omega,e,\bb_e}=\hbar a_{e,\bb_e},
 \label{eq:virtual-edge-depths}
% Monograph source: post:Omega-subordination; eq:Omega-edge-depth-assignment.
\end{equation}
where \(E_{\rm u}(G)\) consists of uncut virtual edges. If
\(\alpha:G\to G'\) is a label-preserving graph isomorphism, and
\(\mathcal I_{A\omega\to Nz}\) is a declared pairing of normal
and antiverse support representatives, covariance of the assignment
requires
\begin{align}
 a_{\alpha(e),\bb'_{\alpha(e)}}&=a_{e,\bb_e},
 \label{eq:depth-isomorphism-covariance}\\
 a_{\mathcal I_{A\omega\to Nz}(e),Nz}&=a_{e,A\omega}.
 \label{eq:depth-pairing-covariance}
\end{align}
These conditions concern the declared edge data. Coincident projected
endpoints, including a tadpole, do not set its depth to zero.
Depth additivity is imposed only for a separately declared concatenation
of parameterized virtual segments; inserting a bivalent Feynman vertex
does not by itself impose such a composition law.

% Monograph source: post:Omega-subordination; eq:Omega-lapse-and-depth-dimensions; cor:fixed-order-uniform-UV-package.
For a general declared edge law, write its dimensionful heat measure
as \(\rho^{\rm heat}_e(\dd\lambda)\). At its fixed depth
\(a_e>0\), the change of variable \(\tau=\lambda/a_e^2\)
defines a dimensionless measure \(\rho_e\), with
\[
\begin{aligned}
 \Phi^{\rm edge}_e(q_E^2)
 &=\int e^{-\lambda q_E^2}\rho^{\rm heat}_e(\dd\lambda)\\
 &=\int e^{-a_e^2\tau q_E^2}\rho_e(\dd\tau),\\
 \int\lambda^{-r}\rho^{\rm heat}_e(\dd\lambda)
 &=a_e^{-2r}M_r[\rho_e].
\end{aligned}
\]
The inverse-moment equality is over the positive half-lines, with
absence of an atom at zero imposed separately when required.
The common-depth formulas below use the allowed benchmark
\(a_e=a>0\), \(\rho_e=\rho\), hence
\(\Phi^{\rm edge}_e(q_E^2)=\Ph(a^2q_E^2)\).
One common rescaled measure is a specialization, not a consequence
of the general depth assignment.

For a fixed finite graph, unequal positive depths and edge measures
are handled by retaining \(a_e\) and the corresponding linewise
gap or moments in the estimates. This gives no uniform theorem for
an infinite graph family. Such a theorem additionally needs control
of the effective damping data, routing and graph constants, vertex
growth, and any graph or label sums. With a fixed rescaled kernel
family, a positive uniform lower bound on the controlling depths is
one sufficient ingredient in those estimates.

The canonical branch of Postulate~\(\Omega\) selects
\(\Ph(z)=e^{-\sqrt z}\), hence \(e^{-a|q_E|}\), as an
additional stable--\(1/2\) specialization. The other kernel families
and the Gaussian bubble below are comparison models within the stated
mathematical class. Once a law has been selected, changing it is a
change of theory data, not a finite mass, residue, or coupling matching
condition.

\subsection{Background operator and internal covariance}

Write the total gauge connection as
\(\mathcal A_\mu=\overline{\mathcal A}_\mu+a_\mu\).
In Euclidean signature choose
\begin{equation}
 \mathscr L_{\bb}=-\overline D_E^{\,\mu}\overline D_{E,\mu}
    \succeq0,\qquad
 \Ph(a^2\mathscr L_{\bb})
     =\int e^{-a^2\tau\mathscr L_{\bb}}\rho(\dd\tau)
 \label{eq:positive-operator}
\end{equation}
on a specified self-adjoint domain. The minus sign makes the flat
eigenvalue \(q_E^2\) nonnegative. Writing
\(\overline D^2\) without this sign convention would obscure which
heat semigroup is being used.

In a common flat spectral representation, the covered internal
covariance is \eqref{eq:covariance-intro}. More general background
Hessians, spin terms, and gauge projectors need not commute with
\(\Ph(\mathscr L_{\bb})\); their ordering and domain must come from
the completed action. No general commutation with a gauge projector
is assumed here. The momentum-space theorem below concerns the
displayed flat representation, or any specified background reduction
for which its absolute envelope has been established.

External LSZ attachments receive no direct \(\Ph\) multiplier.
Operational physical cuts use their declared on-shell projectors.
Both are interfaces of the completed theory. Neither interface is
derived by putting an external momentum into
\(\Ph(0)\).

\begin{proposition}[Free source normalization and its interacting limit]
\label{prop:free-source-normalization}
% Monograph source: prop:source-normalized-realization; rem:source-normalized-why.
On a fixed commuting Euclidean field block, let
\(\mathscr K_0\) be a positive self-adjoint injective free kinetic
operator and \(\Delta=\mathscr K_0^{-1}\). Put
\[
 F=\Ph(a^2\mathscr L_{\bb}),\qquad \Delta_\Ph=F\Delta.
\]
Assume that \(F\) is positive and injective and strongly commutes
with \(\mathscr K_0\): their spectral projections commute.
Work on the common reducing spectral core consisting of vectors
whose joint spectral support is bounded above and separated from
zero in each spectral coordinate, taking the union over such
bounds. Every displayed product
and inverse is then defined there. In particular,
\(T:=F^{-1/2}\) is self-adjoint on its spectral domain.
Introduce an auxiliary Wick field \(\varphi\) with covariance
\(\Delta_\Ph\). The notation \(\varphi^{\rm phys}=T\varphi\)
means the smeared-field identity
\(\varphi^{\rm phys}(f)=\varphi(T^\dagger f)\) on a common
admitted test space; it does not apply an unbounded operator to a
pointwise random field. The covariance forms must extend to any
larger test space subsequently used for local observables.
We use the Hermitian covariance pairing for a complex commuting
field; on a compatible real field block the adjoint below is the
transpose with respect to the real source pairing.
Then its free source--source covariance is
\begin{equation}
 \langle\varphi^{\rm phys}
             (\varphi^{\rm phys})^\dagger\rangle_0
 =T\Delta_\Ph T^\dagger
 =T\Delta_\Ph T=\Delta.
 \label{eq:undressed-source-covariance}
\end{equation}
Contractions between auxiliary interaction insertions still use
\(\Delta_\Ph\). However, a mixed source--vertex contraction is
\begin{equation}
 T\Delta_\Ph=F^{1/2}\Delta,
 \qquad
 (T\Delta_\Ph)\Delta^{-1}=F^{1/2}.
 \label{eq:mixed-source-half-factor}
\end{equation}
Thus free source normalization alone does not prove an undressed
interacting external attachment.
\end{proposition}

\begin{proof}
In the common reducing spectral representation,
\(F\), \(\Delta\), and \(T\) are simultaneous spectral
multipliers, and \(T^\dagger=T\). Thus on the stated core,
\[
 T\Delta_\Ph T^\dagger
 =F^{-1/2}F\Delta F^{-1/2}=\Delta,
 \qquad
 T\Delta_\Ph=F^{-1/2}F\Delta=F^{1/2}\Delta.
\]
These are the two asserted covariance identities. The mixed factor
survives ordinary amputation unless the completed endpoint map or
compatible interactions compensate it, or an appropriate on-shell
normalization is imposed separately.
\end{proof}

For a damping kernel with \(\Ph(q)\to0\), \(T\) is generally
unbounded; the domain clause is essential. The proposition supplies
a free Gaussian realization of the source convention. It does not
construct the completed interacting endpoint map, prove the BRST
master equation, or derive the physical cut rule. Spectral commutation
does not say that \(T^\dagger\) preserves compact or positive-time
support, so it is not a locality or reflection-positivity argument. The source couplings
in the completed functional must distinguish these physical source
coordinates from auxiliary Wick variables, or display the corresponding
source map explicitly. For other components of the BRST multiplet the corresponding identity
is algebraic: the right endpoint operator is the transpose or graded
dual prescribed by the declared source pairing. Its domain and
quadratic-form signs must be checked blockwise. No positive bosonic
Gaussian measure is thereby assigned to fermions or ghosts.

\begin{proposition}[A pure field redefinition cancels the damping]
\label{prop:pure-redefinition-cancellation}
In the real commuting block of
Proposition~\ref{prop:free-source-normalization}, suppose that the
entire interaction is chosen as \(V[T\varphi]\), with sources
coupled to \(T\varphi\), and that the quadratic coefficient is
\(\Delta_\Ph^{-1}\). At fixed background, take an invertible
finite-mode realization preserving the quadratic identities below
and include its change-of-variables Jacobian. Its normalized source
functional equals that of the undamped theory with quadratic
coefficient \(\mathscr K_0\) and interaction \(V\).
In particular this construction is not, by itself, an ultraviolet
improvement of that theory.
\end{proposition}

\begin{proof}
Put \(\chi=T\varphi\). Strong commutation gives
\[
 \Delta_\Ph^{-1}=T\mathscr K_0T,
 \qquad
 \tfrac12\langle\varphi,\Delta_\Ph^{-1}\varphi\rangle
   =\tfrac12\langle\chi,\mathscr K_0\chi\rangle.
\]
The interaction and source term become \(V[\chi]\) and
\(J\cdot\chi\). The finite-mode Jacobian is independent of
fields and sources and cancels in the normalized functional.
Equivalently, every internal line of the transformed vertices has
one \(T\) factor at each endpoint, hence
\begin{equation}
 T(q)\,[F(q)\Delta(q)]\,T(q)=\Delta(q).
 \label{eq:internal-damping-cancellation}
\end{equation}
The cancellation holds graph by graph before any ultraviolet limit.
If the corresponding undamped coefficient diverges when the
finite-mode restriction is removed, this reparametrization leaves
that divergence unchanged.
\end{proof}

For example, \(V[\chi]=\lambda\int\chi^4/4!\) produces the
ordinary undamped tadpole after its two internal endpoint factors
are included. The argument is an exact test of this particular
choice of interaction, not a theorem excluding other completions.
It sharpens the distinction between auxiliary Wick contractions and
the full vertex-weighted graph. The monograph's free source identity
is retained; it does not select \(V[T\varphi]\), and the extra
interacting test here is not attributed to that identity.

\subsection{Carrier algebra and directed boundary values}
\label{sec:carrier}

The fixed geometric pole carrier, its algebraic representative, and
the intrinsic GT axis are identified by
\begin{equation}
\begin{aligned}
 \ee&:=\hat p_{0,\bb}=\hat G_{\bb},\qquad \ee\!\cdot\!\ee=1,\\
 \mathfrak A_{\bb}&=\CC[\ee]/(\ee^2-1),\\
 \je&=\ii\ee,\qquad\je^2=-1.
\end{aligned}
 \label{eq:carrier-algebra}
\end{equation}
The geometric norm and the Clifford product \(\ee^2=1\) are separate
structures. The generator is odd in the \emph{carrier} grading and
Grassmann-even; it commutes with field, spinor, and BRST coefficients.
Thus \(\ee^{2n}=1\) and \(\ee^{2n+1}=\ee\), while all ordinary
Grassmann signs are retained independently.

The directed complex line and its real-algebra isomorphism are
\begin{equation}
 \CC_{\ee}=\operatorname{span}_{\RR}\{1,\je\},\qquad
 \embed(x+\ii y)=x+\ii y\ee,\qquad x,y\in\RR.
 \label{eq:directed-embedding}
\end{equation}
The central complex line \(\operatorname{span}_{\RR}\{1,\ii\}\)
is a different subalgebra. An untransported complex phase times a
directed coefficient need not lie in \(\CC_{\ee}\).
For \(A=A_0+A_1\ee\), with \(A_0,A_1\in\CC\), define
\begin{align}
 \chi_\kappa(A)&=A_0+\kappa A_1,&
 \mathsf P_\kappa&=\frac{1+\kappa\ee}{2},&\kappa&=\pm1,
 \nonumber\\
 A&=\sum_{\kappa=\pm1}\mathsf P_\kappa\chi_\kappa(A),&
 \mathsf P_\kappa\mathsf P_{\kappa'}
   &=\delta_{\kappa\kappa'}\mathsf P_\kappa.
 \label{eq:carrier-characters}
\end{align}
The character label \(\kappa\) is independent of the sector and
evolution labels. Carrier-even and carrier-odd parts are respectively
\(A_0\) and \(A_1\ee\). The inverse \(\embed^{-1}\) is defined
only on \(\CC_{\ee}\), where
\(\chi_+\embed(z)=z\) and \(\chi_-\embed(z)=\overline z\).
In particular, the tag map \(z\mapsto z\ee\) is neither
\(\embed\) nor an algebra homomorphism.

With independent complex and Grassmann sources the coefficient algebra is
\begin{equation}
 \mathfrak B_{\bb}
   =\mathfrak A_{\bb}\widehat\otimes_{\CC}\mathfrak G,
 \label{eq:graded-carrier-source-algebra}
\end{equation}
where \(\mathfrak G\) is the declared Grassmann algebra and
\(\mathfrak A_{\bb}\) lies in its graded centre. The characters
extend by the identity on \(\mathfrak G\). The carrier reality
involution is
\[
 (z_+\mathsf P_++z_-\mathsf P_-)^{\ddagger}
    =\overline{z_-}\mathsf P_++\overline{z_+}\mathsf P_-;
 \qquad \operatorname{Fix}(\ddagger)\cap\mathfrak A_{\bb}
    =\CC_{\ee}.
\]
% Monograph source: def:A5-directed-exp-log; eq:A5-postGTL-reality-graded-OQFT-pairings.
Its extension uses the declared antilinear graded field/source
involution \(\star\):
\[
 (A\widehat\otimes\gamma)^\ddagger
     =A^\ddagger\widehat\otimes\gamma^\star,
 \qquad
 (XY)^\ddagger=(-1)^{|X||Y|}Y^\ddagger X^\ddagger .
\]
Here \(X,Y\) are homogeneous and \(|\cdot|\) denotes Grassmann
parity. On the physical conjugate-source contour, the complete
unweighted pairing \(\int\dd^4x\,\mathcal J_A\phi^A\) is even
and \(\ddagger\)-invariant. Its product with
\(a_{{\rm src},\epsilon,\bb}\) is anti-fixed, whereas the full
source term \(\ii a_{{\rm src},\epsilon,\bb}
\int\dd^4x\,\mathcal J_A\phi^A\) in the Lorentzian exponent
is fixed, since \(\ee^\ddagger=-\ee\) and
\(\ii^\ddagger=-\ii\).
Independent source differentiation precedes restriction to a physical
conjugate-source contour; arbitrary complex sources need not give
conjugate character components.

\begin{proposition}[Carrier reality and positive adjoints]
\label{prop:carrier-reality-adjoint}
On the scalar carrier algebra, the reality operation \(\ddagger\)
interchanges the two character representations:
\begin{equation}
 \chi_\kappa(A^\ddagger)
       =\overline{\chi_{-\kappa}(A)}.
 \label{eq:reality-character-interchange}
\end{equation}
There is no normalized complex-linear functional \(\omega\) on the
full unital algebra satisfying
\(\omega(A^\ddagger A)\geq0\) for every \(A\).
In particular, \(\ddagger\) cannot be the adjoint of a unital
representation of this algebra on a nonzero positive Hilbert space.
\end{proposition}

\begin{proof}
The character formula follows from the definition of \(\ddagger\).
Taking \(A=\ee\) gives
\[
 \ee^\ddagger\ee=-\one,\qquad
 \omega(\ee^\ddagger\ee)=-\omega(\one)=-1.
\]
This contradicts the proposed positivity. Likewise a representation
preserving that involution would give
\(\pi(\ee)^\dagger\pi(\ee)=-I\), impossible on a nonzero
positive Hilbert space.
\end{proof}

For comparison, the same scalar algebra admits a different involution,
which we denote by \(\sharp\):
\begin{equation}
 \left(\sum_\kappa z_\kappa\mathsf P_\kappa\right)^\sharp
      :=\sum_\kappa\overline{z_\kappa}\mathsf P_\kappa.
 \label{eq:coefficient-positive-involution}
\end{equation}
It has \(\ee^\sharp=\ee\), makes each character a
\(\sharp\)-preserving map to ordinary complex numbers, and obeys
\(\|A^\sharp A\|_\chi=\|A\|_\chi^2\) for the norm defined
below. If \(\alpha(A_0+A_1\ee)=A_0-A_1\ee\), then
\(A^\ddagger=\alpha(A^\sharp)\). Thus the two involutions
are not interchangeable. Here \(\sharp\) is only an algebraic
comparison; it neither replaces the declared carrier reality nor
constructs a physical inner product.

The physical adjoint \(\dagger\) belongs to the separately admitted
observable representation, with its own involution on physical fields. No positive \(\ddagger\)-preserving
representation of the entire carrier algebra is assumed. In particular,
carrier reality of a generating functional, positivity of the
character norm, and positivity of physical state sums are three
separate statements. This distinction leaves the vectorial pole
prescription and the two-character reconstruction unchanged.

For a specified scalar function and branch, distinguish
\begin{equation}
 F_{\mathfrak A}(A)
   =\sum_\kappa\mathsf P_\kappa F(\chi_\kappa(A)),
 \qquad
 F_{\ee}=\embed\circ F\circ\embed^{-1}.
 \label{eq:carrier-functional-calculus}
% Monograph source: eq:general-fixed-carrier-functional-calculus; eq:general-character-versus-directed-functional-calculus.
\end{equation}
These agree on the directed line only on a conjugation-invariant
domain with \(F(\overline z)=\overline{F(z)}\) for the chosen
branch. Inverses, logarithms, square roots, and kernel functions are
always interpreted in the indicated calculus. Determinants and
functional logarithms are defined on the two character operators and
reconstructed afterward. The coefficient maps extend entrywise to
matrix-valued factors, preserving the order of their matrix products.

Keep the sector, evolution branch, and positive boundary magnitude
separate:
\begin{equation}
\begin{aligned}
 \sigma_s,\sigma_\epsilon&\in\{+1,-1\},\\
 \eta&:=\sigma_{s\epsilon}=\sigma_s\sigma_\epsilon,\\
 \epsvec&:=\varepsilon\hat p_{0,\bb},\qquad\varepsilon\downarrow0^+.
\end{aligned}
 \label{eq:signs}
\end{equation}
The branch label \(\epsilon\) is not the infinitesimal
\(\varepsilon\). The sign \(\eta\) is outside \(\epsvec\).
On a paired N/A carrier class, use the same fixed pole carrier
\(\hat p_{0,Nb_N}=\hat p_{0,Ab_A}=\hat G\); the opposite
canonical mode orientations are recorded by \(\sigma_s\).
They do not reverse an observer's physical clock.

For a real invariant denominator \(D\), the prescribed boundary is
\begin{equation}
 D\longmapsto D+\ii\eta\epsvec,
 \qquad
 (D+\ii\eta\epsvec)^{-1}
   =\frac{D-\ii\eta\varepsilon\ee}{D^2+\varepsilon^2}.
 \label{eq:directed-inverse}
\end{equation}
This is a boundary value in \(\CC_{\ee}\); it does not shift a
real integration component \(k_\parallel\). The uniform pair of
boundary signs is the Feynman/anti-Feynman pair. A retarded or
advanced prescription has additional frequency dependence and is
not defined by replacing one of these names. The pole datum is also
distinct from the positive, branch-neutral A5 quadratic convergence
regulator of the source functional.

Keep finite \(\varepsilon>0\) through the required source
differentiations and graph integrations, and take the common boundary
limit afterward. Exchanging that limit with an integral or derivative
requires a separate convergence argument. In particular, products
containing running units denote this common regulated expression;
they are not arbitrary products of already singular boundary
distributions. A formula for one line does not prove existence of a
full Lorentzian graph boundary.

For a nonzero aligned GT momentum set
\begin{equation}
 k_{0k,\bb}=k_\parallel\ee,\qquad
 \widehat u_{k,\bb}
  =\frac{k_{0k,\bb}}{\norm{k_{0k,\bb}}}
  =\sgn(k_\parallel)\ee.
 \label{eq:running-carrier}
\end{equation}
At \(k_\parallel=0\) this quotient is undefined. We separately
assign the endpoint value \(\widehat u_{0,\bb}=\sigma_s\ee\);
this convention does not replace nonzero running signs. The
monograph's oriented scalar-line prescription is the following
carrier-valued expression. Calling it a covariance of a particular
source field additionally requires the exchange test below:
\begin{equation}
 \widetilde G^\Ph_{\bb,\eta}(k)
 =\lim_{\varepsilon\downarrow0^+}
 \eta\ii\widehat u_{k,\bb}\,
 \frac{\Ph_{\ee}\!\left(-a^2[k^2+\ii\eta\epsvec]\right)}
      {k^2-m^2+\ii\eta\epsvec}.
 \label{eq:directed-propagator}
% Monograph source: prop:vectorial-poles; eq:general-directed-scalar-coefficient.
\end{equation}
Its continuation and existence as a boundary distribution require
the analytic hypotheses of Section~\ref{sec:analytic}. Its identification
with an inverse Hessian or source covariance is a separate condition;
it is not supplied by the vectorial inverse alone. Equivalently, extracting
the complete directed resolvent coefficient first gives
\begin{align}
 R_\eta^\Ph(k)
  &:=\lim_{\varepsilon\downarrow0^+}\embed^{-1}\!\left[
  \frac{\Ph_{\ee}(-a^2[k^2+\ii\eta\epsvec])}
       {k^2-m^2+\ii\eta\epsvec}\right],\nonumber\\
 \widetilde G^\Ph_{\bb,\eta}(k)
  &=\eta\ii\widehat u_{k,\bb}\embed(R_\eta^\Ph(k)).
 \label{eq:scalar-line-coefficient}
\end{align}
Ordinary \(+\ii\eta0\) in a displayed scalar coefficient below
means this extracted boundary; it never replaces the physical
datum \(\ii\eta\epsvec\).

Only on a fixed-orientation ray is
\(\widehat u_{k,\bb}=\sigma_s\ee\). An explicitly declared
\emph{fixed virtual-carrier benchmark} may assign that unit to
every virtual line independently of its Fourier coefficient. In
that specialization,
\begin{equation}
 \widetilde G^\Ph_{\bb,\eta}
    =\embed(\sigma_\epsilon\ii R_\eta^\Ph).
 \label{eq:fixed-carrier-line}
\end{equation}
For the undamped fixed-carrier line the finite-magnitude decomposition
and its boundary are
\begin{align}
 \frac{\eta\ii\sigma_s\ee}{D+\ii\eta\epsvec}
 &=\sigma_s\frac{\varepsilon}{D^2+\varepsilon^2}
   +\sigma_\epsilon\ii\ee\frac{D}{D^2+\varepsilon^2},\nonumber\\
 \widetilde G_{\bb,\eta}
 &\longrightarrow\sigma_s\pi\delta(D)
     +\sigma_\epsilon\ii\ee\,\PV\frac1D.
 \label{eq:carrier-pv-delta}
\end{align}
Already one line has both carrier parities. Counting its explicit
unit numerators therefore does not determine the parity of a graph.
The Euclidean estimates below concern ordinary coefficient
representations; an unrestricted Lorentzian loop also retains every
running carrier sign and the complete branch data.

For finite \(\varepsilon>0\), let
\(R_\eta^\Ph(k;\varepsilon)\) denote the extracted expression
in \eqref{eq:scalar-line-coefficient} before its limit, and put
\(r_k=\sgn(k_\parallel)\) for a nonzero mode. On a
reflection-compatible kernel branch the full line has characters
\begin{equation}
 \chi_\kappa\!\left(\widetilde G^\Ph_{\bb,\eta}
                  (k;\varepsilon)\right)
 =\kappa\eta r_k\ii\,
  \frac{\Ph(-a^2[k^2+\ii\kappa\eta\varepsilon])}
       {k^2-m^2+\ii\kappa\eta\varepsilon}.
 \label{eq:line-character-values}
\end{equation}
Without that reflection hypothesis, compute the characters of the
specified directed function itself. The factors \(r_{q_i}\)
remain inside every unrestricted running-unit loop integral.

\subsection{Source exchange symmetry and oriented lines}
\label{sec:source-exchange}

\begin{proposition}[Exchange test for a source covariance]
\label{prop:source-exchange}
Fix the support, ordering branch, carrier, and source pairing. Let
\(C^{AB}(x,y)\) be a two-point coefficient obtained by two
independent source derivatives of the same even generating
functional, with the source-first convention of
Section~\ref{sec:bv}. Whenever the derivatives exist and graded
commute, it obeys
\begin{equation}
 C^{AB}(x,y)=(-1)^{\epsilon_A\epsilon_B}C^{BA}(y,x).
 \label{eq:source-covariance-reciprocity}
\end{equation}
For a translation-invariant coefficient this becomes
\(C^{AB}(k)=(-1)^{\epsilon_A\epsilon_B}C^{BA}(-k)\).
In particular, a single identical commuting real scalar has
\(C(k)=C(-k)\), in each carrier character and in the reconstructed
carrier algebra.
\end{proposition}

\begin{proof}
Interchanging the two left source derivatives gives precisely the
Koszul sign in \eqref{eq:source-covariance-reciprocity}.
The common even carrier and source-normalization factors do not
change it. Translation invariance turns endpoint interchange into
momentum reversal. The character projections preserve the identity.
\end{proof}

This necessary test distinguishes an ordered virtual datum from a
source-generated covariance. At fixed finite \(\varepsilon>0\),
the denominator and kernel in \eqref{eq:directed-propagator} are
even under \(k\mapsto-k\), whereas
\(\widehat u_{-k}=-\widehat u_k\) away from
\(k_\parallel=0\). Consequently its unrestricted running-unit
expression satisfies
\begin{equation}
 \widetilde G^\Ph_{\bb,\eta}(-k;\varepsilon)
       =-\widetilde G^\Ph_{\bb,\eta}(k;\varepsilon).
 \label{eq:running-line-reversal}
\end{equation}
The endpoint value on a measure-zero hyperplane does not repair
this mismatch. A nonzero expression of this type cannot itself be
the covariance of the single identical commuting scalar in the
proposition. This obstruction is already present at finite boundary
magnitude, so it is not cured by a choice of distributional limit
or by combining the two carrier characters.

An oriented two-field or endpoint-labelled kernel is a different
possibility: it must specify the reversed component and satisfy
\eqref{eq:source-covariance-reciprocity} for the complete source
matrix. No such completion is constructed by merely writing the
running unit. Nor does source differentiation at a fixed branch
reverse that branch automatically. The fixed virtual-carrier scalar
benchmark used in Section~\ref{sec:literal-bubble} is even and
passes this necessary test, but that fact alone proves neither
positivity nor unitarity.

Accordingly the running-unit formulas are retained as the declared
oriented line data of Ref.~\cite{ChiapponeMonograph}, not promoted
here to a proved Gaussian realization for an identical real scalar.
The unrestricted running-unit source/endpoint realization remains
an explicit completion requirement. The test is an additional
consistency analysis, rather than a replacement of the monograph's
vectorial pole prescription. It also prevents importing the
fixed-carrier scalar benchmark into that unconstructed realization.

\subsection{Complete graphs and their character components}
\label{sec:graph-characters}

For every finite product on one fixed carrier,
\begin{equation}
 \prod_{r=1}^{N}A_r
  =\sum_{\kappa=\pm1}\mathsf P_{\kappa}
           \prod_{r=1}^{N}\chi_\kappa(A_r).
 \label{eq:common-character-product}
\end{equation}
The same character acts on every line, vertex, phase and endpoint
factor. There is no independent character sum for each internal
line: the mixed products vanish because
\(\mathsf P_+\mathsf P_-=0\).

For the complete carrier graph \(\vec{\mathcal G}_{s,\epsilon}\), define
\(\mathcal G^{\rm rep}_{s,\epsilon}
:=\chi_+(\vec{\mathcal G}_{s,\epsilon})\) and
\(\mathcal G^{\rm rep}_{s,\bar\epsilon}
:=\chi_-(\vec{\mathcal G}_{s,\epsilon})\), with
\(\sigma_{\bar\epsilon}=-\sigma_\epsilon\).
The opposite-branch notation refers to the negative-character
projection of the same complete functional. Its calculation uses
the projected orderer, Dyson exponent, vertices and adjoints, pole
prescriptions, external-state conventions and all running line signs.
The carrier-valued graph is therefore
reconstructed from these two projections, not from two unrelated
scalar theories. The bar on the branch label is bookkeeping for the
projected data; it is not complex conjugation. In particular,
\(\chi_-(\ii)=\ii\), so central phases are not conjugated merely
by changing characters. Explicitly,
\begin{equation}
 \vec{\mathcal G}_{s,\epsilon}
  =\mathsf P_+\mathcal G^{\rm rep}_{s,\epsilon}
   +\mathsf P_-\mathcal G^{\rm rep}_{s,\bar\epsilon}.
 \label{eq:complete-graph-character-reconstruction}
% Monograph source: eq:general-directed-graph-character-reconstruction; eq:general-directed-graph-parity-reconstruction.
\end{equation}
Its carrier-even and carrier-odd parts are respectively
\begin{equation}
 \frac{\mathcal G^{\rm rep}_{s,\epsilon}
             +\mathcal G^{\rm rep}_{s,\bar\epsilon}}{2},
 \qquad
 \frac{\mathcal G^{\rm rep}_{s,\epsilon}
             -\mathcal G^{\rm rep}_{s,\bar\epsilon}}{2}\,\ee.
 \label{eq:complete-graph-parity}
\end{equation}
Only when the complete reflection relation
\(\mathcal G^{\rm rep}_{s,\bar\epsilon}
=\overline{\mathcal G^{\rm rep}_{s,\epsilon}}\) has been
established may this be compressed to
\(\vec{\mathcal G}_{s,\epsilon}
=\embed(\mathcal G^{\rm rep}_{s,\epsilon})\).
For example, a central factor \(\ii\) has both character values
equal to \(\ii\); reconstructing it as
\(\embed(\ii)=\ii\ee\) would change the element.
The inverse directed map therefore cannot scalarize an arbitrary
mixed carrier expression.

\begin{proposition}[Character estimates for complete graphs]
\label{prop:character-majorants}
For scalar carrier coefficients define
\begin{equation}
 \norm{A}_{\chi}:=\max_{\kappa=\pm1}|\chi_\kappa(A)|.
 \label{eq:character-norm}
\end{equation}
This is a submultiplicative norm. Let \(\nu\) be a positive
integration measure. If a measurable complete graph integrand is
\(g(\xi)=\sum_\kappa\mathsf P_\kappa g_\kappa(\xi)\)
and \(|g_\kappa(\xi)|\leq M_\kappa(\xi)\), with
\(M_\pm\in L^1(\dd\nu)\), then
\begin{equation}
\begin{aligned}
 \int\norm{g(\xi)}_\chi\dd\nu
 &\leq\int(M_++M_-)\dd\nu<\infty,\\
 \int g\,\dd\nu
 &=\sum_\kappa\mathsf P_\kappa\int g_\kappa\,\dd\nu.
\end{aligned}
 \label{eq:character-integral-majorant}
\end{equation}
The same statement holds uniformly in external parameters whenever
the two majorants are uniform. For matrix coefficients replace the
absolute value by a fixed submultiplicative matrix norm.
\end{proposition}

\begin{proof}
The character map identifies \(\mathfrak A_{\bb}\) with
\(\CC\oplus\CC\). Multiplicativity of each character proves
the norm inequality for products. The integral estimate follows from
\(\max(|g_+|,|g_-|)\leq M_++M_-\); the two finite-dimensional
projections commute with the absolutely convergent integral.
\end{proof}

This norm is an absolute-convergence device. Its compatibility with
\(\sharp\) in \eqref{eq:coefficient-positive-involution} does not
turn \(\ddagger\) into a positive adjoint or supply physical
state norms; Proposition~\ref{prop:carrier-reality-adjoint} prevents
that identification on the full carrier algebra.

For formal source or Grassmann expansions, apply this estimate after
forming the complete graph and extracting a specified coefficient.
Coefficient extraction alone is not multiplicative; products retain
the convolution over source degrees and their Koszul signs. Equivalently,
a finite Grassmann truncation admits a submultiplicative coefficient
norm before integration.

This proposition transfers the UV bounds below to the full carrier
graph when both representatives satisfy them. On real integration
domains the running units contribute signs of modulus at most one,
so they preserve an established absolute envelope. Holomorphy still
requires the hypotheses of Theorem~\ref{thm:holomorphy} in each
character's declared complex variables. In particular, \(\embed\)
is a real-algebra map: holomorphy in a directed argument is not an
unstated identification of \(\ii\) with \(\je\).


\section{Ultraviolet finiteness without divergent subtraction}
\label{sec:uv}

This section is Euclidean. It defines a nonnegative comparison
integral and proves its convergence under explicit hypotheses,
without assuming a Lorentzian Wick rotation.
The connection to BRST is that the vertices whose growth is estimated
must be the vertices of the completed action in
Section~\ref{sec:brst}. BRST invariance is not needed to prove an
estimate for an already specified scalar majorant.

\subsection{Graph data and infrared control}

Fix \(\bb\), a selected evolution branch, and a finite connected
Euclidean virtual graph \(G\) with
\(E_{\rm u}(G)=\{1,\ldots,I\}\) and \(L\geq1\)
independent loop momenta. Every line counted by \(I\) is an
\emph{uncut virtual covariance}; physical external attachments and
on-shell cut projectors are not factors in this comparison integral.
Write
\begin{equation}
\begin{gathered}
 q_i(k,p)=\sum_{\ell=1}^L C_{i\ell}k_\ell+R_i(p),\\
 k\in\RR^{4L},\qquad p\in\mathcal K,\qquad\rank C=L.
\end{gathered}
 \label{eq:routing}
\end{equation}
where \(\mathcal K\) is a compact set of real Euclidean external data.
The routing shifts \(R_i(p)\) are affine-linear and bounded on \(\mathcal K\).
All matrix entries, fields, and orientations are held on their
declared fixed support. Assume measurable graph factors, and suppose
that the full product of vertex and numerator factors obeys
\begin{equation}
 \abs{N_G(k,p)}\leq C_{G,\mathcal K}\bigl(1+\norm{k}^D\bigr),
 \qquad D\geq0.
 \label{eq:numerator-bound}
\end{equation}
Norms in polynomial inequalities are expressed in a fixed reference
momentum unit; the constants restore the dimensions. Nonlocal
completion vertices are included in \eqref{eq:numerator-bound}.
A formal power series for a vertex does not establish this bound for
its untruncated sum.

Use nonnegative measurable Schwinger envelopes
\begin{align}
 C_i(q)&=\int_0^\infty E_i(\alpha)e^{-\alpha q^2}\dd\alpha,
 \label{eq:IR-envelope}\\
 \mathcal I_i&:=\int_0^\infty E_i(\alpha)\dd\alpha<\infty,
 &\mathcal E_i&:=\sup_{0<\alpha\leq\alpha_0}E_i(\alpha)<\infty ,
 \label{eq:IR-conditions}
\end{align}
for a fixed \(\alpha_0>0\). A positive mass gives
\(E_i(\alpha)=e^{-m_i^2\alpha}\) and
\(C_i(q)=(q^2+m_i^2)^{-1}\).
If a massless theory is considered, \eqref{eq:IR-envelope} is an
explicit infrared comparison prescription; an infrared-finite graph
can instead be treated with its own established IR estimate.
The theorem does not assert that all massless graphs are finite,
or that an auxiliary IR prescription can be removed. In particular,
selected massless denominators in
Theorem~\ref{thm:optimized-selected-lines} use a separate direct
line-integral argument and do not retroactively satisfy
\eqref{eq:IR-conditions}.

For clarity, first give every line the same kernel
\eqref{eq:bernstein}. The finite graph majorant is
\begin{equation}
 \mathcal A_G(p)=\int_{\RR^{4L}}
 \frac{\dd^{4L}k}{(2\pi)^{4L}}\,
 \abs{N_G(k,p)}
 \prod_{i=1}^I \Ph(a^2q_i^2)C_i(q_i).
 \label{eq:graph-majorant}
\end{equation}
Finiteness of this positive integral bounds every matrix component
or any fixed finite-dimensional tensor norm of the graph with these
declared denominators and IR prescriptions. This is a \emph{graphwise}
claim at fixed support and finite loop order, not a bound on the number
of graphs or their sum. An IR-modified covariance is not presumed to
majorize an unregulated massless covariance.

\subsection{Global fixed-graph theorem}

For \(r>0\), define the inverse moment
\begin{equation}
 M_r[\rho]=\int_{(0,\infty)}\tau^{-r}\rho(\dd\tau).
 \label{eq:inverse-moment}
\end{equation}
The condition that no atom is present at zero is imposed
\emph{separately}: the integral defining $M_r$ excludes that point.
For example, $M_r[\delta_0]=0$ for every $r>0$, even though the
associated multiplier is $\Ph(z)\equiv1$. For $d\geq0$, the
corresponding undamped massless comparison has
\[
 \int_{|q|<R}|q|^{d-2}\dd^4q
    =\frac{2\pi^2}{d+2}R^{d+2}\longrightarrow\infty.
\]
Thus finiteness of the inverse moment as defined here does not imply
the no-atom hypothesis. The latter is a separate premise of the
convergence statements below. Equivalently, for every $r>0$ the two
requirements can be recorded as one extended-valued integral,
\begin{equation}
 \begin{gathered}
 \rho(\{0\})=0,\qquad M_r[\rho]<\infty\\
 \Longleftrightarrow\quad
 \int_{[0,\infty)}\tau^{-r}\rho(\dd\tau)<\infty.
 \end{gathered}
 \label{eq:extended-inverse-moment}
\end{equation}
Here the integrand on the right takes the extended value
$0^{-r}:=+\infty$. That integral is \emph{not} the definition
of $M_r$ in \eqref{eq:inverse-moment}. This formulation makes the independent
atom condition explicit in a machine-checkable theorem statement.

\begin{theorem}[Global Euclidean convergence]
\label{thm:global-uv}
Assume \eqref{eq:routing}--\eqref{eq:IR-conditions}, finite
\(\Mzero\), and fixed \(a>0\). Either of the following is sufficient
for \(\sup_{p\in\mathcal K}\mathcal A_G(p)<\infty\):
\begin{description}[leftmargin=1cm]
\item[(G)] There is \(\tau_*>0\) with
\(\supp\rho\subset[\tau_*,\infty)\).
\item[(M)] \(\rho(\{0\})=0\) and
\(M_{1+D/2}[\rho]<\infty\).
\end{description}
The joint absolute majorant in loop momenta, Schwinger parameters,
and Laplace--Stieltjes parameters is integrable. All Euclidean
integration-order exchanges used in the proof are therefore valid.
No UV-divergent subtraction is needed to define this graph.
This conclusion concerns the stated Euclidean comparison integral;
it does not establish a physical spectral representation or the
validity of an undressed operational cut.
\end{theorem}

A more conservative sufficient condition is
\begin{equation}
 \int_{(0,1]}\tau^{-\beta}\rho(\dd\tau)<\infty,\qquad
 \beta>1+\frac D2,\qquad\rho(\{0\})=0.
 \label{eq:conservative-moment}
\end{equation}
It implies (M), but is not necessary for this proof. The endpoint
formulation requires the \emph{actual} inverse moment
$M_{1+D/2}[\rho]$ to be finite; borderline power-law behavior
alone does not establish that finiteness.

\begin{proof}
Introduce independent parameters \(\tau_i,\alpha_i\) and set
\begin{equation}
 w_i=\alpha_i+a^2\tau_i,\qquad W=\diag(w_1,\ldots,w_I),
 \qquad A=C^TWC .
 \label{eq:width-matrix}
\end{equation}
Tonelli's theorem applies to the nonnegative comparison integrand
even before its finiteness has been proved. For almost every
positive width vector, \(A\succ0\). Completing the square in all
four Euclidean components gives
\begin{equation}
 \sum_iw_iq_i^2
  =(k+\bmu)^T(A\otimes\one_4)(k+\bmu)+\Delta(w,p),
 \qquad \Delta(w,p)\geq0 .
 \label{eq:complete-square}
\end{equation}

The affine shift is uniformly bounded. For one Euclidean component,
let
\[
 \mu=(C^TWC)^{-1}C^TWR.
\]
If \(\mathfrak B\) is the finite family of invertible \(L\)-row
submatrices \(C_T\), then
\begin{equation}
 \mu=\sum_{T\in\mathfrak B}\pi_T(w)C_T^{-1}R_T,\qquad
 \pi_T(w)=
 \frac{(\det C_T)^2\prod_{i\in T}w_i}
      {\sum_{U\in\mathfrak B}(\det C_U)^2\prod_{i\in U}w_i}.
 \label{eq:shift-mixture}
\end{equation}
The coefficients are nonnegative and sum to one.
\ref{sec:appendix-gaussian} proves this identity directly.
Consequently \(\norm{\bmu}\leq C_{G,\mathcal K}\) even in nonuniform
small-width limits.

After the change of variables
\(y=(A^{1/2}\otimes\one_4)(k+\bmu)\), the Jacobian is
\((\det A)^{-2}\). The Gaussian moments of orders \(0\) and \(D\), together with the
bounded shift, give the following estimate. For \(D=0\), only the
zeroth moment is needed:
\begin{equation}
 \int_{\RR^{4L}}\dd^{4L}k\,\abs{N_G(k,p)}
                 e^{-\sum_iw_iq_i^2}
 \leq C_{G,\mathcal K}\frac{1+\norm{A^{-1}}^{D/2}}{(\det A)^2}.
 \label{eq:gaussian-bound}
\end{equation}
The constant term on the right is necessary. Neither lower-degree
numerator terms nor a nonzero bounded affine shift can be bounded
solely by \(\norm{A^{-1}}^{D/2}\) when this quantity tends to zero.

Choose any invertible \(L\)-row submatrix \(C_S\) and let
\(s_S>0\) be its smallest singular value. Cauchy--Binet and the
Rayleigh quotient give the explicit routing-dependent bounds
\begin{equation}
\begin{aligned}
 \det A&\geq |\det C_S|^2\prod_{i\in S}w_i,\\
 \norm{A^{-1}}&\leq\frac{s_S^{-2}}{\min_{i\in S}w_i}.
\end{aligned}
 \label{eq:chord-bounds}
\end{equation}
Indeed, the determinant contains the indicated positive minor, and
\(v^TAv\geq(\min_{i\in S}w_i)\norm{C_Sv}^2\).
In a fundamental chord basis, \(C_S\) is a signed identity and
both routing constants are one. For a general fixed routing they
are absorbed into the graph constant below.
Using
\((\min_iw_i)^{-r}\leq\sum_iw_i^{-r}\), \(r\geq0\),
the loop-integrated majorant is at most
\begin{equation}
 C_{G,\mathcal K}\prod_{i=1}^IE_i(\alpha_i)
  \prod_{i\in S}(\alpha_i+a^2\tau_i)^{-2}
 \left[1+\sum_{j\in S}
             (\alpha_j+a^2\tau_j)^{-D/2}\right].
 \label{eq:parameter-majorant}
\end{equation}
This inequality holds throughout the positive parameter orthant;
it is not restricted to simultaneous radial scaling.

Every nonchord Schwinger integral contributes
\(\mathcal I_i<\infty\). For a chord, split its integral at
\(\alpha_0\). For \(\gamma>1\) and \(\kappa>0\),
\begin{align}
 \int_0^\infty
  \frac{E_i(\alpha)}{(\alpha+\kappa)^\gamma}\dd\alpha
 &\leq\frac{\mathcal E_i}{\gamma-1}
           \kappa^{1-\gamma}
       +\alpha_0^{-\gamma}\mathcal I_i .
 \label{eq:alpha-envelope-bound}
\end{align}
The first term bounds \(0<\alpha\leq\alpha_0\); the second bounds
the remaining tail. In particular, one must not discard
\(\int_0^\infty E_i(\alpha)\dd\alpha\) for an undamped nonchord
parameter and then claim uniformity at \(m_i=0\).

Let \(h_r(\tau)=1+\tau^{-r}\). Applying
\eqref{eq:alpha-envelope-bound} with
\(\gamma=2\) or \(2+D/2\) gives the remaining majorant
\begin{equation}
 C_{G,\mathcal K,a,E}\left[
   \prod_{i\in S}h_1(\tau_i)
  +\sum_{j\in S}h_{1+D/2}(\tau_j)
          \prod_{i\in S\setminus\{j\}}h_1(\tau_i)
 \right].
 \label{eq:final-tau-majorant}
\end{equation}
Each nonchord \(\tau\)-integration contributes only \(\Mzero\).
Under (G), every factor in \eqref{eq:final-tau-majorant} is bounded
on the support of \(\rho^{\otimes I}\). Under (M), both
\(\int h_1\dd\rho\) and
\(\int h_{1+D/2}\dd\rho\) are finite. Independence of the line
measures factorizes these finite integrals. The positive majorant is
therefore integrable, uniformly on \(\mathcal K\); Fubini now applies to the
signed or tensor graph.
\end{proof}

\begin{remark}[Strength and limits of the estimate]
The global bound already controls all partial, nested, and overlapping
UV regions for the graph class in Theorem~\ref{thm:global-uv}.
It deliberately discards some useful nonchord decay and is generally
not a necessary convergence criterion. Its constants can diverge as
\(a\downarrow0\), as an unregulated positive mass tends to zero, or
as the graph family becomes infinite. No such uniform limit is asserted.
\end{remark}

\begin{remark}[Endpoint moments need not imply stronger moments]
\label{rem:strict-endpoint-example}
For a fixed graph degree put \(r=1+D/2\geq1\) and consider
\begin{equation}
 \rho_r(\dd\tau)
 =\mathbf1_{(0,\mathrm e^{-1})}(\tau)
   \frac{\tau^{r-1}}{\log^2(1/\tau)}\dd\tau.
 \label{eq:endpoint-measure-example}
\end{equation}
This is a finite positive measure without an atom at zero. The
substitution \(u=\log(1/\tau)\) gives
\begin{equation}
 \begin{aligned}
  M_r[\rho_r]&=\int_1^\infty\frac{\dd u}{u^2}=1,\\
  M_{r+\delta}[\rho_r]
     &=\int_1^\infty\frac{e^{\delta u}}{u^2}\dd u=\infty
       \quad(\delta>0).
 \end{aligned}
 \label{eq:endpoint-measure-moments}
\end{equation}
Normalizing its total mass to one changes neither finiteness statement.
The support reaches zero, so there is no positive gap. Nevertheless,
its bounded support gives, for every \(\lambda>0\),
\[
 \int e^{\lambda\tau}\rho_r(\dd\tau)
       \leq e^{\lambda/\mathrm e}\rho_r((0,\mathrm e^{-1}))<\infty.
\]
Its Laplace transform is therefore entire. This example separates
near-zero inverse-moment control from the large-depth condition used
for local holomorphy; neither should be confused with a positive gap.
Thus the endpoint hypothesis genuinely extends the stronger sufficient
condition \eqref{eq:conservative-moment}. It is not claimed to be
necessary for a particular graph, nor does this example provide one
law covering unbounded numerator degrees.
\end{remark}

For example, consider the scalar four-dimensional one-loop bubble
with numerator one, two positive masses, fixed real Euclidean
external momentum, and two power kernels
\(\Ph(z)=(1+z)^{-\nu}\). Its radial UV tail is proportional to
\(r^{-1-4\nu}\), so it is UV convergent for every \(\nu>0\).
The sufficient global criterion with \(D=0\) asks for \(\nu>1\).
Thus the moment condition is sufficient and need not be a sharp
convergence threshold for an individual graph.

\begin{corollary}[Damping a full-rank set of loop lines]
\label{cor:damped-basis}
Retain the routing, numerator, and infrared hypotheses of
Theorem~\ref{thm:global-uv}, and let \(S\subset E_{\rm u}(G)\) be a set of \(L\) uncut
virtual lines for which \(C_S\) is invertible. For each \(i\in S\), suppose
\(\Ph_i\) has a finite positive representing measure \(\rho_i\)
with either a positive gap or
\[
 \rho_i(\{0\})=0,\qquad M_{1+D/2}[\rho_i]<\infty.
\]
On every remaining line, it is enough that the multiplier be measurable
and bounded on the nonnegative real axis; the identity multiplier is allowed.
Then the graph with these line-dependent multipliers is absolutely
convergent, uniformly on \(\mathcal K\).
\end{corollary}

\begin{proof}
Bound the multipliers outside \(S\) by their finite suprema.
In the Schwinger proof use
\(w_i=\alpha_i+a^2\tau_i\) for \(i\in S\) and
\(w_i=\alpha_i\) otherwise. The invertible minor \(C_S\) supplies
the determinant and inverse-matrix bounds, with fixed routing
constants. Equation~\eqref{eq:final-tau-majorant} then contains
only the measures on \(S\). Its factors are integrable by the
individual gap or moment assumptions, while each remaining
Schwinger integral contributes \(\mathcal I_i<\infty\).
A direct proof with explicit one-line integral bounds is given in
\ref{sec:direct-loop-proof}.
\end{proof}

The bounded multipliers outside \(S\) are a freedom in this mathematical
comparison criterion. In applying it to a selected model, all uncut
virtual lines still carry the kernels declared by that model; one may
bound some of those kernels by constants. The allowed identity multiplier
does not turn a physical cut projector into a Euclidean covariance,
or authorize changing the placement rule. A cut removes the cut-line
covariance and its kernel together. The UV estimates concern the
remaining virtual integrations on their established kinematic domain;
they do not prove convergence of the on-shell phase-space integral or
remove infrared singularities.

This corollary gives a direct sufficient coverage test: its selected
damped lines resolve all independent loop directions. It is not a
necessary rank condition for convergence of every graph. Their number alone is
insufficient; the relevant condition is the rank of their routing
rows. The forest criterion below permits useful bounds even when
this simple full-rank sufficient condition is unavailable.

\paragraph{Real momentum limits}
A uniform bound on the integrated absolute value does not by itself
justify dominated convergence. To obtain a common pointwise majorant
for real external limits, use the selected line momenta as integration
variables. Put
\[
 \ell=q_S,\qquad
 k=(C_S^{-1}\otimes\one_4)(\ell-R_S(p)).
\]
The Jacobian is the constant \(|\det C_S|^{-4}\). Bound each
unselected line by its multiplier supremum times \(\mathcal I_i\).
The complete transformed integrand is then bounded independently of
\(p\in\mathcal K\) by
\begin{equation}
 \mathcal M_S(\ell)=C\Bigl(1+\sum_{j\in S}|\ell_j|^D\Bigr)
       \prod_{j\in S}\Ph_j(a_j^2|\ell_j|^2)C_j(\ell_j).
 \label{eq:real-parameter-majorant}
\end{equation}
For \(D=0\), the parenthesis is simply a constant. The one-line
moments in \ref{sec:direct-loop-proof}, with the corresponding fixed
\(a_j>0\), show that \(\mathcal M_S\in L^1(\RR^{4L})\).
Thus, if the transformed numerator and unselected line factors are
continuous in the real external data for almost every \(\ell\),
the graph is continuous and its real external limits pass through
this representation integral. This continuity assumption is not
automatic for arbitrary measurable comparison multipliers. Derivatives
require their own dominating functions; complex continuation and
Lorentzian boundary limits still use Section~\ref{sec:analytic}.
No common pre-Gaussian complex majorant is inferred from a uniform
bound on the integrated absolute values.


\subsection{Optimized selected-line coverage, including massless denominators}
\label{sec:optimized-coverage}

The global theorem uses a single numerator degree and integrable
Schwinger envelopes on every line. Some graphs admit a more economical
proof in selected line momenta. This also permits a massless denominator
on a \emph{selected} line, without assigning a fictitious photon mass.

\begin{theorem}[Line-specific inverse moments on a full-rank routing]
\label{thm:optimized-selected-lines}
Fix the finite routing and compact real external set of
\eqref{eq:routing}. Choose $S$ with $|S|=L$ and invertible $C_S$, and
write $\ell_j=q_j$ for $j\in S$. On selected line $j$, take the actual
nonnegative comparison factor
\begin{equation}
 C_j^{\Ph}(\ell_j)
  :=\frac{\Ph_j(a_j^2|\ell_j|^2)}{|\ell_j|^2+m_j^2},
 \qquad a_j>0,\quad m_j\geq0,
 \label{eq:optimized-selected-factor}
\end{equation}
with a finite positive representing measure $\rho_j$.
Assume the \emph{complete product} of the numerator and all
unselected virtual-line factors, expressed in these variables, obeys
\begin{equation}
 \left|N_G(k(\ell,p),p)\prod_{i\notin S} C_i^{\rm full}(q_i)\right|
 \leq\sum_{\alpha\in\mathcal A}c_\alpha
       \prod_{j\in S}(1+|\ell_j|^{d_{j\alpha}}),
 \label{eq:anisotropic-complete-envelope}
\end{equation}
uniformly in $p\in\mathcal K$, with $\mathcal A$ finite,
$c_\alpha\geq0$ and $d_{j\alpha}\geq0$. Thus an unselected
massless denominator must be controlled as part of the \emph{full}
envelope; it is not assumed bounded merely because its multiplier is.
Set $d_j=\max_\alpha d_{j\alpha}$ and suppose
\begin{equation}
 \rho_j(\{0\})=0,\qquad
 M_{1+d_j/2}[\rho_j]<\infty\quad(j\in S).
 \label{eq:optimized-moment-budget}
\end{equation}
Then the specified complete graph is absolutely convergent,
uniformly over $\mathcal K$, including mixed soft and UV regions
in the selected variables. The following one-line integral does
\emph{not} include the Fourier factor $(2\pi)^{-4}$:
\begin{equation}
 \begin{aligned}
 Q_{d,j}(m_j)
  &:=\int_{\RR^4}\frac{|q|^d\Ph_j(a_j^2|q|^2)}{q^2+m_j^2}
         \dd^4q\\
  &\leq\pi^2\Gamma(1+d/2)a_j^{-2-d}
          M_{1+d/2}[\rho_j].
 \end{aligned}
 \label{eq:optimized-one-line-moment}
\end{equation}
For $m_j=0$ this is an equality, possibly infinite. Under the
\emph{separate no-atom hypothesis}, finiteness of the displayed
inverse moment is necessary and sufficient for that particular
nonnegative weighted one-line integral. It remains a sufficient,
not generally necessary, condition for the complete graph.
\end{theorem}
\begin{proof}
For $q\ne0$ and $m_j\geq0$,
$(q^2+m_j^2)^{-1}\leq|q|^{-2}$. By Tonelli and four-dimensional
polar integration,
\[
 \begin{aligned}
 Q_{d,j}(m_j)&\leq2\pi^2\int\rho_j(\dd\tau)
   \int_0^\infty r^{d+1}e^{-a_j^2\tau r^2}\dd r\\
 &=\pi^2\Gamma(1+d/2)a_j^{-2-d}
                   M_{1+d/2}[\rho_j].
 \end{aligned}
\]
Equality holds for a massless selected denominator. Near $q=0$ the
radial integrand is proportional to $r^{d+1}$ at fixed positive
Laplace depth, so no additional selected-line infrared singularity
is created for $d\geq0$. An atom at $\tau=0$ would instead contribute
an undamped constant and make the massless one-line integral diverge.
It is excluded by hypothesis, not by the moment definition.

The real affine change
$k=(C_S^{-1}\otimes\one_4)(\ell-R_S(p))$ has fixed Jacobian
$|\det C_S|^{-4}$. Insert \eqref{eq:anisotropic-complete-envelope} and integrate
the selected factors independently. With the graph's stated
$(2\pi)^{-4L}$ Fourier convention, the explicit bound is
\begin{equation}
 \mathcal A_G(p)\leq
 \frac{|\det C_S|^{-4}}{(2\pi)^{4L}}
 \sum_{\alpha\in\mathcal A}c_\alpha
 \prod_{j\in S}\bigl[Q_{0,j}(m_j)+Q_{d_{j\alpha},j}(m_j)\bigr].
 \label{eq:anisotropic-integrated-bound}
\end{equation}
Here $\mathcal A_G$ denotes the absolute comparison graph with
the selected factors \eqref{eq:optimized-selected-factor} and
the complete remaining product bounded in
\eqref{eq:anisotropic-complete-envelope}.
Finite mass and the assumed order $1+d_j/2$ imply every smaller
required positive inverse moment. This is an integrable pointwise
majorant in independent selected coordinates, controlling their
partial and simultaneous ultraviolet limits. The theorem assumes
rather than proves the requisite estimate for unselected factors and
completed vertices.
\end{proof}

\begin{corollary}[An explicit selected-line numerical tail]
\label{cor:optimized-selected-tail}
Under Theorem~\ref{thm:optimized-selected-lines}, if for some
$\delta>0$ the moments
$M_{1+(d_j+\delta)/2}[\rho_j]$ are finite, then
\begin{equation}
 \begin{aligned}
 &\int_{|q|>R}\frac{|q|^d\Ph_j(a_j^2|q|^2)}{q^2+m_j^2}
  \dd^4q\\
 &\quad\leq \frac{\pi^2\Gamma(1+(d+\delta)/2)}
                       {R^\delta a_j^{2+d+\delta}}
       M_{1+(d+\delta)/2}[\rho_j],
 \quad 0\leq d\leq d_j.
 \end{aligned}
 \label{eq:optimized-tail-bound}
\end{equation}
Replacing one factor at a time in
\eqref{eq:anisotropic-integrated-bound} and taking the union over
$\{|\ell_j|>R\}$ gives an $O(R^{-\delta})$ tail for the complete
specified graph in the selected variables.
\end{corollary}
\begin{proof}
On $|q|>R$, $|q|^d\leq R^{-\delta}|q|^{d+\delta}$.
Apply \eqref{eq:optimized-one-line-moment} and the finite-product
bound. The union estimate follows from positivity. For each summand in
\eqref{eq:anisotropic-integrated-bound}, apply the tail estimate
separately to its constant and $d_{j\alpha}$ terms; the remaining
line integrals retain their own moment budgets. The bound concerns
\emph{selected-coordinate} truncations, not arbitrary independently
routed cutoffs on divergent constituent integrals.
\end{proof}

For illustration, a massless selected line with
$\Ph_j(z)=(1+z/b_j)^{-\nu_j}$ has the exact moment
\[
 Q_{d,j}(0)=\frac{\pi^2\Gamma(1+d/2)b_j^{1+d/2}}
    {a_j^{2+d}}
 \frac{\Gamma(\nu_j-1-d/2)}{\Gamma(\nu_j)},
 \quad \nu_j>1+d/2.
\]
This particular integral diverges at or below the stated threshold.
If two selected rows require degrees $0$ and $4$, they need
$M_1[\rho_1]$ and $M_3[\rho_2]$, respectively, rather than
$M_3$ on both rows. The improvement relies on the anisotropic
\emph{complete} envelope; decomposing a numerator after discarding
unbounded nonselected factors is not valid.

\paragraph{A two-loop certificate with massless selected lines}
\label{par:two-loop-certificate}
As a checkable example, let \(k_1,k_2\in\RR^4\), take the routing
\(q_1=k_1\), \(q_2=k_2\), \(q_3=k_1-k_2+p\), and let the third
line have a massive denominator \(q_3^2+1\) and a multiplier
bounded by one. The selected routing \(C_{\{1,2\}}=I_2\) has full
rank. In fixed reference momentum units consider the nonnegative
comparison graph
\begin{equation}
 \begin{aligned}
 \mathcal A_{\rm ex}(p)=\int
 \frac{\dd^4k_1\dd^4k_2}{(2\pi)^8}
 &\frac{1+|k_1|^2|k_2|^4}{k_1^2k_2^2(q_3^2+1)}\,\\[-2pt]
 &\quad\times(1+k_1^2)^{-3}(1+k_2^2)^{-4}.
 \end{aligned}
 \label{eq:two-loop-selected-example}
\end{equation}
The denominators are interpreted almost everywhere, and \(p\)
can range over any real compact set. Bounding the third line by one
and applying the exact one-line identity yields
\begin{equation}
 \begin{aligned}
 \mathcal A_{\rm ex}(p)
 &\leq\frac{Q_{0,1}(0)Q_{0,2}(0)
              +Q_{2,1}(0)Q_{4,2}(0)}{(2\pi)^8}\\
 &=\frac{1}{768\pi^4},
 \end{aligned}
 \label{eq:two-loop-selected-certificate}
\end{equation}
using \(Q_{0,1}=Q_{2,1}=\pi^2/2\) and
\(Q_{0,2}=Q_{4,2}=\pi^2/3\). To verify the kernel hypotheses
independently, the representing Gamma laws are
\[
 \rho_\nu(\dd\tau)
 =\frac{\tau^{\nu-1}e^{-\tau}}{\Gamma(\nu)}\dd\tau,
 \qquad
 M_r[\rho_\nu]=\frac{\Gamma(\nu-r)}{\Gamma(\nu)}
 \quad(0<r<\nu).
\]
They have no atom at zero, and the highest selected-line moments are
\(M_2[\rho_3]=1/2\) and \(M_3[\rho_4]=1/6\).
The four displayed $Q$ values follow directly from
\eqref{eq:optimized-one-line-moment} with unit depth: the two
selected propagators contribute their exact radial integrals,
whereas the third line contributes only the bound
$(q_3^2+1)^{-1}\leq1$.
Thus one concrete two-loop comparison is uniformly bounded despite
two massless selected propagators. This example tests the
mathematical certificate and its Fourier normalization; it is not
claimed to be a graph generated by a completed O--QED action.

\begin{proposition}[Complex and real external-momentum certificates for the two-loop example]
\label{prop:two-loop-complex-certificate}
In \eqref{eq:two-loop-selected-example}, interpret the third
factor as the analytic denominator
$[(k_1-k_2+p)^2+1]^{-1}$ for complex $p$, with
$q^2=\sum_{\mu=1}^4q_\mu^2$. Denote this complex integral by
$\mathcal B_{\rm ex}(p)$. It is holomorphic on the tube
$\mathcal T_1=\{p=x+\ii y\in\CC^4:|y|<1\}$, and
\begin{equation}
 |\mathcal B_{\rm ex}(x+\ii y)|
 \leq\frac{1}{768\pi^4(1-|y|^2)}.
 \label{eq:two-loop-complex-bound}
\end{equation}
For every real $p_0$, $0<r<1/2$, and multi-index
$\alpha\in\mathbb N_0^4$, its external derivatives satisfy
\begin{equation}
 |\partial_p^\alpha\mathcal B_{\rm ex}(p_0)|
 \leq\frac{\alpha!\,r^{-|\alpha|}}
 {768\pi^4(1-4r^2)}.
 \label{eq:two-loop-complex-derivative}
\end{equation}
For real external momenta $p,p'\in\RR^4$ one also has the
global stability estimate
\begin{equation}
 \begin{aligned}
 \norm{\nabla_p\mathcal A_{\rm ex}(p)}
 &\leq\frac{9}{8\sqrt3}\,\frac1{768\pi^4},\\
 |\mathcal A_{\rm ex}(p)-\mathcal A_{\rm ex}(p')|
 &\leq\frac{9}{6144\sqrt3\,\pi^4}\norm{p-p'}.
 \end{aligned}
 \label{eq:two-loop-real-Lipschitz}
\end{equation}
For real $p$, $\mathcal B_{\rm ex}(p)=\mathcal A_{\rm ex}(p)$.
The gradient is also globally Lipschitz. For $p,p',h\in\RR^4$,
\begin{equation}
 \begin{aligned}
 &\norm{\nabla\mathcal A_{\rm ex}(p)
             -\nabla\mathcal A_{\rm ex}(p')}\leq
             \frac{\norm{p-p'}}{384\pi^4},\\
 &|\mathcal A_{\rm ex}(p+h)-\mathcal A_{\rm ex}(p)
        -\nabla\mathcal A_{\rm ex}(p)\cdot h|\leq
             \frac{\norm{h}^2}{768\pi^4}.
 \end{aligned}
 \label{eq:two-loop-real-second-order}
\end{equation}
These are coefficient-level real and complex stability bounds, not
assertions about Lorentzian cuts or scattering amplitudes. The tube
can be checked explicitly here because all external dependence
resides in the massive third denominator under the selected routing.
\end{proposition}
\begin{proof}
Write $p=x+\ii y$ and $u=k_1-k_2+x\in\RR^4$. Then
\[
 \Re[(u+\ii y)^2+1]=|u|^2+1-|y|^2
 \geq1-|y|^2>0.
\]
The absolute integrand is therefore bounded by
$(1-|y|^2)^{-1}$ times the same product of the two selected
massless-line factors and its declared numerator. The already
computed one-line integrals yield
\eqref{eq:two-loop-complex-bound} and give a common integrable
majorant on every compact subset of $\mathcal T_1$.
Theorem~\ref{thm:holomorphy} proves joint holomorphy. On a
polydisc of coordinate radius $r$ about real $p_0$ one has
$|\Im p|^2\leq4r^2$; the multivariable Cauchy estimate
then proves \eqref{eq:two-loop-complex-derivative}.
For real $u$, the denominator has the explicit derivative bound
\[
 \norm{\nabla_u(1+|u|^2)^{-1}}
 =\frac{2|u|}{(1+|u|^2)^2}
 \leq\frac9{8\sqrt3},
\]
with equality for $|u|=1/\sqrt3$. The remaining, nonnegative
selected-line weight is independent of $p$ and has integral bounded
by \eqref{eq:two-loop-selected-certificate}. Dominated
differentiation proves the gradient estimate in
\eqref{eq:two-loop-real-Lipschitz}; integration along the real
line segment from $p$ to $p'$ proves its Lipschitz bound.
For $g(u)=(1+|u|^2)^{-1}$, moreover,
\[
 \nabla^2g(u)=-\frac{2I}{(1+|u|^2)^2}
             +\frac{8uu^T}{(1+|u|^2)^3}.
\]
Its tangential and radial eigenvalues are respectively
$-2/(1+t)^2$ and $(6t-2)/(1+t)^3$, with $t=|u|^2$.
The tangential eigenvalue has absolute value at most $2$; the radial
eigenvalue ranges from $-2$ at $t=0$ to a positive maximum $1/2$
at $t=1$. Thus $\norm{\nabla^2g(u)}_{\rm op}\leq2$.
The same integrable selected-line majorant controls the integrated
Hessian, giving
$\norm{\nabla^2\mathcal A_{\rm ex}(p)}_{\rm op}
\leq1/(384\pi^4)$. The fundamental theorem of calculus along
real segments and its integral Taylor remainder prove
\eqref{eq:two-loop-real-second-order}. For each fixed pair of loop momenta, $g(k_1-k_2+p)$ and its
first two derivatives tend to zero as $|p|\to\infty$.
Their globally bounded derivatives are dominated by the same
integrable selected-line weight; dominated convergence therefore
implies that $\mathcal A_{\rm ex}(p)$ and its first two
external derivatives vanish in that limit. This argument uses the
special external-momentum placement in this example and is not a
general consequence of graphwise UV finiteness.
\end{proof}

\subsection{Residual damping after nonpolynomial vertex growth}
\label{sec:residual-damping}

The polynomial bound in \eqref{eq:numerator-bound} is sufficient,
not necessary. The following test compares the \emph{complete}
vertex product with the damping in the same graph. It does not
construct the vertices required by a Ward identity.

\begin{proposition}[Residual-coercivity criterion]
\label{prop:residual-coercivity}
Retain the fixed routing and compact external set in
\eqref{eq:routing}, and assume \(0\leq C_i(q)\leq\mathcal I_i<\infty\).
Either of the following conditions gives uniform absolute convergence
of the graph.

\textit{Quadratic case.} Suppose
\(0\leq\Ph_i(a_i^2q_i^2)\leq B_i^{\rm env}e^{-d_i|q_i|^2}\), with
\(d_i>0\), and the complete numerator satisfies
\begin{equation}
 |N_G(k,p)|\leq C(1+\norm{k}^D)
            \exp\!\left(\sum_i v_i|q_i|^2\right),
 \label{eq:vertex-quadratic-budget}
\end{equation}
where \(v_i\geq0\) are fixed and \(D\geq0\). Put
\begin{equation}
 H_{\rm res}:=C^T\diag(d_i-v_i)C.
 \label{eq:residual-Gram-matrix}
\end{equation}
It is sufficient that \(H_{\rm res}\succ0\); individual
\(d_i-v_i\) need not all be positive. A Laplace gap supplies
\(d_i=a_i^2\tau_{i,*}\).

\textit{Linear case.} Suppose instead
\(0\leq\Ph_i(a_i^2q_i^2)\leq B_i^{\rm env}e^{-a_i|q_i|}\) and replace
\(v_i|q_i|^2\) in \eqref{eq:vertex-quadratic-budget} by
\(v_i|q_i|\). Put \(h_i=a_i-v_i\). The residual coefficients may
have either sign, provided
\begin{equation}
 \gamma_h:=\min_{\substack{\omega\in(\RR^4)^L\\\norm{\omega}=1}}
     \sum_i h_i\left|\sum_{\ell=1}^L C_{i\ell}\omega_\ell\right|>0.
 \label{eq:signed-linear-coercivity}
\end{equation}
In particular, nonnegative \(h_i\) and full loop rank of the rows
with \(h_i>0\) suffice. Here \(B_i^{\rm env}>0\) are finite
upper-envelope constants, not necessarily the masses \(\Ph_i(0)\)
of the representing measures. The growth coefficients have
inverse-momentum-squared units in the quadratic case and
inverse-momentum units in the linear case.
All constants and the positive coercivity margin are uniform on the
external set. The asserted conclusion is Euclidean and does not
supply a Lorentzian continuation.
\end{proposition}

\begin{proof}
The bounded covariances and finite \(B_i^{\rm env}\) can be absorbed
into the constant. In the quadratic case, the remaining exponent is
\(-\sum_i(d_i-v_i)|(Ck)_i+R_i(p)|^2\). Its quadratic part is
\(-k^T(H_{\rm res}\otimes\one_4)k\); its linear and constant
coefficients are bounded on \(\mathcal K\). If
\(\lambda=\lambda_{\min}(H_{\rm res})>0\), Young's inequality gives
\begin{equation}
 \sum_i(d_i-v_i)|q_i|^2
       \geq\tfrac12\lambda\norm{k}^2-C_{\mathcal K}.
 \label{eq:residual-coercive-bound}
\end{equation}
An integrable polynomial times a Gaussian therefore bounds the
complete graph.

For either sign of \(h_i\), the triangle inequalities imply
\[
 h_i|(Ck)_i+R_i|\geq h_i|(Ck)_i|-|h_i|\,|R_i|.
\]
Homogeneity and \eqref{eq:signed-linear-coercivity} therefore give
\begin{equation}
 \sum_i h_i|q_i|\geq\gamma_h\norm{k}
       -\sum_i|h_i|\sup_{p\in\mathcal K}|R_i(p)|.
 \label{eq:signed-linear-affine-bound}
\end{equation}
The remaining polynomial times \(e^{-\gamma_h\norm{k}}\) is
integrable. For a directly computable sufficient test, choose an
invertible \(C_S\) with \(h_*:=\min_{i\in S}h_i>0\). Then
\begin{equation}
 \gamma_h\geq h_*s_{\min}(C_S)
                  -\sum_{i:h_i<0}|h_i|\,\norm{C_i}_2,
 \label{eq:signed-linear-basis-certificate}
\end{equation}
where \(C_i\) is a routing row. The positive rows bound the selected
block norm, and each negative row has operator norm \(\norm{C_i}_2\)
on \((\RR^4)^L\). If all \(h_i\geq0\), the negative sum vanishes
and the previous full-rank sufficient test follows. Both cases give
common pointwise bounds, uniform on \(\mathcal K\).
For example, one loop with two equal routing rows and residuals
\((h_1,h_2)=(2,-1)\), in fixed inverse-momentum units, has
\(\gamma_h=1\). Thus a negative linewise residual need not destroy
convergence of the complete graph. The actual O--QED vertex product
must still establish the stated envelope; this example supplies no
new vertex or contact completion.
\end{proof}

A vanishing coercivity margin does not imply divergence: denominator
power counting or sector estimates may still suffice. Conversely,
rapid line decay alone does not establish a positive residual margin.
For example, cancelling both endpoint factors in
\eqref{eq:internal-damping-cancellation} leaves no such gain.
The criterion allows some growing vertices without assuming that the
Ward-compatible completion actually meets its hypotheses.

\subsection{A forest criterion for weaker linewise hypotheses}
\label{sec:forests}

When the available decay does not meet
Corollary~\ref{cor:damped-basis}, retaining the individual UV regions
can yield weaker sufficient conditions. Failure of that full-rank
criterion does not imply divergence. The relevant principle is the
all-region convergence criterion underlying Euclidean power counting
\cite{Weinberg1960}. Separate scaling checks are sufficient only when
their constants are uniformly controlled in the remaining scales.

\begin{definition}[Admissible sector envelope]
\label{def:sector-envelope}
An admissible finite sector decomposition for \(G\) covers the
ultraviolet region of loop-momentum space up to a null set, with
finite overlap; the remaining bounded region is covered by the
stated infrared estimate. Each sector parametrization includes its
Jacobian, and the transformed integrand has the stated bound.
The decomposition consists of finitely many sectors
\(\mathfrak h\) with scale variables
\(R_j\geq1\), \(j=1,\ldots,r_{\mathfrak h}\), and an absolute
ordinary envelope, including the measure, of the form
\begin{equation}
 C_{G,\mathcal K,\mathfrak h}
 \prod_{j=1}^{r_{\mathfrak h}}
     R_j^{\omega_0(\mathfrak h,j)-1}\dd R_j\,
 \dd\nu_{\mathfrak h},
 \qquad \int\dd\nu_{\mathfrak h}<\infty .
 \label{eq:sector-envelope}
\end{equation}
Angular, lower-scale, and infrared variables are controlled in
\(\dd\nu_{\mathfrak h}\). The degrees \(\omega_0\) include every
completed-action vertex and insertion. These are simultaneous
bounds on the whole sector, not pointwise bounds with constants
depending uncontrollably on other \(R_j\).
\end{definition}

Let a covered line \(e\) have a finite positive representing
measure \(\rho_e\), with \(\rho_e(\{0\})=0\), and suppose
\(M_{\beta_e,e}:=M_{\beta_e}[\rho_e]<\infty\) for an exponent
\(\beta_e>0\).
For \(z>0\),
\begin{equation}
 \Ph_e(z)\leq
 \left(\frac{\beta_e}{\mathrm e}\right)^{\beta_e}
 M_{\beta_e,e}\,z^{-\beta_e},
 \label{eq:moment-power-bound}
\end{equation}
because \(u^{\beta_e}e^{-u}\leq
(\beta_e/\mathrm e)^{\beta_e}\).
A sufficient simultaneous routing bound in sector \(\mathfrak h\)
is
\begin{equation}
 \abs{q_e(k,p)}\geq c_{e,\mathfrak h}
       \prod_jR_j^{\chi_{ej}},\qquad
 c_{e,\mathfrak h}>0,\quad \chi_{ej}\geq0.
 \label{eq:simultaneous-routing}
\end{equation}
It gives total decay
\begin{equation}
 \delta_\Ph(\mathfrak h,j)=2\sum_e\beta_e\chi_{ej},
 \qquad
 \omega_\Ph(\mathfrak h,j)=
       \omega_0(\mathfrak h,j)-\delta_\Ph(\mathfrak h,j).
 \label{eq:forest-degree}
\end{equation}
Only covered uncut internal lines appear in this sum. The simultaneous
bound or an equivalent explicit sector factorization prevents the
same gain from being counted twice for independent scales.

For the gap route, the appropriate condition is coercivity on each
scaled loop subspace. A sufficient form is
\begin{equation}
 \sum_{e\in E_{\rm ctrl}}\tau_{e,*}\abs{q_e(k,p)}^2
   \geq c_{\mathfrak h}\norm{k_{\rm scaled}}^2-C_{\mathcal K},
 \qquad c_{\mathfrak h}>0,
 \label{eq:gap-coercivity}
\end{equation}
with a uniform bound in the other sector variables.
The controlling kernel product then has Gaussian decay in every
scaled direction. One gapped line can control a subspace only if its
routing detects that entire subspace.

\begin{theorem}[Conditional forest convergence]
\label{thm:forest}
Assume an admissible finite sector decomposition and the stated
infrared control. In each sector, suppose the covered kernel product
has a simultaneous bound that either gives
\(\omega_\Ph(\mathfrak h,j)<0\) for every residual scale,
or gives Gaussian coercive control of that scale with the remaining
scales still integrable. Then the graph is absolutely UV convergent
on \(\mathcal K\).
\end{theorem}

\begin{proof}
In a pure moment sector, \eqref{eq:sector-envelope} becomes bounded
by
\[
 C\prod_jR_j^{\omega_\Ph(\mathfrak h,j)-1}\dd R_j\,
 \dd\nu_{\mathfrak h}.
\]
Each radial integral equals
\(-1/\omega_\Ph(\mathfrak h,j)<\infty\), and
\(\nu_{\mathfrak h}\) has finite mass. A coercive Gaussian replaces
the corresponding power test by an integrable Gaussian moment.
The simultaneous hypotheses permit these estimates to be used
together. Sum the finitely many sector majorants.
\end{proof}

The criterion applies once a finite simultaneous sector envelope has
been established; it does not construct that envelope for an arbitrary
completed action. Proper subgraphs and overlapping regions are covered
through their sector scales. Checking only the overall degree of the
graph is insufficient. On a fixed flat graph with all line kernels gapped,
a full-rank chord set automatically supplies global coercivity, as
already captured by Theorem~\ref{thm:global-uv}.

A commonly tempting determinant shortcut is false:
\[
 \mathcal U_G(w)\not\geq c\left(\sum_iw_i\right)^L
 \quad\hbox{uniformly on the whole positive cone}.
\]
For the sunset graph,
\(\mathcal U=w_1w_2+w_1w_3+w_2w_3\), while
\((w_1,w_2,w_3)=(1,\delta,\delta)\) makes
\(\mathcal U\to0\) as \(\delta\downarrow0\).
The chord minor \eqref{eq:chord-bounds} and sector envelopes avoid
this error.

\subsection{All orders, insertions, and finite matching}

\begin{corollary}[Finite-order meaning of all orders]
\label{cor:all-orders}
Fix a declared completed action and its kernel data on the selected
support and evolution branch, a scale \(a>0\), and an infrared
prescription. Suppose every finite loop graph \(G\) with
\(L(G)\geq1\), generated by this fixed set of data, satisfies
Theorem~\ref{thm:global-uv}, Corollary~\ref{cor:damped-basis},
Proposition~\ref{prop:residual-coercivity}, or
Theorem~\ref{thm:forest}, with the corresponding complete-numerator
and infrared bounds. Then every such loop graph is UV finite.
Tree coefficients contain no loop integration; they are included only
on external domains where their vertices and internal exchange factors
are defined, with any poles treated by the declared boundary prescription.
The loop estimates do not remove tree-level kinematic singularities.

For a finite collection of these graphs and a common compact external
set, their bounds can be chosen uniformly over the collection.
If residual field or support labels are summed, the summed amplitudes
additionally require finite support or an absolutely summable family
of graph majorants. This gives no uniform bound in perturbative
order and no convergence theorem for the infinite graph sum.
\end{corollary}

A gap, or inverse moments of every positive order in the fixed kernel
family, can accommodate every finite polynomial numerator degree in
the globally covered graph class. A single finite inverse moment does not automatically
cover an unbounded tower of derivative vertices. Infinite-valence
or nonlocal vertices also require control of graph multiplicity at a
fixed perturbative order.

External-momentum derivatives and composite insertions are covered
only when their \emph{complete differentiated integrands} have an
appropriate common majorant on the claimed domain. Derivatives of a
kernel can change both numerator growth and infrared behaviour;
replacing them by an assumed polynomial-degree cost without an
actual bound is insufficient. BRST and Nielsen insertions must pass
the same test.

For the covered graphs, mass, residue, and coupling conditions are
finite matching conditions on already finite amplitudes. They remain
physical choices of normalization and parameter determination.
UV finiteness does not make amplitudes independent of \(a\), of the
kernel family, or of the completed vertices. It removes the need
for UV-divergent subtraction on the stated domain.
\section{Exact BRST and Slavnov--Taylor identities of the completed action}
\label{sec:brst}

\subsection{Quantum BRST and background covariance}

Let \(G_{\rm gauge}\) be a compact gauge group with structure
constants \(f^{abc}\), and let \(\BR\) be a left odd derivation.
For homogeneous Lie-algebra-valued fields we use
\[
 [X,Y]^a=f^{abc}X^bY^c,\qquad
 (D_\mu[\mathcal A]X)^a
       =\partial_\mu X^a+g f^{abc}\mathcal A_\mu^bX^c.
\]
Grassmann parities are carried by the coefficients, so that
\([X,Y]=-(-1)^{\epsilon_X\epsilon_Y}[Y,X]\).
For \(\mathcal A=\overline{\mathcal A}+a\), take
\begin{align}
 \BR a_\mu^a&=(D_\mu[\mathcal A]c)^a,
 &\BR c^a&=-\frac g2 f^{abc}c^bc^c,
 \label{eq:brst-gauge}\\
 \BR\bar c^a&=B^a,&
 \BR B^a&=0,&
 \BR\overline{\mathcal A}_\mu^a&=0 .
 \label{eq:brst-doublet}
\end{align}
Matter fields transform in their declared gauge representations.
All quantities carry the same suppressed fixed label \(\bb\).
The field \(B\) is the Nakanishi--Lautrup auxiliary field, distinct
from the support label \(b_s\).

The graded Jacobi identity gives \(\BR^2c=0\); on the connection,
\begin{equation}
 \BR^2a_\mu
 =D_\mu(\BR c)+g[\BR a_\mu,c]
 =-\frac g2D_\mu[c,c]+g[D_\mu c,c]=0 .
 \label{eq:nilpotence}
\end{equation}
Nilpotence on the doublet is immediate. The representation identity
gives the matter-field version. These are algebraic statements;
they do not use the ultraviolet kernel.

Since the background is quantum-BRST inert,
\([\BR,\overline D_E]=0\), and functional calculus on a common
domain gives
\begin{equation}
 \BR\Ph(a^2\mathscr L_{\bb})=0.
 \label{eq:kernel-inert}
\end{equation}
Under a background gauge transformation the situation is different:
\begin{equation}
 \delta_\omega\overline{\mathcal A}_\mu=\overline D_\mu\omega,
 \qquad
 \delta_\omega a_\mu=g[a_\mu,\omega].
 \label{eq:background-transform}
\end{equation}
Ghosts, \(B\), and matter transform homogeneously in their respective
representations. The operator in \eqref{eq:positive-operator}
transforms by conjugation. Its functional calculus is therefore
background covariant. Background covariance and quantum BRST
invariance are compatible but distinct requirements
\cite{Abbott1981,BRST1976,Tyutin1975}.

\subsection{Why a completed interaction is required}
\label{sec:completion-required}

Suppose, in a fixed momentum representation,
\[
 S_\Ph(p)=F(p^2)S_0(p),\qquad
 S_\Ph^{-1}(p)=F(p^2)^{-1}S_0^{-1}(p).
\]
A Ward-compatible photon vertex must satisfy
\begin{equation}
 q_\mu\Gamma_\Ph^\mu(p+q,p)
   =S_\Ph^{-1}(p+q)-S_\Ph^{-1}(p).
 \label{eq:dressed-WTI}
\end{equation}
The bare vertex satisfies the same equation with \(S_0^{-1}\).
For nonconstant \(F\), the two right sides are generally different.
Thus \eqref{eq:kernel-inert} cannot prove the interacting Ward
identity for a propagator-only modification.

We assume a completed gauge-fixed action
\begin{equation}
 S_{\rm gf}^{\rm comp}
   =S_{\rm inv}^{\rm comp}+\BR\Psi,\qquad
 \BR S_{\rm inv}^{\rm comp}=0,
 \label{eq:completed-action}
\end{equation}
with all vertices and source interfaces generated consistently.
The existence of such an action with the desired propagators,
physical interfaces, and sufficient UV bounds is substantive input.

The first consistency conditions can be seen in a formal pure-gauge
expansion around the trivial background. Consider
\[
 S_{\rm inv}^{\rm comp}=S_2+gS_3+g^2S_4+\cdots,\qquad
 \BR=\BR_0+g\BR_1,
\]
where \(S_n\) is homogeneous of degree \(n\) in the gauge field,
\(\BR_0a_\mu=\partial_\mu c\), and
\(\BR_1a_\mu=[a_\mu,c]\). BRST invariance requires
\begin{equation}
 \BR_0S_2=0,\qquad
 \BR_0S_3=-\BR_1S_2,\qquad
 \BR_0S_4=-\BR_1S_3.
 \label{eq:quadratic-vertex-hierarchy}
\end{equation}
For two actions with this same BRST transformation,
\(\BR_0\Delta S_3=-\BR_1\Delta S_2\).
A changed quadratic kernel can retain the old cubic vertex only if
\(\BR_1\Delta S_2=0\). This hierarchy constrains the interaction
vertices; their momentum growth must additionally satisfy the UV
estimates for that same completed action.

\subsection{A covariant action and its interaction vertices}

A useful illustration shows why the interaction terms cannot be
ignored. Let \(H(z)=\Ph(z)^{-1}\) on the positive spectrum where
\(\Ph>0\), and formally define, on a common specified domain,
\begin{equation}
 S_{\rm inv}^{H}
 =\frac14\int\dd^4x\,
 \bigl\langle F_{\mu\nu},
 H(a^2[-D[\mathcal A]^2])F_{\mu\nu}\bigr\rangle_{\rm color}.
 \label{eq:covariant-H-action}
\end{equation}
Here \(\langle\ ,\,\rangle_{\rm color}\) is an invariant positive
color pairing. Covariance of \(F_{\mu\nu}\) and of \(D[\mathcal A]\)
implies gauge invariance. Around a flat zero background, its
transverse quadratic kernel is \(p^2H(a^2p^2)\), giving covariance
\(\Ph(a^2p^2)P_T/p^2\).

In the invariant action, \(H\) is evaluated on the full covariant
Laplacian and transforms by conjugation under quantum gauge
transformations. In the gauge-fixing sector below, the same scalar
function is evaluated on the background operator
\(\mathscr L_{\bb}\), which is quantum-BRST inert.

For the gauge-fixing function
\(\mathcal F=\overline D^\mu a_\mu\), consider
\begin{equation}
 \Psi_H=\int\dd^4x\,
 \bar c\,H(a^2\mathscr L_{\bb})
        \left(\mathcal F-\frac{\zeta}{2}B\right),
 \qquad \zeta>0.
 \label{eq:hyper-gauge-fermion}
\end{equation}
Since \(\BR H=0\),
\begin{equation}
 \BR\Psi_H
 =\int\dd^4x\left[
 B H\mathcal F-\frac{\zeta}{2}BHB
             -\bar c\,H\,\BR\mathcal F\right].
 \label{eq:hyper-gauge-expansion}
\end{equation}
With the convention \(\BR\bar c=B\), a literal Euclidean Gaussian
uses the usual imaginary \(B\)-contour, or its equivalent real
auxiliary variable. Eliminating \(B\) yields
\[
 S_{\rm gauge\ fixing}^{H}
       =\frac1{2\zeta}\langle\mathcal F,H\mathcal F\rangle .
\]
If \(\BR\mathcal F=-M c\), the ghost kernel is \(HM\), with
ordered inverse on the common invertible domain after gauge
zero modes have been removed:
\begin{equation}
 (HM)^{-1}=M^{-1}\Ph .
 \label{eq:ghost-ordering}
\end{equation}
It can be written \(\Ph M^{-1}\) only when these operators commute
on the relevant domain. The auxiliary Gaussian determinant
\(\det(\zeta H)^{-1/2}\) is retained in the measure or in a
declared compatible normalization; it need not be independent
of the background.

This is the Gaussian determinant in the chosen coefficient
representation. Its contribution to a carrier-valued logarithm is
formed from the two branch determinants, with their declared
logarithm branches, and the central idempotents. Multiplying the
scalar determinant or its logarithm by \(\ee\) is not that
reconstruction.

Expanding \(H(-D[\mathcal A]^2)\) in
\eqref{eq:covariant-H-action} generates interaction terms.
For an exponentially decreasing \(\Ph\), its inverse can grow
exponentially, and the induced vertices can compensate the
propagator suppression. This illustration establishes a
gauge-invariant algebraic construction; it does not establish the
UV hypotheses for its full vertex expansion, the special external
source interface, or unitarity. It exposes the precise additional
estimates required of a successful completion.

\subsection{BV fields, source signs, and the master identity}
\label{sec:bv}

Use a complete field set \(\phi^A\), including \(a,c,\bar c,B\)
and matter, with Grassmann parity \(\epsilon_A\).
Introduce antifields \(\phi_A^*\) of parity \(\epsilon_A+1\)
and ghost number \(-1-\operatorname{gh}\phi^A\).
The nonminimal pair \((\bar c,\bar c^*)\) is retained.
Our antibracket convention is
\begin{equation}
 (F,G)=\int\dd^4x\sum_A
 \left[
 \frac{\delta_rF}{\delta\phi^A}
 \frac{\delta_lG}{\delta\phi_A^*}
 -
 \frac{\delta_rF}{\delta\phi_A^*}
 \frac{\delta_lG}{\delta\phi^A}
 \right].
 \label{eq:antibracket}
\end{equation}
Left and right derivatives are not interchangeable on odd
coordinates. For a closed off-shell BRST algebra, choose
\begin{equation}
 \Sigma=S_{\rm gf}^{\rm comp}
       +\int\dd^4x\sum_A(-1)^{\epsilon_A}
                   \phi_A^*\,\BR\phi^A.
 \label{eq:bv-action}
\end{equation}
With this convention, \(\BR=(\Sigma,\,\cdot\,)\) on fields.
The displayed gauge and nonminimal terms are
\[
\begin{aligned}
 \Sigma={}&S_{\rm gf}^{\rm comp}
  +\int\dd^4x\,\Bigl[
       a^{*a\mu}(D_\mu c)^a
       -c^{*a}\BR c^a\\
   &\hspace{3.8em}-\bar c^{*a}B^a
       +\hbox{matter antifield terms}\Bigr].
\end{aligned}
\]
Gauge invariance and nilpotence yield
\((\Sigma,\Sigma)=0\).
For more general gauge systems the assumed completed master action
includes the additional antifield terms needed by its algebra.

The physical Lorentzian functional and the Euclidean coefficient
functional have distinct source conventions. To state the dictionary,
first choose the completed independent physical source coordinates;
any auxiliary-to-physical field map is included in the source pairing
and the transformed action. The BRST transformations are pushed
forward with that field map, and the antifield map preserves the BV
pairing. A nonlocal source normalization does not automatically
preserve an undeformed nonlinear gauge transformation. A real field has one source coordinate,
a genuinely complex pair has independent conjugate source coordinates,
and Dirac and ghost pairs retain their declared odd ordering. The
physical conjugation contour is imposed only after differentiation.

% Source conditions: lem:ST-stability-internal-only-early; thm:BRST-Nielsen.
% The following linear cotangent-lift identities are derived in the stated BV convention.
For example, let a Grassmann-even, parity- and ghost-number-preserving
linear map \(T\) be invertible on the admitted reduced domain,
with \(\BR T=0\). Write \(\psi^A\) for auxiliary coordinates
and \(\phi^A=T^A{}_B\psi^B\) for the physical source coordinates.
The source and antifield transformations are
\[
 \mathcal J^{\rm aux}_B=\mathcal J_A T^A{}_B,\qquad
 \psi_B^*=\phi_A^*T^A{}_B,\qquad
 \BR\phi^A=T^A{}_B\BR\psi^B .
\]
Repeated labels here include the spacetime integrations. The nonzero
entries of \(T\) are Grassmann even, have ghost number zero,
and connect fields of equal parity and ghost number.
At fixed background and other data determining
\(T\), the antifield transformation is the cotangent lift:
\[
 \sum_A\phi_A^*\,\delta\phi^A
  =\sum_B\psi_B^*\,\delta\psi^B,\qquad \delta T=0.
\]
It preserves the canonical field--antifield pairing. The same
transformation preserves the source pairing and the signed
antifield term:
\[
 \sum_B(-1)^{\epsilon_B}\psi_B^*\BR\psi^B
    =\sum_A(-1)^{\epsilon_A}\phi_A^*\BR\phi^A .
\]
Thus the canonical antibracket and its declared signs are
preserved character by character. For the corresponding
finite-dimensional change of variables, the measure factor is
\(\mathcal D\psi=\operatorname{Ber}(T^{-1})\mathcal D\phi\);
its continuum counterpart requires the declared determinant and
measure prescription. It may be absorbed only into a compatible
normalization.
For a nonlinear map one needs its graded canonical lift and
Jacobian on the reduced domain; the free covariance identity alone
does not provide them.

On the fixed carrier algebra, with antifields suppressed in this
display, the Lorentzian source template is
\begin{equation}
\begin{split}
 \vct Z^L_{\epsilon,\bb}[\mathcal J]
  =\int\mathcal D\phi\;
  \operatorname{Exp}_{\mathfrak A,\bb}\Bigg[
   &\ii u_{\epsilon,\bb}S_{{\rm sc},\epsilon,\bb}[\phi]
    +\ii a_{{\rm src},\epsilon,\bb}
       \int\dd^4x\,\mathcal J_A\phi^A\\
   &-\frac{\varepsilon}{2}\mathcal Q_{\rm B}[\phi]
  \Bigg].
 \label{eq:physical-L-source-functional}
% Monograph source: eq:A5-postGTL; def:A5-directed-exp-log.
\end{split}
\end{equation}
Here \(S_{{\rm sc},\epsilon,\bb}\) denotes the complete
gauge-fixed coefficient
\(S_{{\rm inv},{\rm sc}}^{\rm comp}+\BR\Psi_{\rm sc}\),
including the ghost sector; only antifield couplings are suppressed
in this display. The exponential is evaluated componentwise in the
central idempotents. The optional convergence form \(\mathcal Q_{\rm B}\)
is nonnegative on its declared commuting field block; fermionic and
ghost quadratic forms remain in the BRST action. The positive
convergence term is branch neutral. It is distinct from the ordered
Green-function datum \(\ii\eta\epsvec\); an exact regulated
matrix inverse requires specification of this datum with the
complete Lorentzian Hessian before inversion. If a convergence
regulator is retained in a BRST
identity, it must satisfy the stated BRST-invariance or BRST-exactness
hypothesis; positivity of \(\mathcal Q_{\rm B}\) alone does not
establish that hypothesis.

% Monograph source: def:A5-directed-exp-log; eq:A5-physical-carrier-involution.
Choose a normalization \(\vct{\mathcal N}\) independent of
sources and, when present, antifields, with nonzero character bodies.
It must also respect each background symmetry and extended BRST
transformation used in the corresponding functional identity.
On a source neighbourhood where both normalized character bodies
are nonzero, choose local logarithm branches. For an even character
\(z_\kappa=b_\kappa+n_\kappa\), with nonzero scalar body
\(b_\kappa\) and positive Grassmann degree \(n_\kappa\), this means
\[
 \log z_\kappa=\log b_\kappa+
  \sum_{r\geq1}\frac{(-1)^{r+1}}{r}
       \left(\frac{n_\kappa}{b_\kappa}\right)^r .
\]
The nilpotent sum terminates in a finite Grassmann algebra and is
otherwise a formal source-degree series. When a compatible vacuum
normalization gives \(\vct Z^L[0]\vct{\mathcal N}^{-1}=1\), use
the local logarithms through \(1\). On a physical conjugate-source
contour, the normalization and branches are chosen so that
\[
 \vct{\mathcal N}^{\ddagger}=\vct{\mathcal N},\qquad
 \operatorname{Log}_{\mathfrak A,\bb}(X^\ddagger)
  =\bigl(\operatorname{Log}_{\mathfrak A,\bb}X\bigr)^\ddagger .
\]
The physical connected generator and its mean fields are
\begin{equation}
 \vct{\mathcal W}^L_{\epsilon,\bb}
   =-\ii\operatorname{Log}_{\mathfrak A,\bb}
       (\vct Z^L_{\epsilon,\bb}\vct{\mathcal N}^{-1}),
 \qquad
 \bar\phi^A
  =a_{{\rm src},\epsilon,\bb}^{-1}
       \frac{\delta_l\vct{\mathcal W}^L_{\epsilon,\bb}}
            {\delta\mathcal J_A}.
 \label{eq:physical-connected-source-rule}
\end{equation}
For ordered source insertions on this same support,
\begin{equation}
\begin{aligned}
 \langle\phi^{A_1}\cdots\phi^{A_n}\rangle_c
  &=\ii^{1-n}a_{{\rm src},\epsilon,\bb}^{-n}
    \frac{\delta^n\vct{\mathcal W}^L_{\epsilon,\bb}}
     {\delta\mathcal J_{A_1}\cdots\delta\mathcal J_{A_n}},\\
 a_{{\rm src},\epsilon,\bb}^{-n}
   &=\eta^{-n}\ee^{\,n}\quad({\rm OQFT}).
\end{aligned}
 \label{eq:physical-source-derivative-parity}
\end{equation}
The rightmost derivative acts first, and the fields on the left
are in the displayed order. All source terms use source-first
coordinates: rewriting a field-first odd source term includes its
Koszul minus sign in the definition of that source coordinate.
For example, \(\bar\eta_\psi\psi+\bar\psi\eta_\psi\) becomes
\(\mathcal J_\psi\psi+\mathcal J_{\bar\psi}\bar\psi\) with
\(\mathcal J_\psi=\bar\eta_\psi\) and
\(\mathcal J_{\bar\psi}=-\eta_\psi\), and likewise for ghosts.
Their carrier power is even for even \(n\) and odd for odd \(n\),
independently of Grassmann parity. On a reduced source domain with
locally invertible graded source-to-mean-field map, the Lorentzian
Legendre transform is
\begin{equation}
 \vct\Gamma^L_{\epsilon,\bb}
  =\vct{\mathcal W}^L_{\epsilon,\bb}
   -a_{{\rm src},\epsilon,\bb}
        \int\dd^4x\,\mathcal J_A\bar\phi^A .
 \label{eq:physical-L-Legendre}
\end{equation}
All products and logarithms here belong to the graded,
source-complexified carrier algebra.

% Monograph source: eq:A5-total-physical-real-spaces; eq:A5-directed-connected-generator-algebra.
If the action, measure, integration contour and source pairing
respect the total \(\ddagger\) involution, then on the physical
conjugate-source contour
\[
 (\vct Z^L\vct{\mathcal N}^{-1})^\ddagger
       =\vct Z^L\vct{\mathcal N}^{-1},\qquad
 (\vct{\mathcal W}^L)^\ddagger=-\vct{\mathcal W}^L,\qquad
 (\vct\Gamma^L)^\ddagger=-\vct\Gamma^L .
\]
Thus the physical connected and effective functionals belong to
the Grassmann-even anti-fixed subspace of \(\mathfrak B_{\bb}\).
Their bosonic scalar bodies lie in
\(\operatorname{span}_{\RR}\{\ii,\ee\}\); this is a reality
condition, not a statement of homogeneous carrier parity.
Independent complex-source calculations still take place in the
full graded carrier algebra.

For the algebraic BRST proof below we work in one admitted
Euclidean character representation, with its source and antifield
coordinates normalized consistently. Thus the unbold symbols
\(\Sigma,Z,W,\Gamma\) in this subsection denote coefficient
functionals. Before any admitted continuation, the Lorentzian source
coefficient has character value
\(\chi_\kappa(a_{{\rm src},\epsilon,\bb})
 =\kappa\eta\) in O--QFT; this sign is part of the map from
\(\mathcal J\) to the coefficient source, not a second sector
fold. The corresponding Euclidean source functional and Legendre
transform are
\begin{align}
 Z[J,\phi^*]
 &=\frac1{\mathcal N}
     \int\mathcal D\phi\,
        \exp\left[-\Sigma[\phi,\phi^*]
                    +\int\dd^4x\,J_A\phi^A\right],
 \label{eq:Z-functional}\\
 W&=\log Z,\qquad
 \varphi^A=\frac{\delta_lW}{\delta J_A},\nonumber\\
 \Gamma[\varphi,\phi^*]&=\int\dd^4x\,J_A\varphi^A-W.
 \label{eq:Legendre}
\end{align}
The graded source pairing is written source first and used
consistently; \(J\) is eliminated at fixed antifields.
The sources have \(\epsilon(J_A)=\epsilon_A\) and
\(\operatorname{gh}(J_A)=-\operatorname{gh}(\phi^A)\), making
\(J_A\phi^A\) even and of ghost number zero.
The normalization \(\mathcal N\) is source and antifield
independent, background-gauge invariant, and chosen invariant
under the gauge-parameter BRST extension used below. A vacuum normalization is allowed
when it has these properties. The field- and antifield-dependent
tree action is \(\Sigma\); its irrelevant vacuum constant is
fixed by \(\mathcal N\).
Equation~\eqref{eq:directed-action} restores the regulator-free
physical action terms that have already been shown to be purely
carrier odd. A complete regulated or loop-corrected effective
functional is instead reconstructed from both characters, using one
signature and one Legendre convention for both representatives:
\begin{equation}
 \vct\Gamma_{\epsilon,\bb}
   =\mathsf P_+\Gamma^{\rm rep}_{\epsilon,\bb}
    +\mathsf P_-
                  \Gamma^{\rm rep}_{\bar\epsilon,\bb},
 \qquad \sigma_{\bar\epsilon}=-\sigma_\epsilon.
 \label{eq:effective-character-reconstruction}
% Monograph source: eq:Omni-connected-effective-carrier-restoration.
\end{equation}
The two representatives are the simultaneous character projections of
one completed functional, including its field and antifield coordinates.
The opposite-branch notation retains the projected source signs,
ordering, adjoint and endpoint rules, and logarithm branches. It is
not obtained by changing only the sign of an isolated denominator,
nor by choosing an unrelated solution of a scalar master equation.
Its even and odd carrier parts are respectively the half-sum and
the half-difference times \(\ee\). The Euclidean proof below
applies to the Euclidean coefficient convention. A physical
Lorentzian identity additionally requires an admitted continuation
that transports the source and antifield normalization and the
Legendre convention consistently. In general neither
\(u_{\epsilon,\bb}\Gamma^{\rm rep}_{\epsilon,\bb}\) nor one
transported complex coefficient reconstructs this functional.

\begin{theorem}[Conditional exact Slavnov--Taylor identity]
\label{thm:ST}
For the closed off-shell system and master action
\eqref{eq:bv-action}, assume an invariant
anomaly-free measure, and a common domain on which the BRST change
of variables, source differentiations, and Legendre transform are
defined without field-space boundary terms. Then the fixed-plane
effective action satisfies
\begin{equation}
 \mathcal S(\Gamma_{\bb})
      :=\frac12(\Gamma_{\bb},\Gamma_{\bb})=0
 \label{eq:ST-master}
\end{equation}
as a formal perturbative identity at every admitted loop order.
If the same completed graphs and required insertions satisfy
Section~\ref{sec:uv}, the identity holds for their finite
coefficients without UV-divergent symmetry-restoring subtraction.
\end{theorem}

\begin{proof}
Perform the infinitesimal change of variables
\(\phi^A\mapsto\phi^A+\theta\BR\phi^A\), with constant odd
\(\theta\), in \eqref{eq:Z-functional}.
The measure and the completed action contribute no breaking.
Since the sources are inert, the left-BRST variation of the source
term is
\[
 \BR(J_A\phi^A)=(-1)^{\epsilon_A}J_A\BR\phi^A .
\]
From \eqref{eq:bv-action},
\(\delta_l\Sigma/\delta\phi_A^*
   =(-1)^{\epsilon_A}\BR\phi^A\).
The change of variables therefore gives the functional identity
\begin{equation}
 \int\dd^4x\sum_A J_A
       \frac{\delta_l Z}{\delta\phi_A^*}=0,
 \qquad
 \int\dd^4x\sum_A J_A
       \frac{\delta_l W}{\delta\phi_A^*}=0.
 \label{eq:Z-W-identity}
\end{equation}
The Legendre transform with the same ordered pairing gives
\(\delta_r\Gamma/\delta\varphi^A=J_A\) and
\(\delta_l\Gamma/\delta\phi_A^*
      =-\delta_lW/\delta\phi_A^*\).
For even \(\Gamma\), the two terms of
\eqref{eq:antibracket} combine into twice the first term.
Substitution yields \eqref{eq:ST-master}.
Expanding that identity in the loop parameter produces the
Slavnov--Taylor hierarchy. The UV conclusion uses the independent
absolute bounds for these very graphs, rather than a cancellation
of divergent terms.
\end{proof}

% Monograph source: thm:BRST-Nielsen; prop:STexact.
% The regulator-insertion signs follow from eq:bv-action and eq:Legendre.
For completeness, an antifield-independent Euclidean regulator
\(R_\varepsilon[\phi]\) inserted as
\(\exp[-\Sigma-R_\varepsilon+J\cdot\phi]\) gives, with the
same invariant measure and ordinary Legendre transform,
\begin{equation}
 \int\dd^4x\,J_A
       \frac{\delta_l W_\varepsilon}{\delta\phi_A^*}
       =-\langle\BR R_\varepsilon\rangle,\qquad
 \frac12(\Gamma_\varepsilon,\Gamma_\varepsilon)
       =\left.\langle\BR R_\varepsilon\rangle
          \right|_{J=J(\varphi)} .
 \label{eq:regulated-ST-breaking}
\end{equation}
Here no regulator term is subtracted from the Legendre transform.
Indeed, the change of variables gives
\(\langle\BR(J\cdot\phi)\rangle
  =\langle\BR R_\varepsilon\rangle\), and the antifield and
Legendre signs then give the displayed equations. A BRST-inert
regulator, or an exact regulator \(R_\varepsilon=\BR\Upsilon_\varepsilon\),
has zero insertion. Otherwise its removal requires a limit in
which this insertion vanishes; an explicit factor of
\(\varepsilon\) alone does not prove that limit. The same
statement holds separately in both carrier characters.

The full carrier-valued identity requires the same hypotheses in
both character components. Project every field, antifield, source,
and coefficient simultaneously, and write
\(\Gamma_{\kappa}
  :=\chi_\kappa(\vct\Gamma)\).
Because the carrier coefficients are Grassmann even and central,
the character maps commute with the declared graded derivatives.
The carrier-valued Slavnov functional is
\begin{equation}
 \mathcal S_{\mathfrak A,\bb}(\vct\Gamma)
  :=\sum_{\kappa=\pm1}\mathsf P_\kappa
        \frac12(\Gamma_{\kappa},\Gamma_{\kappa})
  =0.
 \label{eq:carrier-ST}
% Monograph source: thm:BRST-Nielsen; eq:A5-algebra-valued-Slavnov-functional.
\end{equation}
Each bracket uses its projected canonical field and antifield
coordinates, with exactly the ordering of \eqref{eq:antibracket}.
For a retained regulator this statement additionally requires a
BRST-inert or BRST-exact regulator, and its removal must leave no
anomalous Jacobian or field-space boundary term. Establishing the
identity for only one scalar representative does not determine the
other character component.

When the linear nonminimal dependence is kept as
\(\Gamma=\Gamma_{\rm red}-\int\bar c^*B\), the master equation
is equivalently
\begin{equation}
 \frac12(\Gamma_{\rm red},\Gamma_{\rm red})_{\rm min}
       +\int\dd^4x\,B^a
             \frac{\delta_l\Gamma_{\rm red}}{\delta\bar c^a}=0 .
 \label{eq:reduced-ST}
\end{equation}
The minimal bracket includes gauge, ghost, and matter antifields.
A bracket containing only the minimal pairs cannot by itself
represent the displayed \(B\,\delta\Gamma/\delta\bar c\) term.

Restoring \(\hbar\) as a formal loop-counting parameter, write
\(\Gamma=\Sigma+\sum_{\ell\geq1}\hbar^\ell\Gamma_\ell\).
The linearized operator
\(\mathcal B_\Sigma F=(\Sigma,F)\) obeys
\(\mathcal B_\Sigma^2=0\), and
\begin{equation}
 \mathcal B_\Sigma\Gamma_\ell
 =-\frac12\sum_{r=1}^{\ell-1}
                 (\Gamma_r,\Gamma_{\ell-r}).
 \label{eq:loop-ST}
\end{equation}
The effective action satisfies the nonlinear Slavnov--Taylor identity.
BRST-closed physical representatives need not be BRST-exact.

\subsection{Anomalies and the quantum action principle}

The anomaly-free measure hypothesis is essential. UV convergence
does not cancel an anomalous Jacobian or prove that a formal
field-space integration has no boundary term. Algebraic
renormalization provides a powerful cohomological organization
\cite{PiguetSorella,Slavnov1972,Taylor1971,ZinnJustin1975},
but its locality and quantum-action-principle hypotheses must be
established for the admitted nonlocal completion and insertions.
They are not inherited merely because the undeformed local theory
possessed them.

Theorem~\ref{thm:ST} uses the explicitly assumed invariant
completed functional. It is not an existence proof of a
nonperturbative measure, a proof that a chiral representation is
anomaly free, or a proof of global gauge fixing. These tasks are
part of constructing the completion to which the theorem applies.

\subsection{Background Ward identity}

Assume the complete action, measure, domains, and antifield/source
transformations are covariant under
\eqref{eq:background-transform}. A background change of variables
then yields
\begin{equation}
 \mathcal W_\omega\Gamma_{\bb}=0.
 \label{eq:background-Ward}
\end{equation}
For a functional \(F\), the generator is most unambiguously
specified by its variations:
\begin{align}
 \mathcal W_\omega F
 =\int\dd^4x\bigg\{&
  (\overline D_\mu\omega)^a
      \frac{\delta F}{\delta\overline{\mathcal A}_\mu^a}
      \nonumber\\
 &+\sum_A(\delta_\omega\phi^A)
               \frac{\delta_lF}{\delta\phi^A}
  +\sum_A(\delta_\omega\phi_A^*)
               \frac{\delta_lF}{\delta\phi_A^*}\bigg\}.
 \label{eq:Ward-generator}
\end{align}
Here \(\phi^A\) is the complete quantum field set, and the
antifields transform in the corresponding dual representations.
When \(\mathcal W_\omega\) acts on \(\Gamma\), its quantum-field
arguments are the mean fields \(\varphi^A\).
Cyclicity of an invariant trace shows covariance of a kernel
insertion, but invariance of the whole functional is required for
\eqref{eq:background-Ward}.

\subsection{Nielsen identity and the meaning of gauge independence}
\label{sec:Nielsen}

Let the gauge parameter \(\zeta\) enter only through the
gauge-fixing fermion. The derivative below is taken at fixed
canonical mean fields and antifields, fixed support and carrier,
and fixed kernel and physical-source normalization maps. A
\(\zeta\)-dependent source or field redefinition would contribute
its own variation and is not included in this fixed-map formula.
Introduce the BRST doublet
\[
 \BR\zeta=\chi,\qquad \BR\chi=0,\qquad\chi^2=0.
\]
Assume the extended completed functional and its insertions obey
the extended identity, in the convention of
\eqref{eq:antibracket},
\begin{equation}
 \frac12(\Gamma,\Gamma)+\chi\,\partial_\zeta\Gamma=0.
 \label{eq:extended-ST}
\end{equation}
Set
\(\Xi_\zeta=(\partial_l\Gamma/\partial\chi)|_{\chi=0}\).
Extracting the coefficient of \(\chi\) gives
\begin{equation}
 \partial_\zeta\Gamma\big|_{\chi=0}
       =(\Gamma,\Xi_\zeta)\big|_{\chi=0}
       =\mathcal B_\Gamma\Xi_\zeta\big|_{\chi=0}.
 \label{eq:Nielsen}
\end{equation}
For example, at the classical level
\(\Xi_\zeta=\partial_\zeta\Psi\) in a compatible gauge-fixing
construction, so the variation is BRST exact.
This is the Nielsen identity in the declared convention
\cite{Nielsen1975}.

The corresponding carrier-valued Nielsen identity is read in the
same componentwise sense:
\begin{equation}
 \partial_\zeta\vct\Gamma\big|_{\chi=0}
  =\sum_{\kappa=\pm1}\mathsf P_\kappa
       (\Gamma_{\kappa},\Xi_{\zeta,\kappa})\big|_{\chi=0},
 \qquad
 \Xi_{\zeta,\kappa}
  :=\chi_\kappa(\vct\Xi_\zeta).
 \label{eq:carrier-Nielsen}
% Monograph source: thm:BRST-Nielsen.
\end{equation}
The background Ward identity is lifted with the same simultaneous
projection of all variables. These lifts preserve the coefficient
sign conventions; they do not assign a fixed carrier parity to a
nonlinear functional or to its loop corrections.

Gauge independence concerns admissible, \(\zeta\)-independent
BRST-closed observables, or families representing the same physical
cohomology classes under the Nielsen flow. For these observables and
physical LSZ states, the physical normalizations must be fixed and
the necessary spectral and insertion limits controlled. Off-shell
Green functions generally depend on \(\zeta\).
Varying the physical kernel scale \(a\) or the choice of \(\rho\)
changes the dynamics; neither is a gauge-parameter variation
covered by \eqref{eq:Nielsen}.
\section{Analyticity, positivity, and the physical cut interface}
\label{sec:analytic}

\subsection{Complete monotonicity and analytic continuation}

Equation~\eqref{eq:bernstein} gives a holomorphic function on
\(\Re z>0\). For a compact subset of this half-plane, derivatives
may be taken under the integral because
\(\tau^n e^{-c\tau}\) is bounded for every \(c>0\).
If
\begin{equation}
 \int e^{\epsilon_\rho\tau}\rho(\dd\tau)<\infty
 \quad\hbox{for some }\epsilon_\rho>0,
 \label{eq:exponential-moment}
\end{equation}
the same argument extends holomorphy to
\(\Re z>-\epsilon_\rho\).
Compact support of \(\rho\), or positive exponential moments of
every order, gives an entire Laplace transform. A general
completely monotone function need not extend to a slit plane
\cite{Bernstein1929,Widder1941}.

\begin{proposition}[A gapped counterexample to automatic slit analyticity]
\label{prop:meromorphic-counterexample}
The normalized measure and kernel
\begin{equation}
 \rho=\sum_{n=1}^\infty2^{-n}\delta_n,\qquad
 \Ph(z)=\frac1{2e^z-1}
 \label{eq:meromorphic-counterexample}
\end{equation}
satisfy the gap condition with \(\tau_*=1\), but their complex
continuation has poles away from the negative real axis.
\end{proposition}

\begin{proof}
For \(z\geq0\), the positive geometric series
\[
 \Ph(z)=\sum_{n\geq1}2^{-n}e^{-nz}
\]
proves complete monotonicity and \(\Ph(0)=1\). The gap bound is
\(\Ph(z)\leq e^{-z}\). The denominator vanishes at
\[
 z=-\log2+2\pi\ii n,\qquad n\in\mathbb Z.
\]
For \(n\neq0\), these singularities lie off the negative real axis.
\end{proof}

Even the elementary power kernel
\(\Ph(z)=(1+z/b)^{-\nu}\) has a singularity at \(z=-b\):
a pole for positive integer \(\nu\), and generally a branch point.
Strong Euclidean decay therefore does not imply a pole-free
Lorentzian continuation.

A second distinction concerns growth. Define the directed kinetic
invariant on the fixed plane by
\begin{equation}
 z_{\eta,\bb}(k):=k^2+\ii\eta\epsvec,
 \qquad \epsvec=\varepsilon\ee,
 \qquad \eta=\sigma_s\sigma_\epsilon.
 \label{eq:directed-kinetic-invariant}
\end{equation}
The continuation of \(e^{-tq_E^2}\) is the directed function
\begin{equation}
 \ExpDir\!\left(tz_{\eta,\bb}(k)\right)
 =\embed\!\left[e^{t(k^2+\ii\eta\varepsilon)}\right].
 \label{eq:heat-continuation}
\end{equation}
Its extracted coefficient grows as \(e^{tk^2}\) on real
timelike directions. It is not an oscillatory semigroup factor.
Entireness alone does not justify discarding a large contour arc.

\subsection{Local holomorphy from a common majorant}

Take independent external complex momenta \(z\) after imposing
momentum conservation. In an all-incoming convention,
\(p_n=-\sum_{\ell<n}p_\ell\). It is inconsistent to impose
\(\Im p_\ell^0>0\) on every one of these \(n\) legs.

\begin{theorem}[Holomorphy from a common graph majorant]
\label{thm:holomorphy}
Let \((\Xi,\nu)\) be the positive, sigma-finite measure space of a graph
representation on a connected open set \(\Omega\subset\CC^N\).
For each \(z\), assume \(J(\cdot;z)\) is measurable, and outside
one \(\nu\)-null set assume \(J(\xi;\cdot)\) is holomorphic on
\(\Omega\). Suppose that for every \(U\Subset\Omega\) there is
\(M_U\in L^1(\dd\nu)\) with
\begin{equation}
 |J(\xi;z)|\leq M_U(\xi),\qquad z\in U.
 \label{eq:complex-majorant}
\end{equation}
The bound holds outside a common null set for all \(z\in U\).
Then \(\mathcal M(z)=\int J(\xi;z)\dd\nu(\xi)\) is jointly
holomorphic on \(\Omega\). Local contour integrals and
derivatives may be exchanged with the representation integrals.
\end{theorem}

\begin{proof}
Fubini permits integration over a small contour in any one complex
variable before integration in \(\xi\). Cauchy's theorem makes
that contour integral vanish. Dominated convergence supplies
continuity and local boundedness. Morera's theorem gives separate
holomorphy, and Hartogs' theorem gives joint holomorphy.
Cauchy's formula on a slightly larger compact polydisc supplies an
integrable derivative bound.
\end{proof}


\begin{corollary}[Cauchy budgets for external derivatives]
\label{cor:complex-derivative-budget}
Under Theorem~\ref{thm:holomorphy}, suppose a closed polydisc
$\overline{\mathbb D}(z_0;r_1,\ldots,r_N)$ is inside the admitted
complex domain and one $M\in L^1(\dd\nu)$ majorizes the complete
integrand there. For a multi-index $\alpha\in\mathbb N_0^N$,
\begin{equation}
 |\partial_z^\alpha\mathcal M(z_0)|
  \leq\alpha!\prod_{j=1}^N r_j^{-\alpha_j}
                  \int_\Xi M(\xi)\dd\nu(\xi).
 \label{eq:complex-Cauchy-budget}
\end{equation}
The estimate is for the stated complex-domain representation; a
uniform bound on an \emph{integrated real} Euclidean amplitude is not
a substitute for its pointwise complex majorant. No LSZ, retarded, or
physical boundary conclusion is obtained from this inequality alone.
\end{corollary}
\begin{proof}
The common majorant permits the graph integration and the
multivariable Cauchy contour integrals to be interchanged. Estimate
the resulting contour integrals by $\int M$ and the product of the
inverse radii.
\end{proof}

% Monograph source: lem:wick; eq:external-energy-majorant-bound; rem:two-forest-complex-majorant.
\begin{lemma}[A sufficient massive complex neighbourhood]
\label{lem:massive-neighbourhood}
Fix a compact set \(\mathcal K\) of real Euclidean external data
and use the routing of \eqref{eq:routing}, with momentum conservation
already imposed. Let every Schwinger factor be
\(E_i(\alpha)=e^{-m_i^2\alpha}\), with \(m_i>0\).
Suppose that on an open complex neighbourhood \(\mathcal U_0\)
of \(\mathcal K\), the full numerator has the finite expansion
\[
 N_G(k,z)=\sum_{|\gamma|\leq D}n_\gamma(z)k^\gamma,
 \qquad n_\gamma\in\mathcal O(\mathcal U_0),
\]
with this same loop-momentum degree \(D\) throughout
\(\mathcal U_0\). Assume the gap or inverse-moment hypothesis of
Theorem~\ref{thm:global-uv} for that \(D\), and
\eqref{eq:exponential-moment}. Then there is an open complex
neighbourhood \(\mathcal U\) of \(\mathcal K\), with
\(\overline{\mathcal U}\Subset\mathcal U_0\), on which the
pre-Gaussian Schwinger--Laplace graph has one common pointwise
\(L^1\) majorant. Its integral is holomorphic on each connected
component of \(\mathcal U\).
\end{lemma}

\begin{proof}
Choose finitely many real centres \(p^{(r)}\in\mathcal K\)
and sufficiently small complex balls around them whose closures lie
in \(\mathcal U_0\). The finite holomorphic coefficient family is
bounded on their union, so
\( |N_G(k,z)|\leq C(1+\norm{k}^{D})\) there with the same \(D\).
For a ball of radius \(\delta\), write
\[
 q_i(k,z)=q_i(k,p^{(r)})+h_i(z),\qquad
 |h_i(z)|\leq C_G\delta.
\]
The first term is real. With \(q^2=\sum_\mu q_\mu^2\),
the complex bilinear square, one has
\[
 \Re\bigl(q_i(k,z)^2\bigr)
 \geq\tfrac12|q_i(k,p^{(r)})|^2-C'_G\delta^2.
\]
Consequently a fixed-centre pointwise envelope is
\[
 C(1+\norm{k}^{D})\prod_i
 e^{-\frac12(\alpha_i+a^2\tau_i)|q_i(k,p^{(r)})|^2}
 e^{-\alpha_i(m_i^2-C'_G\delta^2)}
 e^{a^2C'_G\delta^2\tau_i}.
\]
Shrink \(\delta\) so that
\(C'_G\delta^2<\frac12\min_i m_i^2\) and
\(\lambda:=a^2C'_G\delta^2<\epsilon_\rho\).
The reduced masses remain positive, and the tilted measure
\(\rho_\lambda(\dd\tau):=e^{\lambda\tau}\rho(\dd\tau)\)
has finite mass. It has the same support and no new atom at zero.
For every inverse moment \(M_r[\rho]\) used in the UV estimate,
\[
 M_r[\rho_\lambda]
 \leq e^\lambda M_r[\rho]
       +\int_{(1,\infty)}e^{\lambda\tau}\rho(\dd\tau)
 <\infty.
\]
Thus the global Gaussian proof, with the widths multiplied by
\(1/2\), makes this envelope integrable in \(k,\alpha,\tau\).
The finite sum over the covering centres is a common majorant on
their union. The integrand is holomorphic there for every finite
parameter point, so Theorem~\ref{thm:holomorphy} applies on each
connected component.
\end{proof}

This is a local statement near Euclidean data. It does not produce
a full physical sheet or a scattering crossing path.
For stable kernels such as \(e^{-\sqrt z}\), the representing
measure has no positive exponential moment, so
Lemma~\ref{lem:massive-neighbourhood} does not apply. Their explicit
square-root continuation can still be studied on a declared
domain, but its branch geometry and graph majorant require an
independent argument. This limitation can be substantive rather
than merely a failure of a sufficient proof: the literal massive
stable-half scalar bubble is not \(C^6\) at zero Euclidean
external momentum, as proved in
Proposition~\ref{prop:stable-bubble-nonanalytic}.

The complex variables in Theorem~\ref{thm:holomorphy} and
Lemma~\ref{lem:massive-neighbourhood} are ordinary coefficient
variables. The Lorentzian boundary value is fixed separately by
\(\ii\eta\epsvec\); it is not derived by this Euclidean
holomorphy argument. For an invariant coefficient function
\(F^{\rm sheet}\), its directed continuation and extracted
boundary value are
\begin{align}
 F^{\rm dir}_{\bb}
   &:=\embed\circ F^{\rm sheet}\circ\embed^{-1},
       \nonumber\\
 F^{\rm coeff}_{\eta,\bb}(s)
   &:=\lim_{\varepsilon\downarrow0}
       \embed^{-1}\!\left[
       F^{\rm dir}_{\bb}(s+\ii\eta\epsvec)\right].
 \label{eq:analytic-function-directed-interface}
\end{align}
The entire physical graph is then reconstructed by
\eqref{eq:complete-graph-character-reconstruction}.
The same distinction applies to multiple invariant variables.

\subsection{Candidate poles, branches, and contour deformations}

Suppose a continued line numerator \(f\) is holomorphic near
\(D(k)=k^2-m^2=0\). If \(f\) is finite and nonzero there,
multiplication by \(f\) preserves the simple pole and multiplies
its residue by \(f|_{D=0}\). A zero can soften or cancel the pole.
A kernel singularity can add a pole, increase an existing order,
or introduce a new branch surface. Consequently, the precise
inertness statement is
\begin{quote}
The ordinary denominator family is unchanged. Pole preservation
holds locally where the numerator is holomorphic, finite, and
nonzero; the singularities and growth of that numerator require
their own analysis.
\end{quote}

The ordinary denominator Landau equations remain candidate
conditions associated with those denominators
\cite{Landau1959}. They are not the complete singularity
analysis when the kernel or completed vertices have additional
singular surfaces.

A permissible deformation of a loop contour requires an analytic
integrand between the contours, no crossed singularities or
pinches, and sufficient bounds on boundary arcs. A holomorphic
Euclidean parameter integral can instead define a continuation
without such a deformation, as in
Section~\ref{sec:gaussian-bubble}.
These two constructions should not be conflated.
Compact real-momentum bounds do not, on their own, establish
tempered physical boundary values.

\subsection{Reflection positivity is an independent requirement}
\label{sec:OS}

Let \(\theta(x_E^0,\mathbf x)=(-x_E^0,\mathbf x)\), and let
\(\Theta\) include reflection and conjugation. The ungraded formulas
in this subsection concern the Grassmann-even, gauge-invariant
observable algebra. Reflection positivity on that algebra requires
\begin{equation}
 \sum_{i,j}\overline{c_i}c_j
    \langle\Theta F_i\,F_j\rangle_E\geq0
 \label{eq:OS}
\end{equation}
for all finite families of observables smeared at
\(x_E^0>0\). Ordinary positivity of a Euclidean convolution
kernel is different. In particular, \(e^{t\Delta_E}\) spreads
positive-time support across \(x_E^0=0\). Moving this heat
operator onto test functions does not preserve the OS test space.

\begin{proposition}[Heat damping does not preserve OS positivity]
\label{prop:OS-counterexample}
For \(t>0\) and \(\omega>0\), the following time covariance
fails reflection positivity:
\begin{equation}
 C_t(u)=\int_{\RR}\frac{\dd p}{2\pi}\,
       \frac{e^{\ii pu}e^{-tp^2}}{p^2+\omega^2}
 \label{eq:heat-time-covariance}
\end{equation}
\end{proposition}

\begin{proof}
All time moments exist. Set
\[
 c_0=\int\frac{\dd p}{2\pi}
           \frac{e^{-tp^2}}{p^2+\omega^2}>0,\qquad
 c_2=\int\frac{\dd p}{2\pi}
           \frac{p^2e^{-tp^2}}{p^2+\omega^2}>0 .
\]
Then \(C_t(u)=c_0-\tfrac12c_2u^2+O(u^4)\).
Positive-time probes at \(u=\epsilon,2\epsilon\) would require
the Hankel matrix
\[
 \begin{pmatrix}
 C_t(2\epsilon)&C_t(3\epsilon)\\
 C_t(3\epsilon)&C_t(4\epsilon)
 \end{pmatrix}
\]
to be positive semidefinite. Its determinant is
\begin{equation}
 -c_0c_2\epsilon^2+O(\epsilon^4)<0
 \label{eq:OS-negative-determinant}
\end{equation}
for sufficiently small \(\epsilon>0\).
Smooth positive-time approximations to these point probes preserve
the violation by continuity.
\end{proof}

At fixed spatial momentum, a four-dimensional heat-damped scalar
covariance reduces to \eqref{eq:heat-time-covariance} times a
positive constant, with
\(\omega^2=\mathbf p^2+m^2\).
Spatial wave packets localized near that momentum give the same
obstruction. Thus the failure is not an artifact of a nonnormalizable
single spatial mode.

There is a related spectral restriction. If a nonzero physical
scalar two-point coefficient has an unsubtracted positive spectral
representation
\begin{equation}
 C_{\rm phys}(x)=\int_{[0,\infty)}
          \frac{\rho_{\rm phys}(\dd\mu^2)}{x+\mu^2},
 \qquad \rho_{\rm phys}\geq0,\quad x>0,
 \label{eq:physical-spectral}
\end{equation}
then monotone convergence implies
\begin{equation}
 \lim_{x\to\infty}xC_{\rm phys}(x)
       =\rho_{\rm phys}([0,\infty))\in(0,\infty].
 \label{eq:spectral-obstruction}
\end{equation}
A covariance \(\Ph(a^2x)/(x+m^2)\) with
\(\Ph(a^2x)\to0\) cannot itself be that nonzero positive physical
spectral covariance. Auxiliary virtual-line data must therefore
be distinguished from the full positive physical two-point
function. The broader role of spectral positivity in Euclidean
reconstruction is established in
Refs.~\cite{OS1973,OS1975}; rational-propagator obstructions are
also discussed in Ref.~\cite{Arici2018}.

\begin{assumption}[Physical OS completion]
\label{ass:OS-completion}
The completed Schwinger functions of the admitted Grassmann-even,
gauge-invariant observable fields obey reflection positivity
\eqref{eq:OS}, Euclidean covariance, permutation symmetry, and
\(S_0=1\). We assume the sufficient regularity condition \(E0'\)
of Ref.~\cite{OS1975}, Section~IV.1, Eq.~(4.1): for one integer
\(r_{\rm OS}\geq1\) and constants \(C_{\rm OS},\beta_{\rm OS}>0\),
\begin{equation}
 \begin{aligned}
  |S_n(f)|&\leq\sigma_n\,\mathfrak p^{\rm OS}_{n r_{\rm OS}}(f),\\
  0<\sigma_n&\leq C_{\rm OS}(n!)^{\beta_{\rm OS}},\qquad n\geq1.
 \end{aligned}
 \label{eq:OS-hierarchy-growth}
\end{equation}
Here \(S_n\) is the full \(n\)-point Schwinger distribution,
not a connected cumulant or a coefficient of the coupling expansion.
The space \(\mathscr S_0(\RR^{4n})\) consists of Schwartz functions
whose derivatives of every order vanish on coincident-point diagonals;
\(\mathfrak p^{\rm OS}_j\) is the seminorm in Ref.~\cite{OS1975},
Section~II, Eq.~(2.1). Use self-adjoint observable components, or
include their adjoints and the corresponding conjugation in reflection.
For a finite multiplet, impose the conditions jointly on all mixed
correlations, not only on each component separately. Clustering is
additionally assumed when vacuum clustering or uniqueness is invoked.
\end{assumption}

These are assumptions on the completed correlation functions,
not consequences of convergence at each perturbative order.
The integer \(n\) in \eqref{eq:OS-hierarchy-growth} counts
observable insertions, not loops, and the displayed constants are
independent of \(n\). A finite bound for each graph does not supply
this hierarchy estimate or reflection positivity of a full correlation
functional. Nor is positivity a coefficientwise requirement: for
real \(g\), the matrix
\[
 \begin{pmatrix}1&g\\g&g^2\end{pmatrix}
 =\begin{pmatrix}1\\g\end{pmatrix}\begin{pmatrix}1&g\end{pmatrix}
 \succeq0
\]
has a first-order truncation with determinant \(-g^2\).
This elementary example illustrates the truncation issue; it is not
itself a Schwinger hierarchy or an OS counterexample.
The applicable bosonic OS reconstruction then provides a positive
physical Hilbert space and tempered Wightman distributions with
the spectrum condition and local commutativity on the admitted
observable algebra \cite{OS1973,OS1975,StreaterWightman}.
In particular, reconstructed observables obey
\[
 [\mathcal O_{\bb}(x),\mathcal O'_{\bb}(y)]=0
 \quad\hbox{for }(x-y)^2<0.
\]
If physical odd fields are included, their reconstruction needs the
corresponding spin, graded reflection and graded permutation rules.
The locality statement then uses the graded commutator, not an
ordinary commutator for two odd fields. The ungraded assumption above
does not supply that fermionic reconstruction.
No such reconstruction is obtained from positive Laplace weights
alone. Clustering, global distributional bounds, and positivity
must concern the completed Schwinger functions; the absence of a
direct external multiplier does not prove them.
\subsection{External attachments and interacting LSZ residues}
\label{sec:LSZ}

The external-leg placement convention means that the LSZ attachment
and amputation operators are not themselves multiplied by \(\Ph\).
It does not mean that the interacting two-point function is the
same as in the undeformed theory.

The following single-sheet transport describes a resolved scalar
two-point channel whose completed carrier coefficient belongs to
the directed subalgebra. Equivalently, its two complete character
coefficients obey the reflection relation specified below. For a
general mixed carrier two-point function, pole extraction is performed
in both characters before reconstruction; choosing one sheet does
not impose that reflection condition on the other character.

For a chosen completed-theory sheet, let
\(\mathscr R^{(2),\rm sheet}_{\eta,\bb}(z)\) be the
ordinary meromorphic two-point coefficient with its propagator
numerator removed. Here \(m_{{\rm phys},\bb}\) is the
momentum-dimension shell parameter of Section~\ref{sec:framework};
with explicit units it is the physical mass times \(\cGT\). Assume that a neighbourhood of the physical
mass contains precisely one simple pole:
\begin{equation}
 \mathscr R^{(2),\rm sheet}_{\eta,\bb}(z)
  =\frac{Z_{\bb}}{z-m_{\rm phys,\bb}^2}
            +H^{\rm sheet}_{\eta,\bb}(z),
 \qquad 0<|Z_{\bb}|<\infty,
 \label{eq:full-pole}
\end{equation}
where \(H^{\rm sheet}_{\eta,\bb}\) is holomorphic across the
pole. This is a sheet statement before a boundary distribution
is taken. Set
\begin{align}
 \mathscr R^{(2),\rm dir}_{\eta,\bb}
   &:=\embed\circ
          \mathscr R^{(2),\rm sheet}_{\eta,\bb}
                   \circ\embed^{-1},\nonumber\\
 \mathcal R^{(2),\rm coeff}_{\eta,\bb}(p)
   &:=\lim_{\varepsilon\downarrow0}
       \embed^{-1}\!\left[
       \mathscr R^{(2),\rm dir}_{\eta,\bb}
              (p^2+\ii\eta\epsvec)\right].
 \label{eq:LSZ-explicit-coefficient-extraction}
% Monograph source: prop:LSZ; eq:LSZ-directed-two-point-coefficient.
\end{align}
Under this directed two-point hypothesis, and only after a compatible
source/endpoint channel has been specified, the oriented boundary
template is
\begin{align}
 \vec G^{(2)}_{\eta,\bb}(p)
 &=\eta\ii\widehat u_p\,
       \embed\!\left(
           \mathcal R^{(2),\rm coeff}_{\eta,\bb}(p)\right)
       \nonumber\\
 &=\embed\!\left(\sigma_\epsilon\ii
           \mathcal R^{(2),\rm coeff}_{\eta,\bb}(p)\right)
       \quad\text{when }\widehat u_p=\sigma_s\ee.
 \label{eq:LSZ-physical-directed-kernel}
\end{align}
The first line is only an oriented component or raywise pole template.
It is not, by itself, the covariance of an identical real scalar:
Proposition~\ref{prop:source-exchange} still applies. For that field,
the even fixed-carrier specialization passes this necessary test but
requires its own completed source pairing; it must not be substituted
silently for the running-unit model. An endpoint-labelled realization of the
first line must supply its reversed component and verify reciprocity
for the complete source matrix. The assigned zero-mode unit does not
supply that missing matrix.
In particular, its directed coefficient pole part is
\(\embed(Z_{\bb})[p^2-m_{\rm phys,\bb}^2
+\ii\eta\epsvec]^{-1}\). The complete coefficient is
embedded before multiplication by \(\eta\ii\widehat u_p\).

For completeness, write \(r_p=\sgn(p_\parallel)\) at a
nonzero mode and let
\(\mathscr R^{(2),\rm sheet}_{\kappa,\eta,\bb}\) be the
two independently determined, numerator-stripped character
coefficients of a resolved scalar channel. At finite
\(\varepsilon>0\), the reconstruction within this oriented,
numerator-factorized ansatz is
\begin{equation}
 \vec G^{(2)}_{\eta,\bb}(p;\varepsilon)
  =\sum_{\kappa=\pm1}\mathsf P_\kappa\,
     \kappa\eta\ii r_p\,
     \mathscr R^{(2),\rm sheet}_{\kappa,\eta,\bb}
                       (p^2+\ii\kappa\eta\varepsilon).
 \label{eq:LSZ-two-character-reconstruction}
% Monograph source: eq:general-directed-graph-character-reconstruction; prop:LSZ.
\end{equation}
Each character has its own Laurent data on its declared pole
neighbourhood. Equality of their pole locations and the required
relation between their residues are properties of the physical
completion. If
\begin{equation}
 \mathscr R^{(2),\rm sheet}_{-,\eta,\bb}(\overline z)
   =\overline{
      \mathscr R^{(2),\rm sheet}_{+,\eta,\bb}(z)},
 \label{eq:LSZ-character-reflection}
\end{equation}
then \eqref{eq:LSZ-two-character-reconstruction} reduces to the
single-sheet directed formula above. Otherwise both character
coefficients are retained. The common physical boundary limit is
taken only after the admitted reconstruction and integrations.

For an arbitrary completed two-point source matrix, the genuinely
general formula is instead
\begin{equation}
 \vec C^{AB}_{\bb}(p)
   =\sum_{\kappa=\pm1}\mathsf P_\kappa C^{AB}_{\kappa,\bb}(p),
 \label{eq:LSZ-general-source-matrix}
\end{equation}
where each coefficient includes its own full numerator and obeys
\(C^{AB}_{\kappa,\bb}(p)
=(-1)^{\epsilon_A\epsilon_B}C^{BA}_{\kappa,\bb}(-p)\).
Neither scalar dependence on \(p^2\) nor extraction of a common
running sign is automatic. LSZ reduction additionally requires that
the pole belong to the specified physical source channel and admit
the stated spectral and asymptotic limits. A scalar meromorphic
coefficient alone does not construct those limits.

In the selected scalar source channel, put
\(M^2=m_{\rm phys,\bb}^2\) and write
\(h_{\bb}=H^{\rm sheet}_{\eta,\bb}\). The Laurent form
\eqref{eq:full-pole} has the local reciprocal
\begin{equation}
 \begin{aligned}
  \mathfrak F_{\bb}(z)
    &=\frac{z-M^2}{Z_{\bb}+(z-M^2)h_{\bb}(z)},\\
  \mathfrak F_{\bb}(M^2)&=0,\qquad
  \mathfrak F_{\bb}'(M^2)=Z_{\bb}^{-1}.
 \end{aligned}
 \label{eq:local-Laurent-inverse}
\end{equation}
Because \(Z_{\bb}\ne0\), the denominator stays nonzero in a
sufficiently small neighbourhood. Thus this reciprocal extends
holomorphically through the pole with a simple zero; no global
invertibility is asserted. Define its correction relative to the
reference inverse \(z-m_{0,\bb}^2\) by
\begin{equation}
 \mathfrak F_{\bb}(z)
       :=z-m_{0,\bb}^2-\Sigma^{\rm sheet}_{\Ph,\bb}(z),
 \qquad
 \mathfrak F^{\rm dir}_{\bb}
       :=\embed\circ\mathfrak F_{\bb}\circ\embed^{-1}.
 \label{eq:Dyson-directed-inverse-function}
% Monograph source: eq:Dyson-completed-inverse-coefficient; eq:Dyson-local-pole-expansion.
\end{equation}
Here \(\Sigma^{\rm sheet}_{\Ph,\bb}\) includes any tree-level
kinetic modification as well as loop corrections. It is a loop-only
self-energy when the full tree-level inverse in that source channel
is \(z-m_{0,\bb}^2\).

For a general admitted tree inverse \(K_{{\rm tree},\bb}(z)\),
write instead
\begin{equation}
 \begin{aligned}
  \mathfrak F_{\bb}(z)
     &=K_{{\rm tree},\bb}(z)-\Sigma_{{\rm loop},\bb}(z),\\
  \Sigma^{\rm sheet}_{\Ph,\bb}(z)
     &=z-m_{0,\bb}^2-K_{{\rm tree},\bb}(z)
         +\Sigma_{{\rm loop},\bb}(z).
 \end{aligned}
 \label{eq:tree-versus-loop-inverse}
\end{equation}
Both decompositions refer to the same inverse, not to different
choices of the pole prescription.
The inverse in the full Dyson covariance is
\(\mathfrak F^{\rm dir}_{\bb}
(p^2+\ii\eta\epsvec)^{-1}\), taken in the directed
subalgebra. On the local pole neighbourhood established in
\eqref{eq:local-Laurent-inverse}, the finite matching equations are
\begin{equation}
 m_{\rm phys,\bb}^2-m_{0,\bb}^2
     -\Sigma^{\rm sheet}_{\Ph,\bb}(m_{\rm phys,\bb}^2)=0,
 \qquad
 Z_{\bb}^{-1}
     =1-(\Sigma^{\rm sheet}_{\Ph,\bb})'
                               (m_{\rm phys,\bb}^2).
 \label{eq:finite-pole-matching}
\end{equation}
Equivalently, for the same pole \(M^2\),
\begin{equation}
 \begin{aligned}
  K_{{\rm tree},\bb}(M^2)-\Sigma_{{\rm loop},\bb}(M^2)&=0,\\
  Z_{\bb}^{-1}
    &=K_{{\rm tree},\bb}'(M^2)-\Sigma_{{\rm loop},\bb}'(M^2).
 \end{aligned}
 \label{eq:general-tree-pole-matching}
\end{equation}
For example, a specified scalar tree coefficient
\(f(z)/(z-m_0^2)\), with \(f\) holomorphic and nonzero near
\(m_0^2\), has inverse \((z-m_0^2)/f(z)\) and tree residue
\(f(m_0^2)\). Calling the entire correction in
\eqref{eq:Dyson-directed-inverse-function} ``loop only'' would
miss this contribution. This algebraic example does not identify
that coefficient with an independently completed physical source field.
In a general mixed carrier channel the inverse coefficient and
these pole-matching equations are defined for each character
separately. Their one-sheet form here uses the preceding directed
two-point hypothesis.
All quantities in this last equation are coefficient Laurent
data on the selected sheet. Within an admitted channel they determine the pole part of
\eqref{eq:LSZ-physical-directed-kernel}, not an independent
scalar replacement or a proof of source reciprocity.

Internally damped loops and compatible vertices can change
\(\Sigma_\Ph\). A unit-residue convention requires an admissible
finite field normalization that preserves the physical source
pairing and adjoint. In a positive Hermitian scalar channel the
isolated spectral weight is \(Z_{\bb}>0\), and rescaling by
\(Z_{\bb}^{-1/2}\) supplies it. The weaker algebraic condition
\(0<|Z_{\bb}|<\infty\) alone does not establish that physical
normalization \cite{LSZ1955}; neither does \(\Ph(0)=1\).

\begin{proposition}[Positive residues and field--source duality]
\label{prop:positive-residue-matrix}
Consider a finite multiplet \(\varphi_A\) of admitted physical scalar
fields at an isolated stable one-particle mass. Use the adjoint-paired
two-point matrix, with endpoints \(\varphi_A\) and
\(\varphi_B^\dagger\). After removing the conventional propagator
and state-normalization factors, suppose its residue comes from the
positive physical one-particle resolution. Then
\begin{equation}
 Z^{\rm phys}_{AB}
       =\sum_\alpha v_{A\alpha}\overline{v_{B\alpha}},
 \qquad Z^{\rm phys}\succeq0,
 \label{eq:positive-pole-Gram-matrix}
\end{equation}
with finite or convergent overlap sums. A constant field-basis map
\(\varphi'=T\varphi\), with
\((\varphi')^\dagger=\varphi^\dagger T^\dagger\), gives
\(Z'=T Z^{\rm phys}T^\dagger\). An invertible such map can
normalize the residue to the identity on its range, but cannot make
a singular residue the identity on the whole field space.
\end{proposition}

\begin{proof}
The overlaps \(v_{A\alpha}\) are the vacuum-to-one-particle
matrix elements in the chosen normalization; wave packets may be
used before removing common momentum distributions. For every finite
vector \(c\),
\[
 c^\dagger Z^{\rm phys}c
       =\sum_\alpha\left|\sum_A\overline{c_A}v_{A\alpha}\right|^2
       \geq0.
\]
The field map acts on the two endpoints and gives the asserted
congruence. Write the finite-dimensional spectral decomposition as
\(Z^{\rm phys}=\sum_{\lambda>0}\lambda P_\lambda\), and set
\(P_0=P_{\ker Z^{\rm phys}}\) and
\(P_{\mathcal R}=I-P_0\). Define
\begin{equation}
 \begin{aligned}
  T_{\mathcal R}&=\sum_{\lambda>0}\lambda^{-1/2}P_\lambda,\\
  T_{\rm inv}&=T_{\mathcal R}+P_0,\\
  T_{\rm inv}Z^{\rm phys}T_{\rm inv}^\dagger&=P_{\mathcal R}.
 \end{aligned}
 \label{eq:invertible-residue-normalization}
\end{equation}
The last identity also holds with \(T_{\mathcal R}\), but only
\(T_{\rm inv}\) retains an invertible map on the entire field space.
Invertible congruence preserves rank, so the right side can be the
full identity only when \(Z^{\rm phys}\) is nonsingular.
\end{proof}

For this commuting field block, the dual source change is distinct.
With an invertible \(T\) and the Hermitian pairing it is
\begin{equation}
 \begin{aligned}
  J'&=(T^{-1})^\dagger J,\\
  J'^\dagger\varphi'+\varphi'^\dagger J'
    &=J^\dagger\varphi+\varphi^\dagger J.
 \end{aligned}
 \label{eq:physical-field-source-duality}
\end{equation}
A real bilinear pairing uses the inverse transpose instead. Thus
\(Z\mapsto TZT^\dagger\) describes the field change with its
compatible dual sources, not an unspecified reparametrization of \(J\).
A null residue direction has no overlap with this one-particle mass;
it need not be a zero field or have vanishing correlations elsewhere.
For a mixed field matrix, the scalar inverse formulas above apply only
after an admitted channel reduction; entrywise reciprocals are not a
matrix inverse.

This test presupposes a positive physical channel; it does not construct
one from a virtual kernel or impose positivity on raw carrier characters
with their ordering phases. Its adjoint is the physical \(\dagger\),
not the carrier \(\ddagger\) of
Proposition~\ref{prop:carrier-reality-adjoint}.

The massive internal shell samples the directed numerator boundary
\(\lim_{\varepsilon\downarrow0}
\Ph_{\ee}(-a^2[m^2+\ii\eta\epsvec])\).
Its ordinary coefficient is
\(\Ph(-a^2[m^2+\ii\eta0])\) only after extraction.
It need not equal \(\Ph(0)\), and its two sides need not agree.

Even tree exchange graphs can depend on the kernel through their
internal lines and completed vertices. No blanket equality of
tree amplitudes or interacting residues with the undeformed theory
is asserted.

\subsection{Operational cuts and analytic discontinuities}
\label{sec:cut-disc}

The operational cut is prescribed for the entire oriented virtual
line. Its interpretation as a physical completeness projector requires
the admitted source/endpoint channel of Section~\ref{sec:LSZ}. On that
normalized cut-state interface,
\begin{equation}
\begin{aligned}
 &\Cut_{{\rm phys},\bb}\!\left[
  \eta\ii\widehat u_k\,
  \frac{\Ph_{\ee}(-a^2[k^2+\ii\eta\epsvec])}
       {D(k)+\ii\eta\epsvec}\right]\\
 &\qquad:=2\pi\delta(D(k))\,\theta(\eta\ee\!\cdot k),\\
 &D(k)=k^2-m^2.
\end{aligned}
 \label{eq:physical-cut-map}
% Monograph source: eq:stress-cutkosky; eq:operational-directed-line-cut-replacement.
\end{equation}
This operation removes the virtual line, its ordering
coefficient, its mode unit and its kernel together, replacing them
by the normalized locked completeness projector. It is not a
distributional identity for the uncut propagator. On a canonical
aligned ray the extracted uncut coefficient is
\(\sigma_\epsilon\ii R^\Ph_\eta\), as in
\eqref{eq:scalar-line-coefficient}; it is not obtained by
discarding the sector factor in that equation.

For comparison with analytic boundary values, define
\begin{equation}
 R^{\rm coeff}_{\pm,\bb}(k)
 :=\lim_{\varepsilon\downarrow0}
       \embed^{-1}\!\left[
       f^{\rm dir}_{\pm,\bb}(k;\varepsilon)
        (D(k)\mathbin{\pm}\ii\epsvec)^{-1}\right].
 \label{eq:cut-directed-coefficient-extraction}
\end{equation}
Define the corresponding directed boundary coefficients by
\(\mathcal R^{\rm dir}_{\pm,\bb}
:=\embed(R^{\rm coeff}_{\pm,\bb})\).
Only after this extraction write
\(R^{\rm coeff}_{\pm,\bb}=f_\pm/(D\pm\ii0)\).
Their scalar jump is
\(\Disc^{\rm coeff}R=R^{\rm coeff}_+-R^{\rm coeff}_-\).
The corresponding jump of the directed coefficient is
\begin{equation}
 \Disc^{\rm dir}R
 :=\mathcal R^{\rm dir}_{+,\bb}
       -\mathcal R^{\rm dir}_{-,\bb}
 =\embed(\Disc^{\rm coeff}R).
 \label{eq:directed-discontinuity-reconstruction}
\end{equation}
It is distinct from reversing a complete graph's evolution
branch, which also reverses its ordering and Dyson data.
The auxiliary coefficient identity
\((D+\ii0)^{-1}=\PV(1/D)-\ii\pi\delta(D)\)
does not supply the on-shell orientation selector.

Let the continued numerator have boundary values \(f_+\) and
\(f_-\), smooth near a regular shell \(D=0\) with \(\dd D\ne0\).
Assume also that the regulated numerators converge as multipliers of
the resolvent boundary distributions; local \(C^1\) convergence on
compact shell charts is sufficient. The identities below are local
there (with test functions of compact support); a global identity
requires control of other singularities and infinity. Pointwise
numerator limits alone do not justify multiplying two distributional
limits. The ordinary
coefficient jump is
\begin{align}
 \frac{f_+}{D+\ii0}-\frac{f_-}{D-\ii0}
  &=(f_+-f_-)\PV\frac1D
     -\ii\pi(f_++f_-)\big|_{D=0}\delta(D).
 \label{eq:numerator-disc}
% Monograph source: crit:cut-vs-disc; eq:cut-disc-general-boundary-numerator.
\end{align}
To verify the stated sufficient multiplier condition, use a regular
shell chart \((y,\xi)\) with \(y=D\), and absorb its smooth
Jacobian into a compactly supported test function. If the numerator
error tends to zero in \(C^1\), its tested product can be written
\(g_\varepsilon(y,\xi)=g_\varepsilon(0,\xi)\chi(y)
+y r_\varepsilon(y,\xi)\), where \(\chi=1\) near zero and
both coefficients tend uniformly to zero on a fixed compact set.
The pairing of \(\chi\) with \((y\pm\ii\varepsilon)^{-1}\)
is bounded, and \(|y/(y\pm\ii\varepsilon)|\leq1\).
Both error terms therefore vanish. Applying the ordinary resolvent
boundary formula to the remaining fixed smooth numerators proves
\eqref{eq:numerator-disc} on the chart.
If \(f_+=f_-=f\), this reduces to
\begin{equation}
 \Disc^{\rm coeff}\frac{f}{D+\ii0}
       =-2\pi\ii\,f|_{D=0}\delta(D).
 \label{eq:weighted-pole-disc}
\end{equation}
For the directed coefficient, the same formula is transported in
its entirety:
\begin{equation}
 \Disc^{\rm dir}R
 =\embed\!\left[
 (f_+-f_-)\PV\frac1D
   -\ii\pi(f_++f_-)|_{D=0}\delta(D)\right].
 \label{eq:directed-numerator-discontinuity}
\end{equation}
If the common numerator is real, the continuous-numerator
specialization is
\(-2\pi\je f|_{D=0}\delta(D)\).
For complex numerators one must retain the complete embedding.
Neither expression discards the mode numerator or the remaining
graph factors; those enter the full reconstruction
\eqref{eq:complete-graph-character-reconstruction}.

Thus even an entire numerator can change the pole weight.
A numerator branch jump adds the principal-value contribution in
\eqref{eq:numerator-disc}. Removing the damper by a prescribed
\(\Cut_{\rm phys}\) operation does not remove these terms from
\(\Disc_{\rm an}\).

For the stable-half kernel, form the full directed argument and
transported functions before taking a boundary value:
\begin{equation}
\begin{aligned}
 W_{\eta,\bb}(k)&:=-(k^2+\ii\eta\epsvec),\\
 f^{\rm dir}_{a,\eta,\bb}(k;\varepsilon)
   &:=\ExpDir\!\left[-a\SqrtDir(W_{\eta,\bb}(k))\right].
\end{aligned}
 \label{eq:stable-full-directed-kernel}
% Monograph source: eq:stable-half-full-directed-argument; eq:stable-half-directed-boundary-values.
\end{equation}
Here \(\SqrtDir:=\embed\circ\sqrt{\,\cdot\,}
\circ\embed^{-1}\) uses the principal scalar square-root
sheet, and \(\ExpDir:=\embed\circ\exp\circ\embed^{-1}\).
For \(k^2>0\), their directed boundary is
\begin{equation}
\begin{aligned}
 &\lim_{\varepsilon\downarrow0}
 f^{\rm dir}_{a,\eta,\bb}(k;\varepsilon)\\
 &\qquad=\cos(a\sqrt{k^2})+\eta\je\sin(a\sqrt{k^2}),\\
 &\lim_{\varepsilon\downarrow0}
 \embed^{-1}\!\left(f^{\rm dir}_{a,\eta,\bb}(k;\varepsilon)\right)
       =e^{\ii\eta a\sqrt{k^2}}.
\end{aligned}
 \label{eq:stable-directed-boundary}
\end{equation}
The square root of the positive real number \(k^2\) in this
boundary formula is an ordinary nonnegative real root.
For \(m>0\), the shell is regular and lies away from the branch
point. The ordinary coefficient jump and its directed reconstruction
are
\begin{align}
 \Disc^{\rm coeff}R_a
  &=2\ii\theta(k^2)\sin(a\sqrt{k^2})\PV\frac1D
       -2\pi\ii\cos(am)\delta(D),\nonumber\\
 \Disc^{\rm dir}R_a
  &=2\je\theta(k^2)\sin(a\sqrt{k^2})\PV\frac1D
       -2\pi\je\cos(am)\delta(D).
 \label{eq:stable-disc}
\end{align}
The full line remains
\(\widetilde G^\Ph_{F,\bb}
=\eta\ii\widehat u_k\mathcal R^{\rm dir}_{\eta,\bb}\).
Thus the displayed jump concerns its directed denominator/kernel
coefficient, not a graph obtained by omitting that numerator.
At \(m=0\), put
\(z_{\eta,\bb}=k^2+\ii\eta\epsvec\). On the selected
directed sheet,
\begin{equation}
 \frac{\ExpDir[-a\SqrtDir(-z_{\eta,\bb})]}
              {z_{\eta,\bb}}
 =\frac1{z_{\eta,\bb}}
      +\frac{a}{\SqrtDir(-z_{\eta,\bb})}+O(1).
 \label{eq:massless-branch}
\end{equation}
The branch point therefore prevents an isolated meromorphic pole
from following merely from \(\Ph(0)=1\). The cone tip is not
covered by the regular massive-shell discontinuity formula.

\begin{assumption}[Cut--discontinuity compatibility]
\label{ass:cut-matching}
For the completed physical amplitudes, the prescribed on-shell
cut operation agrees with the analytic discontinuity, including
the actual pole residues, completed subamplitudes, kernel branch
terms, and contour contributions. Unphysical states cancel under
the admitted BRST construction, and the remaining state sum has
positive physical norm.
\end{assumption}

This assumption can be established by an action-level unitarity
argument with its requisite analytic conditions, or by an explicit
matching calculation. It cannot be substituted for by
\eqref{eq:kernel-inert} or a positive Laplace measure.
The standard Cutkosky theorem identifies discontinuities under its
analytic and Feynman-integral hypotheses \cite{Cutkosky1960}.
Section~\ref{sec:gaussian-bubble} provides an exact example in which
the naive unweighted cut fails this test.

% Monograph source: prop:LSZ; eq:LSZ-plane-wise-operational-cut-sum; crit:cut-vs-disc.
For the optical relation, assume in addition that the physical
scattering operator is unitary on the chosen positive-norm asymptotic
state space and that the intermediate-state resolution is complete,
including every open physical channel. Writing
\(S^{\rm phys}=\one+\ii T^{\rm phys}\), one obtains
\begin{equation}
 (S^{\rm phys})^\dagger S^{\rm phys}=\one
 \quad\Longrightarrow\quad
 -\ii\bigl[T^{\rm phys}-(T^{\rm phys})^\dagger\bigr]
       =(T^{\rm phys})^\dagger T^{\rm phys}.
 \label{eq:physical-unitarity-identity}
\end{equation}
Here \(\dagger\) is the adjoint on that physical Hilbert space,
not the coefficient reality operation \(\ddagger\).
Proposition~\ref{prop:carrier-reality-adjoint} rules out identifying
these on the full carrier algebra. Neither a character norm nor a
reality relation can replace the positive physical state resolution.
Take a forward matrix element and insert the complete physical
state resolution between the two operators on the right. This
argument first uses normalizable wave packets. Passing to
momentum-space matrix-element coefficients and factoring out the
common momentum-conservation distribution then gives
\begin{equation}
 2\Im\mathcal M^{\rm phys,coeff}_{\bb}(i\to i)
  =\sum_n\int\dd\Pi_{n,\bb}\,
       |\mathcal M^{\rm phys,coeff}_{\bb}(i\to n)|^2.
 \label{eq:optical}
\end{equation}
Here \(\dd\Pi_{n,\bb}\) includes the physical state normalization,
momentum conservation and identical-particle counting factors.
In perturbation theory both sides, including the squared amplitudes,
are expanded to the same order in all couplings. Identifying the state sum with an analytic
discontinuity
additionally requires Assumption~\ref{ass:cut-matching} and the
appropriate analytic boundary relations.

Alternatively, the same forward optical identity follows from a
normalized analytic matching calculation if the physical amplitude
has the Hermitian-analytic boundary relation
\[
 \mathcal M_-^{\rm phys,coeff}(i\to i)
   =\overline{\mathcal M_+^{\rm phys,coeff}(i\to i)}
\]
and its matched discontinuity factorizes as
\[
\begin{aligned}
 &\mathcal M_+^{\rm phys,coeff}(i\to i)
  -\mathcal M_-^{\rm phys,coeff}(i\to i)\\
 &\qquad=\ii\sum_n\int\dd\Pi_{n,\bb}\,
       |\mathcal M_+^{\rm phys,coeff}(i\to n)|^2.
\end{aligned}
\]
The subscripts \(+\) and \(-\) here denote the two boundary
approaches to the selected external-invariant cut after physical
representation and the admitted coefficient extraction. They do
not denote the carrier characters \(\chi_\pm\). A relation
between carrier components alone does not establish this
Hermitian analyticity. Without it, an imaginary entire addition
can change the imaginary part of a one-loop amplitude while
leaving its discontinuity and its tree-level cut unchanged.

A carrier-valued graph, including its cut graph, is reconstructed
from both complete character calculations by
\eqref{eq:complete-graph-character-reconstruction}; its carrier-odd
part is not defined to be an absorptive probability.

Here the amplitudes are those of the completed theory, with
physical residues and state normalization. Plane diagonality
restricts the intermediate states to the same support in the
absence of support-changing interactions. It does not supply
unitarity by itself, and it does not replace completed on-shell
subamplitudes by undeformed ones.

\subsection{Conditional crossing and locality}

Local holomorphy in Theorem~\ref{thm:holomorphy} is a statement
about a graph representation. Physical crossing relates boundary
values in different scattering channels. To invoke a standard
crossing result one needs the applicable physical hypotheses:
locality of reconstructed observables, spectral support,
appropriate stable asymptotic particles, LSZ limits, and a
connected analytic domain joining the channels with controlled
boundary growth. Kernel and vertex singularities must be
included in that domain analysis.

For two-particle scattering with the required stable particles and a
strictly positive physical mass gap, Bros, Epstein, and Glaser construct
a connected on-shell continuation domain \cite{BrosEpsteinGlaser1965}.
Their gap concerns all nonvacuum physical states, not just the external
particles \cite[p.~243]{BrosEpsteinGlaser1965}. The same paper explicitly
separates the additional identification of on-shell boundary values
with the physical amplitudes and postpones its proof
\cite[p.~241]{BrosEpsteinGlaser1965}. Here that identification is part
of the separately assumed scattering and boundary-limit hypotheses;
it is not inferred from existence of the analytic domain alone.
A massless physical channel lies outside this invocation. Neither the
cited domain result nor formal permutation symmetry of a Euclidean
integrand proves crossing for an arbitrary damped completion.

The implication justified here is therefore
\[
 \begin{gathered}
 \text{completed physical reconstruction and LSZ hypotheses}\\
 \text{an admissible channel domain}\\
 \text{and identification of its physical boundary values}\\
 \Longrightarrow\quad
 \text{the corresponding physical crossing relation}.
 \end{gathered}
\]
This statement is separate from the fixed-graph UV theorem and
does not assert a massless scattering construction.

\section{Kernel families, coordinate forms, and worked examples}
\label{sec:examples}

\subsection{Admissible Euclidean kernel families}
\label{sec:kernel-families}

The families below use the dimensionless variable \(z=a^2q_E^2\).
Their suitability for Euclidean UV estimates does not automatically
supply the independent continuation or physical hypotheses.

\paragraph{A single heat depth}
For \(\tau_*>0\),
\begin{equation}
 \Ph(z)=e^{-\tau_*z},\qquad\rho=\delta_{\tau_*}.
 \label{eq:heat-family}
\end{equation}
This satisfies the gap route, and all inverse moments are finite.
It is entire, but its timelike continuation has the growth in
\eqref{eq:heat-continuation}.

\paragraph{Power kernels}
For dimensionless \(b>0\) and \(\nu>0\),
\begin{equation}
 \Ph(z)=(1+z/b)^{-\nu},\qquad
 \rho(\dd\tau)=\frac{b^\nu}{\Gamma(\nu)}
                 \tau^{\nu-1}e^{-b\tau}\dd\tau .
 \label{eq:power-family}
\end{equation}
For \(0<\beta<\nu\),
\begin{equation}
 M_\beta=b^\beta\frac{\Gamma(\nu-\beta)}{\Gamma(\nu)},
 \label{eq:gamma-inverse-moments}
\end{equation}
and the inverse moment diverges for \(\beta\geq\nu\).
Thus \(\nu>1+D/2\) suffices for the global theorem.
Weaker graph-specific conditions may suffice, as the bubble
example shows. The singularity at \(z=-b\) must be retained in a
Lorentzian analysis.

\paragraph{Finite positive exponential mixtures}
For \(w_j>0\), \(\sum_jw_j=1\), and \(\tau_j>0\),
\begin{equation}
 \Ph(z)=\sum_{j=1}^Nw_je^{-\tau_jz},\qquad
 \rho=\sum_{j=1}^Nw_j\delta_{\tau_j}.
 \label{eq:finite-mixture}
\end{equation}
The gap is \(\min_j\tau_j\). These are exponential mixtures,
not rational functions or Pad\'e approximants.

\paragraph{Stretched exponentials}
For \(0<r<1\),
\begin{equation}
 \Ph(z)=e^{-z^r}
 \label{eq:stable-family}
\end{equation}
has a positive stable representing measure. It has no gap, but
every inverse moment is finite:
\begin{align}
 M_\beta
 &=\frac1{\Gamma(\beta)}
     \int_0^\infty z^{\beta-1}\Ph(z)\dd z
   =\frac{\Gamma(\beta/r)}{r\Gamma(\beta)},
 \qquad \beta>0.
 \label{eq:stable-inverse-moments}
\end{align}
The first equality follows from
\(\tau^{-\beta}
 =\Gamma(\beta)^{-1}\int_0^\infty z^{\beta-1}e^{-\tau z}\dd z\)
and Tonelli. The restriction to orders below \(r\) concerns
positive moments of the stable variable, not inverse moments.
At \(r=1\), the measure reduces to a point mass.

For the monograph's canonical \(r=\tfrac12\) specialization,
\begin{equation}
 \rho(\dd\tau)=\frac1{2\sqrt\pi}\,
     \tau^{-3/2}e^{-1/(4\tau)}\dd\tau,\qquad
 M_\beta=\frac{4^\beta\Gamma(\beta+\tfrac12)}{\sqrt\pi}.
 \label{eq:stable-half-density}
\end{equation}
The line multiplier is \(e^{-a|q_E|}\), with
superpolynomial large-momentum decay. The large-\(\tau\) tail
of \eqref{eq:stable-half-density} excludes any positive
exponential moment. Its UV usefulness and its more delicate
continuation therefore coexist.
The O--QED application uses the separately specified normalized massive
variant in Section~\ref{sec:OQED-moments}; its positive auxiliary scale
changes the large-parameter tail and local analytic hypotheses without
creating a positive Laplace support gap.

An ungapped measure does not acquire a gap by mixing it with a
positive atom away from zero: its original near-zero support
remains. Translating the entire measure by \(\tau_*>0\), or
truncating and renormalizing it below that point, does create a
gap. Translation corresponds to
\(\Ph(z)\mapsto e^{-\tau_*z}\Ph(z)\).

The distinction between inverse moments and exponential moments is
summarized in Table~\ref{tab:kernel-comparison}. Write
\[
 \mathcal E_\rho=\left\{\epsilon>0:
      \int e^{\epsilon\tau}\rho(\dd\tau)<\infty\right\}.
\]
Inverse moments enter the UV estimate; exponential moments give a
sufficient local continuation condition. The continuation of an
amplitude also uses the graph-level hypotheses of
Section~\ref{sec:analytic}.

\begin{table*}[t]
\centering
\small
\setlength{\tabcolsep}{3pt}
\renewcommand{\arraystretch}{1.2}
\caption{Representative kernel properties. Here
\(\tau_*,b,\nu>0\), \(\beta>0\), \(0<r<1\), and \(\mu>0\)
in the normalized massive row.}
\label{tab:kernel-comparison}
\begin{tabularx}{\textwidth}{@{}>{\raggedright\arraybackslash}p{0.20\textwidth}
 >{\raggedright\arraybackslash}p{0.28\textwidth}
 >{\raggedright\arraybackslash}p{0.14\textwidth}
 >{\raggedright\arraybackslash}X@{}}
\toprule
Kernel \(\Ph(z)\) & Inverse moment \(M_\beta\) &
\(\mathcal E_\rho\) & Complex continuation\\
\midrule
\(e^{-\tau_*z}\) & \(\tau_*^{-\beta}\), all \(\beta>0\) &
\((0,\infty)\) & Entire\\
\((1+z/b)^{-\nu}\) &
\(\dfrac{b^\beta\Gamma(\nu-\beta)}{\Gamma(\nu)}\),
\(0<\beta<\nu\) & \((0,b)\) &
Pole or branch point at \(z=-b\)\\
\(e^{-z^r}\) &
\(\dfrac{\Gamma(\beta/r)}{r\Gamma(\beta)}\), all \(\beta>0\) &
\(\varnothing\) & Branch point at \(z=0\)\\
\(e^{\mu-\sqrt{z+\mu^2}}\) &
\eqref{eq:OQED-inverse-moments}, all \(\beta>0\) &
\((0,\mu^2]\) & Branch point at \(z=-\mu^2\); no support gap\\
\((2e^z-1)^{-1}\) &
\(\sum_{n\geq1}2^{-n}n^{-\beta}\), all \(\beta>0\) &
\((0,\log2)\) & Poles at \(-\log2+2\pi\ii n\),
\(n\in\mathbb Z\)\\
\bottomrule
\end{tabularx}
\end{table*}

\subsection{Coordinate-space covariance and coincidence limit}

For a finite positive line measure and \(m\geq0\), the bounded
multiplier defines the following Fourier transform initially as a
tempered distribution. Under the coincidence hypotheses stated below,
it is an ordinary continuous covariance function:
\begin{align}
 C_{\Ph,E}(x-y)
 &=\int\rho(\dd\tau)
     \int\frac{\dd^4q_E}{(2\pi)^4}\,
       \frac{e^{\ii q_E\cdot(x-y)}
                 e^{-a^2\tau q_E^2}}{q_E^2+m^2}
 \label{eq:coordinate-covariance}\\
 &=\int\rho(\dd\tau)\int_0^\infty\dd\alpha\,
     \frac{e^{-\alpha m^2}}
          {[4\pi(\alpha+a^2\tau)]^2}\nonumber\\
 &\qquad\times\exp\!\left[-\frac{|x-y|^2}
                         {4(\alpha+a^2\tau)}\right].
 \label{eq:coordinate-heat-form}
\end{align}
The second equality uses the elementary Schwinger representation
\cite{Schwinger1951}. It is a heat-kernel average of the standard
Euclidean covariance. At coincidence, for \(\tau>0\),
\[
 \int_0^\infty
 \frac{e^{-\alpha m^2}}{(\alpha+a^2\tau)^2}\dd\alpha
       \leq\frac1{a^2\tau}.
\]
Therefore
\begin{equation}
 C_{\Ph,E}(0)\leq
       \frac{M_1[\rho]}{16\pi^2a^2},
 \label{eq:coincidence-bound}
\end{equation}
provided there is no atom at zero and \(M_1<\infty\).
For \(m=0\), equality holds. In the stable-half case,
\(M_1=2\), giving
\(C_{\Ph,E}(0)=1/(8\pi^2a^2)\) for the massless scalar
comparison covariance. This finite coincidence limit is a
Euclidean statement, not a reflection-positivity proof.

\subsection{The literal two-line scalar bubble}
\label{sec:literal-bubble}

Fix one support and \(m>0\). The following calculation is a
Euclidean coefficient benchmark with the explicitly fixed virtual
carrier convention. Its displayed bubble function omits the
propagator numerators, vertices, and overall Dyson coefficient.
These data are restored below. For a running-unit Lorentzian
bubble with momenta \(k\) and \(P-k\), the two line factors
\(\sgn(k_\parallel)\) and
\(\sgn(P_\parallel-k_\parallel)\) must instead remain in the
integral. The scalar radial computation does not evaluate that
different integrand.

With one multiplier on each uncut line, the Euclidean scalar bubble is
\begin{equation}
 B_{\Ph,E}(P_E)=
 \int\frac{\dd^4k_E}{(2\pi)^4}\,
 \frac{\Ph(a^2k_E^2)\Ph(a^2(k_E+P_E)^2)}
      {(k_E^2+m^2)((k_E+P_E)^2+m^2)} .
 \label{eq:two-line-bubble}
\end{equation}
Introduce independent \(\tau_1,\tau_2\), independent
\(\alpha_1,\alpha_2\), and \(w_i=\alpha_i+a^2\tau_i\).
Completing the square gives
\[
\begin{aligned}
 &w_1k_E^2+w_2(k_E+P_E)^2\\
 &\quad=(w_1+w_2)\left(k_E+\frac{w_2}{w_1+w_2}P_E\right)^2
       +\frac{w_1w_2}{w_1+w_2}P_E^2.
\end{aligned}
\]
The exact parameter formula is
\begin{align}
 B_{\Ph,E}(P_E)
 =\frac1{16\pi^2}
 &\int\rho(\dd\tau_1)\rho(\dd\tau_2)
       \int_0^\infty\dd\alpha_1\dd\alpha_2 \nonumber\\
 &\times\frac{
 e^{-m^2(\alpha_1+\alpha_2)}
 e^{-\frac{w_1w_2}{w_1+w_2}P_E^2}
 }{(w_1+w_2)^2}.
 \label{eq:two-line-parameters}
\end{align}
This is the actual linewise rule. A single factor attached after
combining and shifting denominators is a different effective
model unless a separate identity establishes equivalence.

In the normal ascending scalar reference convention
\(\sigma_s=\sigma_\epsilon=+1\), use a real \(\phi^4\)
coupling, \(S=1+\ii T\), and tree coefficient
\(\mathcal M_{\rm tree}=-\lambda\). The one-loop
coefficient in this benchmark is
\begin{equation}
\begin{aligned}
 &\mathcal M^{\rm coeff}_{+,\bb}(\sMand,\tMand,\uMand)\\
 &\quad=\mathcal P^{(4)}_{\bb\mid{\rm ext}}
 \left[-\lambda+\frac{\lambda^2}{2}
    \sum_{c=1}^{3}B^{\rm coeff}_{\Ph,+}(v_c)\right]+O(\lambda^3).
\end{aligned}
 \label{eq:scalar-amplitude}
\end{equation}
Here \(v_1=\sMand\), \(v_2=\tMand\), and \(v_3=\uMand\);
the three channel contributions are counted with their multiplicity.
The function \(B^{\rm coeff}_{\Ph,+}\) is the extracted boundary
of the declared continuation of \eqref{eq:two-line-bubble}, as in
\eqref{eq:analytic-function-directed-interface}. This extraction is
performed on the admitted path for each channel; it is not an
independent complex displacement of three constrained invariants. This is a specified scalar benchmark,
not the full gauge completion. The factor \(1/2\) belongs to
the graph symmetry factor and must not be inserted into the
definition of \(B_\Ph\).

Let \(\mathcal M^{\rm coeff}_{-,\bb}\) denote the complete
opposite-branch reference with its own boundary values. For this
real-coupling scalar benchmark with the reflected endpoint and
vertex conventions,
\(\mathcal M^{\rm coeff}_{-,\bb}
=\overline{\mathcal M^{\rm coeff}_{+,\bb}}\). Hence
\begin{align}
 \vec{\mathcal M}_{\bb}
 &=\mathsf P_+\mathcal M^{\rm coeff}_{+,\bb}
      +\mathsf P_-\mathcal M^{\rm coeff}_{-,\bb}
   =\embed(\mathcal M^{\rm coeff}_{+,\bb}),\nonumber\\
 \vec{\mathcal A}_{\bb}
 &=\mathsf P_+\ii\mathcal M^{\rm coeff}_{+,\bb}
      -\mathsf P_-\ii\mathcal M^{\rm coeff}_{-,\bb}
   =\je\vec{\mathcal M}_{\bb}.
 \label{eq:scalar-benchmark-Dyson-reconstruction}
\end{align}
Here \(\vec{\mathcal A}\) is the Dyson transition coefficient;
its opposite branch includes the opposite Dyson sign. Thus a real
tree value \(-\lambda\) gives
\(\vec{\mathcal M}_{\rm tree}=-\lambda\) and
\(\vec{\mathcal A}_{\rm tree}=-\lambda\je\).
If \(\mathcal M_+=A+\ii B\), then
\(\vec{\mathcal M}=A+\je B\) but
\(\vec{\mathcal A}=-B+\je A\). The two quantities have
different parity decompositions. Outside this explicitly reflected
benchmark, use the two complete representatives in
\eqref{eq:complete-graph-character-reconstruction} directly.

\subsection{Exact Gaussian bubble and its discontinuity}
\label{sec:gaussian-bubble}

Choose \(\Ph(z)=e^{-\tau_*z}\) and put \(t=a^2\tau_*>0\).
Then \(w_i=\alpha_i+t\).
Set \(W=w_1+w_2\) and \(x=w_2/W\). The integration region is
\[
\begin{gathered}
 0<x<1,\qquad W\geq W_{\min}(x):=\frac{t}{\min(x,1-x)},\\
 \dd w_1\dd w_2=W\,\dd W\dd x .
\end{gathered}
\]
Equation~\eqref{eq:two-line-parameters} becomes
\begin{equation}
 B_{t,E}(P_E)
   =\frac{e^{2tm^2}}{16\pi^2}
        \int_0^1\dd x\int_{W_{\min}(x)}^\infty
        \frac{\dd W}{W}
             e^{-W[m^2+x(1-x)P_E^2]} .
 \label{eq:gaussian-W}
\end{equation}
Initially for \(\Re z>0\), set
\(E_1(z)=\int_1^\infty e^{-zu}\dd u/u\), and choose its
continuation from the Euclidean region, with the principal branch
and its standard logarithmic jump \cite{DLMF6}. Define
\begin{equation}
\begin{aligned}
 \mathsf S_{\eta,\bb}&:=\sMand+\ii\eta\epsvec,\\
 \Delta^{\rm dir}_{\eta,\bb}(x)
       &:=m^2-x(1-x)\mathsf S_{\eta,\bb},\\
 \EoneDir&:=\embed\circ E_1\circ\embed^{-1}.
\end{aligned}
 \label{eq:Gaussian-directed-channel}
\end{equation}
The directed boundary bubble and its extracted coefficient are
\begin{align}
 \mathcal B^{\rm dir}_{t,\eta,\bb}(\sMand)
 &=\lim_{\varepsilon\downarrow0}
   \frac{e^{2tm^2}}{16\pi^2}
      \int_0^1\EoneDir\!\left(
          W_{\min}(x)\Delta^{\rm dir}_{\eta,\bb}(x)
                                  \right)\dd x,
          \nonumber\\
 B^{\rm coeff}_{t,\eta}(\sMand)
 &:=\embed^{-1}\!\left[
          \mathcal B^{\rm dir}_{t,\eta,\bb}(\sMand)\right]
   =\lim_{\varepsilon\downarrow0}
                    B_t(\sMand+\ii\eta\varepsilon),
 \label{eq:Gaussian-directed-boundary-and-coefficient}
\end{align}
where the ordinary analytic sheet function is
\begin{equation}
\begin{gathered}
 B_t(z)=\frac{e^{2tm^2}}{16\pi^2}
    \int_0^1 E_1\!\left(
       W_{\min}(x)[m^2-x(1-x)z]\right)\dd x,\\
 z\in\CC\setminus[4m^2,\infty).
\end{gathered}
 \label{eq:gaussian-E1}
\end{equation}
This representation is locally uniformly convergent on the
stated cut plane. Indeed,
\(W_{\min}(x)x(1-x)=t\max(x,1-x)\leq t\).
For \(z\) in a compact set and \(x\) near either endpoint,
the real part of the \(E_1\) argument therefore tends uniformly
to \(+\infty\), and its positive-half-plane integral gives
exponential endpoint control. On the remaining compact
\(x\)-interval, the argument avoids the \(E_1\) cut whenever
\(z\notin[4m^2,\infty)\). The physical boundary has only
integrable logarithms at the roots of
\(m^2-x(1-x)\sMand\), including the double root at threshold.
Thus the directed boundary values follow from this evaluated
Euclidean function. This argument does not assume a contour
rotation of the growing unparametrized Lorentzian graph.

For \(\sMand>4m^2\), let
\(\beta(\sMand)=\sqrt{1-4m^2/\sMand}\) denote the positive
real phase-space root. The directed branch identity is
\[
\begin{aligned}
 \lim_{\varepsilon\downarrow0}\Bigl[&\EoneDir\!\left(W[-d-\ii\eta c\epsvec]\right)\\
       &-\EoneDir\!\left(W[-d+\ii\eta c\epsvec]\right)\Bigr]
       =2\pi\eta\je,
\end{aligned}
\]
with \(W,d,c>0\).
Here \(W\) has inverse-momentum-squared units, \(d\) has
momentum-squared units, and \(c\) is dimensionless; the
\(E_1\) arguments are dimensionless. On the open cut interval
use \(W=W_{\min}(x)\), \(c=x(1-x)\), and
\(d=c\sMand-m^2>0\). Consequently
\begin{align}
 \Disc^{\rm dir}_\eta\mathcal B_t
 &:=\mathcal B^{\rm dir}_{t,\eta,\bb}
       -\mathcal B^{\rm dir}_{t,-\eta,\bb}
   =\eta\je\,\frac{e^{2tm^2}}{8\pi}
                \beta(\sMand)\theta(\sMand-4m^2),
        \nonumber\\
 \Disc^{\rm coeff}_\eta B_t
 &:=B^{\rm coeff}_{t,\eta}-B^{\rm coeff}_{t,-\eta}
   =\eta\ii\,\frac{e^{2tm^2}}{8\pi}
                \beta(\sMand)\theta(\sMand-4m^2).
 \label{eq:gaussian-disc}
\end{align}
The ordinary absorptive comparison coefficient is recovered as
\begin{equation}
 \frac{\Disc^{\rm dir}_\eta\mathcal B_t}{2\eta\je}
   =\frac{\Disc^{\rm coeff}_\eta B_t}{2\eta\ii}
   =\frac{e^{2tm^2}}{16\pi}
              \beta(\sMand)\theta(\sMand-4m^2).
 \label{eq:Gaussian-directed-absorptive-coefficient}
\end{equation}
This jump concerns the numerator-stripped bubble. Its vertices,
symmetry factor and Dyson coefficient are restored through
\eqref{eq:scalar-benchmark-Dyson-reconstruction}.

With the standard on-shell one-particle measures, for
future-directed \(P\), \(\eta=+1\), and
\(P^2=\sMand>4m^2\), the two-body
phase-space integral is
\begin{equation}
\begin{aligned}
 \int\dd\Phi_2(P)
 &=\int\prod_{i=1}^2\frac{\dd^3q_i}{(2\pi)^3\,2\omega_i}\,
        (2\pi)^4\delta^{(4)}(P-q_1-q_2)\\
 &=\frac{\beta(\sMand)}{8\pi},\qquad
 \omega_i=\sqrt{\mathbf q_i^2+m^2}.
\end{aligned}
 \label{eq:two-body-phase-space}
\end{equation}
The \(\omega_i>0\) here are intrinsic on-shell momentum
coefficients in the scalar comparison measure, not replacements
for carrier-valued geo-momentum charges.
For the opposite orientation the invariant measure is
\begin{equation}
 \dd\Pi_{1,\eta}(q)
  :=\frac{\dd^4q}{(2\pi)^4}\,2\pi\delta(q^2-m^2)
       \theta(\eta\ee\!\cdot q),
 \label{eq:oriented-one-particle-measure}
\end{equation}
and the two-body integral with \(P=q_1+q_2\) is
\(\beta(\sMand)\theta(\eta\ee\!\cdot P)
\theta(\sMand-4m^2)/(8\pi)\).
An outgoing cut uses \(q_1=k\) and \(q_2=P-k\), so both
selectors have the same declared physical orientation. Evolution-branch
reversal changes the selectors together with the boundary prescription.
The scalar comparisons that follow use \(\eta=+1\).

Consequently, a prescribed undressed unit-vertex bubble cut would
give \(\ii\beta/(8\pi)\), whereas
\eqref{eq:gaussian-disc} gives
\(e^{2tm^2}\ii\beta/(8\pi)\).
They differ even though \(\Ph(0)=1\), the kernel is entire, and
there is a strict Laplace-depth gap. The missing factor must be
accounted for by the actual physical matching and completion if
this operational cut is to represent an analytic discontinuity.

The familiar positive-boundary scalar benchmark normalizations are
\[
 \Disc B_{\rm std}=\frac{\ii\beta}{8\pi},\qquad
 \Im B_{\rm std}=\frac{\beta}{16\pi},\qquad
 \Im\mathcal M_{\rm bubble,std}
      =\frac{\lambda^2\beta}{32\pi}.
\]
The last expression includes the \(1/2\) symmetry factor.
For the Gaussian benchmark it acquires the factor \(e^{2tm^2}\)
in the same convention. The tree term \(-\lambda\) is real;
the factor \(\ii\) in the Dyson coefficient \(\ii\mathcal M\)
is not a tree-level absorptive part.

\subsubsection*{Finite shell normalization and the remaining amplitude}

The next proposition concerns the ordinary analytic sheet functions
\(B_t(z)\), \(b_t(z)\), \(b_{\rm std}(z)\), and \(R_t(z)\).
Their physical directed boundary coefficients are obtained by the
explicit transport below; these coefficient identities do not
replace the full graph reconstruction.

The Gaussian example also permits an explicit normalization test.
Multiplying each scalar line by \(e^{-tm^2}\) makes its massive
pole coefficient unity. For the bubble this gives the comparison
normalization
\begin{equation}
 \widehat B_t(\sMand)=e^{-2tm^2}B_t(\sMand),\qquad
 b_t(\sMand)=\widehat B_t(\sMand)-\widehat B_t(0).
 \label{eq:normalized-bubble}
\end{equation}
Equation~\eqref{eq:scalar-amplitude} continues to use the original
\(B_t\). Define the finite ordinary comparison function by
\begin{equation}
 b_{\rm std}(\sMand)
 =\frac1{16\pi^2}\int_0^1\dd x\,
   \log\!\frac{m^2}{m^2-x(1-x)\sMand},
 \label{eq:subtracted-standard-bubble}
\end{equation}
with the logarithm continued from \(\sMand<4m^2\).

\begin{proposition}[Entire remainder after finite normalization]
\label{prop:entire-bubble-remainder}
For \(m>0\) and \(t>0\), on
\(\CC\setminus[4m^2,\infty)\),
\begin{equation}
 b_t(\sMand)=b_{\rm std}(\sMand)+R_t(\sMand),
 \label{eq:bubble-entire-decomposition}
\end{equation}
where the remainder is entire and has the representation
\begin{equation}
 R_t(\sMand)
 =\frac1{16\pi^2}\int_0^1\dd x
   \int_0^{W_{\min}(x)}\frac{\dd W}{W}\,
       e^{-m^2W}\bigl[1-e^{Wx(1-x)\sMand}\bigr].
 \label{eq:entire-bubble-remainder}
\end{equation}
It vanishes at zero, is real on the real axis, and satisfies
\begin{equation}
 |R_t(\sMand)|
   \leq\frac{3t|\sMand|}{64\pi^2}e^{t|\sMand|}.
 \label{eq:entire-remainder-bound}
\end{equation}
Thus \(b_t\to b_{\rm std}\) uniformly on compact subsets of the
cut plane as \(t\downarrow0\).
\end{proposition}

\begin{proof}
For \(\Re\sMand<4m^2\), the proper-time form of
\eqref{eq:subtracted-standard-bubble} has lower limit \(W=0\).
The normalized difference in \eqref{eq:normalized-bubble} has lower
limit \(W_{\min}(x)\). Subtracting the two convergent integrals
gives \eqref{eq:entire-bubble-remainder}.

Set \(c(x)=x(1-x)\). Throughout the remaining integration region,
\(0\leq Wc(x)\leq t\max(x,1-x)\leq t\). Consequently,
\[
 \frac{|1-e^{Wc(x)\sMand}|}{W}
       \leq c(x)|\sMand|e^{t|\sMand|}.
\]
This is integrable on the region because
\(\int_0^1c(x)W_{\min}(x)\dd x=3t/4\).
It supplies a common majorant on every compact \(\sMand\)-set,
so the remainder is entire by Theorem~\ref{thm:holomorphy} and
satisfies \eqref{eq:entire-remainder-bound}.
Analytic continuation proves \eqref{eq:bubble-entire-decomposition}
on the cut plane; the same bound gives compact-uniform convergence.
\end{proof}

For the directed argument \(\mathsf S_{\eta,\bb}\), set
\(\LogDir:=\embed\circ\log\circ\embed^{-1}\), with the
branch continued from the real subthreshold domain. The same
identity reads, before taking its boundary limit,
\begin{align}
 b_{\rm std}^{\rm dir}(\mathsf S_{\eta,\bb})
 &=\frac1{16\pi^2}\int_0^1
    \LogDir\!\left[
       \frac{m^2}{m^2-x(1-x)\mathsf S_{\eta,\bb}}
             \right]\dd x,\nonumber\\
 R_t^{\rm dir}(\mathsf S_{\eta,\bb})
 &=\frac1{16\pi^2}\int_0^1\dd x
       \int_0^{W_{\min}(x)}\frac{\dd W}{W}\,e^{-m^2W}\nonumber\\
 &\quad\times\left[1-\ExpDir\!\left(
          Wx(1-x)\mathsf S_{\eta,\bb}\right)\right],
           \nonumber\\
 b_t^{\rm dir}(\mathsf S_{\eta,\bb})
 &=b_{\rm std}^{\rm dir}(\mathsf S_{\eta,\bb})
                         +R_t^{\rm dir}(\mathsf S_{\eta,\bb}).
 \label{eq:normalized-bubble-directed-transport}
\end{align}
Each superscript \({\rm dir}\) means transport by \(\embed\);
applying \(\embed^{-1}\) gives the preceding coefficient
identities. Since \(R_t\) is entire and real on the real axis,
its real-axis directed boundary is carrier-even and has no jump.
Its contribution to a complete Dyson graph still carries that
graph's vertex and Dyson factors.

The finite normalization removes the extra shell factor from the
discontinuity, while the entire term retains nontrivial momentum
dependence: \(R_t(\sMand)<0\) for \(0<\sMand<4m^2\).
It also changes the real part above threshold. Matching the
discontinuity therefore does not determine the full amplitude.
At every \(t>0\), the subtraction in
\eqref{eq:normalized-bubble} removes an already finite reference
value. The controlled limit concerns that normalized difference;
the unsubtracted value in \eqref{eq:gaussian-zero} still diverges
as the damping is removed.

\begin{corollary}[The Gaussian remainder is not polynomially bounded]
\label{cor:Gaussian-growth-obstruction}
For real \(s>0\), the entire remainder of
Proposition~\ref{prop:entire-bubble-remainder} satisfies
\begin{equation}
 -R_t(s)\geq
 \frac{\log2}{48\pi^2}e^{-2tm^2}
       \left(e^{2ts/9}-1\right).
 \label{eq:Gaussian-remainder-growth-lower}
\end{equation}
Thus no addition of a finite polynomial makes \(R_t\), or the
normalized bubble boundary \(b_t\), polynomially bounded along
the positive real axis at fixed \(t>0\).
\end{corollary}

\begin{proof}
For positive \(s\), the integrand of \(-R_t(s)\) in
\eqref{eq:entire-bubble-remainder} is nonnegative. Restrict to the
rectangle \(1/3\leq x\leq2/3\), \(t\leq W\leq2t\),
which lies inside \(W\leq W_{\min}(x)\).
There \(x(1-x)\geq2/9\) and \(e^{-m^2W}\geq e^{-2tm^2}\).
Its \(x\)-length is \(1/3\) and
\(\int_t^{2t}\dd W/W=\log2\), proving the bound.
For the ordinary comparison, direct integration of
\eqref{eq:subtracted-standard-bubble} gives on the upper rim,
for \(s>4m^2\),
\[
 16\pi^2 b_{\rm std}(s+\ii0)
 =2-\beta(s)\log\frac{1+\beta(s)}{1-\beta(s)}
       +\ii\pi\beta(s).
\]
It grows at most logarithmically as \(s\to+\infty\).
It cannot cancel the real negative exponential lower bound, nor
can a finite polynomial.
\end{proof}

Matching the jump therefore does not restore the growth assumptions
of a finite-subtraction dispersion argument. This is a specific
obstruction in the isolated Gaussian benchmark. Contributions of a
completed physical amplitude would have to be analyzed together;
the result does not exclude a cancellation supplied by additional
nonpolynomial dynamics. The limit \(t\downarrow0\) on compact
sets remains valid and does not commute with this large-\(s\) test.

\subsubsection*{A reproducible Euclidean numerical check}

At \(P_E=0\), radial integration gives
\begin{align}
 B_{t,E}(0)
 &=\frac1{16\pi^2}
     \int_0^\infty
       \frac{u\,e^{-2tu}}{(u+m^2)^2}\dd u \nonumber\\
 &=\frac{(1+x)e^xE_1(x)-1}{16\pi^2},
 \qquad x=2tm^2 .
 \label{eq:gaussian-zero}
\end{align}
For \(\tau_*=1\) and \(am=1\),
\begin{equation}
 B_{t,E}(0)=
 0.0005318466124095473\ldots .
 \label{eq:numerical-bubble}
\end{equation}
The same value follows from the independent dimensionless
proper-time integral
\[
 \frac1{16\pi^2}\int_0^\infty
       \frac{v\,e^{-v}}{(v+2)^2}\dd v .
\]
This number specifies the kernel, scales, external momentum,
measure, and absence of coupling or symmetry factors. It is an
off-shell Euclidean check, not a predicted scattering measurement.
\subsection{A massive stable-half bubble need not be analytic at zero}
\label{sec:stable-bubble-origin}

The absence of a positive exponential moment for the canonical
stable-half law cannot in general be bypassed by citing positive
internal masses. Consider the same literal, fixed-carrier scalar
coefficient model as \eqref{eq:two-line-bubble}, now with
\[
 C_{a,m}(q)=\frac{e^{-a|q|}}{q^2+m^2},\qquad a>0,\quad m>0.
\]
Let \(\mathcal C_{a,m}(x)\) be its inverse Fourier transform
with measure \(\dd^4q/(2\pi)^4\), and write
\[
 B_{a,m}(P)=\int\frac{\dd^4q}{(2\pi)^4}
                   C_{a,m}(q)C_{a,m}(q+P).
\]

\begin{proposition}[A zero-momentum nonanalyticity of the stable-half bubble]
\label{prop:stable-bubble-nonanalytic}
The function \(B_{a,m}\) is finite for all real \(P\) and is
\(C^5\). It is not \(C^6\) at \(P=0\); in particular, no
function holomorphic in \(P^2\) on a neighbourhood of zero can
restrict to this bubble there. Its coordinate covariance has the
positive large-distance asymptotic
\begin{equation}
 \mathcal C_{a,m}(x)\sim
        \frac{3a}{4\pi^2m^2}|x|^{-5},\qquad |x|\to\infty.
 \label{eq:stable-covariance-long-tail}
\end{equation}
\end{proposition}

\begin{proof}
The complete Fourier and finite-difference argument is given in
Appendix~\ref{sec:stable-bubble-proof}. In brief, the inverse
transform of \(e^{-a|q|}\) is the positive four-dimensional
Poisson kernel
\(3a/[4\pi^2(a^2+|x|^2)^{5/2}]\). Convolution with the
positive massive free Green function, whose total integral is
\(m^{-2}\), gives \eqref{eq:stable-covariance-long-tail}.
The bubble is the Fourier transform of
\(\mathcal C_{a,m}(x)^2\). This nonnegative function has all
absolute moments of order less than six, but an infinite sixth
moment in any fixed direction. A sign-definite sixth central
difference excludes \(C^6\) regularity at zero.
\end{proof}

This is an infrared regularity obstruction, not an ultraviolet
divergence: all inverse Laplace moments and the fixed-graph UV
estimates remain valid. A positive mass in the denominator does
not remove the long-distance tail introduced by the multiplier's
cusp at \(q=0\). Finite polynomial additions cannot restore the
missing differentiability. A physical completed amplitude may
contain different vertices and compensating contributions; it
must establish any claimed analyticity separately. The proposition
does not identify this one auxiliary bubble with a physical
observable or prove failure of every completion.

The result is a new test of the monograph's canonical kernel,
derived here from its stated Euclidean line rule. It is not a
restatement of the monograph's conditional holomorphy theorem.
The effective single-block model in the next subsection has a
different kernel placement and is not subject to this conclusion
merely because it uses the same scalar profile.

\subsection{An effective radial model}
\label{sec:effective-block}

\begin{remark}[Linewise prescription and channel boundary path]
\label{rem:linewise-versus-channel-boundary}
For real Lorentzian momenta \(k,P\) and finite
\(\varepsilon>0\), the common directed prescription combines
exactly as
\begin{align}
 &\frac1{[k^2-m^2+\ii\eta\epsvec]
             [(P-k)^2-m^2+\ii\eta\epsvec]}
       \nonumber\\
 &\qquad=\int_0^1
       \frac{\dd x}{[(k-xP)^2-\Delta(x)
                                      +\ii\eta\epsvec]^2},
 \nonumber\\
 &\Delta(x)=m^2-x(1-x)P^2.
 \label{eq:exact-directed-Feynman-combination}
% Monograph source: eq:feynman_param_final.
\end{align}
In each character the affine denominator has imaginary part
\(\kappa\eta\varepsilon\), so it has no zero for
\(0\leq x\leq1\). For complex external momenta the same
real-parameter representation is used only where both character
segments remain nonzero, or through a separately justified
analytic continuation.
The linewise infinitesimal remains independent of \(x\).
With \(\sMand=P^2\), continuation of an already evaluated
invariant function through
\(P_E^2\mapsto-\mathsf S_{\eta,\bb}\) gives
\begin{equation}
 \Delta^{\rm dir}_{\eta,\bb}(x)
   =m^2-x(1-x)\mathsf S_{\eta,\bb}
   =\Delta(x)-\ii\eta x(1-x)\epsvec.
 \label{eq:channel-directed-effective-denominator}
% Monograph source: eq:effective-kernel-channel-vectorial-denominator.
\end{equation}
These are distinct finite-\(\varepsilon\) formulas. Their
boundary values can be identified only after the selected
threshold-excluding continuation and its non-tangential limit
have been justified. In particular, the Gaussian closed form is
defined by its Euclidean continuation; no equality with a raw
growing Lorentzian loop integral is assumed.
\end{remark}

For a running-unit graph the real shift \(\ell=k-xP\)
also transforms its numerator signs:
\begin{equation}
 \sgn(k_\parallel)\sgn(P_\parallel-k_\parallel)
 =\sgn(\ell_\parallel+xP_\parallel)
   \sgn((1-x)P_\parallel-\ell_\parallel).
 \label{eq:running-signs-after-parameter-shift}
% Monograph source: eq:raw-kernel-under-feynman-translation; eq:running-unit-versus-fixed-virtual-carrier.
\end{equation}
Consequently the denominator combination alone does not turn
that graph into a radial fixed-carrier integral.

For comparison, one may declare an effective single virtual block
after combining denominators:
\begin{equation}
\begin{aligned}
 B_{{\rm eff},E}(P_E)
   &=\int_0^1\dd x\int\frac{\dd^4\ell_E}{(2\pi)^4}
       \frac{\Ph_{\rm eff}(a^2\ell_E^2)}
            {[\ell_E^2+\Delta_E(x)]^2},\\
 \Delta_E(x)&=m^2+x(1-x)P_E^2.
\end{aligned}
 \label{eq:effective-model}
\end{equation}
The measure \(\rho_{\rm eff}\) belongs to this model. It is not
automatically the measure for a product of the two line kernels
in \eqref{eq:two-line-bubble}.

For the explicit choice \(\Ph_{\rm eff}(z)=e^{-\sqrt z}\),
define
\begin{equation}
 J(u;a)=\int_0^\infty r^3e^{-ur^2-ar}\dd r,
 \qquad u>0,\quad a>0.
 \label{eq:J-radial}
\end{equation}
A Schwinger parameter for the squared denominator gives
\begin{equation}
 B_{{\rm eff},E}(P_E)
   =\frac1{8\pi^2}\int_0^1\dd x
        \int_0^\infty\dd u\,
              u\,e^{-u\Delta_E(x)}J(u;a).
 \label{eq:effective-J}
\end{equation}
The corresponding directed external-invariant continuation,
with \(\Delta^{\rm coeff}_{\eta,\bb}(x)
:=\embed^{-1}[\Delta^{\rm dir}_{\eta,\bb}(x)]\),
is initially defined where
\(\Re\Delta^{\rm coeff}_{\eta,\bb}(x)>0\):
\begin{equation}
 \mathcal B^{\rm dir}_{{\rm eff},\eta,\bb}(\sMand;\varepsilon)
  =\frac1{8\pi^2}\int_0^1\dd x\int_0^\infty\dd u\,
       u J(u;a)\ExpDir\!\left[
                  -u\Delta^{\rm dir}_{\eta,\bb}(x)\right].
 \label{eq:effective-directed-J}
\end{equation}
Its ordinary coefficient is obtained by \(\embed^{-1}\);
the physical boundary additionally takes \(\varepsilon\downarrow0\)
when that limit exists on the prescribed domain.
Beyond that directly convergent domain it denotes the continuation
of the established function, whenever the required domain and
majorant exist, not the same positive-real proper-time integral.
The roots of positive real \(u\) in \(J(u;a)\) remain ordinary
real roots; the directed argument is the continued invariant
\(\Delta^{\rm dir}_{\eta,\bb}\).

The exact radial function is
\begin{equation}
 \boxed{\;
 J(u;a)=\frac{
 1+v^2-\sqrt\pi\,v(v^2+\tfrac32)e^{v^2}\erfc(v)}
 {2u^2},
 \qquad v=\frac{a}{2\sqrt u}.
 \;}
 \label{eq:J-closed}
\end{equation}
It follows by applying \(-\partial_a^3\) to
\[
 \int_0^\infty e^{-ur^2-ar}\dd r
 =\frac{\sqrt\pi}{2\sqrt u}\,
      e^{a^2/(4u)}\erfc\!\left(\frac{a}{2\sqrt u}\right).
\]
The minus sign comes from differentiating \(e^{-ar}\) three times.
The useful endpoint checks are
\begin{equation}
 J(u;0)=\frac1{2u^2},\qquad
 J(0;a)=\frac6{a^4}.
 \label{eq:J-limits}
\end{equation}
The second is the limit \(u\downarrow0\) at fixed \(a>0\);
it makes the small-\(u\) integrand of
\eqref{eq:effective-J} integrable.
For numerical evaluation one may use
\(\erfcx(v)=e^{v^2}\erfc(v)\), with the integral
\eqref{eq:J-radial} or a controlled endpoint expansion when
subtraction causes loss of precision.

The alternative
\[
 J_M(u;a,m)
    =\int_0^\infty r^3e^{-ur^2-a\sqrt{r^2+m^2}}\dd r
\]
is a different massive profile. It is not a change of variables
in \eqref{eq:J-radial}. No unproved Tricomi or Appell
reconstruction is needed for either integral.
The original linewise kernel, this effective block, and a massive
radial profile must have distinct model labels.

\subsection{Gauge and ghost covariance coefficients}

For a flat massless Euclidean gauge benchmark with covariant
parameter \(\zeta>0\), the quadratic coefficients may be chosen as
\begin{equation}
 D_{\mu\nu,E}^{\Ph}(q)
   =\frac{\Ph(a^2q^2)}{q^2}
         \bigl(P_{T,\mu\nu}+\zeta P_{L,\mu\nu}\bigr),
 \qquad
 C_{{\rm gh},E}^{\Ph}(q)=\frac{\Ph(a^2q^2)}{q^2},
 \label{eq:gauge-covariances}
\end{equation}
where \(P_{L,\mu\nu}=q_\mu q_\nu/q^2\) and
\(P_T=\delta-P_L\).
The action and ghost ordering that realize these coefficients
must still be specified, as in
Section~\ref{sec:brst}. In particular, transverse damping cannot
be deduced from BRST-exact gauge fixing alone.

For Lorentzian coefficients put
\begin{equation}
 Z_{\eta,\bb}(k):=k^2+\ii\eta\epsvec,
 \qquad
 F^{\rm dir}_{\eta,\bb}(k):=
              \Ph_{\ee}(-a^2Z_{\eta,\bb}(k)).
 \label{eq:gauge-directed-kinetic-data}
\end{equation}
The following displays use the monograph's resolved scalar-pole
covariance prescription at finite \(\varepsilon>0\); their
admitted boundary distributions take the common limit
\(\varepsilon\downarrow0^+\). They are not universal
finite-\(\varepsilon\) identities for the inverse of an
arbitrary regulated matrix Hessian. In the massless quadratic
convention the directed gauge and ghost lines are
\begin{align}
 \vec D^{\Ph}_{\mu\nu;\eta,\bb}(k)
  &=\eta\ii\widehat u_k\,
      F^{\rm dir}_{\eta,\bb}(k)
       \left[\frac{g_{\mu\nu}}{Z_{\eta,\bb}(k)}
        -(1-\zeta)\frac{k_\mu k_\nu}
                           {Z_{\eta,\bb}(k)^2}\right],
                   \nonumber\\
 \vec C^{\Ph}_{{\rm gh};\eta,\bb}(k)
  &=\eta\ii\widehat u_k\,
      \frac{F^{\rm dir}_{\eta,\bb}(k)}
           {Z_{\eta,\bb}(k)}.
 \label{eq:physical-directed-gauge-ghost-lines}
% Monograph source: eq:general-directed-gauge-coefficient; eq:general-directed-ghost-coefficient.
\end{align}
The gauge tensor is symmetric in its two Lorentz indices before
multiplication by the running unit. Consequently its unrestricted
running-unit version has the same exchange obstruction as the real
scalar: it is not an inverse-Hessian covariance for a single real
commuting gauge field without an additional compatible endpoint
realization. The fixed-carrier benchmark does not have this oddness
obstruction. Ghost/antighost and Dirac fields instead require their
complete graded, distinct-field source matrices.

The tensor and ghost coefficients inside the oriented numerator
are extracted by \(\embed^{-1}\) on their directed scalar
entries. Their full lines are reconstructed as in
\eqref{eq:scalar-line-coefficient}, or character by character
when they form part of a mixed carrier graph. The longitudinal
double pole for \(\zeta\ne1\) is already present before the
kernel is inserted.
The O--QED application in \eqref{eq:OQED-directed-lines} retains the
photon's overall minus sign from its declared Maxwell action; it is not
identified with an alternative quadratic-sign template merely because
the absolute Euclidean denominators agree.

For reference, the corresponding massive tensor coefficient is
\begin{equation}
 \mathcal R^{\rm dir}_{\mu\nu;\eta,\bb}(k)
 =\frac{F^{\rm dir}_{\eta,\bb}(k)}
           {Z_{\eta,\bb}(k)-m^2}
    \left[g_{\mu\nu}-(1-\zeta)
       \frac{k_\mu k_\nu}{Z_{\eta,\bb}(k)-\zeta m^2}\right].
 \label{eq:massive-directed-gauge-coefficient}
\end{equation}
The ghost denominator is
\(Z_{\eta,\bb}(k)-\zeta m^2\). These coefficients specify
only the quadratic part; a massive gauge theory still requires
its compatible Higgs/Goldstone, ghost and interaction completion.

For a Lorentzian Dirac line use
\(\{\gamma^\mu,\gamma^\nu\}=2g^{\mu\nu}\one\) and
\(\slashed k=\gamma^\mu k_\mu\). In the monograph's
entrywise coefficient convention, define
\begin{equation}
\begin{aligned}
 R^{D,\Ph}_\eta(k;\varepsilon)
   &:=(\slashed k+m)R^\Ph_\eta(k;\varepsilon),\\
 \vec S^{\Ph}_{\eta,\bb}(k;\varepsilon)
   &:=\eta\ii\widehat u_k\,
                     \embed(R^{D,\Ph}_\eta(k;\varepsilon)).
\end{aligned}
 \label{eq:physical-directed-Dirac-line}
\end{equation}
This rationalized Dirac coefficient prescribes the scalar pole
boundary. For example, without the kernel and with
\(D=k^2-m^2\),
\begin{equation}
 (\slashed k-m)
     \frac{\slashed k+m}{D+\ii\eta\varepsilon}
   =\frac{D}{D+\ii\eta\varepsilon}\,\one.
 \label{eq:Dirac-finite-boundary-check}
\end{equation}
Thus it is not the inverse of the unmodified Dirac matrix at
finite \(\varepsilon\). When an exact regulated matrix
inverse or a Ward identity at finite regulator is required, the
completed Hessian and its boundary datum must be specified and
inverted together in each character. Its identification with a
resolved scalar-pole boundary requires the corresponding limit
argument.
% Monograph source: eq:Omni-CC-Hessian-Green-vectorial-boundary-firewall; eq:background-commuting-virtual-covariance.
The complete spinor coefficient, including its matrix entries,
is embedded consistently. If a different central matrix
realization is retained, use its two character representatives
and \eqref{eq:complete-graph-character-reconstruction}; do not
apply \(\embed^{-1}\) to a mixed-algebra matrix.

\subsection{QED vacuum polarization and the Ward identity}
\label{sec:QED}

Use canonically normalized Euclidean matrices
\(\{\gamma^\mu,\gamma^\nu\}=2\delta^{\mu\nu}\), and set
\begin{equation}
\begin{aligned}
 S_0(k)&=\frac{-\ii\slashed k+m}{k^2+m^2},&
 S_0^{-1}(k)&=\ii\slashed k+m,\\
 \Gamma_0^\mu&=\ii\gamma^\mu,&
 F_k&=\Ph(a^2k^2).
\end{aligned}
 \label{eq:QED-conventions}
\end{equation}
The two-line graph with bare vertices is
\begin{equation}
 \Pi_{\rm bv}^{\mu\nu}(Q)
 =-e^2\int\frac{\dd^4k}{(2\pi)^4}\,
 F_kF_{k+Q}
 \Tr\!\left[\Gamma_0^\mu S_0(k)
                 \Gamma_0^\nu S_0(k+Q)\right].
 \label{eq:bare-QED-loop}
\end{equation}
An overall plane binder
\(\mathcal P^{(2)}_{\bb\mid{\rm ext}}\) is understood.
The factors \(\ii\) are included in \(\Gamma_0^\mu\).
Under an otherwise identical canonical complex identification of
the mass, gamma matrices and binder, this tensor is the negative
of the monograph's displayed \(\gamma\)-vertex loop with the
same prefactor \(-e^2\): the two vertices contribute
\(\ii^2=-1\) \cite{ChiapponeMonograph}.
% Monograph source: eq:Pi-Omni-full.
The Ward contraction and the benchmark below use
\eqref{eq:QED-conventions} throughout. Comparison with a physical
current correlator also retains its declared tensor and
Wick-rotation normalizations.

The calculation uses Euclidean coefficients and a fixed virtual
carrier. Running-unit Lorentzian graphs retain their mode signs
and require a separate evaluation.

Contract with \(Q_\mu\), using
\(Q_\mu\Gamma_0^\mu=S_0^{-1}(k+Q)-S_0^{-1}(k)\).
Cyclicity of the trace gives
\[
 Q_\mu\Pi_{\rm bv}^{\mu\nu}
 =-e^2\int\frac{\dd^4k}{(2\pi)^4}F_kF_{k+Q}
  \Tr\{\Gamma_0^\nu[S_0(k)-S_0(k+Q)]\}.
\]
Assume that the original tensor integrand is absolutely integrable
componentwise at the fixed real Euclidean \(Q\). Its contraction
is initially kept under one integral: the two summands need not
be separately integrable. For the finite positive completely
monotone kernels considered here, Lemma~\ref{lem:Ward-common-cutoff}
justifies the translation on a common bounded domain, proves that
the boundary error vanishes, and proves absolute integrability of
the final remainder. The lemma treats \(m>0\) and also covers
the unchanged locally integrable massless free expressions. It yields
\begin{equation}
 Q_\mu\Pi_{\rm bv}^{\mu\nu}(Q)
 =-e^2\int\frac{\dd^4k}{(2\pi)^4}
 F_k\,[F_{k+Q}-F_{k-Q}]
           \Tr[\Gamma_0^\nu S_0(k)].
 \label{eq:Ward-remainder}
\end{equation}
This remainder is generally nonzero.

\begin{proposition}[Nontransversality under an exact one-loop integrability test]
\label{prop:QED-general-obstruction}
Let \(F(u)=\Ph(a^2u)\) for a nonzero finite positive completely
monotone law, let \(m\geq0\), and assume
\begin{equation}
 \mathcal Q_F:=\int_0^\infty F(u)^2\dd u<\infty.
 \label{eq:QED-Ltwo-condition}
\end{equation}
Then \eqref{eq:bare-QED-loop} is absolutely convergent for every
real \(Q\), is continuous in \(Q\), and obeys
\begin{equation}
\begin{aligned}
 \Pi_{\rm bv}^{\mu\nu}(0)&=\kappa_F\delta^{\mu\nu},\\
 Q_\mu\Pi_{\rm bv}^{\mu\nu}(Q)
      &=\kappa_F Q^\nu+o(|Q|),\\
 \kappa_F
 &=\frac{e^2}{8\pi^2}\int_0^\infty
     \frac{u(u+2m^2)}{(u+m^2)^2}F(u)^2\dd u.
\end{aligned}
 \label{eq:general-Ward-obstruction}
\end{equation}
Here \(\kappa_F>0\) when \(e\ne0\).
For \(m=0\), the rational factor in the integral is one for
\(u>0\). Conversely, absolute convergence of the original
nontrivial tensor at any fixed real \(Q\) implies
\eqref{eq:QED-Ltwo-condition}, by the radial-tail argument in
Lemma~\ref{lem:Ward-common-cutoff}. No differentiability at
\(u=0\) is assumed.
\end{proposition}

\begin{proof}
Choose Hermitian Euclidean gamma matrices and put
\(S_F(k)=F(k^2)S_0(k)\). Then
\[
 \operatorname{tr}[S_F(k)^\dagger S_F(k)]
       =\frac{4F(k^2)^2}{k^2+m^2}.
\]
Equation~\eqref{eq:QED-Ltwo-condition} makes all matrix entries
of \(S_F\) square-integrable, including at \(k=0\) when
\(m=0\). Cauchy--Schwarz bounds the tensor product at every
translate \(Q\); continuity follows from continuity of
translations in \(L^2\). Direct evaluation of the trace at zero
and spherical averaging give
\[
 \Pi_{\rm bv}^{\mu\nu}(0)
 =2e^2\delta^{\mu\nu}\int\frac{\dd^4k}{(2\pi)^4}
          \frac{k^2+2m^2}{(k^2+m^2)^2}F(k^2)^2,
\]
which is \eqref{eq:general-Ward-obstruction}. Its integrand is
positive. The stated small-\(Q\) relation uses only continuity.
For the converse, the trace and monotonicity estimates in
\eqref{eq:Ward-radial-tail} imply
\(\int^\infty rF(r^2)^2\dd r<\infty\). Boundedness of
\(F\) near zero completes \eqref{eq:QED-Ltwo-condition}.
\end{proof}

The exact one-loop condition also has a direct measure test. Tonelli's
theorem applied to the two independent line parameters gives
\begin{equation}
 \mathcal Q_F=\frac1{a^2}
   \int_{[0,\infty)^2}
       \frac{\rho(\dd\tau)\rho(\dd\sigma)}{\tau+\sigma}.
 \label{eq:QED-convolution-moment}
\end{equation}
The equality holds in \([0,\infty]\), with the integrand infinite
at \((0,0)\). If finite, it excludes an atom of \(\rho\) at
zero and equals \(a^{-2}M_1[\rho*\rho]\), where \(\rho*\rho\)
is the convolution measure. It need not require \(M_1[\rho]<\infty\).
Indeed, integrating \(e^{-a^2(\tau+\sigma)u}\) over \(u>0\)
proves the formula directly. This test is for the two-line QED graph;
it is not an all-orders replacement for the moment criteria in
Section~\ref{sec:uv}.

An integration by parts also gives
\begin{equation}
 \kappa_F=-\frac{e^2}{4\pi^2}\int_0^\infty
                 \frac{u^2F(u)F'(u)}{u+m^2}\dd u.
 \label{eq:general-Ward-obstruction-derivative}
\end{equation}
Indeed, complete monotonicity gives \(F'\leq0\) for \(u>0\),
and the boundary term
\(u^2F(u)^2/(u+m^2)\) vanishes at both endpoints.
At infinity this follows from integrability and monotonicity of
\(F^2\); near zero it follows from boundedness of \(F\).
One first integrates on a compact positive interval and then uses
the nonnegative integrands to pass to the endpoints. This proves
convergence of the derivative formula even for a cusp at zero.

For example, in the massless canonical stable-half comparison,
\begin{equation}
 F(u)=e^{-a\sqrt u},\qquad
 \kappa_F=\frac{e^2}{16\pi^2a^2}.
 \label{eq:stable-half-Ward-coefficient}
\end{equation}
For \(F(u)=(1+a^2u)^{-\nu}\) the condition is exactly
\(\nu>1/2\), and at \(m=0\),
\(\kappa_F=e^2/[8\pi^2a^2(2\nu-1)]\).
These are finite nontransverse coefficients of the declared bare
comparison tensor, not physical photon masses. The general remainder
is \(o(|Q|)\); a cubic error requires additional smoothness.

For a sharper remainder, suppose \(k\mapsto F(k^2)\) is \(C^3\).
Require the first through third \(Q\)-derivatives of the complete
weighted expression
\[
 F_k F((k+Q)^2)\Tr[\Gamma_0^\nu S_0(k)]
\]
to have a common \(L^1(\dd^4k)\) majorant for \(Q\) near zero.
This includes the numerator weight and any infrared behaviour, not
just derivatives of the scalar profile alone. Then use
\(\Tr[\Gamma_0^\nu S_0(k)]=4k^\nu/(k^2+m^2)\) and
\[
 F((k+Q)^2)-F((k-Q)^2)
       =4(k\cdot Q)F'(k^2)+O(|Q|^3).
\]
Rotational symmetry then gives
\begin{equation}
 Q_\mu\Pi_{\rm bv}^{\mu\nu}(Q)
 =-4e^2Q^\nu
   \int\frac{\dd^4k}{(2\pi)^4}
       \frac{k^2F(k^2)F'(k^2)}{k^2+m^2}
       +O(|Q|^3).
 \label{eq:Ward-smallQ}
\end{equation}
For the Gaussian kernel \(F(k^2)=e^{-tk^2}\), with \(m>0\),
put \(x_t=2tm^2\). The linear coefficient has the closed form
\begin{align}
 \kappa_t
 &:=-4e^2\int\frac{\dd^4k}{(2\pi)^4}
       \frac{k^2F(k^2)F'(k^2)}{k^2+m^2} \nonumber\\
 &=\frac{e^2t}{4\pi^2}\int_0^\infty
       \frac{u^2e^{-2tu}}{u+m^2}\dd u \nonumber\\
 &=\frac{e^2}{16\pi^2t}
       \left[1-x_t+x_t^2e^{x_t}E_1(x_t)\right].
 \label{eq:gaussian-Ward-coefficient}
\end{align}
The last equality uses
\(u^2/(u+m^2)=u-m^2+m^4/(u+m^2)\).
The positive integral proves \(\kappa_t>0\) for \(e\neq0\), so
\begin{equation}
 Q_\mu\Pi_{\rm bv}^{\mu\nu}(Q)
       =\kappa_tQ^\nu+O(|Q|^3),\qquad
 \Pi_{\rm bv}^{\mu\nu}(0)=\kappa_t\delta^{\mu\nu}.
 \label{eq:gaussian-Ward-violation}
\end{equation}
The second equality follows by rotational covariance and the smooth
expansion at \(Q=0\), or directly by integrating the zero-momentum
trace by parts. This finite nontransverse contribution is a specific
term that the completed functional must cancel; it does not determine
a physical photon mass.

\subsubsection*{A compatible longitudinal vertex}

Let the completed inverse coefficient propagator be
\[
 S_\Ph^{-1}(p)=\ii A(p^2)\slashed p+B(p^2).
\]
For \(S_\Ph=F S_0\), one has \(A=F^{-1}\) and
\(B=mF^{-1}\).
Define the symmetric average and divided differences by
\[
\begin{aligned}
 \overline A(p',p)&=\tfrac12[A(p'^2)+A(p^2)],\\
 A^{[1]}(p',p)&=\frac{A(p'^2)-A(p^2)}{p'^2-p^2},\\
 B^{[1]}(p',p)&=\frac{B(p'^2)-B(p^2)}{p'^2-p^2}.
\end{aligned}
\]
A longitudinal vertex of Ball--Chiu type \cite{BallChiu1980} is
\begin{equation}
\begin{aligned}
 \Gamma_L^\mu(p',p)
 ={}&\ii\overline A(p',p)\gamma^\mu\\
 &+(p'+p)^\mu\Bigl[
     \tfrac{\ii}{2}A^{[1]}(p',p)(\slashed p'+\slashed p)\\
 &\hspace{6.4em}+B^{[1]}(p',p)\Bigr].
\end{aligned}
 \label{eq:longitudinal-vertex}
\end{equation}
Contracting with \(q=p'-p\) gives exactly
\[
 q_\mu\Gamma_L^\mu
       =S_\Ph^{-1}(p')-S_\Ph^{-1}(p).
\]
When the functions are differentiable, equal-invariant difference
quotients are interpreted by their derivative limits.
An arbitrary transverse addition remains possible.

Equation~\eqref{eq:longitudinal-vertex} repairs this one Ward
identity. It does not specify the entire theory. Contact terms,
higher photon vertices, and the non-Abelian generalization must
come from the completed action, and their growth must be included
in the UV estimates. In particular, the inverse \(F^{-1}\) in
this vertex may counteract line suppression.

\begin{proposition}[An exact norm test for a Ward-compatible auxiliary vertex]
\label{prop:Ward-vertex-growth}
Let \(S_\Ph^{-1}(p)=F(p^2)^{-1}(\ii\slashed p+m)\),
with \(F>0\), \(F(0)>0\), \(m\geq0\), and Hermitian Euclidean
gamma matrices. Write \(A_R=F(R^2)^{-1}\), \(A_0=F(0)^{-1}\).
Any vertex satisfying \eqref{eq:dressed-WTI} obeys the exact identity
\begin{equation}
 \left\|n_\mu\Gamma_\Ph^\mu(Rn,0)\right\|^2
       =A_R^2+\frac{m^2}{R^2}(A_R-A_0)^2,
 \label{eq:Ward-vertex-exact-norm}
\end{equation}
for a unit vector \(n\) and \(R>0\). In particular,
\begin{equation}
 \left\|n_\mu\Gamma_\Ph^\mu(Rn,0)\right\|
 \geq\frac1R\left[
       \frac{\sqrt{R^2+m^2}}{F(R^2)}-\frac{m}{F(0)}
                      \right].
 \label{eq:Ward-vertex-growth-lower}
\end{equation}
The norm is the matrix operator norm. For a superpolynomially
decreasing \(F(R^2)\), no globally polynomially bounded vertex
can satisfy this identity in the stated auxiliary field coordinates.
\end{proposition}

\begin{proof}
Set \(p'=Rn\), \(p=0\) in the Ward identity and divide by
\(R\). The contraction is
\(\ii A_R\slashed n+m(A_R-A_0)\one/R\).
Since \(\slashed n\) is Hermitian and \((\slashed n)^2=\one\),
its adjoint times itself is
\([A_R^2+m^2(A_R-A_0)^2/R^2]\one\). This proves the exact
norm. The reverse triangle inequality also gives
\eqref{eq:Ward-vertex-growth-lower}. A transverse addition leaves
this contraction unchanged.
\end{proof}

For the Gaussian and stable-half kernels the lower bound grows
respectively like \(e^{tR^2}\) and \(e^{aR}\). This first test
uses a growing photon momentum \(q=Rn\); by itself it is not a
large-loop-momentum estimate at fixed external photon momentum.
The following diagonal test addresses that distinct regime.

\paragraph{Soft photon, large fermion momentum}
Suppose additionally that \(F\) is differentiable for positive
arguments and the vertex is continuous at \((p',p)=(p,p)\)
for each \(p\ne0\). Put \(A(u)=F(u)^{-1}\). Dividing the
Ward identity with \(p'=p+hn\) by \(h\), then taking
\(h\to0\) for every unit \(n\), fixes
\begin{equation}
\begin{aligned}
 \Gamma_\Ph^\mu(p,p)
 &=\frac{\partial S_\Ph^{-1}(p)}{\partial p_\mu}\\
 &=\ii A(p^2)\gamma^\mu
    +2p^\mu A'(p^2)(\ii\slashed p+m).
\end{aligned}
 \label{eq:soft-Ward-vertex}
\end{equation}
A continuous transverse addition has zero diagonal limit, by the
same directional argument. At \(p=Rn\), write \(A'_R=A'(R^2)\).
For any unit \(v\) orthogonal to \(n\), Hermiticity of the
Euclidean gamma matrices gives
\begin{equation}
\begin{aligned}
 \|n_\mu\Gamma_\Ph^\mu(p,p)\|^2
   &=(A_R+2R^2A'_R)^2+4m^2R^2(A'_R)^2,\\
 \|v_\mu\Gamma_\Ph^\mu(p,p)\|&=A_R.
\end{aligned}
 \label{eq:soft-Ward-norms}
\end{equation}
This proves nonpolynomial vertex growth for the same kernels with
photon momentum fixed at zero and fermion momentum unbounded,
provided the stated diagonal regularity holds. It is still a
vertex test, not a lower bound on a complete loop integral.
The full product of vertices and covariances must be estimated;
individual vertex growth does not prove graph divergence or exclude
cancellation by the required contact graphs. Cancellations or changed
field coordinates require their own estimates.
Proposition~\ref{prop:residual-coercivity} gives explicit sufficient
conditions when nonpolynomial vertex growth leaves a positive damping
margin. This obstruction to a particular sufficient hypothesis is not a proof
that every completed graph diverges. Together with
Proposition~\ref{prop:pure-redefinition-cancellation}, it gives two
explicit tests of the interaction completion instead of assuming
that source normalization or a single Ward identity supplies it.

\paragraph{The stable-half vertex at zero momentum}
The equal-invariant derivative prescription in
\eqref{eq:longitudinal-vertex} only applies where those derivatives
exist. For \(F(p^2)=e^{-a|p|}\), \(a>0\), and \(m>0\),
no Ward-compatible auxiliary vertex can be jointly continuous at
\((p',p)=(0,0)\). Indeed, for every unit \(n\), the Ward identity
forces
\begin{equation}
 \lim_{R\downarrow0}n_\mu\Gamma_\Ph^\mu(Rn,0)
       =\ii\slashed n+am\one.
 \label{eq:stable-half-vertex-origin-limit}
\end{equation}
A continuous vertex would make the left side
\(n_\mu\Gamma_\Ph^\mu(0,0)\), which is odd in \(n\).
The nonzero scalar term \(am\one\) is even in \(n\), a
contradiction. The same obstruction follows from the trace alone.
Thus the difference quotients do not silently supply a regular
zero-momentum completion for this massive cusp profile. This is a
bounded directional discontinuity test, not an assertion of an
infinite vertex norm: the squared contracted norm in
\eqref{eq:Ward-vertex-exact-norm} tends to \(1+a^2m^2\).
For \(m=0\) this particular continuity obstruction
is absent; no higher-regularity conclusion is inferred. These
statements concern the specified auxiliary inverse propagator, not
a separately defined physical current vertex.

The preceding equations are ordinary Euclidean coefficient
calculations. For the full carrier graph, let
\(S_\kappa\), \(\Gamma_\kappa^\mu\), and
\(\Pi_\kappa^{\mu\nu}\) be the complete character
representatives supplied consistently by the completed functional,
with all factors in each character using its complete branch data. Their reconstructions are
\begin{equation}
 \vec S=\sum_{\kappa=\pm1}\mathsf P_{\kappa}S_\kappa,
 \qquad
 \vec\Gamma^\mu=\sum_{\kappa=\pm1}\mathsf P_{\kappa}
                                  \Gamma_\kappa^\mu,
 \qquad
 \vec\Pi^{\mu\nu}=\sum_{\kappa=\pm1}\mathsf P_{\kappa}
                                  \Pi_\kappa^{\mu\nu}.
 \label{eq:QED-character-reconstruction}
\end{equation}
One character is used throughout each traced loop product.
If the inverse propagator exists in both characters, then
\(\vec S^{-1}=\sum_\kappa\mathsf P_{\kappa}S_\kappa^{-1}\).
Consequently, the pair of coefficient Ward identities
\begin{equation}
 q_\mu\Gamma_\kappa^\mu(p',p)
        =S_\kappa^{-1}(p')-S_\kappa^{-1}(p)
 \quad(\kappa=\pm1)
 \label{eq:QED-character-Ward-identities}
\end{equation}
reconstructs exactly the carrier-algebra identity
\(q_\mu\vec\Gamma^\mu
=\vec S^{-1}(p')-\vec S^{-1}(p)\).
Likewise, transversality of the completed tensor must hold in
each character. A scalar coefficient check with omitted running
line signs does not establish this identity for the unrestricted
Lorentzian graph.

For the full completed tensor, write
\begin{equation}
 \Pi_{\rm full}^{\mu\nu}
      =\Pi_{\rm bv}^{\mu\nu}+\Pi_{\rm comp}^{\mu\nu}.
 \label{eq:Pi-complete}
\end{equation}
Here \(\Pi_{\rm comp}\) includes the changes in vertices and all
additional graphs required by the completed functional.
Its Ward identity gives
\(Q_\mu\Pi_{\rm full}^{\mu\nu}=0\).
For \(Q\ne0\), if the completed plane also has the stated
Euclidean covariance and no additional background tensors, then
\begin{equation}
 \Pi_{\rm full}^{\mu\nu}(Q)
    =(Q^2\delta^{\mu\nu}-Q^\mu Q^\nu)
                       \Pi_{\rm full}(Q^2).
 \label{eq:Pi-transverse}
\end{equation}
The scalar function in this equation cannot be computed from
\eqref{eq:bare-QED-loop} alone. The decomposition is not a proof of
regularity at \(Q=0\): a transverse scalar form factor may have a
singularity there. If the tensor is continuous at zero, transversality
forces \(\Pi_{\rm full}^{\mu\nu}(0)=0\). If it is analytic
there, the scalar coefficient has the corresponding regular extension.
These properties require the completed graph sum and its limits.

\subsubsection*{Parameter and cut bookkeeping}

In the two-line Schwinger representation, define
\[
 w_i=\alpha_i+a^2\tau_i,\qquad
 W=w_1+w_2,\qquad x=\frac{w_2}{W}.
\]
The Gaussian exponent after translating the loop variable is
\[
 -m^2(\alpha_1+\alpha_2)-x(1-x)WQ^2.
\]
If \(\alpha_1,\alpha_2,\tau_1,\tau_2\) are still independent
integration variables, \(x\) is this function of them, not an
additional independent Feynman parameter. Changing variables to
\((W,x)\) requires the corresponding Jacobian and lower limits.
The tensor numerator and completion terms must be retained through
that change of variables.

As a reference for a conventional unit-vertex cut with physical
normalization, ordinary QED gives
\begin{equation}
 \Im\Pi^{\rm coeff}_{{\rm std},+}(\sMand)
  =\frac{e^2}{12\pi}
    \left(1+\frac{2m^2}{\sMand}\right)
       \sqrt{1-\frac{4m^2}{\sMand}}\,
       \theta(\sMand-4m^2).
 \label{eq:QED-reference-spectrum}
\end{equation}
Let \(\Pi_{\rm std}^{\rm sheet}\) denote the ordinary
Lorentzian current-correlation scalar in the stated normalization,
and set
\begin{equation}
\begin{aligned}
 \Pi^{\rm dir}_{{\rm std},\eta,\bb}(\sMand)
  &:=\lim_{\varepsilon\downarrow0}
    (\embed\circ\Pi^{\rm sheet}_{\rm std}
                   \circ\embed^{-1})(\mathsf S_{\eta,\bb}),\\
 \Pi^{\rm coeff}_{{\rm std},\eta}
       &:=\embed^{-1}(\Pi^{\rm dir}_{{\rm std},\eta,\bb}).
\end{aligned}
 \label{eq:QED-directed-spectrum-interface}
\end{equation}
For the reflected ordinary reference,
\begin{equation}
 \Disc^{\rm dir}_\eta\Pi_{\rm std}
 =2\eta\je\,
                  \Im\Pi^{\rm coeff}_{{\rm std},+}(\sMand).
 \label{eq:QED-directed-reference-jump}
\end{equation}
The square root in the spectral density is the nonnegative real
phase-space root. This reference scalar omits the full graph's
Dyson factors; the completed polarization and scattering graphs
use \eqref{eq:QED-character-reconstruction}.

This is a normalization benchmark for the ordinary current
matrix element. The completed damped theory has its own
on-shell subamplitudes, residues, and possible kernel branch
contributions. Equality with \eqref{eq:QED-reference-spectrum}
requires their matching; it does not follow from deleting
\(\Ph\) on the cut lines. Resummation of
\(\Pi_{\rm full}\) likewise requires analyzing its actual zeros
and singularities rather than presuming an unchanged pole set.

\section{Application to the fixed-depth O--QED test theory}
\label{sec:OQED-application}

We now apply the preceding criteria to the standalone, support-diagonal
O--QED theory specified in Ref.~\cite{ChiapponeIgnition}, rather than
identifying the bare QED benchmark with a completed gauge theory.
The source supplies a local/reference sector action, a fixed-depth
internal-line rule, and a conditional vertex/contact completion. Its
Eqs.~(48)--(53), (66), and (75)--(83) are the inputs used here. The
inverse moments and graph estimates below are derived for those inputs.
The source does not give an evaluated magnetized completion or its
all-graph growth bounds; a conditional application cannot silently
supply those missing objects.

\subsection{Source action, normalization, and Abelian BRST}
\label{sec:OQED-action}

Fix one support \(\bb=(s,b_s)\), one carrier, and one ordering branch.
The source's coordinate \(G_{0\bb}\) is the centred carrier coordinate
\(X_0\) used here. Its standalone N/A action fold is not the HOUQFT
metric-copy fold. Suppressing the fixed label within the density, write
its conjugation-symmetrized matter block as
\begin{equation}
\begin{aligned}
 \mathcal L_{\rm mat}
 =\tfrac12\bigl[&\bar\Psi(\ii\gamma^\mu D_\mu-m_e)\Psi\\
 &+\overline{\Psi^*}(-\ii\gamma^{\mu *}\bar D_\mu-m_e)\Psi^*\bigr],\\
 D_\mu&=\partial_\mu+\ii e A_\mu,\qquad
 \bar D_\mu=\partial_\mu-\ii e A_\mu .
\end{aligned}
\label{eq:OQED-source-matter}
\end{equation}
Here \(\Psi^*\) denotes the source's componentwise conjugate, not an
undeclared replacement by a charge-conjugate spinor. On the displayed
ascending action branch its sector assembly is
\begin{equation}
\begin{aligned}
 \vec S_{\rm ref}^{\rm lock}
 &=\hat p_0\sum_{s\in\{N,A\}}\sigma_s
       \int\dd^4x_{\bb_s}\,\ell_{\rm GT}^{K-1}
       \mathcal L_{{\rm ref},\bb_s},\\
 \mathcal L_{\rm ref}
 &=\mathcal L_{\rm mat}-\tfrac14F_{\mu\nu}F^{\mu\nu}
       -\frac{(\partial\cdot A)^2}{2\alpha_{\rm gf}} .
\end{aligned}
\label{eq:OQED-source-action}
\end{equation}
The free Abelian ghost term suppressed in this physical-amplitude
display is retained in the BRST functional. The source requires
\(\alpha_{\rm gf}\ne0\); its Feynman-gauge specialization is
\(\alpha_{\rm gf}=1\). The carrier factor is inserted in assembling
the action, not treated as an extra scalar coupling. On another fixed
ordering branch use the action-unit convention
\eqref{eq:action-branch-unit} together with its source convention.
The proofs below are fibrewise: the displayed assembly is not a sum
used to cancel divergent N/A coefficients.

The scale \(\ell_{\rm GT}\) normalizes the locked measure. It is neither
the virtual propagation depth nor a plasma response length. Passing to
the canonical coefficient variables of \eqref{eq:QED-conventions}
requires a consistent field/source normalization and the matched charge
\(e\); it is not accomplished by dropping a dimensional prefactor from
only the kinetic term. The conjugation display alone also does not fix
an extra factor of two in a fermion loop. Physical species multiplicities
and the independent integration variables belong to the completed
source and measure record. Throughout the tests below the flat vacuum
coefficient has one four-component Dirac species. Additional diagonal
species contribute a sum of closed species traces, each with its own
charge product \(e_f^2\), mass, kernel, and matching contacts. For
\(N_f\) identical species this is \(N_f\) times the one-species loop,
not a factor \(N_f\) at each vertex. A flavour-mixing theory instead
retains the full ordered flavour matrices inside that trace.

The gauge convention in the source is
\(A_\mu\mapsto A_\mu-\partial_\mu\alpha\),
\(\Psi\mapsto e^{\ii e\alpha}\Psi\). For a left odd derivation a
compatible Abelian differential is
\begin{equation}
\begin{aligned}
 \BR A_\mu&=-\partial_\mu c,&\quad \BR c&=0,\\
 \BR\Psi&=\ii e c\Psi,&\BR\bar\Psi&=\ii e\bar\Psi c,\\
 \BR\Psi^*&=-\ii e c\Psi^*,&
 \BR\overline{\Psi^*}&=-\ii e\overline{\Psi^*}c,\\
 \BR\bar c&=B,&\BR B&=0.
\end{aligned}
\label{eq:OQED-Abelian-BRST}
\end{equation}
The right-ordered ghost in \(\BR\bar\Psi\) is intentional: the graded
Leibniz rule gives \(\BR(\bar\Psi\Psi)=0\). Also
\(\BR(D_\mu\Psi)=\ii e cD_\mu\Psi\), \(\BR F_{\mu\nu}=0\),
and \(\BR^2=0\). These identities follow directly by cancellation of
\(\partial_\mu c\) terms and by \(c^2=0\).
For example, in Euclidean coefficient variables put
\(\mathcal F=-\partial\cdot A\) and
\(\Psi_{\rm gf}=\int\bar c\,H(-\partial_E^2)
(\mathcal F-\zeta B/2)\), where \(H\) is field independent and
\(\zeta>0\). Then
\[
 \BR\Psi_{\rm gf}=\int\bigl[BH\mathcal F-\tfrac\zeta2BHB
                  +\bar cH(-\partial_E^2)c\bigr].
\]
The same auxiliary-field contour qualification as in
Section~\ref{sec:brst} gives positive Euclidean gauge fixing. The
Abelian ghosts remain free on this linear-gauge branch. This verifies
the differential and gauge-fixing sector, not the unspecified nonlocal
matter/contact completion. The BV theorem applies to that completion
only with its measure, source maps, and insertions established.

\subsection{The specified kernel passes every finite inverse-moment test}
\label{sec:OQED-moments}

Use the following dictionary, without equating independent scales:
\[
 a:=a_{\rm ker}^{(p)}>0,\qquad M:=M_{\rm ker}\geq0,\qquad
 \ell_\Omega=\hbar a,\qquad \mu:=aM.
\]
The electron pole parameter \(m_e\) remains in its denominator;
\(M\) is the auxiliary kernel scale. With a dimensionful nonnegative
kinetic variable \(u\), the source's normalized representative is
\begin{equation}
\begin{aligned}
 F_{a,M}(u)&:=e^{-a h_M(u)},&
 h_M(u)&:=\sqrt{u+M^2}-M,\\
 \Ph_\mu(z)&:=e^{\mu-\sqrt{z+\mu^2}},&
 F_{a,M}(u)&=\Ph_\mu(a^2u).
\end{aligned}
\label{eq:OQED-normalized-kernel}
\end{equation}
In particular \(F_{a,M}(0)=1\). For \(M>0\) this is not the same
profile as either \(e^{-a\sqrt u}\) or the unnormalized
\(e^{-a\sqrt{u+M^2}}\).

\begin{proposition}[Explicit O--QED kernel certificate]
\label{prop:OQED-moment-certificate}
The dimensionless Laplace measure of \(\Ph_\mu\) is the probability
measure
\begin{equation}
 \rho_\mu(\dd\tau)=\frac{e^\mu}{2\sqrt\pi}\,
 \tau^{-3/2}\exp\!\left(-\mu^2\tau-\frac1{4\tau}\right)\dd\tau.
\label{eq:OQED-dimensionless-measure}
\end{equation}
It has no atom at zero and no positive support gap, including when
\(\mu>0\). Every inverse moment is finite. For \(r>0\),
\begin{equation}
 M_r[\rho_\mu]=
 \begin{cases}
 \displaystyle\frac{e^\mu}{\sqrt\pi}(2\mu)^{r+1/2}
                      K_{r+1/2}(\mu),&\mu>0,\\[2pt]
 \displaystyle\frac{4^r\Gamma(r+1/2)}{\sqrt\pi},&\mu=0.
 \end{cases}
\label{eq:OQED-inverse-moments}
\end{equation}
Here \(K_\nu\) is the modified Bessel function, not the number of GT
directions. In particular,
\begin{equation}
\begin{aligned}
 M_1&=2(1+\mu),\\
 M_2&=4(\mu^2+3\mu+3),\\
 M_3&=8(\mu^3+6\mu^2+15\mu+15).
\end{aligned}
\label{eq:OQED-first-moments}
\end{equation}
For \(\mu>0\) its positive exponential moments exist precisely for
\(0<b\leq\mu^2\), with
\begin{equation}
 \int e^{b\tau}\rho_\mu(\dd\tau)
       =\exp\!\left(\mu-\sqrt{\mu^2-b}\right).
\label{eq:OQED-exponential-moment}
\end{equation}
There is no positive exponential moment when \(\mu=0\).
\end{proposition}

\begin{proof}
Rescale the source's dimensionful heat parameter \(\lambda=a^2\tau\)
in its Eq.~(66). The stable-half subordination identity gives
\(\int e^{-z\tau}\rho_\mu(\dd\tau)=\Ph_\mu(z)\), and evaluation
at zero proves normalization. The density is strictly positive on
\((0,\infty)\); tempering its large-\(\tau\) tail creates no gap
at small \(\tau\). The factor \(e^{-1/(4\tau)}\) makes every inverse
moment integrable at zero, and the finite measure tail controls infinity.
Using the standard integral representation of \(K_\nu\)
\cite[Eq.~10.32.10]{DLMF10}, or its elementary change of variable,
\[
 \int_0^\infty \tau^{-r-3/2}e^{-\mu^2\tau-1/(4\tau)}\dd\tau
       =2(2\mu)^{r+1/2}K_{r+1/2}(\mu).
\]
For \(\mu=0\), substitute \(v=1/(4\tau)\). The three integer
moments follow by half-integer evaluation, or by integrating
\(z^{r-1}\Ph_\mu(z)/\Gamma(r)\) after setting
\(y=\sqrt{z+\mu^2}-\mu\). Finally, replacing \(\mu^2\) by
\(\mu^2-b\) in the density proves
\eqref{eq:OQED-exponential-moment}. At \(b=\mu^2\) the tail is
\(\tau^{-3/2}\), still integrable; beyond it the integral diverges.
At that endpoint its first differentiated moment diverges, so finiteness
there does not extend holomorphy across the branch point.
\end{proof}

The dimensionful inverse moment is
\begin{equation}
 \int\lambda^{-r}\rho^{\rm heat}_{a,M}(\dd\lambda)
       =a^{-2r}M_r[\rho_{aM}].
\label{eq:OQED-dimensional-moment}
\end{equation}
Thus the actual condition \(M_{1+D/2}<\infty\) is verified for every
finite numerator degree, without inferring it from normalization or
from absence of an atom. For \(a\geq a_->0\) and \(aM\leq\mu_+\), the density also gives
\[
 a^{-2r}M_r[\rho_{aM}]
 \leq a_-^{-2r}e^{\mu_+}\frac{4^r\Gamma(r+1/2)}{\sqrt\pi}.
\]
This supplies explicit uniform kernel constants on the corresponding
compact parameter sets. No uniform damping-removal limit follows.
A useful separate bound is
\begin{equation}
 e^{-a\sqrt u}\leq F_{a,M}(u)
       \leq\min\{1,e^{aM}e^{-a\sqrt u}\},\qquad u\geq0.
\label{eq:OQED-linear-exponential-bounds}
\end{equation}
It follows from \(\sqrt u-M\leq h_M(u)\leq\sqrt u\).
Increasing \(M\) at fixed \(a,u>0\) weakens the suppression; it does
not improve the large-momentum exponential rate.


\subsection{Positive-depth stability and singular depth limits}
\label{sec:OQED-depth-limits}

The UV theorem fixes positive depth. We now distinguish three
mathematically different questions: controlled differentiation at
positive depth, large-depth localization, and the singular removal
of damping. They concern the
positive Euclidean scalar line and bubble comparisons, not a
calculation of the completed polarization tensor.

\begin{proposition}[Depth derivatives and one-line asymptotics]
\label{prop:OQED-depth-stability}
For $F_{a,M}$ of \eqref{eq:OQED-normalized-kernel} and every
integer $n\geq0$,
\begin{equation}
 \partial_a^nF_{a,M}(u)=(-h_M(u))^nF_{a,M}(u),\qquad
 h_M(u)=\sqrt{u+M^2}-M.
 \label{eq:OQED-depth-derivatives}
\end{equation}
For $a\in[a_-,a_+]$, $0<a_-\leq a_+$, and $0\leq M\leq M_+$,
\begin{equation}
 |\partial_a^nF_{a,M}(q^2)|
 \leq e^{a_+M_+}|q|^n e^{-a_-|q|}.
 \label{eq:OQED-depth-uniform-envelope}
\end{equation}
In particular, such derivatives preserve an established complete
graph bound when the additional polynomial numerator degree and
infrared factors are controlled. Also, for $u>0$,
$\partial_MF_{a,M}(u)>0$.

Define a Fourier-normalized one-line comparison for $d\geq0$:
\begin{equation}
 \mathscr Q_d(a,M,m)
 :=\int_{\RR^4}\frac{\dd^4q}{(2\pi)^4}
   \frac{|q|^dF_{a,M}(q^2)}{q^2+m^2},
 \quad m\geq0.
 \label{eq:OQED-one-line-depth}
\end{equation}
For every fixed $m,M\geq0$, it is finite at $a>0$, and
\begin{equation}
 \lim_{a\downarrow0}a^{d+2}\mathscr Q_d(a,M,m)
 =\frac{\Gamma(d+2)}{8\pi^2}.
 \label{eq:OQED-zero-depth-one-line}
\end{equation}
The convergence is uniform for $(m,M)$ in bounded nonnegative sets.
For fixed $m>0$, the two different strong-depth limits are
\begin{equation}
 \begin{aligned}
 \lim_{a\to\infty}a^{d+4}\mathscr Q_d(a,0,m)
   &=\frac{\Gamma(d+4)}{8\pi^2m^2},\\
 \lim_{a\to\infty}a^{(d+4)/2}\mathscr Q_d(a,M,m)
   &=\frac{(2M)^{(d+4)/2}\Gamma((d+4)/2)}
           {16\pi^2m^2},\quad M>0.
 \end{aligned}
 \label{eq:OQED-strong-depth-one-line}
\end{equation}
These asymptotics hold in their stated regimes, not uniformly over
arbitrary joint limits of $a$ and $M$.
\end{proposition}
\begin{proof}
The derivatives follow by differentiating the exponential at fixed
$M$. The inequalities
$0\leq h_M(q^2)\leq|q|$ and
$F_{a,M}(q^2)\leq e^{a_+M_+}e^{-a_-|q|}$ prove the uniform bound.
Also $\partial_Mh_M(u)=M/\sqrt{u+M^2}-1<0$ for $u>0$.

Four-dimensional radial integration gives
\[
 \mathscr Q_d=\frac1{8\pi^2}\int_0^\infty
 \frac{r^{d+3}e^{-a(\sqrt{r^2+M^2}-M)}}{r^2+m^2}\dd r.
\]
For the zero-depth limit, set $r=v/a$ and multiply by $a^{d+2}$.
The resulting denominator is $v^2+a^2m^2$ and the exponent is
$-\sqrt{v^2+a^2M^2}+aM$. It converges pointwise to $-v$, with
an integrable common upper bound proportional to
$v^{d+1}e^{-v}$ for $a$ sufficiently small and bounded masses.
Dominated convergence gives \eqref{eq:OQED-zero-depth-one-line},
including $m=0$. For $M=0$ and $a\to\infty$, the same $r=v/a$
substitution, now multiplied by $a^{d+4}$, leaves the limiting
massive denominator $m^2$ and the moment $\Gamma(d+4)$.
For fixed $M>0$, set $r=v/\sqrt a$; then
$a(\sqrt{M^2+v^2/a}-M)\to v^2/(2M)$.
The radial Gaussian integral gives the second limit.
For the last limit, the rescaled exponent is
$v^2/[\sqrt{M^2+v^2/a}+M]$. For all $a\geq1$ it is
at least $v^2/(2M+v)$, which gives an integrable dominator
proportional to $v^{d+3}e^{-v^2/(2M+v)}$ after bounding the
massive denominator by $m^{-2}$. This also justifies the earlier
changes of limit.
\end{proof}


\begin{corollary}[An explicit O--QED exponential momentum tail]
\label{cor:OQED-exponential-line-tail}
For $a>0$, $M,m\geq0$, $d\geq0$ and $R>0$, the
Fourier-normalized line integral of
\eqref{eq:OQED-one-line-depth} obeys
\begin{equation}
 \begin{aligned}
 &\int_{|q|>R}\frac{\dd^4q}{(2\pi)^4}
     \frac{|q|^dF_{a,M}(q^2)}{q^2+m^2}\\
 &\qquad\leq \frac{e^{aM}}{8\pi^2a^{d+2}}
                     \Gamma(d+2,aR),
 \end{aligned}
 \label{eq:OQED-exponential-line-tail}
\end{equation}
where $\Gamma(s,x)=\int_x^\infty
 t^{s-1}e^{-t}\dd t$ is the upper incomplete gamma
function. For each fixed positive depth the right side is
exponentially small as $R\to\infty$. Selected-line graph tails
inherit finite products and union bounds only when the
\emph{complete} anisotropic numerator and unselected-line
estimate in Theorem~\ref{thm:optimized-selected-lines}
has been verified. These are Euclidean coefficient bounds,
not contour-arc or physical-cut estimates.
\end{corollary}
\begin{proof}
Use $F_{a,M}(r^2)\leq e^{aM}e^{-ar}$ and
$(r^2+m^2)^{-1}\leq r^{-2}$ for $r>0$. Four-dimensional
radial integration bounds the tail by
$e^{aM}(8\pi^2)^{-1}\int_R^\infty r^{d+1}e^{-ar}\dd r$,
which is the stated incomplete-gamma expression.
\end{proof}

For every $M\geq0$, the literal two-line scalar comparison
used below is already defined for real Euclidean $P$ by
\begin{equation}
 B_{a,M,m}(P):=\int\frac{\dd^4q}{(2\pi)^4}
 \frac{F_{a,M}(q^2)F_{a,M}((q+P)^2)}
 {(q^2+m^2)((q+P)^2+m^2)},\quad m>0.
 \label{eq:OQED-literal-depth-bubble}
\end{equation}
This real Euclidean definition also covers $M=0$; the stronger
holomorphic-tube statement below assumes $M>0$.


\begin{proposition}[Depth monotonicity of the literal scalar bubble]
\label{prop:OQED-bubble-depth-monotone}
Fix $m>0$, $M\geq0$, and real Euclidean $P$. Put
$H_{M,P}(q):=h_M(q^2)+h_M((q+P)^2)$. For each $a>0$,
$a\mapsto B_{a,M,m}(P)$ is completely monotone:
\begin{equation}
 \begin{aligned}
 &(-1)^n\partial_a^n B_{a,M,m}(P)\\
 &\quad=\int\frac{\dd^4q}{(2\pi)^4}
 \frac{H_{M,P}(q)^n e^{-aH_{M,P}(q)}}
 {(q^2+m^2)((q+P)^2+m^2)}>0
 \end{aligned}
 \label{eq:OQED-bubble-depth-derivatives}
\end{equation}
for every integer $n\geq0$. For $a\geq a_->0$,
$0\leq M\leq M_+$, and any real $P$,
\begin{equation}
 |\partial_a^nB_{a,M,m}(P)|\leq
 c_{n,a_-}B_{a_-/2,M_+,m}(0),
 \label{eq:OQED-bubble-depth-uniform}
\end{equation}
where $c_{0,a_-}=1$ and
$c_{n,a_-}=(2n/(\mathrm e\,a_-))^n$ for $n\geq1$.
The bubble is holomorphic in the \emph{depth} $a$ on
$\Re a>0$ at fixed real $P$. For fixed $a>0$ it is
nondecreasing in the auxiliary mass $M$ (strictly for
$P$ in the stated real domain). None of these statements
implies analyticity in complex external momentum when $M=0$,
or reflection positivity of a physical two-point function.
\end{proposition}
\begin{proof}
The two line multipliers equal $e^{-aH_{M,P}}$ with
$H_{M,P}\geq0$ and $H_{M,P}>0$ almost everywhere. On
$\Re a\geq a_->0$, derivatives are dominated because
$H^n e^{-aH}\leq c_{n,a_-}e^{-a_-H/2}$.
The resulting positive integrand is integrable; its integral
is bounded by $c_{n,a_-}B_{a_-/2,M,m}(P)$.
Cauchy--Schwarz for the real even line covariance gives
$B_{a_-/2,M,m}(P)\leq B_{a_-/2,M,m}(0)$.
Since $\partial_Mh_M(u)<0$ for $u>0$, the latter is at
most $B_{a_-/2,M_+,m}(0)$. The same majorant proves
holomorphy in the stated depth half-plane. Positivity
of the derivatives and monotonicity in $M$ follow directly
from the respective pointwise inequalities.
\end{proof}

\begin{theorem}[Controlled crossover for $aM$ fixed]
\label{thm:OQED-strong-depth-crossover}
Fix $m>0$ and $d\geq0$, and consider $M=\mu/a$ with fixed
$\mu\geq0$ as $a\to\infty$. Define
\begin{equation}
 \mathcal H_d(\mu):=\int_0^\infty v^{d+3}
                  e^{-\sqrt{v^2+\mu^2}+\mu}\dd v.
 \label{eq:OQED-crossover-H}
\end{equation}
Then, uniformly for $\mu$ in compact subsets of $[0,\infty)$,
\begin{equation}
 \begin{aligned}
 \lim_{a\to\infty}a^{d+4}\mathscr Q_d(a,\mu/a,m)
    &=\frac{\mathcal H_d(\mu)}{8\pi^2m^2},\\
 \lim_{a\to\infty}a^4B_{a,\mu/a,m}(0)
    &=\frac{3+6\mu+4\mu^2}{64\pi^2m^4}.
 \end{aligned}
 \label{eq:OQED-strong-depth-crossover}
\end{equation}
Here $B_{a,M,m}$ is the literal real Euclidean comparison
\eqref{eq:OQED-literal-depth-bubble}. Furthermore $\mathcal H_0(\mu)=6+6\mu+2\mu^2$,
$\mathcal H_d(0)=\Gamma(d+4)$, and
\begin{equation}
 \mathcal H_d(\mu)\sim\tfrac12(2\mu)^{(d+4)/2}
          \Gamma((d+4)/2)\qquad(\mu\to\infty).
 \label{eq:OQED-large-mu-crossover}
\end{equation}
The last limit concerns the crossover coefficient; it is not a
uniform joint $a,M$ asymptotic over unbounded $aM$.
\end{theorem}
\begin{proof}
Set $r=v/a$ and $M=\mu/a$. The exponent becomes
$\sqrt{v^2+\mu^2}-\mu$, and the denominator tends to $m^2$,
respectively $m^4$. For $\mu$ bounded above by $\mu_+$,
$e^{\mu-\sqrt{v^2+\mu^2}}\leq e^{\mu_+}e^{-v}$,
which supplies a common integrable radial majorant.
The limiting bubble integral is
\[
 \frac1{8\pi^2m^4}\int_0^\infty v^3
  e^{-2(\sqrt{v^2+\mu^2}-\mu)}\dd v.
\]
With $t=\sqrt{v^2+\mu^2}-\mu$ one has
$v^3\dd v=(t^3+3\mu t^2+2\mu^2t)\dd t$;
its integral equals $3/8+3\mu/4+\mu^2/2$.
The same substitution with $e^{-t}$ gives $\mathcal H_0$.
Finally rescale $v=\sqrt{2\mu}\,z$ in $\mathcal H_d$;
the exponent tends to $-z^2$ and is dominated by a function
proportional to $e^{-c\min(z^2,z)}$ for $\mu$ bounded below.
This yields \eqref{eq:OQED-large-mu-crossover}.
\end{proof}

\subsection{Fixed-graph application and the remaining vertex condition}
\label{sec:OQED-graph-application}

In the flat coefficient prototype, with this representative assigned to
the designated photon and fermion lines, use
\begin{equation}
\begin{aligned}
 S_{a,M}(q)&=F_{a,M}(q^2)
                 \frac{-\ii\slashed q+m_e}{q^2+m_e^2},\\
 D_{a,M,E}^{\mu\nu}(q)&=
           \frac{F_{a,M}(q^2)}{q^2}\delta^{\mu\nu}
                 \quad(\alpha_{\rm gf}=1).
\end{aligned}
\label{eq:OQED-Euclidean-lines}
\end{equation}
These are auxiliary uncut coefficients, not a redefinition of the
source-normalized physical two-point functions. There is one kernel
per designated uncut line. If dressed lines are used, no additional
copy of their product is inserted into the graph. The same rule applies
to contacts and to graphs with no internal photon. The Lorentzian
source convention, at finite boundary magnitude before coefficient
extraction and the admitted common limit, retains
\begin{equation}
\begin{aligned}
 z_{\eta,\bb}(q)&=q^2+\ii\eta\epsvec,\\
 F^{\rm dir}_{a,M}(q)&=
    \ExpDir\!\left[aM-a\SqrtDir(M^2-z_{\eta,\bb}(q))\right],\\
 \vec S(q)&=\eta\ii\widehat u_q\,
    \embed\!\left[(\slashed q+m_e)
      \embed^{-1}\!\left(\frac{F^{\rm dir}_{a,M}(q)}
                               {z_{\eta,\bb}(q)-m_e^2}\right)\right],\\
 \vec D^{\mu\nu}(q)&=-\eta\ii\widehat u_q\,
       \frac{F^{\rm dir}_{a,M}(q)}{z_{\eta,\bb}(q)}g^{\mu\nu}.
\end{aligned}
\label{eq:OQED-directed-lines}
\end{equation}
The photon minus sign is the convention of the source's Maxwell action
and Eq.~(53); it is not erased by using an absolute Euclidean estimate.
The spinor embedding is the entrywise convention of
\eqref{eq:physical-directed-Dirac-line}. These remain resolved
scalar-pole templates, with the same finite-regulator and source-exchange
qualifications as Sections~\ref{sec:carrier} and~\ref{sec:LSZ}.

\begin{theorem}[O--QED application with explicit kernel hypotheses]
\label{thm:OQED-fixed-graph}
Fix the support, branch, source/endpoint conventions, and positive
line depths of the O--QED record. Suppose a finite completed loop graph
has full-rank flat routing and the infrared hypotheses of
Theorem~\ref{thm:global-uv}. On a full-rank set of loop lines use
\(F_{a_i,M_i}\), with \(a_i>0\), \(M_i\geq0\), and bound all
other multipliers as in Corollary~\ref{cor:damped-basis}.
If the full completed numerator has a finite polynomial bound of degree
\(D_G\), the graph is absolutely Euclidean convergent, uniformly on
the admitted compact external set. The kernel part of this assertion
is established by \eqref{eq:OQED-inverse-moments}, not left as an
unevaluated assumption.

Alternatively, assume bounded covariance envelopes and a complete
numerator bound
\begin{equation}
 |N_G(k,p)|\leq C(1+\norm{k}^{D_G})
          \exp\!\left(\sum_i v_i|q_i|\right),\qquad v_i\geq0.
\label{eq:OQED-vertex-budget}
\end{equation}
In this second alternative, every exponentially bounded line has its
specified \(a_i>0\); for a merely bounded unselected multiplier set
\(a_i=v_i=0\) in the bound, without changing its model definition.
It is sufficient that the signed residuals \(h_i=a_i-v_i\) obey
\eqref{eq:signed-linear-coercivity}; a positive right side in
\eqref{eq:signed-linear-basis-certificate} is one explicit certificate.
The simpler nonnegative-residual, full-loop-rank test remains sufficient.
These bounds concern the product of all vertices and numerators in the
graph, including contact vertices. The constants \(e^{a_iM_i}\) belong
to the envelope, not to a renormalization of the kernel measure.

If either condition, or the verified forest alternative, holds for
all graphs and insertions of one fixed completed action, and the
independent anomaly-free BRST hypotheses hold for the same action in
both characters, its admitted coefficients are graphwise UV finite and
satisfy the identities of Section~\ref{sec:brst} at every finite order.
This last implication does not certify an unspecified completion.
\end{theorem}

\begin{proof}
For each selected line the required moment is finite:
\[
 a_i^{-2-D_G}M_{1+D_G/2}[\rho_{a_iM_i}]<\infty.
\]
Corollary~\ref{cor:damped-basis} therefore applies directly. In the second case
\eqref{eq:OQED-linear-exponential-bounds} leaves
\(\exp[-\sum_i(a_i-v_i)|q_i|]\) times a fixed constant
\(\exp(\sum_i a_iM_i)\). The linear residual-coercivity proof then
supplies an integrable polynomial times \(e^{-c\norm{k}}\).
The character norm transfers the two coefficient estimates to the
carrier graph. The finite-order and BRST conclusions use their respective
corollary and theorem with the same complete graph data.
\end{proof}

For the conventional local trilinear comparison vertices, the fermion
numerators alone give \(D_G\leq I_f\), with \(I_f\) the number
of internal fermion lines. Additional polynomial vertex derivatives
increase that bound by their total degree. The polarization bubble,
electron self-energy, electron--photon vertex, and four-photon fermion
box have respectively the sufficient benchmark degrees
\(2,1,2,4\), hence moment orders \(2,3/2,2,3\).
All four required moments are available. Together with the local
infrared analysis below, this proves convergence of these specified
comparison integrands; it does not claim that unchanged local vertices
satisfy the dressed Ward identity. The massless photon prevents a
direct invocation of the globally bounded-covariance hypothesis until
that separate infrared analysis is supplied.

\paragraph{Massless photons and infrared control}
The physical photon denominator in \eqref{eq:OQED-Euclidean-lines}
is massless. A positive-mass comparison is not substituted for it inside
a Ward identity. In the vacuum one-loop self-energy and vertex graphs
with \(m_e>0\), choose the photon momentum as the loop variable.
At its origin the only photon singularity is \(|q|^{-2}\), locally
integrable with the four-dimensional measure; the remaining massive
fermion factors are bounded there on compact real external sets.
The polarization and fermion box have only massive internal fermions.
Their polynomial-vertex benchmark integrals are therefore finite
without adding a photon mass. This off-shell Euclidean statement does
not remove on-shell soft-photon singularities or prove uniform
infrared bounds at higher orders.

A useful extension avoids assuming individually bounded denominators.
Let \((\Xi,\nu)\) be the positive measure space of all loop or spectral
variables, and let \(R:\Xi\to[0,\infty)\) be a measurable radial size
in reference momentum units. Let the nonnegative measurable function
\(W_G\) contain the rational, matrix-norm, and polynomial factors left
after extracting a residual exponential. Suppose, with constants
\(C,C_0,c>0\), \(N\geq0\) uniform on the external domain,
\begin{equation}
\begin{aligned}
 |I_G(\xi,p)|&\leq C e^{-cR(\xi)}W_G(\xi,p),\quad c>0,\\
 \int_{R(\xi)\leq R_0}W_G(\xi,p)\dd\nu(\xi)
    &\leq C_0(1+R_0)^N,\qquad R_0\geq1.
\end{aligned}
\label{eq:OQED-spectral-volume-certificate}
\end{equation}
Then the graph is absolutely convergent: splitting into
\(n\leq R<n+1\) bounds it by
\(CC_0\sum_{n\geq0}e^{-cn}(n+2)^N<\infty\).
The series beginning at \(n=\lfloor R_*\rfloor\) also bounds the
numerical tail beyond \(R_*\). Local infrared integrability is part of the
second inequality, not inferred from exponential UV suppression.
This also gives an explicit sufficient extension to a specified
thermal integration/sum or Landau-level representation, provided its
multiplicities, transition matrix elements, and completed vertex growth
are included in \(W_G\). It does not infer those estimates from the
flat vacuum propagator.

\subsection{An evaluated one-loop certificate and its Ward limitation}
\label{sec:OQED-one-loop}

The exact one-loop integrability condition can be evaluated without
special functions. For every \(p>0\), the substitution
\(y=\sqrt{u+M^2}-M\), \(u=y^2+2My\), gives
\begin{equation}
 \int_0^\infty F_{a,M}(u)^p\dd u
   =2\int_0^\infty(y+M)e^{-pa y}\dd y
   =\frac{2}{p^2a^2}+\frac{2M}{pa}.
\label{eq:OQED-kernel-Lp}
\end{equation}
Consequently \(\mathcal Q_{F_{a,M}}=1/(2a^2)+M/a\) and
Proposition~\ref{prop:QED-general-obstruction} applies to the source's
vacuum prototype with bare coefficient vertices. In particular,
\begin{equation}
\begin{aligned}
 \kappa_{a,M,m_e}
 &=\frac{e^2}{8\pi^2}\int_0^\infty
  \frac{u(u+2m_e^2)}{(u+m_e^2)^2}
  e^{-2a(\sqrt{u+M^2}-M)}\dd u,\\
 0<\kappa_{a,M,m_e}
 &\leq\frac{e^2}{16\pi^2a^2}(1+2aM)\quad(e\ne0).
\end{aligned}
\label{eq:OQED-explicit-Ward-obstruction}
\end{equation}
Equality holds for \(m_e=0\) and is strict for \(m_e>0\).
Indeed, the rational factor is
\(1-m_e^4/(u+m_e^2)^2\), between zero and one. Thus
\(\Pi_{\rm bv}^{\mu\nu}(0)=\kappa_{a,M,m_e}\delta^{\mu\nu}\)
and \(P_\mu\Pi_{\rm bv}^{\mu\nu}(P)
=\kappa_{a,M,m_e}P^\nu+o(|P|)\). The finite bare comparison
is not transverse. Neither normalization at \(u=0\) nor choosing
\(M>0\) removes this obstruction.

\begin{center}
\small
\begin{tabular}{@{}crr@{}}
\toprule
\(aM\) & \(a^2\mathcal Q_F\) & \(\kappa/e^2\), \(a=m_e=1\)\\
\midrule
0 & 0.5 & 0.003774745738\\
0.5 & 1.0 & 0.008196730399\\
1 & 1.5 & 0.013357748315\\
2 & 2.5 & 0.024602484262\\
\bottomrule
\end{tabular}
\end{center}
These are dimensionless coefficient benchmarks in one reference
momentum unit, not calibrated plasma inputs or physical photon masses.
The tabulated values were checked using both the positive integral and
its integration-by-parts form
\eqref{eq:general-Ward-obstruction-derivative}. Values with \(M\leq m_e\)
are valid Euclidean comparisons, not certificates of an isolated
Lorentzian shell below the kernel branch point.

The normalized massive kernel also makes the distinction between
regularity and ultraviolet growth concrete. For
\(A(u)=F_{a,M}(u)^{-1}\),
\begin{equation}
 A'(u)=\frac{a}{2\sqrt{u+M^2}}A(u),\qquad M>0.
\label{eq:OQED-inverse-derivative}
\end{equation}
Hence the Ball--Chiu-type vertex of
\eqref{eq:longitudinal-vertex}, with \(B=m_eA\), has regular
divided differences at the origin and diagonal value
\(\Gamma^\mu(0,0)=\ii\gamma^\mu\). This removes the specific
massive zero-momentum cusp obstruction of the \(M=0\) kernel.
It does not bound the vertex polynomially at large momentum:
\(A(R^2)\sim e^{a(R-M)}\).

\begin{proposition}[A Ward-compatible two-vertex bubble fails the graphwise test]
\label{prop:OQED-Ward-bubble-obstruction}
Use the flat vacuum coefficient \(S_{a,M}=F_{a,M}S_0\),
\(a>0\), \(m_e>0\), \(M\geq0\), and \(e\ne0\).
Assume the vertex satisfies the dressed Ward identity and is continuous
at \((p',p)=(k,k)\) for every \(k\ne0\). Consider precisely the
two-vertex term of the source's polarization skeleton at \(P=0\),
without its contact contribution. Its trace over external Euclidean
indices is not absolutely integrable. More precisely, if
\[
 I_W^{\mu\nu}(k):=-e^2\Tr\bigl[
       \Gamma^\mu(k,k)S_{a,M}(k)
       \Gamma^\nu(k,k)S_{a,M}(k)\bigr],
\]
then, for \(u=k^2>0\), \(\Omega=\sqrt{u+M^2}\),
\begin{equation}
\begin{aligned}
 \sum_{\mu=1}^4I_W^{\mu\mu}(k)=e^2\biggl[
 &\frac{8(u+2m_e^2)}{(u+m_e^2)^2}
 -\frac{8a u}{\Omega(u+m_e^2)}\\
 &-\frac{4a^2u}{u+M^2}\biggr]
 \longrightarrow-4e^2a^2 .
\end{aligned}
\label{eq:OQED-Ward-bubble-trace}
\end{equation}
Thus this particular two-vertex graph cannot meet a graphwise
absolute-convergence theorem on an external set containing zero.
\end{proposition}

\begin{proof}
The soft Ward identity \eqref{eq:soft-Ward-vertex} fixes the diagonal
vertex; a continuous transverse addition has zero diagonal limit.
Multiplying by the specified propagator gives the exact cancellation
\begin{equation}
 \Gamma^\mu(k,k)S_{a,M}(k)
       =\ii\gamma^\mu S_0(k)+\frac{a k^\mu}{\Omega}\one.
\label{eq:OQED-Ward-line-cancellation}
\end{equation}
Let \(G_0^\mu=\ii\gamma^\mu S_0\). The Euclidean Clifford trace
identities give
\[
 \sum_\mu\Tr[G_0^\mu G_0^\mu]
     =-\frac{8(u+2m_e^2)}{(u+m_e^2)^2},\qquad
 \Tr G_0^\mu=\frac{4k^\mu}{u+m_e^2}.
\]
The cross terms in the squared product contribute
\(8a u/[\Omega(u+m_e^2)]\), and its scalar-square terms contribute
\(4a^2u/(u+M^2)\). Multiplication by \(-e^2\) proves
\eqref{eq:OQED-Ward-bubble-trace}. Its nonzero asymptote implies
failure of absolute integrability of at least one tensor component.
For an explicit radial exhaustion,
\begin{equation}
 \int_{|k|<R}\frac{\dd^4k}{(2\pi)^4}
       \sum_\mu I_W^{\mu\mu}(k)
       =-\frac{e^2a^2}{8\pi^2}R^4+O(R^3).
\label{eq:OQED-Ward-bubble-exhaustion}
\end{equation}
This radius only displays the asymptotic failure; it is not a new
physical cutoff. The locally integrable \(M=0\) endpoint does not
affect the ultraviolet conclusion.
\end{proof}

\begin{corollary}[The obstruction persists at nonzero transfer]
\label{cor:OQED-fixed-transfer-obstruction}
Retain \(S_{a,M}=F_{a,M}S_0\), \(a>0\), \(m_e>0\),
\(M\geq0\), and \(e\ne0\), but no longer assume diagonal
continuity. Fix real \(P\ne0\), let \(k'=k+P\), and require both
vertices of the two-vertex bubble to satisfy
\(P\cdot\Gamma=S_{a,M}^{-1}(k')-S_{a,M}^{-1}(k)\) in its common
external orientation. Define its complete two-vertex coefficient by
\[
 \begin{aligned}
 I_W^{\mu\nu}(k,P):=-e^2\Tr\bigl[&\Gamma^\mu(k',k)S_{a,M}(k)\\
                    &\Gamma^\nu(k,k')S_{a,M}(k')\bigr].
 \end{aligned}
\]
This graph is not componentwise absolutely integrable. Specifically,
\begin{equation}
\begin{aligned}
 P_\mu P_\nu I_W^{\mu\nu}(k,P)
 =-e^2\Tr\bigl[&S_{a,M}^{-1}(k')S_{a,M}(k)\\
               &+S_{a,M}^{-1}(k)S_{a,M}(k')-2\one\bigr],
\end{aligned}
\label{eq:OQED-double-Ward-contraction}
\end{equation}
and along \(k=rn\), \(|n|=1\), \(r\to\infty\),
\begin{equation}
 P_\mu P_\nu I_W^{\mu\nu}(rn,P)
       \longrightarrow-8e^2\bigl[\cosh(a n\cdot P)-1\bigr].
\label{eq:OQED-fixed-transfer-obstruction}
\end{equation}
\end{corollary}

\begin{proof}
Use the two vertices \(\Gamma^\mu(k',k)\) and
\(\Gamma^\nu(k,k')\). Their Ward identities have the same sign when
both are contracted with \(P\); the second follows by reversing its
own transfer. Expanding the product of the two inverse differences
and using the ordered inverse identities gives
\eqref{eq:OQED-double-Ward-contraction}.
For a finite-momentum check put \(u=k^2\), \(u'=k'^2\),
\(b=k\cdot k'+m_e^2\), and
\(d_\Omega=\sqrt{u'+M^2}-\sqrt{u+M^2}\). Direct Clifford traces give
\begin{equation}
\begin{split}
 P_\mu P_\nu I_W^{\mu\nu}(k,P)
 =-4e^2\biggl[ b\biggl(&\frac{e^{a d_\Omega}}{u+m_e^2}
                  +\frac{e^{-a d_\Omega}}{u'+m_e^2}\biggr)-2\biggr].
\end{split}
\label{eq:OQED-double-Ward-finite-trace}
\end{equation}
No factors \(F\) are discarded in this identity. Along \(k=rn\),
\(d_\Omega=n\cdot P+O(r^{-1})\) uniformly on the unit sphere,
and \(b/(u+m_e^2)=1+O(r^{-1})\), as does its primed partner.
This proves the stated limit with a uniform \(O(r^{-1})\) error. It is strictly negative
on an angular set of nonzero measure, so its absolute radial integral
diverges. Absolute integrability of every tensor component would imply
integrability of this fixed contraction, a contradiction. Transverse
vertex additions do not change either Ward contraction.
\end{proof}

This is stronger than noting that the vertices grow. It is an explicit
failure for the specified vacuum prototype and the two-vertex term in
the source's conditional skeleton. Satisfying these Ward contractions
alone does not establish that the vertices are functional derivatives
of one completed action. It does not compute the contact contribution,
prove divergence of its sum with that contribution, or exclude a
physical completion with a different admitted source/field realization.
If finiteness relies on combining the bubble and contacts, that combined
integrand must be formed first and supplied with its own absolute
majorant. A cancellation between separately undefined integrals is not
a graphwise certificate. Nor does the Ward identity determine every
transverse or contact term. The source already leaves those objects
in its completion record; the present calculation identifies a specific
requirement they must address.

\subsection{Contact completion must be tested before integration}
\label{sec:OQED-contact-test}

Ref.~\cite{ChiapponeIgnition}, Eq.~(78), explicitly retains a contact
sector but does not evaluate it. We first state the mathematical
requirement for combining it with the bubble. A vanishing integrated
Ward contraction is not, by itself, a bound on the absolute integrand.

\begin{lemma}[Finite common-routing sums]
\label{lem:OQED-grouped-convergence}
At a fixed perturbative order, suppose a finite set of action-generated
terms has been placed on one measure space \((\Xi,\nu)\), with all
routing Jacobians, signs, and multiplicities retained. Let
\[
 J_{\mathcal B}(\xi,p)=\sum_{j\in\mathcal B}c_jJ_j(\xi,p)
\]
be their pointwise sum. Suppose the chosen common finite-domain
prescription represents that action coefficient, and there is an
integrable envelope \(M_{\mathcal B}\) such that
\[
 |J_{\mathcal B}(\xi,p)|\leq M_{\mathcal B}(\xi),\qquad p\in\mathcal K.
\]
Then its common-domain limit exists absolutely and is independent of
exhaustion. If the same property holds locally on a complex domain
and the summed integrand is holomorphic there almost everywhere, the
summed coefficient is holomorphic. These conclusions do not assert
absolute convergence of the separate terms.
\end{lemma}
\begin{proof}
Integrate the sum before taking the exhaustion limit. Dominated
convergence gives the absolute limit and its uniform tail bound; the
complex conclusion follows from Theorem~\ref{thm:holomorphy}. No exchange
with separately divergent integrals is used.
\end{proof}

This is a criterion for a declared finite sum, not an application of
Theorem~\ref{thm:global-uv} to a divergent constituent. Separate changes
of loop routing inside divergent terms are not licensed: any boundary
terms induced by such changes must be controlled in the same common
domain. All subintegration regions, not only simultaneous large
momenta, must satisfy the envelope at higher loops. At one loop a
local integrable envelope together with
\(\lVert J_{\mathcal B}(k,p)\rVert\leq C(1+|k|)^{-4-\delta}\),
\(\delta>0\), outside a fixed ball is one simple sufficient condition.
Cancellation after angular averaging alone is not an absolute bound.

For example, if the contact contribution to the double contraction
has a continuous uniform angular limit at \(k=rn\), absolute
integrability of the complete one-loop sum forces that limit to be
\[
 \lim_{r\to\infty}P_\mu P_\nu I_{\rm contact}^{\mu\nu}(rn,P)
       =8e^2[\cosh(a n\cdot P)-1].
\]
Otherwise a nonzero limiting sum on an open angular set would reproduce
the divergent radial integral in
Corollary~\ref{cor:OQED-fixed-transfer-obstruction}. This cancellation
of the leading term is necessary under the stated limit assumption,
not sufficient: slower-decaying remainders can still be nonintegrable.

\paragraph{A separately specified contact test}
To examine that distinction explicitly, impose an additional assumption
not supplied by the PDF: in a flat vacuum coefficient calculation with
field-independent measure, a constant Abelian background enters the
quadratic fermion symbol only by the shift
\[
 \mathcal D[\mathcal A](k)=\mathcal D(k+e\mathcal A),\qquad
 \mathcal D(k)=F_{a,M}(k^2)^{-1}(\ii\slashed k+m_e).
\]
Here \(\mathcal A\) is the constant background on the noncompact flat
chart, and the two external insertions carry zero momentum. This
specifies that constant-background response, not the nonconstant-field
covariantization or its full transverse vertices. Fix \(a,M,m_e>0\) and \(e\ne0\).
Using the same bubble sign as \eqref{eq:bare-QED-loop}, the corresponding
contact contribution from a second variation of the trace logarithm is
\(e^2\Tr[\mathcal D_{,\mu\nu}S]\), with \(S=\mathcal D^{-1}\)
and commas denoting momentum derivatives. No second multiplier is
inserted on this already dressed line. This definition fixes the
relative bubble/contact sign for the following auxiliary test; it is
not inferred from the Ward identity alone.

\begin{proposition}[A contact cancellation that is not absolute convergence]
\label{prop:OQED-constant-contact}
In this constant-background realization the combined coefficient is
\begin{equation}
\begin{aligned}
 J_{\rm cb}^{\mu\nu}(k)
 &=e^2\Tr[\mathcal D_{,\mu\nu}S
           -\mathcal D_{,\mu}S\mathcal D_{,\nu}S]\\
 &=e^2\partial_\mu\partial_\nu\log\det\mathcal D(k).
\end{aligned}
\label{eq:OQED-contact-Hessian}
\end{equation}
It has the exact tensor form, with \(u=k^2\), \(\Omega=\sqrt{u+M^2}\),
\begin{equation}
\begin{split}
 \frac{J_{\rm cb}^{\mu\nu}}{e^2}
 =4a\left(\frac{\delta^{\mu\nu}}{\Omega}
               -\frac{k^\mu k^\nu}{\Omega^3}\right)
   +\frac{4\delta^{\mu\nu}}{u+m_e^2}
               -\frac{8k^\mu k^\nu}{(u+m_e^2)^2}.
\end{split}
\label{eq:OQED-contact-explicit}
\end{equation}
The leading constant of the two-vertex term cancels, but the summed
tensor is still not componentwise absolutely integrable. In fact,
\begin{equation}
\begin{split}
 \sum_\mu J_{\rm cb}^{\mu\mu}(k)
 =e^2\left[
 \frac{4a(3u+4M^2)}{(u+M^2)^{3/2}}
       +\frac{8(u+2m_e^2)}{(u+m_e^2)^2}\right]>0,
\end{split}
\label{eq:OQED-contact-trace}
\end{equation}
and its radial exhaustion is exactly
\begin{equation}
 \int_{|k|<R}\frac{\dd^4k}{(2\pi)^4}
            \sum_\mu J_{\rm cb}^{\mu\mu}(k)
  =\frac{e^2}{2\pi^2}\left[
        \frac{aR^4}{\sqrt{R^2+M^2}}+\frac{R^4}{R^2+m_e^2}\right].
\label{eq:OQED-contact-exhaustion}
\end{equation}
\end{proposition}
\begin{proof}
Differentiate \(S\mathcal D=\one\), so
\(\partial_\nu S=-S\mathcal D_{,\nu}S\), and use cyclicity of the
finite matrix trace. This proves \eqref{eq:OQED-contact-Hessian}.
In four Dirac components,
\[
 \log\det\mathcal D(k)
   =4a(\sqrt{u+M^2}-M)+2\log[(u+m_e^2)/p_{\rm ref}^2]
\]
up to a constant. The fixed momentum unit \(p_{\rm ref}>0\) only makes
the logarithm dimensionless; all derivatives are independent of it.
Twice differentiating gives \eqref{eq:OQED-contact-explicit}, whose
trace is \eqref{eq:OQED-contact-trace}. The divergence theorem on a
finite ball, with sphere area \(2\pi^2R^3\), gives
\eqref{eq:OQED-contact-exhaustion}. Its leading term is
\(e^2aR^3/(2\pi^2)\), proving the asserted failure. The determinant
here is only a finite matrix identity at fixed momentum; no determinant
of an unbounded field operator is being defined by it.
\end{proof}

Subtracting the separately typed ordinary symbol
\(\mathcal D_{\rm ref}=\ii\slashed k+m_e\) in this same raw test
removes the second term of \eqref{eq:OQED-contact-trace}, not the
nonintegrable first term. Nor may the integral be set to zero merely
because \eqref{eq:OQED-contact-Hessian} is a total derivative: its
surface term is displayed explicitly. A symmetry-preserving prescription
for an otherwise divergent functional is a different construction from
an absolutely convergent Lebesgue integral. This calculation neither
assigns a physical photon mass nor changes the fixed-depth kernel rule.
It tests the additional constant-background, trace-log, and measure
assumptions just stated. Additional action terms, measure determinants,
or source maps need their own derivation and common-integrand estimate;
they are not silently excluded or supplied by this test.

\begin{corollary}[Changing the radial damper does not repair this raw contact test]
\label{cor:OQED-general-contact-obstruction}
Retain the constant-background shift, trace-log sign, one-species
normalization, and field-independent measure of
Proposition~\ref{prop:OQED-constant-contact}. Replace \(F_{a,M}\) by
any positive radial multiplier \(F(u)\) that is \(C^2\) for
\(u>u_0>0\) and satisfies \(F(u)\to0\) as \(u\to\infty\).
Set \(h_F=-\log F\), keep \(m_e>0\) and \(e\ne0\), and subtract
the ordinary symbol's contact-plus-bubble coefficient \emph{before}
integration. Then the raw difference is not componentwise absolutely
integrable on the exterior of a ball. No monotonicity or rate of decay
of \(F\) is needed for this assertion.
\end{corollary}
\begin{proof}
At fixed momentum, \(\mathcal D_F=F^{-1}\mathcal D_{\rm ref}\)
has four Dirac components. Subtracting its finite matrix logarithmic
derivatives therefore gives exactly
\begin{equation}
 \Delta J_F^{\mu\nu}(k)
     =4e^2\partial_\mu\partial_\nu h_F(k^2).
 \label{eq:OQED-general-contact-difference}
\end{equation}
Choose \(R_0^2>u_0\). On each finite annulus the divergence theorem
and the retained Fourier normalization give
\begin{equation}
\begin{split}
 \int_{R_0<|k|<R}\frac{\dd^4k}{(2\pi)^4}
             \sum_\mu\Delta J_F^{\mu\mu}(k)
 =\frac{e^2}{\pi^2}\bigl[&R^4h_F'(R^2)\\
                      &-R_0^4h_F'(R_0^2)\bigr].
\end{split}
 \label{eq:OQED-general-contact-flux}
\end{equation}
Indeed, the outward radial derivative of \(4e^2h_F(k^2)\) is
\(8e^2R h_F'(R^2)\), and the sphere area is \(2\pi^2R^3\).
If every tensor component were absolutely integrable outside \(R_0\),
the trace integrals on the left would be bounded independently of \(R\).
Hence \(|u^2h_F'(u)|\leq C\) for all sufficiently large \(u\).
It would follow that \(\int^\infty |h_F'(u)|\dd u<\infty\), so
\(h_F\) had a finite limit. This contradicts
\(h_F(u)=-\log F(u)\to+\infty\).
\end{proof}

The distinction between trace cancellation and absolute tensor control
can also be read off pointwise. In the Euclidean spectral matrix norm,
define, for real \(u>u_0\),
\[
 \mathcal H_F(u):=\max\{|h_F'(u)|,
                         |h_F'(u)+2u h_F''(u)|\}.
\]
The radial and three tangential eigenvalues of
\(\Delta J_F=8e^2h_F'I+16e^2h_F''kk^T\) give the exact identity
\begin{equation}
 \int_{|k|>R_0}\frac{\dd^4k}{(2\pi)^4}
            \|\Delta J_F(k)\|_2
   =\frac{e^2}{2\pi^2}\int_{R_0^2}^\infty u\mathcal H_F(u)\dd u.
 \label{eq:OQED-contact-norm-test}
\end{equation}
This identity holds with either side infinite. Finite-dimensional norm
equivalence makes finiteness equivalent to componentwise absolute
integrability. Since \(\mathcal H_F\geq|h_F'|\), it also gives a
direct alternative proof of the obstruction. Conversely, a nonzero
limit of \(F\) alone is not sufficient: for
\(h_F(u)=1-(1+u)^{-1}\), the multiplier tends to \(e^{-1}\)
and the flux in \eqref{eq:OQED-general-contact-flux} stays bounded,
but \(u|h_F'(u)|\sim u^{-1}\) is not integrable. The trace test
must therefore not be reversed into a tensor-convergence criterion.

This result concerns only the raw constant-background realization above,
including its ordinary-symbol subtraction. It is not a statement about
every contact completion, a renormalized response difference, or an
infinite-dimensional determinant. It shows why changing the radial
profile alone cannot supply the missing common-integrand bound while
those assumptions are kept. Field-dependent measure contributions,
different endpoint realizations, or additional action terms must be
constructed and estimated, not presumed absent or presumed to cancel it.

\subsection{Local analytic gain and the directed pole test}
\label{sec:OQED-analytic}

For \(M>0\), the positive exponential moments in
\eqref{eq:OQED-exponential-moment} satisfy the additional kernel
hypothesis of Lemma~\ref{lem:massive-neighbourhood}. Therefore every
flat positive-mass graph with the lemma's polynomial holomorphic
numerator has a holomorphic neighbourhood of each compact real
Euclidean external set. In particular, for \(m_e>0\) the bare
polarization and the literal two-line scalar benchmark with mass \(m_e\)
are locally holomorphic at zero. For the latter, Euclidean rotational
invariance makes its local convergent expansion a function of \(P_E^2\).
The \(M=0\) nonanalyticity of Proposition~\ref{prop:stable-bubble-nonanalytic}
is consequently not a theorem about this \(M>0\) profile. No uniform
analytic neighbourhood as \(M\downarrow0\) is asserted. Massless
photon graphs, nonpolynomial completed vertices, and a medium require
separate common-majorant arguments.

\begin{proposition}[An explicit holomorphic tube for the scalar comparison]
\label{prop:OQED-bubble-tube}
For \(a,M,m>0\), let
\(C(q)=F_{a,M}(q^2)/(q^2+m^2)\), with the principal square-root branch, and set
\[
 \mathcal T_*=\{P\in\CC^4:|\Im P|<2m_*\},\qquad
 m_*:=\min\{m,M\}.
\]
The literal scalar comparison bubble has the holomorphic continuation
\begin{equation}
 B(P)=\int_{\RR^4}\frac{\dd^4k}{(2\pi)^4}
                     C(k-P/2)C(k+P/2),\qquad P\in\mathcal T_*.
\label{eq:OQED-bubble-tube}
\end{equation}
For every \(0\leq d<2m_*\), its integrand obeys the common bound
\begin{equation}
 |C(k-P/2)C(k+P/2)|
 \leq\frac{e^{2aM}e^{-2a|k|}}{(m^2-d^2/4)^2},
 \qquad |\Im P|\leq d.
\label{eq:OQED-bubble-tube-majorant}
\end{equation}
Near zero write \(B(P)=b(P^2)\), with \(P^2=\sum_\mu P_\mu^2\).
Then \(b(w)\) is holomorphic for \(|w|<4m_*^2\). For
\(0<r<\mathcal R<4m_*^2\), its Taylor coefficients \(b_n\) satisfy
\begin{equation}
\begin{aligned}
 |b_n|&\leq B_{\mathcal R}\mathcal R^{-n},\\
 \sup_{|w|\leq r}\left|b(w)-\sum_{n=0}^N b_n w^n\right|
   &\leq B_{\mathcal R}
          \frac{(r/\mathcal R)^{N+1}}{1-r/\mathcal R},\\
 B_{\mathcal R}&:=\frac{3e^{2aM}}
           {64\pi^2a^4(m^2-\mathcal R/4)^2}.
\end{aligned}
\label{eq:OQED-bubble-Taylor-bound}
\end{equation}
\end{proposition}
\begin{proof}
For real \(P\), a real translation of the absolutely convergent
integral gives the symmetric routing in \eqref{eq:OQED-bubble-tube}.
For complex \(P=x+\ii y\), put \(q_\pm=k\pm x/2\pm\ii y/2\).
For \(\Lambda=m\) or \(M\),
\[
 \Re(q_\pm^2+\Lambda^2)
   =|k\pm x/2|^2+\Lambda^2-|y|^2/4>0.
\]
Both denominators are nonzero and both principal square roots are
holomorphic. For \(\Re z>0\),
\(\Re\sqrt z=\sqrt{(|z|+\Re z)/2}\geq\sqrt{\Re z}\).
Thus \(|F_{a,M}(q_\pm^2)|\leq e^{aM}e^{-a|k\pm x/2|}\),
and \(|q_\pm^2+m^2|\geq m^2-d^2/4\). The triangle inequality
\(|k-x/2|+|k+x/2|\geq2|k|\) proves the envelope. Its integral is
\[
 \int\frac{\dd^4k}{(2\pi)^4}e^{-2a|k|}
                =\frac{3}{64\pi^2a^4}.
\]
Theorem~\ref{thm:holomorphy} proves the tube statement. Real rotational
invariance makes the local power series a series in \(P^2\).
Equivalently, the even holomorphic function \(B(te_1)\) for
\(|t|<2m_*\) descends to \(b(w)\), \(w=t^2\), throughout the
stated disk. On \(|w|=\mathcal R\), choose either square root of
\(w\); the bound applies with \(d=\sqrt{\mathcal R}\).
Cauchy's estimate and the geometric-series tail give
\eqref{eq:OQED-bubble-Taylor-bound}. No complex loop-momentum
translation or Lorentzian contour rotation is used.
\end{proof}

This is a sufficient domain and an explicit remainder for the literal
scalar comparison, not a maximal physical sheet. The same two-line
routing with a polynomial numerator admits this tube by an additional
compact-set polynomial majorant, with the stated positive masses.
The result does not cover the nonintegrable Ward-vertex bubble, the
constant-background contact test, or the unevaluated medium response.
It demonstrates concretely how a positive \(M\) changes the local
analytic problem; the guaranteed tube shrinks to zero as \(M\downarrow0\).

\paragraph{Certified scalar truncation and the first Taylor coefficient}
For the same symmetric routing, let \(B^{[R]}(P)\) be
\eqref{eq:OQED-bubble-tube} restricted to \(|k|\leq R\), \(R\geq0\).
The pointwise envelope gives a uniform absolute tail bound throughout
\(|\Im P|\leq d<2m_*\):
\begin{equation}
\begin{aligned}
 |B(P)-B^{[R]}(P)|
 &\leq\frac{e^{2aM-2aR}}{8\pi^2(m^2-d^2/4)^2}\,T_a(R),\\
 T_a(R)&:=\frac{R^3}{2a}+\frac{3R^2}{(2a)^2}
                   +\frac{6R}{(2a)^3}+\frac6{(2a)^4}.
\end{aligned}
 \label{eq:OQED-bubble-radial-tail}
\end{equation}
Four-dimensional polar integration reduces the omitted envelope to
\(\int_R^\infty r^3e^{-2ar}\dd r=e^{-2aR}T_a(R)\).
The radius is a numerical integration parameter, not a new theory scale.
This bound may be added to a separately certified quadrature error on
the retained ball; a floating-point comparison alone is not such a
certificate. Both errors refer to the same routing and complex domain.

The first two coefficients in \(b(w)=b_0+b_1w+\cdots\) also have
one-dimensional representations. Put
\(f(u)=F_{a,M}(u)/(u+m^2)\). Then
\begin{equation}
\begin{aligned}
 b_0&=\frac1{16\pi^2}\int_0^\infty u f(u)^2\dd u,\\
 b_1&=-\frac1{32\pi^2}\int_0^\infty u^2 f'(u)^2\dd u<0,\\
 f'(u)&=-f(u)\left[\frac{a}{2\sqrt{u+M^2}}+\frac1{u+m^2}\right].
\end{aligned}
 \label{eq:OQED-bubble-first-coefficients}
\end{equation}
To prove the derivative identity, differentiate the real bubble at zero
and integrate by parts, using smoothness at the origin and exponential
decay at infinity. Rotational invariance gives
\[
 \partial_{P_i}\partial_{P_j}B(0)
 =-\int\frac{\dd^4k}{(2\pi)^4}\partial_i C(k)\partial_j C(k)
 =2\delta_{ij}b_1.
\]
Since \(\partial_i C=2k_i f'(k^2)\), spherical averaging yields
\eqref{eq:OQED-bubble-first-coefficients}; strict negativity follows
from \(f'\ne0\). This is the Euclidean invariant coefficient. Under
\(w=-s\), its first invariant derivative changes sign. Neither sign
is a physical damping rate or an evaluation of the completed O--QED
polarization. Together with \eqref{eq:OQED-bubble-Taylor-bound}, these
formulas provide controlled approximation data for the stated scalar
comparison without interchanging an unbounded Taylor series with a loop
integral.

\begin{proposition}[Signed scalar approximation bounds]
\label{prop:OQED-scalar-CM-enclosure}
For the scalar comparison with \(a,M,m>0\), the invariant function
\(b(w)\) admits a finite positive Laplace measure \(\nu\) on
\((0,\infty)\):
\begin{equation}
 \begin{aligned}
  b(w)&=\int_0^\infty e^{-\lambda w}\nu(\dd\lambda),\\
  \int_0^\infty e^{s\lambda}\nu(\dd\lambda)&<\infty,
       \qquad 0<s<4m_*^2.
 \end{aligned}
 \label{eq:OQED-scalar-positive-measure}
\end{equation}
Thus \(b\) is holomorphic for \(\Re w>-4m_*^2\) and, on the
real half-line, \((-1)^n b^{(n)}(w)>0\) for every \(n\geq0\).
Its coefficients \(b_n=b^{(n)}(0)/n!\) alternate in sign. Every
finite truncation \(p_N(w)=\sum_{n=0}^N b_nw^n\) satisfies
\begin{equation}
 0\leq(-1)^{N+1}[b(w)-p_N(w)]
       \leq |b_{N+1}|w^{N+1},\qquad w\geq0.
 \label{eq:OQED-scalar-signed-enclosure}
\end{equation}
In particular, odd truncations are lower bounds and even truncations
are upper bounds on the positive Euclidean invariant axis.
\end{proposition}
\begin{proof}
Use the two independent line parameters in
\eqref{eq:two-line-parameters}, with \(w_i=\alpha_i+a^2\tau_i\)
and \(\rho=\rho_{aM}\). Push forward the positive measure
\[
 \frac{e^{-m^2(\alpha_1+\alpha_2)}}{16\pi^2(w_1+w_2)^2}
 \dd\alpha_1\dd\alpha_2\rho(\dd\tau_1)\rho(\dd\tau_2)
\]
by \(\lambda=w_1w_2/(w_1+w_2)>0\). Its mass is the finite
number \(b_0\), and Tonelli proves the first equality for \(w\geq0\).
Since \(\lambda\leq(w_1+w_2)/4\), multiplication by
\(e^{s\lambda}\) is bounded by reducing the Schwinger mass square
to \(m^2-s/4>0\) and tilting each line measure by
\(e^{sa^2\tau/4}\). For \(s<4m_*^2\), the tilt is strictly
below \((aM)^2\); its mass and required inverse moments remain
finite. The existing Gaussian estimate therefore proves the second
line of \eqref{eq:OQED-scalar-positive-measure}. It supplies all
moments of \(\nu\), a common holomorphic majorant on compact
subsets of the half-plane, and
\[
 (-1)^n b^{(n)}(w)=\int\lambda^n e^{-\lambda w}\nu(\dd\lambda).
\]
Strict positivity follows because the measure is nonzero and supported
on positive \(\lambda\). Finally the integral form of the Taylor
remainder of \(e^{-x}\), for \(x\geq0\), has sign
\((-1)^{N+1}\) and modulus at most \(x^{N+1}/(N+1)!\).
Integrating that finite-order inequality proves
\eqref{eq:OQED-scalar-signed-enclosure} without exchanging an
infinite series and an integral.
\end{proof}

This positive measure is a proper-time measure for the auxiliary
\emph{bubble}; it is not a positive physical mass-spectral measure or
an OS reconstruction. Its half-plane continuation agrees with
\eqref{eq:OQED-bubble-tube} wherever the latter is used, by uniqueness
from the real Euclidean region. The signed inequalities apply at each
finite order even outside the guaranteed Taylor disk; they do not
assert convergence of that series there. They supplement, rather than
replace, the complex-disk estimate \eqref{eq:OQED-bubble-Taylor-bound}.
For a useful independent coefficient check, integration by parts and
four-dimensional rotational averaging also give
\begin{equation}
 \begin{aligned}
 b_2&=\frac1{192\pi^2}\int_0^\infty
             u[2f'(u)+uf''(u)]^2\dd u>0,\\
 b_1^2&\leq2b_0b_2 .
 \end{aligned}
 \label{eq:OQED-scalar-second-coefficient}
\end{equation}
Indeed, \(\partial_{P_1}^4B(0)=\int(\partial_1^2C)^2\dd^4k/(2\pi)^4
=24b_2\). For a radial \(C\), its integrated value is one eighth
of \(\int(\Delta C)^2\dd^4k/(2\pi)^4\); here
\(\Delta C=4[2f'+uf'']\). Boundary terms vanish under the same
positive-mass smoothness and exponential bounds. The last inequality
is Cauchy--Schwarz for \(1,\lambda\) in \(L^2(\nu)\).

At the line level the branch in \eqref{eq:OQED-directed-lines} starts
at the timelike invariant \(q^2=M^2\). If \(0<m_e^2<M^2\), the
numerator is holomorphic and nonzero near the electron shell, but its
coefficient there is
\begin{equation}
 f_{a,M}(m_e^2)
    =e^{a(M-\sqrt{M^2-m_e^2})},
\label{eq:OQED-electron-shell-weight}
\end{equation}
not one. The photon shell at zero instead has unit numerator for
\(M>0\). These are local auxiliary pole statements; physical residues
are those of the completed source matrix. If \(m_e^2=M^2\), the pole
coincides with the kernel branch point and the isolated-meromorphic
pole argument is unavailable. For \(m_e^2>M^2\) the two numerator
boundaries are \(e^{aM}e^{\pm\ii a\sqrt{m_e^2-M^2}}\), and the
cut--discontinuity test retains their branch contribution. The source's
kernel-only inequality \(m_e^2<M^2\) does not exclude medium cuts,
vertex singularities, or competing physical poles.

The internal vectorial datum \(\ii\eta\epsvec\) remains separate
from the retarded response limit \(z=\omega+\ii0^+\) after the
physical readout. Positive Laplace weights do not establish that
readout, reflection positivity, or undressed operational-cut matching.
The spectral obstruction in \eqref{eq:spectral-obstruction} continues
to apply to a raw covariance with this exponentially decaying profile.


\subsection{A finite zero-depth limit after common scalar subtraction}
\label{sec:OQED-scalar-subtraction}

The positive-depth graph theorem does \emph{not} extend the
unsubtracted scalar bubble uniformly to $a=0$. Nevertheless, a
particular \emph{difference formed before integration} has a
finite limit. This provides a useful comparison with ordinary
subtracted scalar field theory, not a kernel-off renormalization
of the Ward-completed gauge polarization.

\begin{theorem}[Subtracted literal bubble at vanishing depth]
\label{thm:OQED-subtracted-zero-depth}
Fix $m>0$ and $M\geq0$. For real Euclidean $P$, $w=P^2\geq0$,
define $C_{a,M,m}(q):=F_{a,M}(q^2)/(q^2+m^2)$ and
\begin{equation}
 D_{a,M,m}(w):=B_{a,M,m}(0)-B_{a,M,m}(P).
 \label{eq:OQED-subtracted-bubble-def}
\end{equation}
At every $a>0$ this is an absolutely convergent \emph{common}
integral and satisfies
\begin{equation}
 0\leq D_{a,M,m}(w)\leq\frac{w}{64\pi^2m^2}.
 \label{eq:OQED-subtraction-uniform-bound}
\end{equation}
As $a\downarrow0$ with $M$ fixed,
\begin{equation}
 \boxed{\displaystyle
 \lim_{a\downarrow0}D_{a,M,m}(w)
   =\frac1{16\pi^2}\int_0^1
      \log\left(1+\frac{x(1-x)w}{m^2}\right)\dd x.}
 \label{eq:OQED-subtraction-limit}
\end{equation}
Convergence is uniform when $w$ and $M$ range over compact
subsets of $[0,\infty)$ and $m$ ranges over a compact subset of
$(0,\infty)$. In addition,
\begin{equation}
 \lim_{a\downarrow0}D_{a,M,m}'(0)
  =\frac1{96\pi^2m^2}.
 \label{eq:OQED-subtraction-slope-limit}
\end{equation}
In contrast, the \emph{unsubtracted} comparison has the
universal finite-part asymptotic
\begin{equation}
 B_{a,M,m}(0)=\frac1{8\pi^2}
  \left[\log\frac1{2am}-\gamma_{\mathrm E}-\frac12\right]
  +o(1),\quad a\downarrow0,
 \label{eq:OQED-unsubtracted-log-depth}
\end{equation}
where $\gamma_{\mathrm E}$ is Euler's constant. In this
normalization the constant inside the logarithm is specified,
not absorbed into an arbitrary reference momentum scale.
The comparison makes no assertion about a jointly completed
O--QED kernel-off gauge functional or about uniform Lorentzian
boundary values.
\end{theorem}
\begin{proof}
For $a>0$, $C_{a,M,m}$ is real, even, and square-integrable.
A real translation gives the manifestly positive \emph{single}
representation
\begin{equation}
 D_{a,M,m}(P^2)=\frac12\int\frac{\dd^4k}{(2\pi)^4}
 [C_{a,M,m}(k-P/2)-C_{a,M,m}(k+P/2)]^2.
 \label{eq:OQED-subtraction-difference-square}
\end{equation}
This is an identity between convergent terms at $a>0$; its
right-hand side remains the preferred representation at the limit.

For the bound and the limit, use the positive dimensionful heat-depth
probability law $\rho^{\rm heat}_{a,M}$ for $F_{a,M}$, defined
as the push-forward of $\rho_{aM}$ under $\tau\mapsto a^2\tau$,
so $F_{a,M}(u)=\int e^{-\lambda u}\rho^{\rm heat}_{a,M}
(\dd\lambda)$. Use it in \eqref{eq:two-line-parameters}. Let
$w_i=\alpha_i+\lambda_i$ and $S=\alpha_1+\alpha_2$.
The proper-time coefficient for $D$ has the nonnegative factor
\[
 \frac{e^{-m^2S}}{(w_1+w_2)^2}
 \left[1-e^{-\frac{w_1w_2}{w_1+w_2}w}\right].
\]
Since $\lambda:=w_1w_2/(w_1+w_2)\leq(w_1+w_2)/4$ and
$1-e^{-\lambda w}\leq\lambda w$, this factor is at most
$w e^{-m^2S}/[4(w_1+w_2)]\leq w e^{-m^2S}/(4S)$.
The integral of $e^{-m^2S}/S$ over the positive $\alpha$ quadrant
is $m^{-2}$. Including $1/(16\pi^2)$ proves the bound.

For each fixed $u\geq0$, $F_{a,M}(u)\to1$ as $a\downarrow0$.
The normalized positive heat-depth measures therefore concentrate at
$\lambda_i=0$: for any $s,\delta>0$,
\[
 \rho^{\rm heat}_{a,M}([\delta,\infty))
 \leq\frac{1-F_{a,M}(s)}{1-e^{-s\delta}}\longrightarrow0.
\]
The preceding $\alpha$-integrable upper bound permits first taking
the heat-measure limit for $\alpha_1,\alpha_2>0$ and then using
dominated convergence in those variables. With both heat depths set
to zero, put $\alpha_1=xS$, $\alpha_2=(1-x)S$. The resulting
positive integral is
\[
 \frac1{16\pi^2}\int_0^1\dd x\int_0^\infty
  \frac{e^{-m^2S}}S
       [1-e^{-Sx(1-x)w}]\dd S,
\]
which evaluates to \eqref{eq:OQED-subtraction-limit}. The same
uniform bound applied to finite differences, or the bounded first
$w$-derivative of the proper-time integrand, gives equicontinuity
and hence compact-uniform convergence. The derivative at $w=0$
converges by the identical domination and equals
$[16\pi^2m^2]^{-1}\int_0^1x(1-x)\dd x$.

The argument is uniform on the stated compact $M,m$ sets:
$1-F_{a,M}(s)\leq a\sqrt{s}$ controls the heat-law
concentration independently of $M$, and $m\geq m_->0$ gives
one common integrable Schwinger envelope. The same domination
controls the first $w$-derivative. The explicit finite part
in \eqref{eq:OQED-unsubtracted-log-depth} is proved in
Proposition~\ref{prop:OQED-zero-depth-finite-part} below.
Thus finiteness is achieved only after the \emph{common}
subtraction; the unsubtracted integral remains divergent.
\end{proof}


\begin{proposition}[Universal logarithmic finite part at zero depth]
\label{prop:OQED-zero-depth-finite-part}
For the literal scalar bubble at $P=0$, with $m>0$, $M\geq0$,
put $s=2am>0$ and
$L_{a,m}:=[\log(1/s)-\gamma_{\mathrm E}-1/2]/(8\pi^2)$.
Then
\begin{equation}
 |B_{a,0,m}(0)-L_{a,m}|
 \leq \frac{(2+\pi/4)s}{8\pi^2},
 \label{eq:OQED-unsubtracted-finitepart-error}
\end{equation}
and
\begin{equation}
 0\leq B_{a,M,m}(0)-B_{a,0,m}(0)
 \leq(e^{2aM}-1)B_{a,0,m}(0).
 \label{eq:OQED-kernel-mass-finitepart-error}
\end{equation}
Consequently \eqref{eq:OQED-unsubtracted-log-depth} holds
uniformly for $M$ in bounded nonnegative sets and $m$ in
compact positive intervals. Its divergent coefficient and
finite constant do not depend on $M$ under this limiting
prescription. It is a coefficient-level statement about the
unsubtracted \emph{auxiliary scalar bubble}, not a physical
counterterm or an interacting QED amplitude.
\end{proposition}
\begin{proof}
At $M=0$, radial substitution $r=mx$ gives
\begin{equation}
 8\pi^2B_{a,0,m}(0)
 =\int_0^\infty\frac{x^3e^{-sx}}{(1+x^2)^2}\dd x
 =I(s)-K(s),
 \label{eq:OQED-massless-kernel-log-integrals}
\end{equation}
where
$I(s)=\int_0^\infty e^{-sx}x/(1+x^2)\dd x$ and
$K(s)=\int_0^\infty e^{-sx}x/(1+x^2)^2\dd x$.
Let
$g(x)=x/(1+x^2)-\mathbf1_{[1,\infty)}(x)/x$.
One checks directly that
$\int_0^\infty g(x)\dd x=0$ and
$\int_0^\infty x|g(x)|\dd x=1$. Hence
\[
 I(s)=E_1(s)+\int_0^\infty e^{-sx}g(x)\dd x,
 \qquad
 |I(s)+\gamma_{\mathrm E}+\log s|\leq2s,
\]
using $|E_1(s)+\gamma_{\mathrm E}+\log s|\leq s$.
Also $K(0)=1/2$ and
$|K(s)-1/2|\leq s\int_0^\infty
 x^2(1+x^2)^{-2}\dd x=\pi s/4$.
Together these inequalities establish
\eqref{eq:OQED-unsubtracted-finitepart-error}.

For general $M$, let
$\delta_M(r):=r+M-\sqrt{r^2+M^2}\in[0,M]$.
Then
$e^{-2ah_M(r^2)}=e^{-2ar}e^{2a\delta_M(r)}$.
Multiplying by the positive massive radial weight and
integrating gives \eqref{eq:OQED-kernel-mass-finitepart-error}.
For bounded $M$, its right side is
$O(aM\log(1/a))$ as $a\downarrow0$, uniformly when
$m$ stays in a compact positive interval. This establishes
the stated universal finite part with no subtraction
of two separately undefined integrals.
\end{proof}

\subsection{The kernel-off comparison is not a new UV theorem}
\label{sec:OQED-reference}

Following Ref.~\cite{ChiapponeIgnition}, Eqs.~(80)--(83), retain the
same state, gauge, source normalization, action-completion prescription,
and renormalization/matching record for
\begin{equation}
 \Delta\Pi_{R,\Phi}(a)
       :=\Pi_{R,\Phi}^{\rm matched}(a)-\Pi_{R,\rm ref}^{\rm ren},
       \qquad \Phi_{\rm off}\equiv1.
\label{eq:OQED-matched-increment}
\end{equation}
The reference is separately defined. Finiteness of the fixed-depth
term does not render an unrenormalized kernel-off QED integral finite.
The shared record fixes physical normalization conditions and the
reference prescription; it does not mean that identical divergent bare
counterterms must be added to an already finite graph. Where the two
completed renormalized tensors and their analytic continuations exist,
their Ward identities imply transversality of their difference. This
linearity argument determines neither its sign nor a spatial power.

The optional real interpolation illustrates the boundary of the UV
hypotheses. To avoid confusing its parameter with the sector label,
write
\begin{equation}
 F_t=(1-t)+tF_{a,M},\qquad
 \rho_t=(1-t)\delta_0+t\rho_{aM},\qquad 0\leq t\leq1.
\label{eq:OQED-interpolation-atom}
\end{equation}
For every \(t<1\), its measure has a nonzero atom at zero and
\(F_t(u)\to1-t\) as \(u\to\infty\). Thus the moment theorem
cannot be applied along the whole path merely because it applies at
\(t=1\). For the bare two-line zero-momentum prototype,
\begin{equation}
 \left.\partial_tF_t(u)^2\right|_{t=0}=2(F_{a,M}(u)-1),\qquad
 \partial_t^2F_t(u)^2=2(F_{a,M}(u)-1)^2.
\label{eq:OQED-interpolation-derivatives}
\end{equation}
The two factors tend to \(-2\) and \(2\), respectively. Multiplying
the undamped QED tensor, whose absolute radial tail is proportional
to \(r\,\dd r\) on an angular set of nonzero measure, makes neither
derivative integrable. The unrenormalized difference factor
\(F_{a,M}^2-1\to-1\) has the same obstruction. These examples do
not disprove differentiability of a separately completed, consistently
renormalized response family. They show why the uniform derivative
majorants assumed in the PDF's optional interpolation audit are
additional hypotheses, not consequences of fixed-depth UV convergence.

In contrast, derivatives with respect to the depth on a compact interval
\(0<a_-\leq a\leq a_+\) satisfy
\[
 |\partial_a^nF_{a,M}(u)|=h_M(u)^nF_{a,M}(u)
              \leq u^{n/2}F_{a_-,M}(u).
\]
Every fixed finite number of line derivatives therefore has a polynomial
cost covered by the verified moments, when the remaining completed
vertex derivatives and infrared bounds are supplied. This is a
positive-depth parameter estimate, not a physical zero-depth limit.

\subsection{From the UV certificate to the medium response}
\label{sec:OQED-medium}

The vacuum line \eqref{eq:OQED-Euclidean-lines} carries no density,
temperature, magnetic spectrum, or distribution function. The source
instead specifies a state \(\mathcal S\), its integration/summation
measure \(\dd\mu_E^{\mathcal S}\), and a nonnegative kinetic
operator/eigenvalue \(\mathcal Q_{E,\mathcal S}\). In the scalar
spectral realization used there,
\begin{equation}
 S_{E,\Phi}(Q;\mathcal S)
   =F_{a,M}(\mathcal Q_{E,\mathcal S}(Q))
                     S_{E,\rm ref}(Q;\mathcal S).
\label{eq:OQED-medium-line}
\end{equation}
A noncommuting background realization instead requires its ordered
operator covariance and common domain. One cannot infer scalar
commutation with spin or magnetic projectors from
\eqref{eq:OQED-medium-line}.
With all arguments and index contractions provided by that realization,
the source's one-loop skeleton is
\begin{equation}
\begin{aligned}
 \Pi_{E,\Phi}^{\mu\nu}(P)
  ={}&-e^2\int\dd\mu_E^{\mathcal S}(Q)\,
    \Tr[\Gamma_\Phi^\mu S_{E,\Phi}(Q)
              \Gamma_\Phi^\nu S_{E,\Phi}(Q+P)]\\
   &+\Pi_{E,\rm contact,\Phi}^{\mu\nu}(P).
\end{aligned}
\label{eq:OQED-completed-skeleton}
\end{equation}
For a magnetic spectral basis, the compact \(Q+P\) notation includes
the actual transition kernels and sums, not a scalar substitution for
Landau-level matrix elements. The UV proof applies after those factors
and the contact terms have been expanded at a specified perturbative
order or bounded as a declared skeleton. A skeleton built from dressed
objects does not by itself have one strict perturbative loop order.

Equation~\eqref{eq:OQED-spectral-volume-certificate} states a concrete
sufficient estimate for the resulting measure, including all discrete
multiplicities. KMS or a separately justified Euclidean state is required
for this route. A genuinely nonequilibrium response needs the real-time
functional stated in the source, with its own estimates; the vacuum
Gaussian proof is not automatically a theorem for that functional.
Similarly, the common holomorphic domain and retarded boundary topology
must be proved for the completed increment, not inferred from a
pointwise expansion of the kernel. At even order that expansion must
retain odd--odd products with vertices, propagators, and state factors.

\paragraph{A quantitative transfer bound}
The ultraviolet analysis can supply convergent representations and
controlled numerical tails. Their relevance to the source's RT
application is especially clear once an error bound has been established
on the same complex continuation domain used by its pole calculation.
Let \(X=\Delta\Pi_{R,\Phi}\), \(D_0=D_{\rm EM,0}\),
\(U=V_{\rm RT,EM}\), and \(V=V_{\rm EM,RT}\). First work with
bounded operators on fixed electromagnetic and retained RT Hilbert
spaces: \(D_0\) is boundedly invertible, \(X,Y\) act on the
electromagnetic space, and \(U,V\) have the indicated opposite channel
directions. Assume \(r_X=\|D_0X\|<1\). The source's exact Schur
correction is
\begin{equation}
 \delta K_X=-U(R_X-D_0)V,\qquad
 R_X=(D_0^{-1}-X)^{-1},
\label{eq:OQED-Schur-increment}
\end{equation}
with \(K_0=D_{\rm RT,0}^{-1}\). For another admissible response
\(Y\), with \(r_Y=\|D_0Y\|<1\), the resolvent identity gives
\begin{equation}
\begin{aligned}
 \delta K_X-\delta K_Y&=-U R_X(X-Y)R_Y V,\\
 \|\delta K_X-\delta K_Y\|
 &\leq\frac{\|U\|\|V\|\|D_0\|^2}
                {(1-r_X)(1-r_Y)}\|X-Y\|.
\end{aligned}
\label{eq:OQED-Schur-error-bound}
\end{equation}
No commutation is used: multiplication of both inverse identities
proves the ordered resolvent formula. An unbounded electromagnetic
inverse can also be admitted on a fixed graph-norm domain
\(\mathcal D_{\rm EM}\). Require both \(D_0^{-1}\) and
\(D_0^{-1}-Y\) to be bounded isomorphisms from
\(\mathcal D_{\rm EM}\) to \(\mathcal H_{\rm EM}\) in these norms, and use
\(X-Y:\mathcal D_{\rm EM}\to\mathcal H_{\rm EM}\),
\(R_Y:\mathcal H_{\rm EM}\to\mathcal D_{\rm EM}\),
\(U:\mathcal D_{\rm EM}\to\mathcal H_{\rm RT}\), and
\(V:\mathcal H_{\rm RT}\to\mathcal H_{\rm EM}\).
Every inverse and product below then uses those fixed domain/codomain
norms; in particular \(R_Y(X-Y)\) is bounded on
\(\mathcal D_{\rm EM}\). A merely formal common dense domain is not
a substitute for these mapping bounds. The retained RT operator
estimates below remain bounded-space statements.

For the pole calculation use a finite retained pencil, or a separately
admitted analytic Fredholm simple-pole realization. Assume the operators
are holomorphic on the same neighbourhood of the algebraically simple
baseline pole \(z_*\), and use the source's \emph{pole} bookkeeping family
\begin{equation}
 K_{X,\lambda_{\rm p}}(z)
       :=K_0(z)+\lambda_{\rm p}\delta K_X(z),\qquad
 z_X(0)=z_*,\quad
 \dot z_X:=\left.\partial_{\lambda_{\rm p}}z_X\right|_0.
 \label{eq:OQED-pole-bookkeeping-family}
\end{equation}
Here \(X\) and its full Schur correction are held fixed while
\(\lambda_{\rm p}\) is varied. This is not the family obtained by
replacing \(X\) by \(\lambda_{\rm p}X\) inside the electromagnetic
resolvent: the latter has first variation \(-UD_0XD_0V\), rather
than \(\delta K_X\). Neither bookkeeping parameter is a new coupling.
Let \(v\) and \(w^\dagger\) be the baseline right/left null vectors
with \(w^\dagger K_0'(z_*)v\ne0\). The pole-family derivatives and
rate coefficients satisfy
\begin{equation}
\begin{aligned}
 \dot z_X&=-\frac{w^\dagger\delta K_X(z_*)v}
                      {w^\dagger K_0'(z_*)v},\qquad
 \Gamma_X^{(1)}=-2\Im\dot z_X,\\
 |\Gamma_X^{(1)}-\Gamma_Y^{(1)}|
 &\leq2\Lambda_*\|\delta K_X(z_*)-\delta K_Y(z_*)\|,\\
 \Lambda_*&=\frac{\|w\|\|v\|}{|w^\dagger K_0'(z_*)v|}.
\end{aligned}
\label{eq:OQED-pole-error-bound}
\end{equation}
These follow by differentiating the null-vector equation and by the
operator norm inequality. The adjoint is the physical represented
adjoint, not the carrier reality operation. The result propagates a
certified response error into the exact-Schur, first-order pole rate;
it does not bound the finite pole-shift remainder. That remainder and
its projected rate accuracy remain separate requirements of the PDF.
A numerical error proved only on real Euclidean momenta cannot be
inserted into \eqref{eq:OQED-Schur-error-bound} at a complex RT pole
without a controlled continuation of the error itself.

\begin{corollary}[Certification from an approximate response]
\label{cor:OQED-aposteriori-Schur}
Suppose only an approximate response \(Y\) and a certified bound
\(\|X-Y\|\leq\epsilon_\Pi\) are available on the common complex domain.
Assume \(R_Y=(D_0^{-1}-Y)^{-1}\) exists there and
\(\|R_Y\|\leq B_Y\), with \(q:=B_Y\epsilon_\Pi<1\), uniformly on
the domain under consideration. Then the exact inverse exists and
\begin{equation}
\begin{aligned}
 \|R_X\|&\leq\frac{B_Y}{1-q},\\
 \|\delta K_X-\delta K_Y\|
       &\leq E_K:=\frac{\|U\|\|V\|B_Y^2\epsilon_\Pi}{1-q},\\
 |\Gamma_X^{(1)}-\Gamma_Y^{(1)}|&\leq2\Lambda_*E_K
                 \quad\text{at the common baseline pole}.
\end{aligned}
\label{eq:OQED-aposteriori-Schur}
\end{equation}
The norms of \(U,V\) are their uniform bounds on that domain, or
their pointwise bounds when only a pointwise statement is made.
For the graph-domain realization, \(B_Y\) is the
\(\mathcal H_{\rm EM}\to\mathcal D_{\rm EM}\) bound and
\(\epsilon_\Pi\) the \(\mathcal D_{\rm EM}\to\mathcal H_{\rm EM}\)
bound. Analyticity of \(X,Y,D_0,U,V\) in the indicated bounded-operator
spaces gives holomorphic inverses and Schur corrections; the norm
inequalities alone do not establish that analytic hypothesis.
\end{corollary}
\begin{proof}
Write \(E=X-Y\). The ordered factorization
\[
 D_0^{-1}-X=(D_0^{-1}-Y)(\one-R_YE)
\]
and the Neumann series give
\(R_X=(\one-R_YE)^{-1}R_Y\). The exact resolvent difference
\(R_X-R_Y=R_X E R_Y\) gives the second estimate, and
\eqref{eq:OQED-pole-error-bound} gives the third. This argument does
not assume a bound on the unknown \(\|D_0X\|\). In particular,
\(B_Y=\|D_0\|/(1-r_Y)\) is available when \(r_Y<1\).
\end{proof}

The strict margin \(\Gamma_Y^{(1)}>2\Lambda_*E_K\) certifies a
positive exact first-order transferred rate on the same represented
branch. It does not certify a positive finite shifted-root rate:
that additional statement needs the independently controlled pole
remainder. Near a null rate the absolute error bound remains meaningful,
whereas a relative error ratio need not. The bound uses the physical
pole functional and is invariant under rescaling either null vector.

\begin{proposition}[A finite-pole enclosure on a retained matrix space]
\label{prop:OQED-finite-pole-enclosure}
On a finite-dimensional retained RT space, define the full pencils
\(K_X=K_0+\delta K_X\) and \(K_Y=K_0+\delta K_Y\).
Suppose both are holomorphic on a neighbourhood of the closed disk
\(\overline{\mathbb D}(z_Y,r)\), \(r>0\), and that
\(\det K_Y\) has exactly one zero there, the simple zero at \(z_Y\),
with no zero on the boundary. Assume a certified bound
\(\|K_X(z)-K_Y(z)\|\leq E_K(z)\) on that boundary and
\begin{equation}
 \sup_{|z-z_Y|=r}\|K_Y(z)^{-1}\|E_K(z)<1.
 \label{eq:OQED-finite-pole-contour-test}
\end{equation}
Then \(\det K_X\) has exactly one zero in the disk, also simple,
and its location \(z_X\) satisfies \(|z_X-z_Y|<r\).
For \(Z\in\{X,Y\}\), use these enclosed endpoint roots to define
the rate-comparison quantities relative to one fixed \(z_*\):
\begin{equation}
\begin{aligned}
 \Gamma_Z^{\rm fin}&:=-2\Im(z_Z-z_*),\\
 |\Gamma_X^{\rm fin}-\Gamma_Y^{\rm fin}|&<2r.
\end{aligned}
 \label{eq:OQED-finite-pole-rate-enclosure}
\end{equation}
Thus \(\Gamma_Y^{\rm fin}>2r\) certifies positivity of the
corresponding \(X\) quantity. Its interpretation as the finite rate of
the selected baseline mode additionally requires the branch-identification
condition below; endpoint proximity alone does not establish it.
\end{proposition}
\begin{proof}
On the boundary, every matrix
\(K_Y+t(K_X-K_Y)=K_Y[\one+tK_Y^{-1}(K_X-K_Y)]\),
\(0\leq t\leq1\), is invertible by the Neumann estimate.
The argument principle applied to its determinant preserves the number
of interior zeros counted with algebraic multiplicity along this
homotopy. That number is one at \(t=0\), so it is one at \(t=1\).
The unique zero is simple and lies strictly inside the disk. The rate
bound follows from \(|\Im(z_X-z_Y)|\leq|z_X-z_Y|\).
\end{proof}

Here \(E_K\) may be supplied by
Corollary~\ref{cor:OQED-aposteriori-Schur}, provided its common-domain
inverse hypotheses exclude electromagnetic poles throughout the disk.
The approximate pencil's root and its isolation must themselves be
certified; an estimated slope or a small residual at one point is not
the contour test. This proposition bounds the \emph{finite} root
difference without a first-order replacement. It is distinct from
\eqref{eq:OQED-pole-error-bound} and does not assert an
infinite-dimensional determinant theorem. Root clusters, branch singularities
intersecting the disk, and unproved continuum reductions require
separate spectral arguments.

\paragraph{Endpoint matching is not baseline-mode tracking}
The proposition matches roots through \(K_Y+t(K_X-K_Y)\) at the
physical endpoint \(\lambda_{\rm p}=1\). This is not the pole family
\eqref{eq:OQED-pole-bookkeeping-family}. A concrete counterexample,
in fixed reference units, is
\[
 \begin{gathered}
 K_0(z)=z^2-1,\qquad K_Z(z)=z^2-q_Z,\\
 q_Y=-1-\ii\delta_{\rm ex},\qquad q_X=-1+\ii\delta_{\rm ex}.
 \end{gathered}
\]
Take \(\delta_{\rm ex}=0.01\), and track each affine family
\(z^2-[1+\lambda_{\rm p}(q_Z-1)]\) from the baseline root \(+1\).
Both paths have simple roots throughout. They end at the principal
roots \(z_Y=\sqrt{q_Y}\) and \(z_X^{\rm tracked}=\sqrt{q_X}\),
near \(-\ii\) and \(+\ii\), respectively. Nevertheless, on the
circle \(|z-z_Y|=r=0.05\),
\[
 \frac{\|K_X-K_Y\|}{|K_Y|}
 \leq\frac{2\delta_{\rm ex}}{r(2|z_Y|-r)}<0.206<1.
\]
The enclosed \(X\) root is therefore \(-\sqrt{q_X}\), not
\(z_X^{\rm tracked}\). With baseline \(+1\), the matched endpoint
has a positive rate-comparison quantity while the tracked \(X\) branch
has the opposite sign. This is an auxiliary matrix-pencil example, not an O--QED model. It
does not contradict the root-count theorem; it shows why physical mode
identification is an additional input.

One sufficient identification certificate uses a fixed positively
oriented simple contour \(\Gamma_0\). Its interior contains the endpoint
disk and exactly one baseline zero, the simple root \(z_*\). Assume
all pencils are holomorphic on a neighbourhood of the closed interior
and \(K_0\) is invertible on the contour. Require
\begin{equation}
 \sup_{z\in\Gamma_0}\|K_0(z)^{-1}\|
       \max\{\|\delta K_X(z)\|,\|\delta K_Y(z)\|\}<1.
 \label{eq:OQED-common-baseline-contour}
\end{equation}
The boundary then stays invertible for the entire two-parameter family
\(K_0+\lambda_{\rm p}[(1-t)\delta K_Y+t\delta K_X]\),
\((\lambda_{\rm p},t)\in[0,1]^2\). The argument principle gives
one simple interior root throughout. Uniqueness identifies both endpoint
roots with the continuations of the same baseline mode. A less restrictive
certified contour family may be used, but it must establish that same
identification, not only endpoint root counts.

\begin{corollary}[A contour certificate for the finite pole remainder]
\label{cor:OQED-pole-Taylor-remainder}
On the same finite retained matrix space, let \(K_0\) and \(\delta K\)
be holomorphic near \(\overline{\mathbb D}(z_*,R_{\rm p})\),
with \(R_{\rm p}>0\).
Suppose \(\det K_0\) has precisely one zero in this disk, the simple
zero \(z_*\), and none on its boundary \(\Gamma_*\). Assume a
certified constant \(0<\vartheta_{\rm p}<1\) satisfies
\begin{equation}
 \sup_{z\in\Gamma_*}\norm{K_0(z)^{-1}\delta K(z)}
                    \leq\vartheta_{\rm p}.
 \label{eq:OQED-complex-pole-margin}
\end{equation}
The pole family \(K_{\lambda_{\rm p}}=K_0+\lambda_{\rm p}\delta K\)
has one simple interior root \(z(\lambda_{\rm p})\), holomorphic for
\(|\lambda_{\rm p}|<\vartheta_{\rm p}^{-1}\), with
\(z(0)=z_*\). Writing \(z(\lambda_{\rm p})-z_*
=\sum_{n\geq1}c_n\lambda_{\rm p}^n\), one has
\begin{equation}
 \begin{aligned}
  |c_n|&\leq R_{\rm p}\vartheta_{\rm p}^{\,n},\\
  \left|z(1)-z_*-\sum_{n=1}^{N}c_n\right|
     &\leq\frac{R_{\rm p}\vartheta_{\rm p}^{\,N+1}}
                    {1-\vartheta_{\rm p}},\\
  |\Gamma^{\rm fin}-\Gamma^{(1)}|
     &\leq\frac{2R_{\rm p}\vartheta_{\rm p}^{\,2}}
                    {1-\vartheta_{\rm p}}.
 \end{aligned}
 \label{eq:OQED-pole-remainder-certificate}
\end{equation}
Here \(N\geq1\), \(c_1=\dot z\), and the rates use the same
baseline and pole bookkeeping parameter as
\eqref{eq:OQED-pole-bookkeeping-family}. The radius and margin are
certification data, not additional theory parameters.
\end{corollary}
\begin{proof}
For every \(|\lambda_{\rm p}|<\vartheta_{\rm p}^{-1}\), the
Neumann estimate keeps the boundary inverse well defined. Homotopy
from zero preserves the number of interior zeros with multiplicity,
so there is exactly one, and it is simple. The argument principle
also gives the root itself:
\begin{equation}
 z(\lambda_{\rm p})=\frac1{2\pi\ii}\oint_{\Gamma_*}
 z\,\Tr\!\left[K_{\lambda_{\rm p}}(z)^{-1}
                    \partial_zK_{\lambda_{\rm p}}(z)\right]\dd z.
 \label{eq:OQED-pole-contour-moment}
\end{equation}
The contour integrand is locally uniformly holomorphic in
\(\lambda_{\rm p}\), so the root is holomorphic throughout that
disk and remains within distance \(R_{\rm p}\) of \(z_*\).
Cauchy's coefficient estimate on \(|\lambda_{\rm p}|=r_{\lambda}<\vartheta_{\rm p}^{-1}\) gives
\(|c_n|\leq R_{\rm p}r_{\lambda}^{-n}\). Letting
\(r_{\lambda}\uparrow\vartheta_{\rm p}^{-1}\) and summing the
geometric tail at \(\lambda_{\rm p}=1\) proves the first two
bounds. The rate bound follows by taking the imaginary part.
No inference from a real-parameter estimate to complex analyticity is
used: the affine pencil and boundary inverse were controlled for
complex \(\lambda_{\rm p}\) first.
\end{proof}

For \(\delta K=\delta K_X\), the unknown exact correction may be
bounded on \(\Gamma_*\) using
\[
 \norm{K_0^{-1}\delta K_X}
 \leq\norm{K_0^{-1}}\bigl(\norm{\delta K_Y}+E_K\bigr).
\]
If this proves \eqref{eq:OQED-complex-pole-margin}, and the same
response certificate controls the baseline point, then
\begin{equation}
 |\Gamma_X^{\rm fin}-\Gamma_Y^{(1)}|
 \leq2\Lambda_*E_K(z_*)+
       \frac{2R_{\rm p}\vartheta_{\rm p}^{\,2}}
                     {1-\vartheta_{\rm p}}.
 \label{eq:OQED-combined-finite-rate-error}
\end{equation}
This is one explicit, possibly conservative route to the previously
separate pole-remainder requirement. A margin below one guarantees a
bound, not necessarily a small bound. In the \(e^{-\ii zt}\)
response convention, for an unstable baseline,
positive \(\Gamma^{\rm fin}\) means reduced growth; strict decay of
this mode requires \(\Gamma^{\rm fin}>2\Im z_*\), not merely
\(\Gamma^{\rm fin}>0\).

\paragraph{A counted root and its observable residue}
A simple determinant zero is a pole of the matrix inverse, not of every
source/readout component. Let \(K(z_0)v=0\),
\(w^\dagger K(z_0)=0\), and
\(d_{\rm p}:=w^\dagger K'(z_0)v\ne0\) at a simple zero.
The local resolvent form is \cite[Theorem~2.4]{Beyn2012}
\begin{equation}
 K(z)^{-1}=\frac{v w^\dagger}{d_{\rm p}(z-z_0)}+H(z),
 \qquad H\text{ holomorphic near }z_0.
 \label{eq:OQED-readout-resolvent}
\end{equation}
For completeness, the adjugate formula gives a rank-one pole
coefficient proportional to \(v w^\dagger\); Jacobi's formula for
\(\partial_z\det K\) fixes its coefficient to \(d_{\rm p}^{-1}\).
For holomorphic input and output maps on that neighbourhood,
\begin{equation}
 \begin{aligned}
 \mathcal O(z)&=L_{\rm out}(z)K(z)^{-1}L_{\rm in}(z),\\
 \mathop{\rm Res}_{z=z_0}\mathcal O
   &=\frac{[L_{\rm out}(z_0)v]
                   [w^\dagger L_{\rm in}(z_0)]}{d_{\rm p}}.
 \end{aligned}
 \label{eq:OQED-observable-pole-residue}
\end{equation}
Both factors must be nonzero for this readout to contain the pole.
For example, \(K=\diag(z-z_0,1)\),
\(L_{\rm out}=(0,1)\), and \(L_{\rm in}=(0,1)^T\) give
\(\mathcal O=1\), despite the simple determinant zero. Thus root
location, continuation of the baseline mode, and its nonzero overlap
with the selected source and readout are distinct certificates. This
qualification changes neither the internal vectorial prescription nor
the pole of the full retained inverse. The maps are physical response
maps, not the carrier characters. The residue in
\eqref{eq:OQED-readout-resolvent} is not in general an idempotent:
for \(K(z)=2(z-z_0)\) it is \(1/2\). It must not be confused
with a normalized spectral projector.

\begin{proposition}[Observable residue error]
\label{prop:OQED-observable-residue-error}
Suppose the finite matrix pencils \(K_X,K_Y\) are holomorphic near
and inside a positively oriented rectifiable simple contour \(\Gamma\),
and each has exactly one interior determinant zero, both simple.
Use the same holomorphic input and output maps in
\(\mathcal O_Z=L_{\rm out}K_Z^{-1}L_{\rm in}\), \(Z=X,Y\),
with no other interior singularities. In compatible operator norms,
assume the following certified uniform boundary bounds
\[
 \begin{gathered}
  \|K_Y^{-1}\|\leq B_\Gamma,\qquad
  \|K_X-K_Y\|\leq\epsilon_K,\qquad
  q_\Gamma:=B_\Gamma\epsilon_K<1,\\
  \|L_{\rm out}\|\leq L_{{\rm out},\Gamma},\qquad
  \|L_{\rm in}\|\leq L_{{\rm in},\Gamma}.
 \end{gathered}
\]
Let \(\ell(\Gamma)\) be the contour length and
\(\mathcal R_Z^{\rm obs}\) the residue of \(\mathcal O_Z\) at
its enclosed zero, allowing this residue to vanish. Then
\begin{equation}
 \begin{aligned}
  \|\mathcal R_X^{\rm obs}-\mathcal R_Y^{\rm obs}\|
      &\leq E_{\rm obs},\\
  E_{\rm obs}&:=\frac{\ell(\Gamma)}{2\pi}
       L_{{\rm out},\Gamma}L_{{\rm in},\Gamma}
       \frac{B_\Gamma^2\epsilon_K}{1-q_\Gamma}.
 \end{aligned}
 \label{eq:OQED-observable-residue-error}
\end{equation}
A certified lower bound
\(\|\mathcal R_Y^{\rm obs}\|\geq\underline r_Y>E_{\rm obs}\)
therefore proves that the exact readout contains its enclosed pole.
\end{proposition}
\begin{proof}
Put \(E=K_X-K_Y\). The boundary Neumann estimate gives
\(\|K_X^{-1}\|\leq B_\Gamma/(1-q_\Gamma)\), and
\[
 K_X^{-1}-K_Y^{-1}=-K_X^{-1}E K_Y^{-1}.
\]
Each observable has at most one simple pole inside \(\Gamma\), so
\(\mathcal R_Z^{\rm obs}=(2\pi\ii)^{-1}
\oint_\Gamma\mathcal O_Z(z)\dd z\). Integrating the ordered
resolvent difference proves \eqref{eq:OQED-observable-residue-error};
the reverse triangle inequality proves the nonvanishing assertion.
The two enclosed pole locations need not coincide.
\end{proof}

The lower bound must include any error in evaluating the approximate
residue. When the input or output maps themselves are approximated,
their errors must also be propagated; the displayed bound assumes common
exact maps. Root counts and baseline-mode identification remain separate
requirements. A nonzero residue certifies visibility in this response
channel, not positivity of that residue, a damping sign, or a completed
physical reconstruction.

\paragraph{Certifying a contour bound from finitely many evaluations}
For a practical sufficient test, use the spectral norm on the finite
matrix space and a circle \(\Gamma\) with nodes \(z_j\). Let every
point of \(\Gamma\) lie within arc length \(h\) of a node. Suppose
\(\underline\sigma_j\) are certified lower bounds for
\(\sigma_{\min}(K_Y(z_j))\), including evaluation and arithmetic
errors, and suppose a bound \(L\geq\sup_\Gamma\|K_Y'(z)\|_2\)
has been proved. For a certified uniform boundary error \(E_*\), put
\begin{equation}
 \begin{aligned}
  d_*&:=\min_j\underline\sigma_j-Lh,\\
  E_*&\geq\sup_\Gamma\|K_X-K_Y\|_2.
 \end{aligned}
 \label{eq:OQED-mesh-contour-margin}
\end{equation}
If \(d_*>E_*\geq0\), then the approximate pencil is nonsingular
on the entire circle and
\(\sup_\Gamma\|K_Y^{-1}\|_2\|K_X-K_Y\|_2\leq E_*/d_*<1\).
Indeed, integration of \(K_Y'\) along the intervening arc bounds
the matrix difference by \(Lh\), and the variational definition of
the smallest singular value gives the same Lipschitz bound for it.
This certifies the boundary inequality, not the interior root count,
which remains a separate hypothesis. Finite sampling without the
between-node bound is insufficient: \(1+(z/r)^n\) equals \(2\) at
\(z_j=r e^{2\pi\ii j/n}\) but vanishes between those nodes.

The remaining root count can, for example, be certified independently by
\begin{equation}
 N_\Gamma=\frac1{2\pi\ii}\oint_\Gamma
          \Tr[K_Y(z)^{-1}K_Y'(z)]\dd z.
 \label{eq:OQED-certified-root-count}
\end{equation}
With holomorphy inside and nonsingularity on the contour, this is the
nonnegative integer number of zeros of \(\det K_Y\), counted with
multiplicity. A certified quadrature value \(\widetilde N\) with
error \(\epsilon_N\) establishes \(N_\Gamma=1\) whenever
\(|\widetilde N-1|+\epsilon_N<1/2\). An uncertified near-integer
floating-point value is not a root count. Nor does a count of one make
a numerical centre an exact zero: if a disk of radius \(r\) about
\(z_c\) encloses \(z_X\), while only
\(|z_Y-z_c|\leq\epsilon_z\) is certified, then the resulting
root-distance bound is \(r+\epsilon_z\), not \(r\). The rate
bound inherits the same additional factor \(2\epsilon_z\).

\paragraph{Error from eliminating unretained modes}
All preceding finite-matrix certificates concern the declared exact
retained pencil. A numerical compression of a larger theory requires
an additional reduction estimate. Write the larger inverse-response
pencil as
\[
 \mathscr K(z)=\begin{pmatrix}A(z)&B(z)\\C(z)&D(z)\end{pmatrix}.
\]
Assume \(D(z)^{-1}\) exists on the common domain. The exact retained
pencil is then \(A-BD^{-1}C\), not merely \(A\). An approximation \(\widetilde A\) which discards that complement
must therefore include, in compatible operator norms,
\begin{equation}
 \|A-BD^{-1}C-\widetilde A\|
 \leq\|A-\widetilde A\|+\|B\|\|D^{-1}\|\|C\|
 \label{eq:OQED-retained-complement-error}
\end{equation}
in its total error budget. The blocks and complementary inverse must
be holomorphic on the same disk neighbourhood, with the required
uniform bounds on its contour. Uniform norm bounds alone do not
supply holomorphy.
Near a singular complement, this reduction and its finite-pole
certificate are unavailable until a larger retained space or another
proved reduction is used. No continuum error estimate is supplied by
matrix quadrature alone.

The same elimination must transport the source and readout maps.
For the compatible block input and output
\(J=(J_r,J_c)^T\) and \(L=(L_r,L_c)\), define
\(S=A-BD^{-1}C\). Wherever the inverses exist, block factorization
gives the exact response identity
\begin{equation}
 \begin{aligned}
  L\mathscr K^{-1}J
     &=L_{\rm eff}S^{-1}J_{\rm eff}+L_cD^{-1}J_c,\\
  L_{\rm eff}&=L_r-L_cD^{-1}C,\\
  J_{\rm eff}&=J_r-BD^{-1}J_c.
 \end{aligned}
 \label{eq:OQED-Schur-source-readout}
\end{equation}
Indeed, solving the complementary row gives
\(x_c=D^{-1}(J_c-Cx_r)\), followed by
\(Sx_r=J_{\rm eff}\); inserting both into the output proves the
formula, with all products in their displayed order. If the blocks,
complementary inverse, sources, and readouts are holomorphic, the last
term is holomorphic and the pole residue comes from the two effective
endpoint maps. Keeping only \(L_r,J_r\) is valid only under conditions
that remove the omitted terms, such as \(L_c=J_c=0\); controlling the
retained pencil alone does not justify that omission. The identity is a
block-response statement, not an additional microscopic O--QED action.

The strict \(H_{k^2}\) status, the low-wavenumber transfer test, and the
finite-band \(H_{T,k^2}\) match in Ref.~\cite{ChiapponeIgnition} are
not consequences of UV finiteness. Neither is incremental passivity or
the sign of the transferred rate. Only after those separate tests can
its reduced law \(\Gamma_{\rm OQED}^{(0)}=e^2k^2
\mathcal K_{r_{\rm resp}}(\ell_{\rm dyn};x,\beta)\) be used.
The laboratory coordinate \(\ell_{\rm dyn}\) is not \(a\),
\(\ell_\Omega\), or \(\ell_{\rm GT}\). No fusion-yield or calibrated
attenuation number is inferred from the kernel certificate.

\begin{table*}[t]
\centering
\small
\setlength{\tabcolsep}{5pt}
\renewcommand{\arraystretch}{1.15}
\caption{O--QED application: what the explicit calculations certify and
what remains to be established for one completed interacting theory.
A failed sufficient criterion is not automatically an impossibility
result for a different completion.}
\label{tab:OQED-certification-status}
\begin{tabularx}{\textwidth}{@{}>{\raggedright\arraybackslash}p{0.21\textwidth}
 >{\raggedright\arraybackslash}p{0.37\textwidth}
 >{\raggedright\arraybackslash}X@{}}
\toprule
Object & Established in this article & Scope or outstanding requirement\\
\midrule
Normalized square-root kernel & All finite inverse moments; positive-depth derivative and explicit selected-line tail estimates & Properties of a specified virtual kernel, not of completed vertices\\
Selected Euclidean graph & Absolute convergence when the full routing, vertex envelope, line moments, and IR bounds meet the stated theorem & These hypotheses are not derived for every graph of the O--QED action\\
Bare QED comparison & A finite tensor with a nonzero small-momentum Ward remainder & Not the Ward-completed polarization or a physical photon mass\\
Ward-compatible auxiliary vertices & The indicated vacuum bubble has a nonintegrable contraction at fixed nonzero transfer & Its separately defined contact-completed sum requires a new common-integrand estimate\\
Constant-background contact test & Leading bubble growth cancels, but the declared trace-log Hessian is still not absolutely integrable & Tests an additional auxiliary realization, not every possible completed contact sector\\
Physical amplitude and response & Conditional source, residue, contour, and error bounds & Positive physical reconstruction, cut matching, and a completed medium response remain separate\\
\bottomrule
\end{tabularx}
\end{table*}

\subsection{Application status}
\label{sec:OQED-status}

For the supplied O--QED profile, normalization, complete monotonicity,
all inverse moments, fixed positive-depth derivatives, their singular
zero-depth and strong-depth limits, and the bare one-loop integrability
test are explicit calculations. The optimized selected-line theorem
also certifies appropriate massless selected denominators provided
the remaining complete graph factors obey its anisotropic bound. For
\(M>0\), the additional local analytic kernel hypothesis is also
verified. The Abelian BRST differential and the source action's sector
and carrier conventions are matched to this article.

The interaction vertices and contacts remain conditional, as do their
graphwise or grouped estimates, the physical source realization, and
the state-dependent reconstruction. The Ward-compatible vacuum bubble in
Proposition~\ref{prop:OQED-Ward-bubble-obstruction} and
Corollary~\ref{cor:OQED-fixed-transfer-obstruction} actually fails the
individual-graph test, despite the verified kernel moments. The additional
constant-background contact test cancels its leading zero-transfer
term but fails the grouped absolute bound as well. Its raw reference
difference remains nonintegrable for every positive radial multiplier
tending to zero under the same shift, trace-log, and measure assumptions.
This is a specified auxiliary realization, not an evaluation of the
PDF's full completion. The scalar tube, its signed real-invariant and
complex Taylor bounds, and its finite zero-depth \emph{subtracted}
comparison are independent results. The first-order and finite-pole
error certificates are conditional bounds on the explicitly stated
analytic and operator domains. Endpoint root enclosure
requires baseline-branch identification; the complex pole-family
margin also supplies an explicit finite-remainder bound. A selected
source/readout channel must have nonzero residue to contain that pole;
\eqref{eq:OQED-observable-residue-error} supplies a sufficient quantitative
test when its contour data are available. Thus the application proves neither unconditional all-orders
finiteness of the unspecified completion nor a generic impossibility
theorem for O--QED. A positive interacting-theory certificate would
have to exhibit \emph{one specified completed action}, its exact
source/vertex/contact Feynman rules, and a common graphwise or
correctly grouped integrable majorant for every claimed diagram.
Only then could the BRST and physical-interface hypotheses be tested
on the same completed construction. The present results supply
quantitative kernel certificates and explicit interaction-level
obstructions against which that future calculation can be checked.
Table~\ref{tab:OQED-certification-status} separates the corresponding
mathematical results from the unresolved theory-level obligations.



\section{Conclusions}
\label{sec:conclusions}

At a fixed positive kernel scale, the Laplace--Stieltjes representation
turns ultraviolet suppression into a family of Gaussian estimates.
Theorem~\ref{thm:global-uv} proves absolute convergence uniformly on
compact Euclidean external sets when the completed numerator and
infrared envelopes satisfy its hypotheses. Its proof controls
nonuniform parameter limits through a bounded affine shift and a
factorized inverse-moment majorant. Theorem~\ref{thm:forest} retains
individual UV scales when weaker linewise damping is sufficient.
Theorem~\ref{thm:optimized-selected-lines} gives line-specific
moment budgets and allows selected massless denominators under a
complete anisotropic envelope; its no-atom hypothesis is independent
of the inverse moments as defined here. Corollary~\ref{cor:optimized-selected-tail}
quantifies ultraviolet numerical tails. The two-loop example additionally
has an explicit global Lipschitz gradient (and hence $C^{1,1}$
regularity), a controlled quadratic real-momentum remainder,
and complex-tube derivative bounds.
Proposition~\ref{prop:residual-coercivity} also admits nonpolynomial
completed numerators when a strictly positive quadratic or linear
damping margin remains. Together these results establish finiteness
of every admitted finite loop graph. Their all-orders scope is
graphwise and does not assert convergence of the infinite series.

The symmetry argument applies to the same graph family. A completed
anomaly-free BRST/BV action yields the Slavnov--Taylor hierarchy with
finite coefficients whenever its vertices and insertions satisfy the
convergence bounds. The background Ward identity expresses background
covariance. The Nielsen identity gives gauge independence for the
fixed or consistently transported physical observables and states of
Section~\ref{sec:Nielsen}, under its normalization and limit hypotheses.
The regulator-free physical action retains its branch-directed
carrier. Euclidean calculations use explicitly normalized
coefficient sources; physical connected and effective functionals
are reconstructed from both carrier characters, with the vectorial
boundary prescription retained before every Lorentzian inversion.
Carrier reality and the positive physical adjoint remain distinct;
character reconstruction alone provides no positive state space.
Source differentiation also imposes graded exchange symmetry. The
unrestricted running-unit scalar line is retained as oriented data;
its identification with an identical real-boson covariance is not
asserted and fails the direct evenness test without additional source
or endpoint structure.

The examples identify concrete tests for an interacting model.
The scalar bubble keeps one multiplier on each internal line and
exhibits the shell weights in its analytic discontinuity. A finite
shell normalization isolates an entire remainder and gives a
controlled comparison with the ordinary subtracted amplitude, but
its nonpolynomial real-axis growth survives finite polynomial matching.
The literal massive stable-half bubble is UV finite while failing
\(C^6\) regularity at zero external momentum. The QED calculation
proves that the bare-vertex loop is nontransverse
under the stated exact one-loop integrability condition. A longitudinal
vertex satisfies the dressed Ward identity wherever its defining
quotients exist. The contracted norm identities then test vertex
growth at high transfer and, under diagonal continuity, in the
soft-photon limit at large fermion momentum. The massive stable-half
profile also forces a directional discontinuity at the zero-momentum
endpoint. These tests connect the covariance to the vertices and
normalizations required by an interacting theory. The kernel families
and radial integrals provide
explicit analytic and numerical checks of the Euclidean construction.

The O--QED application explicitly verifies all inverse moments for the
normalized massive square-root kernel and evaluates its one-loop
integrability coefficient. Its positive-depth derivative bounds,
nonuniform zero-depth limits, complete monotonicity in depth,
exponential numerical tails, and the strong-depth crossover are proved
for the auxiliary scalar comparisons. A positive auxiliary kernel scale removes
the specific zero-momentum cusp obstruction and admits local holomorphy
under the massive graph hypotheses. It does not tame the Ward vertex:
the corresponding two-vertex vacuum polarization term has a nonzero
large-momentum Ward-contraction limit and is not individually integrable.
A separately specified constant-background contact Hessian cancels the
leading term but leaves a nonintegrable total derivative. Its radial
flux proves that changing the vanishing multiplier alone cannot make
the raw ordinary-reference difference absolutely convergent while the
same auxiliary assumptions are retained. A successful completion needs
an absolute bound on each graph or on a precisely defined common-routing
sum, including all ultraviolet regions.
The scalar comparison has an explicit holomorphic tube, a positive
proper-time measure for its invariant dependence, signed finite-order
Taylor enclosures, and numerical tail bounds. Its first coefficients
also have one-dimensional formulas. Forming the literal scalar
zero-momentum subtraction before integration yields a finite
$a\downarrow0$ comparison, although its unsubtracted value diverges.
Its unsubtracted logarithmic finite part is also explicit and
independent of the auxiliary kernel mass in the stated limit.
This does not perform a gauge-theory kernel-off subtraction.
Kernel-off matching and interpolation are not
automatic extensions of the fixed-depth theorem. On a certified common
complex response domain, the typed a posteriori Schur estimate controls
inverse existence and the stated pole-family derivative. A separate
finite-dimensional contour test encloses an endpoint root without
replacing it by that derivative. Control of the complex pole parameter
also bounds its finite Taylor remainder. Baseline-mode identification,
nonzero source/readout residues, certified contour evaluations, and any
omitted-mode error remain separate requirements. The residue error bound
and exact endpoint transport make these response tests explicit.
None of these estimates
determines a hydrodynamic power or supplies the missing medium response.

Local holomorphy is established by a common complex majorant.
Physical reconstruction and crossing additionally require the stated
positivity, spectral, and asymptotic hypotheses. A specific completion
must therefore meet source reciprocity, UV control, and BRST symmetry
together and, when a scattering interpretation is claimed, establish
the required physical state space and cut matching. In the
pure-redefinition construction of
Proposition~\ref{prop:pure-redefinition-cancellation}, the transformed
vertices cancel the apparent line suppression. The convergence and
symmetry conditions must therefore be verified for the same completed
action. The present results provide those tests and their conditional
consequences at fixed scale and finite perturbative order. Convergence
of the full series and integration over unrestricted families of
supports require further estimates.

\begin{samepage}
\section*{Acknowledgements and use of generative AI}

The framework is the author's independent research programme.
Generative AI tools assisted with \LaTeX\ preparation, comparison with
the cited source material, editorial refinement, and mathematical
consistency checks. The author is responsible for the scientific
claims and for reviewing and approving the final manuscript.
\end{samepage}
\appendix
% Keep theorem numbers compact; elsarticle retains its native appendix headings.
\renewcommand{\thetheorem}{\Alph{section}.\arabic{theorem}}

\section{Details of the Gaussian and determinant estimates}
\label{sec:appendix-gaussian}

\subsection{Weighted-basis identity for the affine shift}

Let \(C\) be an \(I\times L\) real matrix of full column rank,
\(W=\diag(w_i)\) with \(w_i>0\), and
\(A=C^TWC\). For one component of the routing vector,
\[
 \mu=A^{-1}C^TWR.
\]
Cauchy--Binet gives
\begin{equation}
 \det A=\sum_{|T|=L}(\det C_T)^2\prod_{i\in T}w_i.
 \label{eq:CB-appendix}
\end{equation}
Terms with singular \(C_T\) vanish.

For component \(j\) of \(\mu\), let \(C^{(j)}\) be \(C\)
with column \(j\) replaced by \(R\).
Cramer's rule and a second Cauchy--Binet expansion give
\begin{align}
 \mu_j
 &=\frac{\det(C^TWC^{(j)})}{\det A}\nonumber\\
 &=\frac{
 \sum_{|T|=L}\det C_T\det C_T^{(j)}
                          \prod_{i\in T}w_i}
 {\det A}.
 \label{eq:Cramer-shift}
\end{align}
For invertible \(C_T\), Cramer's rule in that square subsystem
says
\[
 \det C_T^{(j)}
     =\det C_T\,(C_T^{-1}R_T)_j.
\]
Substitution yields \eqref{eq:shift-mixture}. Since there are
only finitely many such subsystems, compact external data imply
a uniform bound independently of relative widths. The finite convex
representation supplies the uniformity needed when some widths
vanish much faster than others.

\subsection{Gaussian moments with the affine shift retained}

Writing \(n=4L\), let
\(y=(A^{1/2}\otimes\one_4)(k+\bmu)\).
For \(D=0\), the numerator is bounded and the zeroth Gaussian
integral suffices. For \(D>0\), let
\(c_D=\max\{1,2^{D-1}\}\) and
\(B=\sup_{w_i>0,\,p\in\mathcal K}\norm{\bmu(w,p)}<\infty\).
The inequality \((x+y)^D\leq c_D(x^D+y^D)\) gives
\[
 1+\norm{k}^D\leq
 1+c_DB^D+c_D\norm{A^{-1}}^{D/2}\norm{y}^D.
\]
For every real \(D\geq0\), the radial Gaussian moment is
\begin{equation}
 \int_{\RR^n}|y|^D e^{-|y|^2}\dd^ny
   =\frac{\pi^{n/2}\Gamma((n+D)/2)}{\Gamma(n/2)}.
 \label{eq:Gaussian-moment-exact}
\end{equation}
It follows by integrating over the unit sphere and substituting
\(v=|y|^2\). Multiplying by the Jacobian
\((\det A)^{-2}\) proves \eqref{eq:gaussian-bound}.
For \(D=0\), only the zeroth moment is required; the displayed
constant-term convention remains harmless.

\subsection{A direct momentum-space proof of the coverage criterion}
\label{sec:direct-loop-proof}

For Corollary~\ref{cor:damped-basis}, write the bounded multiplier
on an unselected line as \(\Psi_i\), with
\(\sup_{z\geq0}|\Psi_i(z)|\leq B_i\). Its whole line factor obeys
\(\abs{\Psi_i(a^2q_i^2)}C_i(q_i)\leq B_i\mathcal I_i\).
Use the selected momenta as independent variables:
\[
 k=(C_S^{-1}\otimes\one_4)(q_S-R_S(p)),\qquad
 \dd^{4L}k=|\det C_S|^{-4}\dd^{4L}q_S.
\]
The bounded external shift gives
\(1+\norm{k}^D\leq C_{G,\mathcal K,S}
(1+\sum_{j\in S}|q_j|^D)\) for \(D>0\).
For \(D=0\), a constant bound suffices.

For \(0\leq d\leq D\), define the one-line integral
\begin{align}
 Q_{d,j}
 &:=\int_{\RR^4}|q|^d\Ph_j(a^2q^2)C_j(q)\dd^4q \nonumber\\
 &=\pi^2\Gamma(2+d/2)
   \int\rho_j(\dd\tau)\int_0^\infty
   \frac{E_j(\alpha)\dd\alpha}
        {(\alpha+a^2\tau)^{2+d/2}}.
 \label{eq:one-line-moment}
\end{align}
Here \(|q|^0\) means the constant function one. Tonelli's theorem
and the radial Gaussian integral justify the second equality.
Writing \(M_{0,j}=\rho_j([0,\infty))\),
\eqref{eq:alpha-envelope-bound} gives
\begin{align}
 Q_{d,j}\leq\pi^2\Gamma(2+d/2)
 \bigg[&\frac{\mathcal E_j a^{-2-d}}{1+d/2}
           M_{1+d/2}[\rho_j] \nonumber\\
       &+\alpha_0^{-2-d/2}\mathcal I_j M_{0,j}\bigg].
 \label{eq:one-line-moment-bound}
\end{align}
Finite measure mass and the assumed moment make both
\(Q_{0,j}\) and \(Q_{D,j}\) finite. Apart from fixed routing and
Fourier constants, the transformed graph majorant is bounded by
\[
 \left(\prod_{i\notin S}B_i\mathcal I_i\right)
 \left[\prod_{j\in S}Q_{0,j}
   +\sum_{j\in S}Q_{D,j}
        \prod_{h\in S\setminus\{j\}}Q_{0,h}\right].
\]
For \(D=0\), the first product alone is enough. The bound holds
on all loop-momentum space, so it controls partial and simultaneous
UV limits with constants uniform on \(\mathcal K\).
For a positive mass, \(\mathcal I_i=m_i^{-2}\).
The large-radius factor \(r^{D+1}\Ph_j(a^2r^2)\) explains the
inverse-moment order \(1+D/2\) in this sufficient estimate.

\section{Normalization and limiting checks for the examples}
\label{sec:appendix-checks}

\subsection{The two-body phase-space factor}

In the future rest frame \(P=(\sqrt{\sMand},\mathbf0)\),
the spatial delta function in \eqref{eq:two-body-phase-space}
sets \(\mathbf q_2=-\mathbf q_1\).
For equal masses, \(\omega=\sqrt{|\mathbf q_1|^2+m^2}\);
the remaining integral is
\[
 \frac1{4(2\pi)^2}
     \int\frac{4\pi r^2\dd r}{\omega^2}
         \delta(\sqrt{\sMand}-2\omega).
\]
Since \(\dd(2\omega)/\dd r=2r/\omega\), this equals
\[
 \frac{r_*}{8\pi\omega_*}
       =\frac1{8\pi}\sqrt{1-\frac{4m^2}{\sMand}}.
\]
The imaginary part of a Schwarz-reflecting scalar bubble is
half its discontinuity divided by \(\ii\).
A separate graph symmetry factor must then be applied at the
amplitude level.

\subsection{Zero-momentum Gaussian bubble}

Set \(b=2t\) and \(x=bm^2\). Radial integration of
\eqref{eq:two-line-bubble} at zero external momentum gives
\[
 16\pi^2B_{t,E}(0)
   =\int_0^\infty
          \frac{u\,e^{-bu}}{(u+m^2)^2}\dd u.
\]
Let
\[
 I_1=\int_0^\infty\frac{e^{-bu}}{u+m^2}\dd u
       =e^xE_1(x).
\]
Differentiating with respect to \(m^2\),
\[
 I_2=\int_0^\infty\frac{e^{-bu}}{(u+m^2)^2}\dd u
       =\frac1{m^2}-bI_1.
\]
Using \(u/(u+m^2)^2=1/(u+m^2)-m^2/(u+m^2)^2\)
yields
\[
 16\pi^2B_{t,E}(0)
   =I_1-m^2I_2=(1+x)e^xE_1(x)-1.
\]
As \(x\to\infty\), the bracket is \(x^{-2}+O(x^{-3})\).
As \(x\downarrow0\), it diverges logarithmically.
The finite-\(a\) theorem thus does not authorize interchanging
integration with the formal removal of damping.

\subsection{Radial kernel: endpoint and numerical stability}

Expanding \(e^{-ur^2}\) to a fixed finite order in
\eqref{eq:J-radial} gives, as \(u\downarrow0\) at fixed
\(a>0\),
\begin{equation}
 J(u;a)
 =\frac6{a^4}-\frac{120u}{a^6}
        +\frac{2520u^2}{a^8}
        -\frac{60480u^3}{a^{10}}+O(u^4).
 \label{eq:J-endpoint-expansion}
\end{equation}
Each coefficient follows from
\(\int_0^\infty r^n e^{-ar}\dd r=n!/a^{n+1}\).
More generally, set
\[
 P_N(u;a)=\sum_{n=0}^N
       \frac{(-u)^n(2n+3)!}{n!\,a^{2n+4}}.
\]
Taylor's remainder for \(e^{-ur^2}\), integrated against the positive
weight \(r^3e^{-ar}\), gives the explicit enclosure
\begin{equation}
\begin{gathered}
 0\leq(-1)^{N+1}\bigl[J(u;a)-P_N(u;a)\bigr]\\
 \leq\frac{(2N+5)!}{(N+1)!\,a^{2N+6}}u^{N+1},
 \qquad a>0,\quad u\geq0.
\end{gathered}
 \label{eq:J-truncation-bound}
\end{equation}
Odd truncations are lower bounds and even truncations are upper
bounds. In particular,
\[
 0\leq J(u;a)-P_3(u;a)
       \leq\frac{1663200\,u^4}{a^{12}}.
\]
These estimates apply to each finite truncation. The scaling identity
\[
 J(u;a)=a^{-4}j(u/a^2),\qquad
 j(\xi)=\int_0^\infty x^3e^{-x-\xi x^2}\dd x
\]
identifies \(u/a^2\) as the dimensionless endpoint parameter.
The estimates become useful for approximation when that parameter is
small; no interchange with the infinite asymptotic series is needed.

Two dimensionless numerical checks of
\eqref{eq:J-closed} against the defining radial integral are
\begin{center}
\begin{tabular}{@{}ccr@{}}
\toprule
\(u\)&\(a\)&\(J(u;a)\), both evaluations\\
\midrule
\(1\)&\(1\)&\(0.147563809330584\)\\
\(2\)&\(0.4\)&\(0.086614103683222\)\\
\bottomrule
\end{tabular}
\end{center}
The representation by \(\erfcx\) controls exponential overflow,
but subtractive cancellation for large \(v\) should still be
handled with the integral or an appropriate finite endpoint
expansion.

\subsection{A common-domain justification of the Ward remainder}
\label{sec:Ward-cutoff-proof}

\begin{lemma}[Common-domain Ward translation]
\label{lem:Ward-common-cutoff}
Fix \(m>0\), \(a>0\), and \(Q\in\RR^4\), and use
\eqref{eq:QED-conventions}. Let
\[
\begin{gathered}
 F(u)=\Ph(a^2u)=\int_{[0,\infty)}e^{-a^2\tau u}\rho(\dd\tau),\\
 \rho\geq0,\qquad \rho([0,\infty))<\infty.
\end{gathered}
\]
Assume that the original tensor integrand
\[
 H_Q^{\mu\nu}(k)
 =F_kF_{k+Q}
   \Tr[\Gamma_0^\mu S_0(k)\Gamma_0^\nu S_0(k+Q)]
\]
is in \(L^1(\RR^4)\) componentwise. Then
\[
 \mathcal R_Q^\nu(k)
 =F_k[F_{k+Q}-F_{k-Q}]\,
       \Tr[\Gamma_0^\nu S_0(k)]
\]
is also in \(L^1(\RR^4)\), and
\eqref{eq:Ward-remainder} follows without assuming that the two
terms in the contracted difference are separately integrable.
\end{lemma}

\begin{proof}
Put \(T^\nu(k)=4k^\nu/(k^2+m^2)\).
First, absolute integrability of \(H_Q\) forces the radial tail
condition
\begin{equation}
 \int_{r_0}^{\infty}rF(r^2)^2\dd r<\infty
 \quad\hbox{for some }r_0>0.
 \label{eq:Ward-radial-tail}
\end{equation}
Indeed, write \(k=rn\), \(|n|=1\), and fix a diagonal component
\(\mu\), with no sum over \(\mu\). The trace has the uniform
large-\(r\) expansion
\[
 \Tr[\Gamma_0^\mu S_0(k)\Gamma_0^\mu S_0(k+Q)]
 =\frac{4(2n_\mu^2-1)}{r^2}+O_{Q,m}(r^{-3}).
\]
On the fixed angular region \(n_\mu^2\leq1/4\), its modulus is at
least \(c/r^2\) for sufficiently large \(r\).
Since \(F\) is nonnegative and nonincreasing,
\[
 F(r^2)F(|k+Q|^2)\geq F((r+|Q|)^2)^2.
\]
Integration on this angular region, followed by a bounded radial
shift, proves \eqref{eq:Ward-radial-tail}.
For the nonincreasing function \(f(r)=F(r^2)^2\),
\[
 \frac38R^2f(R)
 \leq\int_{R/2}^{R}rf(r)\dd r\longrightarrow0.
\]
Thus \(R^2F(R^2)^2\to0\).

The radius \(R\) below only exhausts the integration domain; it
is not a new physical regulator or kernel parameter.
Now let \(B_R=\{k:|k|<R\}\) and set
\begin{align*}
 I_R^\nu
 &=\int_{B_R}\frac{\dd^4k}{(2\pi)^4}
       F_kF_{k+Q}[T^\nu(k)-T^\nu(k+Q)],\\
 J_R^\nu
 &=\int_{B_R}\frac{\dd^4k}{(2\pi)^4}
       F_k[F_{k+Q}-F_{k-Q}]T^\nu(k),\\
 C_Q^\nu(k)&=F_kF_{k-Q}T^\nu(k).
\end{align*}
All these finite-domain integrals exist. Translating only within
the finite cutoff gives the exact identity
\begin{align}
 I_R^\nu&=J_R^\nu+E_R^\nu,\nonumber\\
 E_R^\nu
 &=\int_{B_R}\frac{\dd^4k}{(2\pi)^4}C_Q^\nu(k)
   -\int_{B_R+Q}\frac{\dd^4k}{(2\pi)^4}C_Q^\nu(k).
 \label{eq:Ward-finite-cutoff}
\end{align}
For sufficiently large \(R>4|Q|\), the symmetric difference of the balls lies in
a shell of volume \(O_Q(R^3)\), and \(|T^\nu(k)|\leq C/R\)
there. Both \(|k|\) and \(|k-Q|\) are at least
\(R-2|Q|\). Monotonicity therefore gives
\begin{equation}
 |E_R^\nu|
 \leq C_QR^2F((R-2|Q|)^2)^2
 \longrightarrow0.
 \label{eq:Ward-boundary-error}
\end{equation}

It remains to establish the absolute integrability of the final
remainder. For \(u>0\), the positive measure representation gives
\begin{equation}
 -uF'(u)
 =\int_{[0,\infty)}(a^2\tau u)e^{-a^2\tau u}\rho(\dd\tau)
 \leq\frac2{\mathrm e}F(u/2),
 \label{eq:Ward-CM-derivative}
\end{equation}
because \(x e^{-x/2}\leq2/\mathrm e\).
The mean-value theorem along \(k+tQ\), \(-1\leq t\leq1\),
then yields, for sufficiently large \(r=|k|>2|Q|\),
\[
 |F_{k+Q}-F_{k-Q}|
 \leq\frac{C_Q}{r}F(r^2/8).
\]
Using \(|T^\nu(k)|\leq4/r\), the absolute radial tail of
\(\mathcal R_Q^\nu\) is bounded by a constant times
\[
 \int^\infty rF(r^2)F(r^2/8)\dd r.
\]
This is finite by Cauchy--Schwarz, \eqref{eq:Ward-radial-tail},
and a constant rescaling of the radial variable in the second
factor. On compact momentum sets, \(m>0\) and \(F(0)<\infty\)
give local integrability. Hence \(J_R^\nu\) converges absolutely.
The algebraic Ward contraction identifies
\(-e^2\lim_{R\to\infty}I_R^\nu\) with
\(Q_\mu\Pi_{\rm bv}^{\mu\nu}\). Taking the limit in
\eqref{eq:Ward-finite-cutoff} proves
\eqref{eq:Ward-remainder}.
\end{proof}

For the unchanged canonical massless free expressions, the same
proof applies: \(T^\nu(k)=O(|k|^{-1})\), and the original free
tensor is bounded locally by
\(C/(|k|\,|k+Q|)\), with the coincident case \(Q=0\) of order
\(|k|^{-2}\). These singularities are locally integrable in four
Euclidean dimensions. An additional infrared prescription must
retain the algebraic Ward contraction and justify the finite-domain
changes of variables; an integrable comparison envelope alone
does not establish that compatibility.

The distinction between the original tensor and its split terms is
material. For \(F(u)=(1+a^2u)^{-\beta}\), \(m>0\), and
\(1/2<\beta\leq3/4\), the tensor has an integrable absolute
radial tail \(r^{1-4\beta}\), while each split term has a
nonintegrable absolute tail \(r^{2-4\beta}\).
The common-cutoff proof includes this range.

\paragraph{Relation to the monograph}
% Monograph source: eq:Pi-Omni-full; lem:ST-stability-internal-only-early.
The monograph distinguishes the bare-vertex tensor from the
completed transverse tensor and derives the latter's Ward
identity from the completed A5/Q1 functional
\cite{ChiapponeMonograph}. The lemma above adds the explicit
finite-cutoff and \(L^1\) justification for the bare-vertex
remainder; it is not a replacement for that action-level argument.
It uses precisely \(\Gamma_0^\mu=\ii\gamma^\mu\) and
\(S_0^{-1}=\ii\slashed k+m\) from
\eqref{eq:QED-conventions}.

\subsection{Stable-half covariance tail and bubble differentiability}
\label{sec:stable-bubble-proof}

Use the inverse Fourier convention of
Section~\ref{sec:stable-bubble-origin}. Subordination by the
stable-half density gives
\begin{align}
 P_a(x)
 &=\int_0^\infty
    \frac{a\,e^{-a^2/(4v)}}{2\sqrt\pi\,v^{3/2}}
    \frac{e^{-|x|^2/(4v)}}{(4\pi v)^2}\dd v\nonumber\\
 &=\frac{3a}{4\pi^2(a^2+|x|^2)^{5/2}}.
 \label{eq:four-dimensional-Poisson-kernel}
\end{align}
The second line follows by substituting
\(y=(a^2+|x|^2)/(4v)\) and using
\(\Gamma(5/2)=3\sqrt\pi/4\). This derives the coefficient
rather than importing a different Fourier normalization.
The positive massive free Green function is
\[
\begin{aligned}
 G_m(y)&=\int_0^\infty e^{-m^2v}
       (4\pi v)^{-2}e^{-|y|^2/(4v)}\dd v,\\
 \int_{\RR^4}G_m(y)\dd^4y&=m^{-2}.
\end{aligned}
\]
It has an exponentially decaying integral tail. For example,
\(c|y|\leq |y|^2/(8v)+2c^2v\) implies
\[
 \int e^{c|y|}G_m(y)\dd^4y
       \leq\frac{4}{m^2-2c^2}<\infty,
       \qquad 0<c<m/\sqrt2.
\]
Fourier multiplication therefore gives the positive convolution
\(\mathcal C_{a,m}=G_m*P_a\).

Put \(r=|x|\). On \(|y|\leq r/2\), the ratio
\(P_a(x-y)/P_a(x)\) is bounded uniformly for large \(r\),
and tends to one for every fixed \(y\). Dominated convergence
against \(G_m(y)\dd^4y\) gives the limiting ratio \(m^{-2}\)
on this region. On \(|y|>r/2\), bound \(P_a(x-y)\) by
\(\|P_a\|_\infty\) and use the exponentially small integral
tail of \(G_m\); division by \(P_a(x)\asymp r^{-5}\)
still leaves zero. This proves
\eqref{eq:stable-covariance-long-tail}. Also
\(\|\mathcal C_{a,m}\|_\infty\leq m^{-2}\|P_a\|_\infty\),
so there is no local coordinate singularity.

The momentum covariance is real, even and square-integrable.
Fourier convolution, with the stated \((2\pi)^{-4}\) convention,
gives the exact identity
\begin{equation}
 B_{a,m}(P)=\int_{\RR^4}e^{\ii P\cdot x}
                         \mathcal C_{a,m}(x)^2\dd^4x.
 \label{eq:stable-bubble-coordinate-square}
\end{equation}
Since \(\mathcal C_{a,m}(x)^2\asymp r^{-10}\) at infinity,
all its absolute moments of orders \(0\leq j<6\) are finite:
the radial tail is \(r^{j-7}\dd r\). Differentiation under
the integral proves \(C^5\) regularity. On an angular set with
\(|x_1|\geq r/2\), however, the sixth moment has the divergent
tail \(\int^\infty\dd r/r\). Thus
\(\int x_1^6\mathcal C_{a,m}(x)^2\dd^4x=\infty\).

For a proof that this forbids \(C^6\), put \(b(h)=B_{a,m}(h e_1)\)
and use the symmetric sixth difference
\[
 d_h^6b:=\sum_{j=0}^6(-1)^j\binom{6}{j}b((3-j)h).
\]
Equation~\eqref{eq:stable-bubble-coordinate-square} yields
\begin{equation}
 -\frac{d_h^6b}{h^6}
   =\int_{\RR^4}\frac{64\sin^6(hx_1/2)}{h^6}
                  \mathcal C_{a,m}(x)^2\dd^4x.
 \label{eq:stable-sixth-difference}
\end{equation}
The integrand is nonnegative and converges pointwise to
\(x_1^6\mathcal C_{a,m}(x)^2\). Fatou's lemma makes the
lower limit of the right side infinite. If \(b\) were \(C^6\)
near zero, its left side would tend to the finite number
\(-b^{(6)}(0)\), a contradiction. Holomorphy in \(P^2\)
around zero would in particular imply analyticity of \(b(h)\),
and is therefore also excluded. This completes the proof of
Proposition~\ref{prop:stable-bubble-nonanalytic}.

\section{Logical dependence of the principal results}
\label{sec:dependence}

The following dependency ledger records the inputs needed to combine the results.
In particular, the convergence and symmetry rows must refer to the
same completed action and the same graph family.
The adjoint and reconstruction assumptions concern physical source
channels and full correlation functions, not just their auxiliary
carrier coefficients.

\begin{description}[style=nextline]
\item[Plane binding.] \textit{Inputs:} Fixed labelled axial fibre, independent A2/GTP momentum binder,
and support-diagonal action; admitted canonical or source representation.
\textit{Conclusion:} Selection rule on the chosen support; positivity and scattering
unitarity require the physical completion.

\item[Carrier reconstruction.] \textit{Inputs:} Fixed pole carrier,
declared vectorial boundary limits, both character projections of
the same complete functional, and compatible functional-calculus branches.
\textit{Conclusion:} Carrier-valued lines and functionals; one complex coefficient suffices
only under the stated reflection or directed-subalgebra conditions.

\item[Reality and adjoints.] \textit{Inputs:} The declared carrier
reality \(\ddagger\); a separately specified physical adjoint
\(\dagger\). \textit{Conclusion:} Reality relations for coefficients
do not establish physical positivity. A positive unital representation
preserving \(\ddagger\) on the entire carrier algebra is excluded by
Proposition~\ref{prop:carrier-reality-adjoint}.

\item[Source covariance.] \textit{Inputs:} Graded source derivatives at fixed branch, source pairing, and
translation-invariant field components.
\textit{Conclusion:} Exchange symmetry of the complete source matrix. An odd oriented
running line is not an identical commuting scalar covariance.

\item[Global graph convergence.] \textit{Inputs:} Full-rank routing, completed numerator bound, gap or sufficient
inverse moments, and the stated IR control.
\textit{Conclusion:} Each finite graph at fixed \(a>0\) and admitted IR data;
singular parameter limits need separate bounds.

\item[Optimized line coverage.] \textit{Inputs:} Full-rank
selected routing, line-specific inverse moments, and an anisotropic
bound covering the complete numerator and unselected lines.
\textit{Conclusion:} Absolute convergence even with selected massless
$1/q^2$ denominators; stronger moments give explicit selected-variable
UV tails. Unselected massless singularities are not excused.

\item[Residual damping.]
\textit{Inputs:} A bound on the full nonpolynomial vertex product and
a strictly positive quadratic or linear damping margin, with the
stated bounded infrared covariances.
\textit{Conclusion:} Uniform absolute Euclidean convergence of the
complete graph. Failure of this sufficient test alone proves no
divergence and no continuation statement.

\item[O--QED kernel application.] \textit{Inputs:} The normalized kernel
\(F_{a,M}\), fixed \(a>0\), \(M\geq0\), and the same completed
vertices and infrared/source conventions as in the claimed amplitude.
\textit{Conclusion:} All inverse moments are verified and the bare
one-loop coefficient is finite; full graphwise finiteness still requires
vertex bounds. The Ward-compatible vacuum two-vertex term fails that
test before contacts. The declared constant-background contact test
also fails absolute integrability; its raw radial reference difference
has the obstruction in Corollary~\ref{cor:OQED-general-contact-obstruction}.
Other completions remain to be specified. A grouped result requires
its own common-routing envelope.
Medium sums require an explicit spectral/infrared bound, and the
kernel-off reference is separately renormalized and matched.

\item[Forest convergence.] \textit{Inputs:} Uniform simultaneous sector envelopes, every active UV scale
controlled, and IR/angle bounds.
\textit{Conclusion:} All verified sector scales; each damping gain is assigned through
a simultaneous routing estimate.

\item[Slavnov--Taylor identity.] \textit{Inputs:} Completed closed BRST/BV action in both characters, compatible
source/antifield maps and source covariances, invariant anomaly-free
measure, and admitted insertions.
\textit{Conclusion:} Formal perturbative identity for the completed functional;
its vertex growth is checked independently.

\item[Nielsen identity.] \textit{Inputs:} Extended BRST functional, admissible gauge-parameter insertion,
fixed kernel/source maps, and compatible physical normalization.
\textit{Conclusion:} Gauge independence for admitted fixed or Nielsen-transported physical
classes; off-shell functions can depend on the gauge parameter.

\item[Local holomorphy.] \textit{Inputs:} Holomorphic graph integrand and a common pointwise complex
\(L^1\) majorant.
\textit{Conclusion:} The declared complex domain; not a holomorphic neighbourhood of
the origin for the literal massive stable-half bubble. Boundary
growth and channel continuation require their own estimates.

\item[Depth limits and scalar subtraction.] \textit{Inputs:} The normalized
square-root kernel at positive depth and the literal scalar bubble,
with $m>0$ for the subtraction and strong-depth limits.
\textit{Conclusion:} Explicit depth derivatives, distinct strong-depth
regimes and crossover, a nonuniform undamped line limit, and a finite
$a\downarrow0$ limit for the \emph{already subtracted} scalar bubble.
The literal bubble is completely monotone in positive depth,
with an explicit uniform derivative bound and an exponentially
decaying O--QED line tail. Its unsubtracted zero-depth limit has
a universal logarithmic finite part with a quantitative error
bound. These are not interacting gauge or retarded-response results.

\item[Scalar approximation.] \textit{Inputs:} A literal scalar bubble
with \(a,M,m>0\) and its positive proper-time measure. \textit{Conclusion:} Holomorphy for
\(\Re w>-4\min(m,M)^2\), signed finite-order bounds for real
\(w\geq0\), and the separately stated complex Taylor and radial-tail
bounds. This is not physical spectral positivity or a completed gauge
response.

\item[Physical reconstruction.] \textit{Inputs:} Completed even-observable
full correlation functions satisfying Assumption~\ref{ass:OS-completion},
including \eqref{eq:OS-hierarchy-growth} uniformly in insertion
number. Odd physical fields require the corresponding graded
reconstruction. Finiteness graph by graph is not this hierarchy
condition or an order-by-order positivity claim.
\textit{Conclusion:} A positive physical reconstruction on the admitted
observable algebra; an individual virtual-line kernel need not be
a physical spectral covariance.

\item[Pole normalization.] \textit{Inputs:} An isolated pole of the
full inverse in a compatible physical source channel; the actual tree
inverse and loop correction; positive one-particle spectral overlaps.
\textit{Conclusion:} Finite normalization on the nonzero residue
subspace, with an invertible field map and its dual source map when
all field coordinates are retained. Neither \(\Ph(0)=1\) nor an
arbitrary nonzero complex Laurent coefficient establishes positivity.

\item[Cut matching and optical relation.] \textit{Inputs:} Normalized physical states and completed subamplitudes; cut--jump
matching. The optical relation also requires unitary scattering
and a complete state resolution, or an equivalent Hermitian-analytic
cut-factorization argument.
\textit{Conclusion:} Completed amplitudes with their actual residues and branch terms;
positive norms and support diagonality alone do not imply the
optical identity.

\item[Gaussian shell normalization.] \textit{Inputs:} The exact fixed-carrier Gaussian bubble and its finite reference
subtraction.
\textit{Conclusion:} The usual normalized jump plus an entire remainder; the remainder
has nonpolynomial positive-real-axis growth.

\item[Response and pole errors.] \textit{Inputs:} A certified complex-domain
response error, compatible operator/domain norms, and an invertible
approximate electromagnetic response with a strict Neumann margin.
The endpoint root enclosure additionally requires the holomorphic
retained matrix pencils and contour isolation in
Proposition~\ref{prop:OQED-finite-pole-enclosure}. Interpreting its root
as the selected physical pole requires baseline-branch identification;
numerical contour samples require between-node bounds, and a compression
of a larger theory requires control of the discarded complement.
\textit{Conclusion:} Schur and first-order pole-family error bounds;
under the respective additional hypotheses, an endpoint root enclosure
and a bound on the selected finite rate. A complex pole-family margin
supplies the explicit Taylor remainder in
\eqref{eq:OQED-pole-remainder-certificate}; an observed response pole
also needs the nonzero input/output residue in
\eqref{eq:OQED-observable-pole-residue}. Its nonvanishing can be
certified by \eqref{eq:OQED-observable-residue-error}; eliminating a
complement also transports the endpoint maps as in
\eqref{eq:OQED-Schur-source-readout}. A real Euclidean error alone
supplies neither conclusion at an unverified complex continuation point.

\item[Physical crossing.] \textit{Inputs:} Local physical reconstruction,
suitable asymptotic poles and LSZ limits, and controlled continuation
between the channels, with the physical boundary-value identification
established independently of the domain result. Invocation of the cited
Bros--Epstein--Glaser domain theorem also requires its stability and
strictly positive physical mass-gap hypotheses.
\textit{Conclusion:} Crossing for the admitted scattering channels
and asymptotic states; the continuation must account for kernel
and vertex singularities.
\end{description}

\section{Relation to the foundational constructions}
\label{sec:monograph-differences}

The locked-plane geometry, directed vectorial boundary prescription,
two-character carrier reconstruction, and original Gaussian routing
proof follow the framework of Ref.~\cite{ChiapponeMonograph}.
The present article states its hypotheses independently and develops
additional sufficient Euclidean convergence criteria: the endpoint
inverse-moment condition, line-specific full-rank coverage with
massless selected denominators, residual coercivity for some
growing vertices, and explicit tail and complex-domain estimates.
The two-loop comparison following
Theorem~\ref{thm:optimized-selected-lines}, together with
Proposition~\ref{prop:two-loop-complex-certificate}, is a
mathematical test of those criteria, not a graph attributed to the
physical action. Its bounds refer to the explicit nonnegative
selected-line integrand, not to an independently completed
O--QED Feynman graph.

Several calculations delimit the use of the auxiliary covariances.
Source exchange symmetry prevents an unrestricted odd running-unit
line from being identified with the covariance of one identical
commuting real field without a compatible endpoint/source matrix.
Carrier reality $\ddagger$ is not the positive physical adjoint.
The Gaussian bubble, massive stable-half bubble, and reflection-
positivity counterexample separately test discontinuities, low-
momentum analyticity, and physical reconstruction. None changes
the original vectorial pole datum or identifies Euclidean damping
with a physical spectral representation.

The action-level BRST/BV and Nielsen identities require a compatible
completed action, source measures, and insertions. The common-domain
bare QED Ward argument in Appendix~\ref{sec:Ward-cutoff-proof}
adds an explicit absolute-integration justification; it does not
replace the monograph's completed-action Ward argument.
The Euclidean $\Gamma_0^\mu=\ii\gamma^\mu$ convention used in
\eqref{eq:QED-conventions} is stated separately from the
$\gamma^\mu$-vertex display in the reference, so their overall
prefactors must be converted before comparison.

For the standalone O--QED example, the action and kernel input come
from Ref.~\cite{ChiapponeIgnition}, especially its
Eqs.~(48)--(53), (66), and (75)--(83).
The inverse moments, depth-limit formulas, scalar comparison bounds,
and explicit Ward-bubble obstruction are calculated here. The
constant-background contact test additionally imposes a shift
covariantization, a trace-log realization, and a field-independent
measure; it is not an evaluation of the source's unspecified full
contact completion. The nonintegrability of these particular
terms neither proves divergence of every possible completion nor
supplies a missing combined-graph estimate.

The pole and readout certificates use the stated simple-root,
Schur-complement, and common-contour assumptions; the local inverse
residue invokes the standard result of Ref.~\cite{Beyn2012}.
Reflection positivity and crossing retain their separate OS and
physical-boundary hypotheses \cite{OS1975,BrosEpsteinGlaser1965}.
Collectively the mathematical results establish conditional
convergence and diagnostic obstructions. A single interacting
theory satisfying all graph, symmetry, source, and physical
reconstruction requirements has not been constructed or excluded.

\begingroup
\raggedright
\begin{thebibliography}{99}

\bibitem{ChiapponeMonograph}
G.~J.~Chiappone,
\emph{Groupoid Hyper--Omniverse Unified Quantum Field Theory},
author-supplied working revision dated 6 November 2026,
used for the A5/Q1 construction and fixed-plane, UV, BRST, and
analyticity conventions. Related Zenodo identifier supplied by the author:
\href{https://doi.org/10.5281/zenodo.21942369}{doi:\nolinkurl{10.5281/zenodo.21942369}}
(exact deposited version/file correspondence not independently verified).
\bibitem{ChiapponeIgnition}
G.~J.~Chiappone,
\emph{Omniverse QED geo--phase damping in the hydrodynamic window:
a first--principles, hypothesis--scoped fixed--depth response framework
for Rayleigh--Taylor damping in MagLIF},
author-supplied 185-page manuscript dated 8 August 2026,
used here for Eqs.~(48)--(53), (66), (75)--(83), and (99)--(118).
Related Zenodo identifier supplied by the author:
\href{https://doi.org/10.5281/zenodo.20371104}{doi:\nolinkurl{10.5281/zenodo.20371104}}
(exact deposited version/file correspondence not independently verified).
\bibitem{DLMF10}
\emph{NIST Digital Library of Mathematical Functions},
Release 1.2.8 (15 September 2026), F.~W.~J. Olver \emph{et al.} (eds.),
\S10.32, ``Integral Representations,'' especially
\href{https://dlmf.nist.gov/10.32.E10}{Eq.~10.32.10}
(accessed 6 November 2026).
\bibitem{Weinberg1960}
S.~Weinberg,
High-energy behavior in quantum field theory,
\emph{Phys.\ Rev.}\ \textbf{118} (1960) 838--849.
\href{https://doi.org/10.1103/PhysRev.118.838}{doi:\nolinkurl{10.1103/PhysRev.118.838}}.
\bibitem{Abbott1981}
L.~F.~Abbott,
The background field method beyond one loop,
\emph{Nucl.\ Phys.\ B} \textbf{185} (1981) 189--203.
\href{https://doi.org/10.1016/0550-3213(81)90371-0}{doi:\nolinkurl{10.1016/0550-3213(81)90371-0}}.
\bibitem{BRST1976}
C.~Becchi, A.~Rouet, and R.~Stora,
Renormalization of gauge theories,
\emph{Ann.\ Phys.}\ \textbf{98} (1976) 287--321.
\href{https://doi.org/10.1016/0003-4916(76)90156-1}{doi:\nolinkurl{10.1016/0003-4916(76)90156-1}}.
\bibitem{Tyutin1975}
I.~V.~Tyutin,
\emph{Gauge Invariance in Field Theory and Statistical Physics
in Operator Formalism}, P.~N.~Lebedev Physical Institute preprint
No.~39 (1975); arXiv reproduction (latest revision, v2,
15 December 2008),
\href{https://arxiv.org/abs/0812.0580v2}{arXiv:0812.0580v2}.
\bibitem{PiguetSorella}
O.~Piguet and S.~P.~Sorella,
\emph{Algebraic Renormalization: Perturbative Renormalization,
Symmetries and Anomalies}, Lecture Notes in Physics Monographs,
vol.~28, Springer, Berlin/Heidelberg (original edition 1995;
eBook published 11 September 2008).
\href{https://doi.org/10.1007/978-3-540-49192-7}{doi:\nolinkurl{10.1007/978-3-540-49192-7}}.
\bibitem{Slavnov1972}
A.~A.~Slavnov,
Ward identities in gauge theories,
\emph{Theor.\ Math.\ Phys.}\ \textbf{10} (1972) 99--104
(English translation of \emph{Teor.\ Mat.\ Fiz.} \textbf{10}
(1972) 153--161).
\href{https://doi.org/10.1007/BF01090719}{doi:\nolinkurl{10.1007/BF01090719}}.
\bibitem{Taylor1971}
J.~C.~Taylor,
Ward identities and charge renormalization of the Yang--Mills field,
\emph{Nucl.\ Phys.\ B} \textbf{33} (1971) 436--444.
\href{https://doi.org/10.1016/0550-3213(71)90297-5}{doi:\nolinkurl{10.1016/0550-3213(71)90297-5}}.
\bibitem{ZinnJustin1975}
J.~Zinn-Justin,
Renormalization of gauge theories,
in H.~Rollnik and K.~Dietz (eds.),
\emph{Trends in Elementary Particle Theory},
Lecture Notes in Physics, vol.~37, Springer, Berlin/Heidelberg
(1975), pp.~1--39.
\href{https://doi.org/10.1007/3-540-07160-1_1}{doi:\nolinkurl{10.1007/3-540-07160-1_1}}.
\bibitem{Nielsen1975}
N.~K.~Nielsen,
On the gauge dependence of spontaneous symmetry breaking in gauge
theories, \emph{Nucl.\ Phys.\ B} \textbf{101} (1975) 173--188.
\href{https://doi.org/10.1016/0550-3213(75)90301-6}{doi:\nolinkurl{10.1016/0550-3213(75)90301-6}}.
\bibitem{Bernstein1929}
S.~Bernstein,
Sur les fonctions absolument monotones,
\emph{Acta Math.}\ \textbf{52} (1929) 1--66.
\href{https://doi.org/10.1007/BF02592679}{doi:\nolinkurl{10.1007/BF02592679}}.
\bibitem{Widder1941}
D.~V.~Widder,
\emph{The Laplace Transform}, Princeton Mathematical Series, vol.~6,
Princeton University Press (original edition 1941;
Princeton Legacy Library eBook, 25 July 2016).
\href{https://doi.org/10.1515/9781400876457}{doi:\nolinkurl{10.1515/9781400876457}}.
\bibitem{Landau1959}
L.~D.~Landau,
On analytic properties of vertex parts in quantum field theory,
\emph{Nucl.\ Phys.}\ \textbf{13} (1959) 181--192.
\href{https://doi.org/10.1016/0029-5582(59)90154-3}{doi:\nolinkurl{10.1016/0029-5582(59)90154-3}}.
\bibitem{OS1973}
K.~Osterwalder and R.~Schrader,
Axioms for Euclidean Green's functions,
\emph{Commun.\ Math.\ Phys.}\ \textbf{31} (1973) 83--112.
\href{https://doi.org/10.1007/BF01645738}{doi:\nolinkurl{10.1007/BF01645738}}.
\bibitem{OS1975}
K.~Osterwalder and R.~Schrader,
Axioms for Euclidean Green's functions II
(with an appendix by Stephen Summers),
\emph{Commun.\ Math.\ Phys.}\ \textbf{42} (1975) 281--305.
\href{https://doi.org/10.1007/BF01608978}{doi:\nolinkurl{10.1007/BF01608978}}.
\bibitem{Arici2018}
F.~Arici, D.~Becker, C.~Ripken, F.~Saueressig,
and W.~D.~van Suijlekom,
Reflection positivity in higher derivative scalar theories,
\emph{J.\ Math.\ Phys.}\ \textbf{59} (2018) 082302.
\href{https://doi.org/10.1063/1.5027231}{doi:\nolinkurl{10.1063/1.5027231}};
\href{https://arxiv.org/abs/1712.04308v1}{arXiv:1712.04308v1}
(latest arXiv version, 12 December 2017).
\bibitem{StreaterWightman}
R.~F.~Streater and A.~S.~Wightman,
\emph{PCT, Spin and Statistics, and All That},
original edition, W.~A. Benjamin (1964); additions and corrections
(1978); Princeton Landmarks in Mathematics and Physics,
Princeton University Press paperback with revised preface and
corrections (2000); electronic edition 11 August 2016.
\href{https://doi.org/10.1515/9781400884230}{doi:\nolinkurl{10.1515/9781400884230}}.
\bibitem{LSZ1955}
H.~Lehmann, K.~Symanzik, and W.~Zimmermann,
Zur Formulierung quantisierter Feldtheorien,
\emph{Nuovo Cimento}\ \textbf{1} (1955) 205--225.
\href{https://doi.org/10.1007/BF02731765}{doi:\nolinkurl{10.1007/BF02731765}}.
\bibitem{Cutkosky1960}
R.~E.~Cutkosky,
Singularities and discontinuities of Feynman amplitudes,
\emph{J.\ Math.\ Phys.}\ \textbf{1} (1960) 429--433.
\href{https://doi.org/10.1063/1.1703676}{doi:\nolinkurl{10.1063/1.1703676}}.
\bibitem{BrosEpsteinGlaser1965}
J.~Bros, H.~Epstein, and V.~Glaser,
A proof of the crossing property for two-particle amplitudes in
general quantum field theory,
\emph{Commun.\ Math.\ Phys.}\ \textbf{1} (1965) 240--264.
\href{https://doi.org/10.1007/BF01646307}{doi:\nolinkurl{10.1007/BF01646307}}.
\bibitem{Schwinger1951}
J.~Schwinger,
On gauge invariance and vacuum polarization,
\emph{Phys.\ Rev.}\ \textbf{82} (1951) 664--679.
\href{https://doi.org/10.1103/PhysRev.82.664}{doi:\nolinkurl{10.1103/PhysRev.82.664}}.
\bibitem{DLMF6}
\emph{NIST Digital Library of Mathematical Functions},
Release 1.2.8 (15 September 2026), F.~W.~J. Olver \emph{et al.} (eds.),
Chapter~6 (exponential integrals: definitions, continuation,
and asymptotics); see
\href{https://dlmf.nist.gov/6.2}{\S6.2},
\href{https://dlmf.nist.gov/6.4}{\S6.4},
and \href{https://dlmf.nist.gov/6.12}{\S6.12}
(accessed 6 November 2026).
\bibitem{Beyn2012}
W.-J.~Beyn,
An integral method for solving nonlinear eigenvalue problems,
\emph{Linear Algebra Appl.}\ \textbf{436} (2012) 3839--3863.
\href{https://doi.org/10.1016/j.laa.2011.03.030}{doi:\nolinkurl{10.1016/j.laa.2011.03.030}};
\href{https://arxiv.org/abs/1003.1580v1}{arXiv:1003.1580v1}
(latest arXiv version, 8 March 2010).
\bibitem{BallChiu1980}
J.~S.~Ball and T.-W.~Chiu,
Analytic properties of the vertex function in gauge theories.~I,
\emph{Phys.\ Rev.\ D}\ \textbf{22} (1980) 2542--2549.
\href{https://doi.org/10.1103/PhysRevD.22.2542}{doi:\nolinkurl{10.1103/PhysRevD.22.2542}}.
\end{thebibliography}
\endgroup
\end{document}
